\chapter{Intel® Embree Overview}\label{intel-embree-overview}

Intel® Embree is a high-performance ray tracing library developed at
Intel, which is released as open source under the
\href{http://www.apache.org/licenses/LICENSE-2.0}{Apache 2.0 license}.
Intel® Embree supports x86 CPUs under Linux, macOS, and Windows; ARM
CPUs on Linux and macOS; as well as Intel® GPUs under Linux and Windows.

Intel® Embree targets graphics application developers to improve the
performance of photo-realistic rendering applications. Embree is
optimized towards production rendering, by putting focus on incoherent
ray performance, high quality acceleration structure construction, a
rich feature set, accurate primitive intersection, and low memory
consumption.

Embree's feature set includes various primitive types such as triangles
(as well quad and grids for lower memory consumption); Catmull-Clark
subdivision surfaces; various types of curve primitives, such as flat
curves (for distant views), round curves (for closeup views), and normal
oriented curves, all supported with different basis functions (linear,
Bézier, B-spline, Hermite, and Catmull Rom); point-like primitives, such
as ray oriented discs, normal oriented discs, and spheres; user defined
geometries with a procedural intersection function; multi-level
instancing; filter callbacks invoked for any hit encountered; motion
blur including multi-segment motion blur, deformation blur, and
quaternion motion blur; and ray masking.

Intel® Embree contains ray tracing kernels optimized for the latest x86
processors with support for SSE, AVX, AVX2, and AVX-512 instructions,
and uses runtime code selection to choose between these kernels. Intel®
Embree contains algorithms optimized for incoherent workloads
(e.g.~Monte Carlo ray tracing algorithms) and coherent workloads
(e.g.~primary visibility and hard shadow rays) as well as supports for
dynamic scenes by implementing high-performance two-level spatial index
structure construction algorithms.

Intel® Embree supports applications written with the Intel® Implicit
SPMD Program Compiler (Intel® ISPC, \url{https://ispc.github.io/}) by
providing an ISPC interface to the core ray tracing algorithms. This
makes it possible to write a renderer that automatically vectorizes and
leverages SSE, AVX, AVX2, and AVX-512 instructions.

Intel® Embree supports Intel GPUs through the
\href{https://www.khronos.org/sycl/}{SYCL} open standard programming
language. SYCL allows to write C++ code that can be run on various
devices, such as CPUs and GPUs. Using Intel® Embree application
developers can write a single source renderer that executes efficiently
on CPUs and GPUs. Maintaining just one code base this way can
significantly improve productivity and eliminate inconsistencies between
a CPU and GPU version of the renderer. Embree supports GPUs based on the
Xe HPG and Xe HPC microarchitecture, which support hardware accelerated
ray tracing do deliver excellent levels of ray tracing performance.

\section{Supported Platforms}\label{supported-platforms}

Embree supports Windows (32-bit and 64-bit), Linux (64-bit), and macOS
(64-bit). Under Windows, Linux and macOS x86 based CPUs are supported,
while ARM CPUs are currently only supported under Linux and macOS (e.g.~
Apple M1). ARM support for Windows experimental.

Embree supports Intel GPUs based on the Xe HPG microarchitecture (Intel®
Arc™ GPU) under Linux and Windows and Xe HPC microarchitecture (Intel®
Data Center GPU Flex Series and Intel® Data Center GPU Max Series) under
Linux.

The code compiles with the Intel® Compiler, Intel® oneAPI DPC++
Compiler, GCC, Clang, and the Microsoft Compiler. To use Embree on the
GPU the Intel® oneAPI DPC++ Compiler must be used. Please see section
\hyperref[compiling-embree]{Compiling Embree} for details on tested
compiler versions.

Embree requires at least an x86 CPU with support for SSE2 or an Apple M1
CPU.

\section{Embree Support and Contact}\label{embree-support-and-contact}

If you encounter bugs please report them via
\href{https://github.com/embree/embree/issues}{Embree's GitHub Issue
Tracker}.

For questions and feature requests please write us at
\href{mailto:embree_support@intel.com}{\nolinkurl{embree\_support@intel.com}}.

To receive notifications of updates and new features of Embree please
subscribe to the \href{https://groups.google.com/d/forum/embree/}{Embree
mailing list}.

\section{Version History}\label{version-history}

\subsection{Embree 4.1.0}\label{embree-4.1.0}

\begin{itemize}
\tightlist
\item
  Added support for Intel® Data Center GPU Max Series.
\item
  Added ARM64 Linux support.
\item
  Added EMBREE\_BACKFACE\_CULLING\_SPHERES cmake option. The new cmake
  option defaults to OFF.
\end{itemize}

\subsection{Embree 4.0.1}\label{embree-4.0.1}

\begin{itemize}
\tightlist
\item
  Improved performance for Tiger Lake, Comet Lake, Cannon Lake, Kaby
  Lake, and Skylake client CPUs by using 256 bit SIMD instructions by
  default.
\item
  Fixed broken motion blur of RTC\_GEOMETRY\_TYPE\_ROUND\_LINEAR\_CURVE
  geometry type.
\item
  Fixed bvh build retry issue for TBB 2020.3
\item
  Added support for Intel® Data Center GPU Flex Series
\item
  Fixed issue on systems without a SYCL platform.
\end{itemize}

\subsection{Embree 4.0.0}\label{embree-4.0.0}

\begin{itemize}
\tightlist
\item
  This Embree release adds support for Intel® Arc™ GPUs through SYCL.
\item
  The SYCL support of Embree is in beta phase. Current functionality,
  quality, and GPU performance may not reflect that of the final
  product. Please read the documentation section ``Embree SYCL Known
  Issues'' for known limitations.
\item
  Embree CPU support in this release as at Gold level, incorporating the
  same quality and performance as previous releases.
\item
  A small number of API changes were required to get optimal experience
  and performance on the CPU and GPU. See documentation section
  ``Upgrading from Embree 3 to Embree 4'' for details.
\item
  rtcIntersect and rtcOccluded function arguments changed slightly.
\item
  RTCIntersectContext is renamed to RTCRayQuery context and most members
  moved to new RTCIntersectArguments and RTCOccludedArguments
  structures.
\item
  rtcFilterIntersection and rtcFilterOcclusion API calls got replaced by
  rtcInvokeIntersectFilterFromGeometry and
  rtcInvokeOccludedFilterFromGeometry API calls.
\item
  rtcSetGeometryEnableFilterFunctionFromArguments enables argument
  filter functions for some geometry.
\item
  RTC\_RAY\_QUERY\_FLAG\_INVOKE\_ARGUMENT\_FILTER ray query flag enables
  argument filter functions for each geometry.
\item
  User geometry callbacks have to return if a valid hit was found.
\item
  Ray masking is enabled by default now as required by most users.
\item
  The default ray mask for geometries got changed from 0xFFFFFFFF to
  0x1.
\item
  Removed ray stream API as rarely used with minimal performance
  benefits over packet tracing.
\item
  Introduced rtcForwardIntersect/rtcForwardOccluded API calls to trace
  tail recursive rays from user geometry callback.
\item
  The rtcGetGeometryUserDataFromScene API call got added to be used in
  SYCL code.
\item
  Added support for user geometry callback function pointer passed
  through ray query context
\item
  Feature flags enable reducing code complexity for optimal performance
  on the GPU.
\item
  Fixed compilation issues for ARM AArch64 processor under Linux.
\item
  Setting default frequency level to SIMD256 for ARM on all platforms.
  This allows using double pumped NEON execution by enabling
  EMBREE\_ISA\_NEON2X in cmake under Linux.
\item
  Fixed missing end caps of motion blurred line segments.
\item
  EMBREE\_ISPC\_SUPPORT is turned OFF by default.
\item
  Embree drops support of the deprecated Intel(R) Compiler. It is
  replaced by the Intel(R) oneAPI DPC++/C++ Compiler on Windows and
  Linux and the Intel(R) C++ Classic Compiler on MacOS (latest tested
  versions is 2023.0.0).
\end{itemize}

\subsection{Embree 3.13.5}\label{embree-3.13.5}

\begin{itemize}
\tightlist
\item
  Fixed bug in bounding flat Catmull Rom curves of subdivision level 4.
\item
  Improved self intersection avoidance for
  RTC\_GEOMETRY\_TYPE\_DISC\_POINT geometry type. Intersections are
  skipped if the ray origin lies inside the sphere defined by the point
  primitive. Self intersection avoidance can get disabled at compile
  time using the EMBREE\_DISC\_POINT\_SELF\_INTERSECTION\_AVOIDANCE
  cmake option.
\item
  Fixed spatial splitting for non-planar quads.
\end{itemize}

\subsection{Embree 3.13.4}\label{embree-3.13.4}

\begin{itemize}
\tightlist
\item
  Using 8-wide BVH and double pumped NEON instructions on Apple M1 gives
  8\% performance boost.
\item
  Fixed binning related crash in SAH BVH builder.
\item
  Added EMBREE\_TBB\_COMPONENT cmake option to define the
  component/library name of Intel® TBB (default: tbb).
\item
  Embree supports now Intel® oneAPI DPC++/C++ Compiler 2022.0.0
\end{itemize}

\subsection{Embree 3.13.3}\label{embree-3.13.3}

\begin{itemize}
\tightlist
\item
  Invalid multi segment motion blurred normal oriented curves are
  properly excluded from BVH build.
\item
  Fixing issue with normal oriented curve construction when center curve
  curvature is very large. Due to this change normal oriented curve
  shape changes slightly.
\item
  Fixed crash caused by disabling a geometry and then detaching it from
  the scene.
\item
  Bugfix in emulated ray packet intersection when EMBREE\_RAY\_PACKETS
  is turned off.
\item
  Bugfix for linear quaternion interpolation fallback.
\item
  Fixed issues with spaces in path to Embree build folder.
\item
  Some fixes to compile Embree in SSE mode using WebAssembly.
\item
  Bugfix for occlusion rays with grids and ray packets.
\item
  We do no longer provide installers for Windows and macOS, please use
  the ZIP files instead.
\item
  Upgrading to Intel® ISPC 1.17.0 for release build.
\item
  Upgrading to Intel® oneTBB 2021.5.0 for release build.
\end{itemize}

\subsection{Embree 3.13.2}\label{embree-3.13.2}

\begin{itemize}
\tightlist
\item
  Avoiding spatial split positions that are slightly out of geometry
  bounds.
\item
  Introduced rtcGetGeometryThreadSafe function, which is a thread safe
  version of rtcGetGeometry.
\item
  Using more accurate rcp implementation.
\item
  Bugfix to rare corner case of high quality BVH builder.
\end{itemize}

\subsection{Embree 3.13.1}\label{embree-3.13.1}

\begin{itemize}
\tightlist
\item
  Added support for Intel® ISPC ARM target.
\item
  Releases upgrade to Intel® TBB 2021.3.0 and Intel® ISPC 1.16.1
\end{itemize}

\subsection{Embree 3.13.0}\label{embree-3.13.0}

\begin{itemize}
\tightlist
\item
  Added support for Apple M1 CPUs.
\item
  RTC\_SUBDIVISION\_MODE\_NO\_BOUNDARY now works properly for
  non-manifold edges.
\item
  CMake target `uninstall' is not defined if it already exists.
\item
  Embree no longer reads the .embree3 config files, thus all
  configuration has to get passed through the config string to
  rtcNewDevice.
\item
  Releases upgrade to Intel® TBB 2021.2.0 and Intel® ISPC 1.15.0
\item
  Intel® TBB dll is automatically copied into build folder after build
  on windows.
\end{itemize}

\subsection{Embree 3.12.2}\label{embree-3.12.2}

\begin{itemize}
\tightlist
\item
  Fixed wrong uv and Ng for grid intersector in robust mode for AVX.
\item
  Removed optimizations for Knights Landing.
\item
  Upgrading release builds to use Intel® oneTBB 2021.1.1
\end{itemize}

\subsection{Embree 3.12.1}\label{embree-3.12.1}

\begin{itemize}
\tightlist
\item
  Changed default frequency level to SIMD128 for Skylake, Cannon Lake,
  Comet Lake and Tiger Lake CPUs. This change typically improves
  performance for renderers that just use SSE by maintaining higher CPU
  frequencies. In case your renderer is AVX optimized you can get higher
  ray tracing performance by configuring the frequency level to simd256
  through passing frequency\_level=simd256 to rtcNewDevice.
\end{itemize}

\subsection{Embree 3.12.0}\label{embree-3.12.0}

\begin{itemize}
\tightlist
\item
  Added linear cone curve geometry support. In this mode a real
  geometric surface for curves with linear basis is rendered using
  capped cones. They are discontinuous at edge boundaries.
\item
  Enabled fast two level builder for instances when low quality build is
  requested.
\item
  Bugfix for BVH build when geometries got disabled.
\item
  Added EMBREE\_BACKFACE\_CULLING\_CURVES cmake option. This allows for
  a cheaper round linear curve intersection when correct internal
  tracking and back hits are not required. The new cmake option defaults
  to OFF.
\item
  User geometries with invalid bounds with lower\textgreater upper in
  some dimension will be ignored.
\item
  Increased robustness for grid interpolation code and fixed returned
  out of range u/v coordinates for grid primitive.
\item
  Fixed handling of motion blur time range for sphere, discs, and
  oriented disc geometries.
\item
  Fixed missing model data in releases.
\item
  Ensure compatibility to newer versions of Intel® oneTBB.
\item
  Motion blur BVH nodes no longer store NaN values.
\end{itemize}

\subsection{Embree 3.11.0}\label{embree-3.11.0}

\begin{itemize}
\tightlist
\item
  Round linear curves now automatically check for the existence of left
  and right connected segments if the flags buffer is empty. Left
  segments exist if the segment(id-1) + 1 == segment(id) and similarly
  for right segments.
\item
  Implemented the min-width feature for curves and points, which allows
  to increase the radius in a distance dependent way, such that the
  curve or points thickness is n pixels wide.
\item
  Round linear curves are closed now also at their start.
\item
  Embree no longer supports Visual Studio 2013 starting with this
  release.
\item
  Bugfix in subdivision tessellation level assignment for non-quad base
  primitives
\item
  Small meshes are directly added to top level build phase of two-level
  builder to reduce memory consumption.
\item
  Enabled fast two level builder for user geometries when low quality
  build is requested.
\end{itemize}

\subsection{Embree 3.10.0}\label{embree-3.10.0}

\begin{itemize}
\tightlist
\item
  Added EMBREE\_COMPACT\_POLYS CMake option which enables double indexed
  triangle and quad leaves to reduce memory consumption in compact mode
  by an additional 40\% at about 15\% performance impact. This new mode
  is disabled by default.
\item
  Compile fix for Intel® oneTBB 2021.1-beta05
\item
  Releases upgrade to Intel® TBB 2020.2
\item
  Compile fix for Intel® ISPC v1.13.0
\item
  Adding RPATH to libembree.so in releases
\item
  Increased required CMake version to 3.1.0
\item
  Made instID member for array of pointers ray stream layout optional
  again.
\end{itemize}

\subsection{Embree 3.9.0}\label{embree-3.9.0}

\begin{itemize}
\tightlist
\item
  Added round linear curve geometry support. In this mode a real
  geometric surface for curves with linear basis is rendered using
  capped cones with spherical filling between the curve segments.
\item
  Added rtcGetSceneDevice API function, that returns the device a scene
  got created in.
\item
  Improved performance of round curve rendering by up to 1.8x.
\item
  Bugfix to sphere intersection filter invocation for back hit.
\item
  Fixed wrong assertion that triggered for invalid curves which anyway
  get filtered out.
\item
  RelWithDebInfo mode no longer enables assertions.
\item
  Fixed an issue in FindTBB.cmake that caused compile error with Debug
  build under Linux.
\item
  Embree releases no longer provide RPMs for Linux. Please use the RPMs
  coming with the package manager of your Linux distribution.
\end{itemize}

\subsection{Embree 3.8.0}\label{embree-3.8.0}

\begin{itemize}
\tightlist
\item
  Added collision detection support for user geometries (see rtcCollide
  API function)
\item
  Passing geomID to user geometry callbacks.
\item
  Bugfix in AVX512VL codepath for rtcIntersect1
\item
  For sphere geometries the intersection filter gets now invoked for
  front and back hit.
\item
  Fixed some bugs for quaternion motion blur.
\item
  RTCRayQueryContext always non-const in Embree API
\item
  Made RTCHit aligned to 16 bytes in Embree API
\end{itemize}

\subsection{New Features in Embree
3.7.0}\label{new-features-in-embree-3.7.0}

\begin{itemize}
\tightlist
\item
  Added quaternion motion blur for correct interpolation of rotational
  transformations.
\item
  Fixed wrong bounding calculations when a motion blurred instance did
  instantiate a motion blurred scene.
\item
  In robust mode the depth test consistently uses tnear \textless= t
  \textless= tfar now in order to robustly continue traversal at a
  previous hit point in a way that guarantees reaching all hits, even
  hits at the same place.
\item
  Fixed depth test in robust mode to be precise at tnear and tfar.
\item
  Added next\_hit tutorial to demonstrate robustly collecting all hits
  along a ray using multiple ray queries.
\item
  Implemented robust mode for curves. This has a small performance
  impact but fixes bounding problems with flat curves.
\item
  Improved quality of motion blur BVH by using linear bounds during
  binning.
\item
  Implemented issue with motion blur builder where number of time
  segments for SAH heuristic were counted wrong due to some numerical
  issues.
\item
  Fixed an accuracy issue with rendering very short fat curves.
\item
  rtcCommitScene can now get called during rendering from multiple
  threads to lazily build geometry. When Intel® TBB is used this causes
  a much lower overhead than using rtcJoinCommitScene.
\item
  Geometries can now get attached to multiple scenes at the same time,
  which simplifies mapping general scene graphs to API.
\item
  Updated to Intel® TBB 2019.9 for release builds.
\item
  Fixed a bug in the BVH builder for Grid geometries.
\item
  Added macOS Catalina support to Embree releases.
\end{itemize}

\subsection{New Features in Embree
3.6.1}\label{new-features-in-embree-3.6.1}

\begin{itemize}
\tightlist
\item
  Restored binary compatibility between Embree 3.6 and 3.5 when
  single-level instancing is used.
\item
  Fixed bug in subgrid intersector
\item
  Removed point query alignment in Intel® ISPC header
\end{itemize}

\subsection{New Features in Embree
3.6}\label{new-features-in-embree-3.6}

\begin{itemize}
\tightlist
\item
  Added Catmull-Rom curve types.
\item
  Added support for multi-level instancing.
\item
  Added support for point queries.
\item
  Fixed a bug preventing normal oriented curves being used unless
  timesteps were specified.
\item
  Fixed bug in external BVH builder when configured for dynamic build.
\item
  Added support for new config flag ``user\_threads=N'' to device
  initialization which sets the number of threads used by Intel® TBB but
  created by the user.
\item
  Fixed automatic vertex buffer padding when using rtcSetNewGeometry API
  function.
\end{itemize}

\subsection{New Features in Embree
3.5.2}\label{new-features-in-embree-3.5.2}

\begin{itemize}
\tightlist
\item
  Added EMBREE\_API\_NAMESPACE cmake option that allows to put all
  Embree API functions inside a user defined namespace.
\item
  Added EMBREE\_LIBRARY\_NAME cmake option that allows to rename the
  Embree library.
\item
  When Embree is compiled as static library, EMBREE\_STATIC\_LIB has no
  longer to get defined before including the Embree API headers.
\item
  Added CPU frequency\_level device configuration to allow an
  application to specify the frequency level it wants to run on. This
  forces Embree to not use optimizations that may reduce the CPU
  frequency below that level. By default Embree is configured to the the
  AVX-heavy frequency level, thus if the application uses solely non-AVX
  code, configuring the Embree device with ``frequency\_level=simd128''
  may give better performance.
\item
  Fixed a bug in the spatial split builder which caused it to fail for
  scenes with more than 2\^{}24 geometries.
\end{itemize}

\subsection{New Features in Embree
3.5.1}\label{new-features-in-embree-3.5.1}

\begin{itemize}
\tightlist
\item
  Fixed ray/sphere intersector to work also for non-normalized rays.
\item
  Fixed self intersection avoidance for ray oriented discs when
  non-normalized rays were used.
\item
  Increased maximal face valence for subdiv patch to 64 and reduced
  stack size requirement for subdiv patch evaluation.
\end{itemize}

\subsection{New Features in Embree
3.5.0}\label{new-features-in-embree-3.5.0}

\begin{itemize}
\tightlist
\item
  Changed normal oriented curve definition to fix waving artefacts.
\item
  Fixed bounding issue for normal oriented motion blurred curves.
\item
  Fixed performance issue with motion blurred point geometry.
\item
  Fixed generation of documentation with new pandoc versions.
\end{itemize}

\subsection{New Features in Embree
3.4.0}\label{new-features-in-embree-3.4.0}

\begin{itemize}
\tightlist
\item
  Added point primitives (spheres, ray-oriented discs, normal-oriented
  discs).
\item
  Fixed crash triggered by scenes with only invalid primitives.
\item
  Improved robustness of quad/grid-based intersectors.
\item
  Upgraded to Intel® TBB 2019.2 for release builds.
\end{itemize}

\subsection{New Features in Embree
3.3.0}\label{new-features-in-embree-3.3.0}

\begin{itemize}
\tightlist
\item
  Added support for motion blur time range per geometry. This way
  geometries can appear and disappear during the camera shutter and time
  steps do not have to start and end at camera shutter interval
  boundaries.
\item
  Fixed crash with pathtracer when using --triangle-sphere command line.
\item
  Fixed crash with pathtracer when using --shader ao command line.
\item
  Fixed tutorials showing a black window on macOS 10.14 until moved.
\end{itemize}

\subsection{New Features in Embree
3.2.4}\label{new-features-in-embree-3.2.4}

\begin{itemize}
\tightlist
\item
  Fixed compile issues with ICC 2019.
\item
  Released ZIP files for Windows are now provided in a version linked
  against Visual Studio 2013 and Visual Studio 2015.
\end{itemize}

\subsection{New Features in Embree
3.2.3}\label{new-features-in-embree-3.2.3}

\begin{itemize}
\tightlist
\item
  Fixed crash when using curves with RTC\_SCENE\_FLAG\_DYNAMIC combined
  with RTC\_BUILD\_QUALITY\_MEDIUM.
\end{itemize}

\subsection{New Features in Embree
3.2.2}\label{new-features-in-embree-3.2.2}

\begin{itemize}
\tightlist
\item
  Fixed intersection distance for unnormalized rays with line segments.
\item
  Removed libmmd.dll dependency in release builds for Windows.
\item
  Fixed detection of AppleClang compiler under MacOSX.
\end{itemize}

\subsection{New Features in Embree
3.2.1}\label{new-features-in-embree-3.2.1}

\begin{itemize}
\tightlist
\item
  Bugfix in flat mode for hermite curves.
\item
  Added EMBREE\_CURVE\_SELF\_INTERSECTION\_AVOIDANCE\_FACTOR cmake
  option to control self intersection avoidance for flat curves.
\item
  Performance fix when instantiating motion blurred scenes. The
  application should best use two (or more) time steps for an instance
  that instantiates a motion blurred scene.
\item
  Fixed AVX512 compile issue with GCC 6.1.1.
\item
  Fixed performance issue with rtcGetGeometryUserData when used during
  rendering.
\item
  Bugfix in length of derivatives for grid geometry.
\item
  Added BVH8 support for motion blurred curves and lines. For some
  workloads this increases performance by up to 7\%.
\item
  Fixed rtcGetGeometryTransform to return the local to world transform.
\item
  Fixed bug in multi segment motion blur that caused missing of
  perfectly axis aligned geometry.
\item
  Reduced memory consumption of small scenes by 4x.
\item
  Reduced temporal storage of grid builder.
\end{itemize}

\subsection{New Features in Embree
3.2.0}\label{new-features-in-embree-3.2.0}

\begin{itemize}
\tightlist
\item
  Improved watertightness of robust mode.
\item
  Line segments, and other curves are now all contained in a single BVH
  which improves performance when these are both used in a scene.
\item
  Performance improvement of up to 20\% for line segments.
\item
  Bugfix to Embree2 to Embree3 conversion script.
\item
  Added support for Hermite curve basis.
\item
  Semantics of normal buffer for normal oriented curves has changed to
  simplify usage. Please see documentation for details.
\item
  Using GLFW and imgui in tutorials.
\item
  Fixed floating point exception in static variable initialization.
\item
  Fixed invalid memory access in rtcGetGeometryTransform for non-motion
  blur instances.
\item
  Improved self intersection avoidance for flat curves. Transparency
  rays with tnear set to previous hit distance do not need curve radius
  based self intersection avoidance as same hit is calculated again. For
  this reason self intersection avoidance is now only applied to ray
  origin.
\end{itemize}

\subsection{New Features in Embree
3.1.0}\label{new-features-in-embree-3.1.0}

\begin{itemize}
\tightlist
\item
  Added new normal-oriented curve primitive for ray tracing of
  grass-like structures.
\item
  Added new grid primitive for ray tracing tessellated and displaced
  surfaces in very memory efficient manner.
\item
  Fixed bug of ribbon curve intersector when derivative was zero.
\item
  Installing all static libraries when EMBREE\_STATIC\_LIB is enabled.
\item
  Added API functions to access topology of subdivision mesh.
\item
  Reduced memory consumption of instances.
\item
  Improved performance of instances by 8\%.
\item
  Reduced memory consumption of curves by up to 2x.
\item
  Up to 5\% higher performance on AVX-512 architectures.
\item
  Added native support for multiple curve basis functions. Internal
  basis conversions are no longer performed, which saves additional
  memory when multiple bases are used.
\item
  Fixed issue with non thread safe local static variable initialization
  in VS2013.
\item
  Bugfix in rtcSetNewGeometry. Vertex buffers did not get properly
  overallocated.
\item
  Replaced ImageMagick with OpenImageIO in the tutorials.
\end{itemize}

\subsection{New Features in Embree
3.0.0}\label{new-features-in-embree-3.0.0}

\begin{itemize}
\tightlist
\item
  Switched to a new version of the API which provides improved
  flexibility but is not backward compatible. Please see ``Upgrading
  from Embree 2 to Embree 3'' section of the documentation for upgrade
  instructions. In particular, we provide a Python script that performs
  most of the transition work.
\item
  User geometries inside an instanced scene and a top-level scene no
  longer need to handle the instID field of the ray differently. They
  both just need to copy the context.instID into the ray.instID field.
\item
  Support for context filter functions that can be assigned to a ray
  query.
\item
  User geometries can now invoke filter functions using the
  rtcFilterIntersection and rtcFilterOcclusion calls.
\item
  Higher flexibility through specifying build quality per scene and
  geometry.
\item
  Geometry normal uses commonly used right-hand rule from now on.
\item
  Added self-intersection avoidance to ribbon curves and lines.
  Applications do not have to implement self-intersection workarounds
  for these primitive types anymore.
\item
  Added support for 4 billion primitives in a single scene.
\item
  Removed the RTC\_MAX\_USER\_VERTEX\_BUFFERS and
  RTC\_MAX\_INDEX\_BUFFERS limitations.
\item
  Reduced memory consumption by 192 bytes per instance.
\item
  Fixed some performance issues on AVX-512 architectures.
\item
  Individual Contributor License Agreement (ICLA) and Corporate
  Contributor License Agreement (CCLA) no longer required to contribute
  to the project.
\end{itemize}

\subsection{New Features in Embree
2.17.5}\label{new-features-in-embree-2.17.5}

\begin{itemize}
\tightlist
\item
  Improved watertightness of robust mode.
\item
  Fixed floating point exception in static variable initialization.
\item
  Fixed AVX512 compile issue with GCC 6.1.1.
\end{itemize}

\subsection{New Features in Embree
2.17.4}\label{new-features-in-embree-2.17.4}

\begin{itemize}
\tightlist
\item
  Fixed AVX512 compile issue with GCC 7.
\item
  Fixed issue with not thread safe local static variable initialization
  in VS2013.
\item
  Fixed bug in the 4 and 8-wide packet intersection of instances with
  multi-segment motion blur on AVX-512 architectures.
\item
  Fixed bug in rtcOccluded4/8/16 when only AVX-512 ISA was enabled.
\end{itemize}

\subsection{New Features in Embree
2.17.3}\label{new-features-in-embree-2.17.3}

\begin{itemize}
\tightlist
\item
  Fixed GCC compile warning in debug mode.
\item
  Fixed bug of ribbon curve intersector when derivative was zero.
\item
  Installing all static libraries when EMBREE\_STATIC\_LIB is enabled.
\end{itemize}

\subsection{New Features in Embree
2.17.2}\label{new-features-in-embree-2.17.2}

\begin{itemize}
\tightlist
\item
  Made BVH build of curve geometry deterministic.
\end{itemize}

\subsection{New Features in Embree
2.17.1}\label{new-features-in-embree-2.17.1}

\begin{itemize}
\tightlist
\item
  Improved performance of occlusion ray packets by up to 50\%.
\item
  Fixed detection of Clang for CMake 3 under MacOSX
\item
  Fixed AVX code compilation issue with GCC 7 compiler caused by
  explicit use of vzeroupper intrinsics.
\item
  Fixed an issue where Clang address sanitizer reported an error in the
  internal tasking system.
\item
  Added fix to compile on 32 bit Linux distribution.
\item
  Fixed some wrong relative include paths in Embree.
\item
  Improved performance of robust single ray mode by 5\%.
\item
  Added EMBREE\_INSTALL\_DEPENDENCIES option (default OFF) to enable
  installing of Embree dependencies.
\item
  Fixed performance regression for occlusion ray streams.
\item
  Reduced temporary memory requirements of BVH builder for curves and
  line segments.
\item
  Fixed performance regression for user geometries and packet ray
  tracing.
\item
  Fixed bug where wrong closest hit was reported for very curvy hair
  segment.
\end{itemize}

\subsection{New Features in Embree
2.17.0}\label{new-features-in-embree-2.17.0}

\begin{itemize}
\tightlist
\item
  Improved packet ray tracing performance for coherent rays by 10-60\%
  (requires RTC\_INTERSECT\_COHERENT flag).
\item
  Improved ray tracing performance for incoherent rays on AVX-512
  architectures by 5\%.
\item
  Improved ray tracing performance for streams of incoherent rays by
  5-15\%.
\item
  Fixed tbb\_debug.lib linking error under Windows.
\item
  Fast coherent ray stream and packet code paths now also work in robust
  mode.
\item
  Using less aggressive prefetching for large BVH nodes which results in
  1-2\% higher ray tracing performance.
\item
  Precompiled binaries have stack-protector enabled, except for
  traversal kernels. BVH builders can be slightly slower due to this
  change. If you want stack-protectors disabled please turn off
  EMBREE\_STACK\_PROTECTOR in cmake and build the binaries yourself.
\item
  When enabling ISAs individually, the 8-wide BVH was previously only
  available when the AVX ISA was also selected. This issue is now fixed,
  and one can enable only AVX2 and still get best performance by using
  an 8-wide BVH.
\item
  Fixed rtcOccluded1 and rtcOccluded1Ex API functions which were broken
  in Intel® ISPC.
\item
  Providing MSI installer for Windows.
\end{itemize}

\subsection{New Features in Embree
2.16.5}\label{new-features-in-embree-2.16.5}

\begin{itemize}
\tightlist
\item
  Bugfix in the robust triangle intersector that rarely caused NaNs.
\item
  Fixed bug in hybrid traversal kernel when BVH leaf was entered with no
  active rays. This rarely caused crashes when used with instancing.
\item
  Fixed bug introduced in Embree 2.16.2 which caused instancing not to
  work properly when a smaller than the native SIMD width was used in
  ray packet mode.
\item
  Fixed bug in the curve geometry intersector that caused rendering
  artefacts for Bézier curves with p0=p1 and/or p2=p3.
\item
  Fixed bug in the curve geometry intersector that caused hit results
  with NaNs to be reported.
\item
  Fixed masking bug that caused rare cracks in curve geometry.
\item
  Enabled support for SSE2 in precompiled binaries again.
\end{itemize}

\subsection{New Features in Embree
2.16.4}\label{new-features-in-embree-2.16.4}

\begin{itemize}
\tightlist
\item
  Bugfix in the ribbon intersector for hair primitives. Non-normalized
  rays caused wrong intersection distance to be reported.
\end{itemize}

\subsection{New Features in Embree
2.16.3}\label{new-features-in-embree-2.16.3}

\begin{itemize}
\tightlist
\item
  Increased accuracy for handling subdivision surfaces. This fixes
  cracks when using displacement mapping but reduces performance at
  irregular vertices.
\item
  Fixed a bug where subdivision geometry was not properly updated when
  modifying only the tessellation rate and vertex array.
\end{itemize}

\subsection{New Features in Embree
2.16.2}\label{new-features-in-embree-2.16.2}

\begin{itemize}
\tightlist
\item
  Fixed bug that caused NULL ray query context in intersection filter
  when instancing was used.
\item
  Fixed an issue where uv's where outside the triangle (or quad) for
  very small triangles (or quads). In robust mode we improved the uv
  calculation to avoid that issue, in fast mode we accept that
  inconsistency for better performance.
\item
  Changed UV encoding for non-quad subdivision patches to allow a
  subpatch UV range of \texttt{{[}-0.5,1.5{[}}. Using this new encoding
  one can use finite differences to calculate derivatives if required.
  Please adjust your code in case you rely on the old encoding.
\end{itemize}

\subsection{New Features in Embree
2.16.1}\label{new-features-in-embree-2.16.1}

\begin{itemize}
\tightlist
\item
  Workaround for compile issues with Visual Studio 2017
\item
  Fixed bug in subdiv code for static scenes when using tessellation
  levels larger than 50.
\item
  Fixed low performance when adding many geometries to a scene.
\item
  Fixed high memory consumption issue when using instances in dynamic
  scene (by disabling two level builder for user geometries and
  instances).
\end{itemize}

\subsection{New Features in Embree
2.16.0}\label{new-features-in-embree-2.16.0}

\begin{itemize}
\tightlist
\item
  Improved multi-segment motion blur support for scenes with different
  number of time steps per mesh.
\item
  New top level BVH builder that improves build times and BVH quality of
  two-level BVHs.
\item
  Added support to enable only a single ISA. Previously code was always
  compiled for SSE2.
\item
  Improved single ray tracing performance for incoherent rays on AVX-512
  architectures by 5-10\%.
\item
  Improved packet/hybrid ray tracing performance for incoherent rays on
  AVX-512 architectures by 10-30\%.
\item
  Improved stream ray tracing performance for coherent rays in
  structure-of-pointers layout by 40-70\%.
\item
  BVH builder for compact scenes of triangles and quads needs
  essentially no temporary memory anymore. This doubles the maximal
  scene size that can be rendered in compact mode.
\item
  Triangles no longer store the geometry normal in fast/default mode
  which reduces memory consumption by up to 20\%.
\item
  Compact mode uses BVH4 now consistently which reduces memory
  consumption by up to 10\%.
\item
  Reduced memory consumption for small scenes (of 10k-100k primitives)
  and dynamic scenes.
\item
  Improved performance of user geometries and instances through BVH8
  support.
\item
  The API supports now specifying the geometry ID of a geometry at
  construction time. This way matching the geometry ID used by Embree
  and the application is simplified.
\item
  Fixed a bug that would have caused a failure of the BVH builder for
  dynamic scenes when run on a machine with more then 1000 threads.
\item
  Fixed a bug that could have been triggered when reaching the maximal
  number of mappings under Linux (\texttt{vm.max\_map\_count}). This
  could have happened when creating a large number of small static
  scenes.
\item
  Added huge page support for Windows and MacOSX (experimental).
\item
  Added support for Visual Studio 2017.
\item
  Removed support for Visual Studio 2012.
\item
  Precompiled binaries now require a CPU supporting at least the SSE4.2
  ISA.
\item
  We no longer provide precompiled binaries for 32-bit on Windows.
\item
  Under Windows one now has to use the platform toolset option in CMake
  to switch to Clang or the Intel® Compiler.
\item
  Fixed a bug for subdivision meshes when using the incoherent scene
  flag.
\item
  Fixed a bug in the line geometry intersection, that caused reporting
  an invalid line segment intersection with primID -1.
\item
  Buffer stride for vertex buffers of different time steps of triangle
  and quad meshes have to be identical now.
\item
  Fixed a bug in the curve geometry intersection code when passed a
  perfect cylinder.
\end{itemize}

\subsection{New Features in Embree
2.15.0}\label{new-features-in-embree-2.15.0}

\begin{itemize}
\tightlist
\item
  Added \texttt{rtcCommitJoin} mode that allows thread to join a build
  operation. When using the internal tasking system this allows Embree
  to solely use the threads that called \texttt{rtcCommitJoin} to build
  the scene, while previously also normal worker threads participated in
  the build. You should no longer use \texttt{rtcCommit} to join a
  build.
\item
  Added \texttt{rtcDeviceSetErrorFunction2} API call, which sets an
  error callback function which additionally gets passed a user provided
  pointer (\texttt{rtcDeviceSetErrorFunction} is now deprecated).
\item
  Added \texttt{rtcDeviceSetMemoryMonitorFunction2} API call, which sets
  a memory monitor callback function which additionally get passed a
  user provided pointer. (\texttt{rtcDeviceSetMemoryMonitorFunction} is
  now deprecated).
\item
  Build performance for hair geometry improved by up to 2×.
\item
  Standard BVH build performance increased by 5\%.
\item
  Added API extension to use internal Morton-code based builder, the
  standard binned-SAH builder, and the spatial split-based SAH builder.
\item
  Added support for BSpline hair and curves. Embree uses either the
  Bézier or BSpline basis internally, and converts other curves, which
  requires more memory during rendering. For reduced memory consumption
  set the \texttt{EMBREE\_NATIVE\_SPLINE\_BASIS} to the basis your
  application uses (which is set to \texttt{BEZIER} by default).
\item
  Setting the number of threads through
  \texttt{tbb::taskscheduler\_init} object on the application side is
  now working properly.
\item
  Windows and Linux releases are build using AVX-512 support.
\item
  Implemented hybrid traversal for hair and line segments for improved
  ray packet performance.
\item
  AVX-512 code compiles with Clang 4.0.0
\item
  Fixed crash when ray packets were disabled in CMake.
\end{itemize}

\subsection{New Features in Embree
2.14.0}\label{new-features-in-embree-2.14.0}

\begin{itemize}
\tightlist
\item
  Added \texttt{ignore\_config\_files} option to init flags that allows
  the application to ignore Embree configuration files.
\item
  Face-varying interpolation is now supported for subdivision surfaces.
\item
  Up to 16 user vertex buffers are supported for vertex attribute
  interpolation.
\item
  Deprecated \texttt{rtcSetBoundaryMode} function, please use the new
  \texttt{rtcSetSubdivisionMode} function.
\item
  Added \texttt{RTC\_SUBDIV\_PIN\_BOUNDARY} mode for handling boundaries
  of subdivision meshes.
\item
  Added \texttt{RTC\_SUBDIV\_PIN\_ALL} mode to enforce linear
  interpolation for subdivision meshes.
\item
  Optimized object generation performance for dynamic scenes.
\item
  Reduced memory consumption when using lots of small dynamic objects.
\item
  Fixed bug for subdivision surfaces using low tessellation rates.
\item
  Hair geometry now uses a new ribbon intersector that intersects with
  ray-facing quads. The new intersector also returns the v-coordinate of
  the hair intersection, and fixes artefacts at junction points between
  segments, at the cost of a small performance hit.
\item
  Added \texttt{rtcSetBuffer2} function, that additionally gets the
  number of elements of a buffer. In dynamic scenes, this function
  allows to quickly change buffer sizes, making it possible to change
  the number of primitives of a mesh or the number of crease features
  for subdivision surfaces.
\item
  Added simple `viewer\_anim' tutorial for rendering key frame
  animations and `buildbench' for measuring BVH (re-)build performance
  for static and dynamic scenes.
\item
  Added more AVX-512 optimizations for future architectures.
\end{itemize}

\subsection{New Features in Embree
2.13.0}\label{new-features-in-embree-2.13.0}

\begin{itemize}
\tightlist
\item
  Improved performance for compact (but not robust) scenes.
\item
  Added robust mode for motion blurred triangles and quads.
\item
  Added fast dynamic mode for user geometries.
\item
  Up to 20\% faster BVH build performance on the second generation
  Intel® Xeon Phi™ processor codenamed Knights Landing.
\item
  Improved quality of the spatial split builder.
\item
  Improved performance for coherent streams of ray packets (SOA layout),
  e.g.~for fast primary visibility.
\item
  Various bug fixes in tessellation cache, quad-based spatial split
  builder, etc.
\end{itemize}

\subsection{New Features in Embree
2.12.0}\label{new-features-in-embree-2.12.0}

\begin{itemize}
\tightlist
\item
  Added support for multi-segment motion blur for all primitive types.
\item
  API support for stream of pointers to single rays
  (\texttt{rtcIntersect1Mp} and \texttt{rtcOccluded1Mp})
\item
  Improved BVH refitting performance for dynamic scenes.
\item
  Improved high-quality mode for quads (added spatial split builder for
  quads)
\item
  Faster dynamic scenes for triangle and quad-based meshes on AVX2
  enabled machines.
\item
  Performance and correctness bugfix in optimization for streams of
  coherent (single) rays.
\item
  Fixed large memory consumption (issue introduced in Embree v2.11.0).
  If you use Embree v2.11.0 please upgrade to Embree v2.12.0.
\item
  Reduced memory consumption for dynamic scenes containing small meshes.
\item
  Added support to start and affinitize Intel® TBB worker threads by
  passing ``\texttt{start\_threads=1,set\_affinity=1}'' to
  \texttt{rtcNewDevice}. These settings are recommended on systems with
  a high thread count.
\item
  \texttt{rtcInterpolate2} can now be called within a displacement
  shader.
\item
  Added initial support for Microsoft's Parallel Pattern Library (PPL)
  as tasking system alternative (for optimal performance Intel® TBB is
  highly recommended).
\item
  Updated to Intel® TBB 2017 which is released under the Apache v2.0
  license.
\item
  Dropped support for Visual Studio 2012 Win32 compiler. Visual Studio
  2012 x64 is still supported.
\end{itemize}

\subsection{New Features in Embree
2.11.0}\label{new-features-in-embree-2.11.0}

\begin{itemize}
\tightlist
\item
  Improved performance for streams of coherent (single) rays flagged
  with \texttt{RTC\_INTERSECT\_COHERENT}. For such coherent ray streams,
  e.g. primary rays, the performance typically improves by 1.3-2×.
\item
  New spatial split BVH builder for triangles, which is 2-6× faster than
  the previous version and more memory conservative.
\item
  Improved performance and scalability of all standard BVH builders on
  systems with large core counts.
\item
  Fixed \texttt{rtcGetBounds} for motion blur scenes.
\item
  Thread affinity is now on by default when running on the latest Intel®
  Xeon Phi™ processor.
\item
  Added AVX-512 support for future Intel® Xeon processors.
\end{itemize}

\subsection{New Features in Embree
2.10.0}\label{new-features-in-embree-2.10.0}

\begin{itemize}
\tightlist
\item
  Added a new curve geometry which renders the sweep surface of a circle
  along a Bézier curve.
\item
  Intersection filters can update the \texttt{tfar} ray distance.
\item
  Geometry types can get disabled at compile time.
\item
  Modified and extended the ray stream API.
\item
  Added new callback mechanism for the ray stream API.
\item
  Improved ray stream performance (up to 5-10\%).
\item
  Up to 20\% faster morton builder on machines with large core counts.
\item
  Lots of optimizations for the second generation Intel® Xeon Phi™
  processor codenamed Knights Landing.
\item
  Added experimental support for compressed BVH nodes (reduces node size
  to 56-62\% of uncompressed size). Compression introduces a typical
  performance overhead of \textasciitilde10\%.
\item
  Bugfix in backface culling mode. We do now properly cull the backfaces
  and not the frontfaces.
\item
  Feature freeze for the first generation Intel® Xeon Phi™ coprocessor
  codenamed Knights Corner. We will still maintain and add bug fixes to
  Embree v2.9.0, but Embree 2.10 and future versions will no longer
  support it.
\end{itemize}

\subsection{New Features in Embree
2.9.0}\label{new-features-in-embree-2.9.0}

\begin{itemize}
\tightlist
\item
  Improved shadow ray performance (10-100\% depending on the scene).
\item
  Added initial support for ray streams (10-30\% higher performance
  depending on ray coherence in the stream).
\item
  Added support to calculate second order derivatives using the
  \texttt{rtcInterpolate2} function.
\item
  Changed the parametrization for triangular subdivision faces to the
  same scheme used for pentagons.
\item
  Added support to query the Embree configuration using the
  \texttt{rtcDeviceGetParameter} function.
\end{itemize}

\subsection{New Features in Embree
2.8.1}\label{new-features-in-embree-2.8.1}

\begin{itemize}
\tightlist
\item
  Added support for setting per geometry tessellation rate (supported
  for subdivision and Bézier geometries).
\item
  Added support for motion blurred instances.
\end{itemize}

\subsection{New Features in Embree
2.8.0}\label{new-features-in-embree-2.8.0}

\begin{itemize}
\tightlist
\item
  Added support for line segment geometry.
\item
  Added support for quad geometry (replaces triangle-pairs feature).
\item
  Added support for linear motion blur of user geometries.
\item
  Improved performance through AVX-512 optimizations.
\item
  Improved performance of lazy scene build (when using Intel® TBB 4.4
  update 2).
\item
  Improved performance through huge page support under linux.
\end{itemize}

\subsection{New Features in Embree
2.7.1}\label{new-features-in-embree-2.7.1}

\begin{itemize}
\tightlist
\item
  Internal tasking system supports cancellation of build operations.
\item
  Intel® ISPC mode for robust and compact scenes got significantly
  faster (implemented hybrid traversal for bvh4.triangle4v and
  bvh4.triangle4i).
\item
  Hair rendering got faster as we fixed some issues with the SAH
  heuristic cost factors.
\item
  BVH8 got slight faster for single ray traversal (improved sorting when
  hitting more than 4 boxes).
\item
  BVH build performance got up to 30\% faster on CPUs with high core
  counts (improved parallel partition code).
\item
  High quality build mode again working properly (spatial splits had
  been deactivated in v2.7.0 due to some bug).
\item
  Support for merging two adjacent triangles sharing a common edge into
  a triangle-pair primitive (can reduce memory consumption and BVH build
  times by up to 50\% for mostly quad-based input meshes).
\item
  Internal cleanups (reduced number of traversal kernels by more
  templating).
\item
  Reduced stack size requirements of BVH builders.
\item
  Fixed crash for dynamic scenes, triggered by deleting all geometries
  from the scene.
\end{itemize}

\subsection{New Features in Embree
2.7.0}\label{new-features-in-embree-2.7.0}

\begin{itemize}
\tightlist
\item
  Added device concept to Embree to allow different components of an
  application to use Embree without interfering with each other.
\item
  Fixed memory leak in twolevel builder used for dynamic scenes.
\item
  Fixed bug in tessellation cache that caused crashes for subdivision
  surfaces.
\item
  Fixed bug in internal task scheduler that caused deadlocks when using
  \texttt{rtcCommitThread}.
\item
  Improved hit-distance accuracy for thin triangles in robust mode.
\item
  Added support to disable ray packet support in cmake.
\end{itemize}

\subsection{New Features in Embree
2.6.2}\label{new-features-in-embree-2.6.2}

\begin{itemize}
\tightlist
\item
  Fixed bug triggered by instantiating motion blur geometry.
\item
  Fixed bug in hit UV coordinates of static subdivision geometries.
\item
  Performance improvements when only changing tessellation levels for
  subdivision geometry per frame.
\item
  Added ray packet intersectors for subdivision geometry, resulting in
  improved performance for coherent rays.
\item
  Reduced virtual address space usage for static geometries.
\item
  Fixed some AVX2 code paths when compiling with GCC or Clang.
\item
  Bugfix for subdiv patches with non-matching winding order.
\item
  Bugfix in ISA detection of AVX-512.
\end{itemize}

\subsection{New Features in Embree
2.6.1}\label{new-features-in-embree-2.6.1}

\begin{itemize}
\tightlist
\item
  Major performance improvements for ray tracing subdivision surfaces,
  e.g.~up to 2× faster for scenes where only the tessellation levels are
  changing per frame, and up to 3× faster for scenes with lots of crease
  features
\item
  Initial support for architectures supporting the new 16-wide AVX-512
  ISA
\item
  Implemented intersection filter callback support for subdivision
  surfaces
\item
  Added \texttt{RTC\_IGNORE\_INVALID\_RAYS} CMake option which makes the
  ray intersectors more robust against full tree traversal caused by
  invalid ray inputs (e.g.~INF, NaN, etc)
\end{itemize}

\subsection{New Features in Embree
2.6.0}\label{new-features-in-embree-2.6.0}

\begin{itemize}
\tightlist
\item
  Added \texttt{rtcInterpolate} function to interpolate per vertex
  attributes
\item
  Added \texttt{rtcSetBoundaryMode} function that can be used to select
  the boundary handling for subdivision surfaces
\item
  Fixed a traversal bug that caused rays with very small ray direction
  components to miss geometry
\item
  Performance improvements for the robust traversal mode
\item
  Fixed deadlock when calling \texttt{rtcCommit} from multiple threads
  on same scene
\end{itemize}

\subsection{New Features in Embree
2.5.1}\label{new-features-in-embree-2.5.1}

\begin{itemize}
\tightlist
\item
  On dual socket workstations, the initial BVH build performance almost
  doubled through a better memory allocation scheme
\item
  Reduced memory usage for subdivision surface objects with crease
  features
\item
  \texttt{rtcCommit} performance is robust against unset ``flush to
  zero'' and ``denormals are zero'' flags. However, enabling these flags
  in your application is still recommended
\item
  Reduced memory usage for subdivision surfaces with borders and
  infinitely sharp creases
\item
  Lots of internal cleanups and bug fixes for both Intel® Xeon® and
  Intel® Xeon Phi™
\end{itemize}

\subsection{New Features in Embree
2.5.0}\label{new-features-in-embree-2.5.0}

\begin{itemize}
\tightlist
\item
  Improved hierarchy build performance on both Intel Xeon and Intel Xeon
  Phi
\item
  Vastly improved tessellation cache for ray tracing subdivision
  surfaces
\item
  Added \texttt{rtcGetUserData} API call to query per geometry user
  pointer set through \texttt{rtcSetUserData}
\item
  Added support for memory monitor callback functions to track and limit
  memory consumption
\item
  Added support for progress monitor callback functions to track build
  progress and cancel long build operations
\item
  BVH builders can be used to build user defined hierarchies inside the
  application (see tutorial \hyperref[bvh-builder]{BVH Builder})
\item
  Switched to Intel® TBB as default tasking system on Xeon to get even
  faster hierarchy build times and better integration for applications
  that also use Intel® TBB
\item
  \texttt{rtcCommit} can get called from multiple Intel® TBB threads to
  join the hierarchy build operations
\end{itemize}

\subsection{New Features in Embree
2.4}\label{new-features-in-embree-2.4}

\begin{itemize}
\tightlist
\item
  Support for Catmull Clark subdivision surfaces (triangle/quad base
  primitives)
\item
  Support for vector displacements on Catmull Clark subdivision surfaces
\item
  Various bug fixes (e.g.~4-byte alignment of vertex buffers works)
\end{itemize}

\subsection{New Features in Embree
2.3.3}\label{new-features-in-embree-2.3.3}

\begin{itemize}
\tightlist
\item
  BVH builders more robustly handle invalid input data (Intel Xeon
  processor family)
\item
  Motion blur support for hair geometry (Xeon)
\item
  Improved motion blur performance for triangle geometry (Xeon)
\item
  Improved robust ray tracing mode (Xeon)
\item
  Added \texttt{rtcCommitThread} API call for easier integration into
  existing tasking systems (Xeon and Intel Xeon Phi coprocessor)
\item
  Added support for recording and replaying all
  \texttt{rtcIntersect}/\texttt{rtcOccluded} calls (Xeon and Xeon Phi)
\end{itemize}

\subsection{New Features in Embree
2.3.2}\label{new-features-in-embree-2.3.2}

\begin{itemize}
\tightlist
\item
  Improved mixed AABB/OBB-BVH for hair geometry (Xeon Phi)
\item
  Reduced amount of pre-allocated memory for BVH builders (Xeon Phi)
\item
  New 64-bit Morton code-based BVH builder (Xeon Phi)
\item
  (Enhanced) Morton code-based BVH builders use now tree rotations to
  improve BVH quality (Xeon Phi)
\item
  Bug fixes (Xeon and Xeon Phi)
\end{itemize}

\subsection{New Features in Embree
2.3.1}\label{new-features-in-embree-2.3.1}

\begin{itemize}
\tightlist
\item
  High quality BVH mode improves spatial splits which result in up to
  30\% performance improvement for some scenes (Xeon)
\item
  Compile time enabled intersection filter functions do not reduce
  performance if no intersection filter is used in the scene (Xeon and
  Xeon Phi)
\item
  Improved ray tracing performance for hair geometry by \textgreater20\%
  on Xeon Phi. BVH for hair geometry requires 20\% less memory
\item
  BVH8 for AVX/AVX2 targets improves performance for single ray tracing
  on Haswell by up to 12\% and by up to 5\% for hybrid (Xeon)
\item
  Memory conservative BVH for Xeon Phi now uses BVH node quantization to
  lower memory footprint (requires half the memory footprint of the
  default BVH)
\end{itemize}

\subsection{New Features in Embree
2.3}\label{new-features-in-embree-2.3}

\begin{itemize}
\tightlist
\item
  Support for ray tracing hair geometry (Xeon and Xeon Phi)
\item
  Catching errors through error callback function
\item
  Faster hybrid traversal (Xeon and Xeon Phi)
\item
  New memory conservative BVH for Xeon Phi
\item
  Faster Morton code-based builder on Xeon
\item
  Faster binned-SAH builder on Xeon Phi
\item
  Lots of code cleanups/simplifications/improvements (Xeon and Xeon Phi)
\end{itemize}

\subsection{New Features in Embree
2.2}\label{new-features-in-embree-2.2}

\begin{itemize}
\tightlist
\item
  Support for motion blur on Xeon Phi
\item
  Support for intersection filter callback functions
\item
  Support for buffer sharing with the application
\item
  Lots of AVX2 optimizations, e.g.~\textasciitilde20\% faster 8-wide
  hybrid traversal
\item
  Experimental support for 8-wide (AVX/AVX2) and 16-wide BVHs (Xeon Phi)
\end{itemize}

\subsection{New Features in Embree
2.1}\label{new-features-in-embree-2.1}

\begin{itemize}
\tightlist
\item
  New future proof API with a strong focus on supporting dynamic scenes
\item
  Lots of optimizations for 8-wide AVX2 (Haswell architecture)
\item
  Automatic runtime code selection for SSE, AVX, and AVX2
\item
  Support for user-defined geometry
\item
  New and improved BVH builders:

  \begin{itemize}
  \tightlist
  \item
    Fast adaptive Morton code-based builder (without SAH-based top-level
    rebuild)
  \item
    Both the SAH and Morton code-based builders got faster (Xeon Phi)
  \item
    New variant of the SAH-based builder using triangle pre-splits (Xeon
    Phi)
  \end{itemize}
\end{itemize}

\subsection{New Features in Embree
2.0}\label{new-features-in-embree-2.0}

\begin{itemize}
\tightlist
\item
  Support for the Intel® Xeon Phi™ coprocessor platform
\item
  Support for high-performance ``packet'' kernels on SSE, AVX, and Xeon
  Phi
\item
  Integration with the Intel® Implicit SPMD Program Compiler (Intel®
  ISPC)
\item
  Instantiation and fast BVH reconstruction
\item
  Example photo-realistic rendering engine for both C++ and Intel® ISPC
\end{itemize}

\chapter{Installation of Embree}\label{installation-of-embree}

\section{Windows Installation}\label{windows-installation}

Embree linked against Visual Studio 2015 are provided as a ZIP file
\href{https://github.com/embree/embree/releases/download/v4.1.0/embree-4.1.0.x64.vc14.windows.zip}{embree-4.1.0.x64.vc14.windows.zip}.
After unpacking this ZIP file, you should set the path to the
\texttt{lib} folder manually to your \texttt{PATH} environment variable
for applications to find Embree.

\section{Linux Installation}\label{linux-installation}

The Linux version of Embree is also delivered as a \texttt{tar.gz} file:
\href{https://github.com/embree/embree/releases/download/v4.1.0/embree-4.1.0.x86_64.linux.tar.gz}{embree-4.1.0.x86\_64.linux.tar.gz}.
Unpack this file using \texttt{tar} and source the provided
\texttt{embree-vars.sh} (if you are using the bash shell) or
\texttt{embree-vars.csh} (if you are using the C shell) to set up the
environment properly:

\begin{Shaded}
\begin{Highlighting}[]
\NormalTok{tar xzf embree}\OperatorTok{{-}}\FloatTok{4.1}\ErrorTok{.0.x86\_64.linux.tar.gz}
\NormalTok{source embree}\OperatorTok{{-}}\FloatTok{4.1}\ErrorTok{.0.x86\_64.linux}\OperatorTok{/}\NormalTok{embree}\OperatorTok{{-}}\NormalTok{vars}\OperatorTok{.}\NormalTok{sh}
\end{Highlighting}
\end{Shaded}

We recommend adding a relative \texttt{RPATH} to your application that
points to the location where Embree (and TBB) can be found,
e.g.~\texttt{\$ORIGIN/../lib}.

\section{macOS Installation}\label{macos-installation}

The macOS version of Embree is also delivered as a ZIP file:
\href{https://github.com/embree/embree/releases/download/v4.1.0/embree-4.1.0.x86_64.macosx.zip}{embree-4.1.0.x86\_64.macosx.zip}.
Unpack this file using \texttt{tar} and source the provided
\texttt{embree-vars.sh} (if you are using the bash shell) or
\texttt{embree-vars.csh} (if you are using the C shell) to set up the
environment properly:

\begin{Shaded}
\begin{Highlighting}[]
\NormalTok{unzip embree}\OperatorTok{{-}}\FloatTok{4.1}\ErrorTok{.0.x64.macosx.zip}\NormalTok{    source embree}\OperatorTok{{-}}\FloatTok{4.1}\ErrorTok{.0.x64.macosx}\OperatorTok{/}\NormalTok{embree}\OperatorTok{{-}}\NormalTok{vars}\OperatorTok{.}\NormalTok{sh}
\end{Highlighting}
\end{Shaded}

If you want to ship Embree with your application, please use the Embree
library of the provided ZIP file. The library name of that Embree
library is of the form \texttt{@rpath/libembree.4.dylib} (and similar
also for the included TBB library). This ensures that you can add a
relative \texttt{RPATH} to your application that points to the location
where Embree (and TBB) can be found, e.g.~\texttt{@loader\_path/../lib}.

\section{Building Embree
Applications}\label{building-embree-applications}

The most convenient way to build an Embree application is through CMake.
Just let CMake find your unpacked Embree package using the
\texttt{FIND\_PACKAGE} function inside your \texttt{CMakeLists.txt}
file:

\begin{Shaded}
\begin{Highlighting}[]
\NormalTok{ FIND\_PACKAGE}\OperatorTok{(}\NormalTok{embree }\DecValTok{4}\NormalTok{ REQUIRED}\OperatorTok{)}
\end{Highlighting}
\end{Shaded}

For CMake to properly find Embree you need to set the
\texttt{embree\_DIR} variable to the folder containing the
\texttt{embree\_config.cmake} file. You might also have to set the
\texttt{TBB\_DIR} variable to the path containing
\texttt{TBB-config.cmake} of a local TBB install, in case you do not
have TBB installed globally on your system, e.g:

\begin{Shaded}
\begin{Highlighting}[]
\NormalTok{cmake }\OperatorTok{{-}}\NormalTok{D embree\_DIR}\OperatorTok{=}\NormalTok{path\_to\_embree\_package}\OperatorTok{/}\NormalTok{lib}\OperatorTok{/}\NormalTok{cmake}\OperatorTok{/}\NormalTok{embree}\OperatorTok{{-}}\FloatTok{4.1}\ErrorTok{.0}\OperatorTok{/}\NormalTok{ \textbackslash{}}
      \OperatorTok{{-}}\NormalTok{D TBB\_DIR}\OperatorTok{=}\NormalTok{path\_to\_tbb\_package}\OperatorTok{/}\NormalTok{lib}\OperatorTok{/}\NormalTok{cmake}\OperatorTok{/}\NormalTok{tbb}\OperatorTok{/}\NormalTok{ \textbackslash{}}
      \OperatorTok{..}
\end{Highlighting}
\end{Shaded}

The \texttt{FIND\_PACKAGE} function will create an \texttt{embree}
target that you can add to your target link libraries:

\begin{Shaded}
\begin{Highlighting}[]
\NormalTok{TARGET\_LINK\_LIBRARIES}\OperatorTok{(}\NormalTok{application embree}\OperatorTok{)}
\end{Highlighting}
\end{Shaded}

For a full example on how to build an Embree application please have a
look at the \texttt{minimal} tutorial provided in the \texttt{src}
folder of the Embree package and also the contained \texttt{README.txt}
file.

\section{Building Embree SYCL
Applications}\label{building-embree-sycl-applications}

Building Embree SYCL applications is also best done using CMake. Please
first get some compatible SYCL compiler and setup the environment as
decribed in sections \hyperref[linux-sycl-compilation]{Linux SYCL
Compilation} and \hyperref[windows-sycl-compilation]{Windows SYCL
Compilation}.

Also perform the setup steps from the previous
\hyperref[building-embree-applications]{Building Embree Applications}
section.

Please also have a look at the \hyperref[minimal]{Minimal} tutorial that
is provided with the Embree release, for an example how to build a
simple SYCL application using CMake and Embree.

To properly compile your SYCL application you have to add additional
SYCL compile flags for each C++ file that contains SYCL device side code
or kernels as described next.

\subsection{JIT Compilation}\label{jit-compilation}

We recommend using just in time compilation (JIT compilation) together
with \hyperref[sycl-jit-caching]{SYCL JIT caching} to compile Embree
SYCL applications. For JIT compilation add these options to the
compilation phase of all C++ files that contain SYCL code:

\begin{Shaded}
\begin{Highlighting}[]
\OperatorTok{{-}}\NormalTok{fsycl }\OperatorTok{{-}}\NormalTok{Xclang }\OperatorTok{{-}}\NormalTok{fsycl}\OperatorTok{{-}}\NormalTok{allow}\OperatorTok{{-}}\NormalTok{func}\OperatorTok{{-}}\NormalTok{ptr }\OperatorTok{{-}}\NormalTok{fsycl}\OperatorTok{{-}}\NormalTok{targets}\OperatorTok{=}\NormalTok{spir64}
\end{Highlighting}
\end{Shaded}

These options enable SYCL two phase compilation (\texttt{-fsycl}
option), enable function pointer support
(\texttt{-Xclang\ -fsycl-allow-func-ptr} option), and just in time (JIT)
compilation only (\texttt{-fsycl-targets=spir64} option).

The following link options have to get added to the linking stage of
your application when using just in time compilation:

\begin{Shaded}
\begin{Highlighting}[]
\OperatorTok{{-}}\NormalTok{fsycl }\OperatorTok{{-}}\NormalTok{fsycl}\OperatorTok{{-}}\NormalTok{targets}\OperatorTok{=}\NormalTok{spir64}
\end{Highlighting}
\end{Shaded}

For a full example on how to build an Embree SYCL application please
have a look at the SYCL version of the \texttt{minimal} tutorial
provided in the \texttt{src} folder of the Embree package and also the
contained \texttt{README.txt} file.

Please have a look at the \hyperref[compiling-embree]{Compiling Embree}
section on how to create an Embree package from sources if required.

\subsection{AOT Compilation}\label{aot-compilation}

Ahead of time compilation (AOT compilation) allows to speed up first
application start up time as device binaries are precompiled. We do not
recommend using AOT compilation as it does not allow the usage of
specialization constants to reduce code complexity.

For ahead of time compilation add these compile options to the
compilation phase of all C++ files that contain SYCL code:

\begin{Shaded}
\begin{Highlighting}[]
\OperatorTok{{-}}\NormalTok{fsycl }\OperatorTok{{-}}\NormalTok{Xclang }\OperatorTok{{-}}\NormalTok{fsycl}\OperatorTok{{-}}\NormalTok{allow}\OperatorTok{{-}}\NormalTok{func}\OperatorTok{{-}}\NormalTok{ptr }\OperatorTok{{-}}\NormalTok{fsycl}\OperatorTok{{-}}\NormalTok{targets}\OperatorTok{=}\NormalTok{spir64\_gen}
\end{Highlighting}
\end{Shaded}

These options enable SYCL two phase compilation (\texttt{-fsycl}
option), enable function pointer support
(\texttt{-Xclang\ -fsycl-allow-func-ptr} option), and ahead of time
(AOT) compilation (\texttt{-fsycl-targets=spir64\_gen} option).

The following link options have to get added to the linking stage of
your application when compiling ahead of time for Xe HPG devices:

\begin{Shaded}
\begin{Highlighting}[]
\OperatorTok{{-}}\NormalTok{fsycl }\OperatorTok{{-}}\NormalTok{fsycl}\OperatorTok{{-}}\NormalTok{targets}\OperatorTok{=}\NormalTok{spir64\_gen}
\OperatorTok{{-}}\NormalTok{Xsycl}\OperatorTok{{-}}\NormalTok{target}\OperatorTok{{-}}\NormalTok{backend}\OperatorTok{=}\NormalTok{spir64\_gen }\StringTok{"{-}device XE\_HPG\_CORE"}
\end{Highlighting}
\end{Shaded}

This in particular configures the devices for AOT compilation to
\texttt{XE\_HPG\_CORE}.

To get a list of all device supported by AOT compilation look at the
help of the device option in ocloc tool:

\begin{Shaded}
\begin{Highlighting}[]
\NormalTok{ocloc compile }\OperatorTok{{-}{-}}\NormalTok{help}
\end{Highlighting}
\end{Shaded}

\chapter{Compiling Embree}\label{compiling-embree}

We recommend using the prebuild Embree packages from
\url{https://github.com/embree/embree/releases}. If you need to compile
Embree yourself you need to use CMake as described in the following.

Do not enable fast-math optimizations in your compiler as this mode is
not supported by Embree.

\section{Linux and macOS}\label{linux-and-macos}

To compile Embree you need a modern C++ compiler that supports C++11.
Embree is tested with the following compilers:

Linux

\begin{itemize}
\tightlist
\item
  Intel® oneAPI DPC++/C++ Compiler 2023.1.0
\item
  Intel® oneAPI DPC++/C++ Compiler 2023.0.0
\item
  oneAPI DPC++/C++ Compiler 2023-04-17
\item
  Clang 5.0.0
\item
  Clang 4.0.0
\item
  GCC 10.0.1 (Fedora 32) AVX512 support
\item
  GCC 8.3.1 (Fedora 28) AVX512 support
\item
  GCC 7.3.1 (Fedora 27) AVX2 support
\item
  GCC 7.3.1 (Fedora 26) AVX2 support
\item
  GCC 6.4.1 (Fedora 25) AVX2 support
\item
  Intel® Implicit SPMD Program Compiler 1.19.0
\item
  Intel® Implicit SPMD Program Compiler 1.18.1
\item
  Intel® Implicit SPMD Program Compiler 1.17.0
\item
  Intel® Implicit SPMD Program Compiler 1.16.1
\item
  Intel® Implicit SPMD Program Compiler 1.15.0
\item
  Intel® Implicit SPMD Program Compiler 1.14.1
\item
  Intel® Implicit SPMD Program Compiler 1.13.0
\item
  Intel® Implicit SPMD Program Compiler 1.12.0
\end{itemize}

macOS x86

\begin{itemize}
\tightlist
\item
  Intel® C++ Classic Compiler 2023.1.0
\item
  Intel® C++ Classic Compiler 2023.0.0
\item
  Apple Clang 12.0.5 (macOS 11.7.1)
\end{itemize}

macOS M1

\begin{itemize}
\tightlist
\item
  Apple Clang 12.0.5 (macOS 11.7.1)
\end{itemize}

Embree supports using the Intel® Threading Building Blocks (TBB) as the
tasking system. For performance and flexibility reasons we recommend
using Embree with the Intel® Threading Building Blocks (TBB) and best
also use TBB inside your application. Optionally you can disable TBB in
Embree through the \texttt{EMBREE\_TASKING\_SYSTEM} CMake variable.

Embree supports the Intel® Implicit SPMD Program Compiler (Intel® ISPC),
which allows straightforward parallelization of an entire renderer. If
you want to use Intel® ISPC then you can enable
\texttt{EMBREE\_ISPC\_SUPPORT} in CMake. Download and install the Intel®
ISPC binaries from
\href{https://ispc.github.io/downloads.html}{ispc.github.io}. After
installation, put the path to \texttt{ispc} permanently into your
\texttt{PATH} environment variable or you set the
\texttt{EMBREE\_ISPC\_EXECUTABLE} variable to point at the ISPC
executable during CMake configuration.

You additionally have to install CMake 3.1.0 or higher and the developer
version of \href{https://www.glfw.org/}{GLFW} version 3.

Under macOS, all these dependencies can be installed using
\href{http://www.macports.org/}{MacPorts}:

\begin{Shaded}
\begin{Highlighting}[]
\NormalTok{sudo port install cmake tbb}\OperatorTok{{-}}\NormalTok{devel glfw}\OperatorTok{{-}}\NormalTok{devel}
\end{Highlighting}
\end{Shaded}

Depending on your Linux distribution you can install these dependencies
using \texttt{yum} or \texttt{apt-get}. Some of these packages might
already be installed or might have slightly different names.

Type the following to install the dependencies using \texttt{yum}:

\begin{Shaded}
\begin{Highlighting}[]
\NormalTok{sudo yum install cmake}
\NormalTok{sudo yum install tbb}\OperatorTok{{-}}\NormalTok{devel}
\NormalTok{sudo yum install glfw}\OperatorTok{{-}}\NormalTok{devel}
\end{Highlighting}
\end{Shaded}

Type the following to install the dependencies using \texttt{apt-get}:

\begin{Shaded}
\begin{Highlighting}[]
\NormalTok{sudo apt}\OperatorTok{{-}}\NormalTok{get install cmake}\OperatorTok{{-}}\NormalTok{curses}\OperatorTok{{-}}\NormalTok{gui}
\NormalTok{sudo apt}\OperatorTok{{-}}\NormalTok{get install libtbb}\OperatorTok{{-}}\NormalTok{dev}
\NormalTok{sudo apt}\OperatorTok{{-}}\NormalTok{get install libglfw3}\OperatorTok{{-}}\NormalTok{dev}
\end{Highlighting}
\end{Shaded}

Finally, you can compile Embree using CMake. Create a build directory
inside the Embree root directory and execute \texttt{ccmake\ ..} inside
this build directory.

\begin{Shaded}
\begin{Highlighting}[]
\NormalTok{mkdir build}
\NormalTok{cd build}
\NormalTok{ccmake }\OperatorTok{..}
\end{Highlighting}
\end{Shaded}

Per default, CMake will use the compilers specified with the \texttt{CC}
and \texttt{CXX} environment variables. Should you want to use a
different compiler, run \texttt{cmake} first and set the
\texttt{CMAKE\_CXX\_COMPILER} and \texttt{CMAKE\_C\_COMPILER} variables
to the desired compiler. For example, to use the Clang compiler instead
of the default GCC on most Linux machines (\texttt{g++} and
\texttt{gcc}), execute

\begin{Shaded}
\begin{Highlighting}[]
\NormalTok{cmake }\OperatorTok{{-}}\NormalTok{DCMAKE\_CXX\_COMPILER}\OperatorTok{=}\NormalTok{clang}\OperatorTok{++} \OperatorTok{{-}}\NormalTok{DCMAKE\_C\_COMPILER}\OperatorTok{=}\NormalTok{clang }\OperatorTok{..}
\end{Highlighting}
\end{Shaded}

Running \texttt{ccmake} will open a dialog where you can perform various
configurations as described below in
\hyperref[cmake-configuration]{CMake Configuration}. After having
configured Embree, press \texttt{c} (for configure) and \texttt{g} (for
generate) to generate a Makefile and leave the configuration. The code
can be compiled by executing make.

\begin{Shaded}
\begin{Highlighting}[]
\NormalTok{make }\OperatorTok{{-}}\NormalTok{j }\DecValTok{8}
\end{Highlighting}
\end{Shaded}

The executables will be generated inside the build folder. We recommend
installing the Embree library and header files on your system. Therefore
set the \texttt{CMAKE\_INSTALL\_PREFIX} to \texttt{/usr} in cmake and
type:

\begin{Shaded}
\begin{Highlighting}[]
\NormalTok{sudo make install}
\end{Highlighting}
\end{Shaded}

If you keep the default \texttt{CMAKE\_INSTALL\_PREFIX} of
\texttt{/usr/local} then you have to make sure the path
\texttt{/usr/local/lib} is in your \texttt{LD\_LIBRARY\_PATH}.

You can also uninstall Embree again by executing:

\begin{Shaded}
\begin{Highlighting}[]
\NormalTok{sudo make uninstall}
\end{Highlighting}
\end{Shaded}

You can also create an Embree package using the following command:

\begin{Shaded}
\begin{Highlighting}[]
\NormalTok{make package}
\end{Highlighting}
\end{Shaded}

Please see the \hyperref[building-embree-applications]{Building Embree
Applications} section on how to build your application with such an
Embree package.

\section{Linux SYCL Compilation}\label{linux-sycl-compilation}

There are two options to compile Embree with SYCL support: The open
source \href{https://github.com/intel/llvm/}{``oneAPI DPC++ Compiler''}
or the
\href{https://www.intel.com/content/www/us/en/developer/articles/tool/oneapi-standalone-components.html\#dpcpp-cpp}{``Intel(R)
oneAPI DPC++/C++ Compiler''}. Other SYCL compilers are not supported.

The ``oneAPI DPC++ Compiler'' is more up-to-date than the ``Intel(R)
oneAPI DPC++/C++ Compiler'' but less stable. The current tested version
of the ``oneAPI DPC++ compiler is

\begin{itemize}
\tightlist
\item
  \href{https://github.com/intel/llvm/releases/download/sycl-nightly\%2F20230417/dpcpp-compiler.tar.gz}{oneAPI
  DPC++ Compiler 2023-04-17}
\end{itemize}

The compiler can be downloaded and simply extracted. The oneAPI DPC++
compiler 2022-12-14 can be set up executing the following command in a
Linux (bash) shell:

\begin{Shaded}
\begin{Highlighting}[]
\NormalTok{wget https}\OperatorTok{:}\CommentTok{//github.com/intel/llvm/releases/download/sycl{-}nightly\%2F20230417/dpcpp{-}compiler.tar.gz}
\NormalTok{tar xzf dpcpp}\OperatorTok{{-}}\NormalTok{compiler}\OperatorTok{.}\NormalTok{tar}\OperatorTok{.}\NormalTok{gz}
\NormalTok{source }\OperatorTok{./}\NormalTok{dpcpp\_compiler}\OperatorTok{/}\NormalTok{startup}\OperatorTok{.}\NormalTok{sh}
\end{Highlighting}
\end{Shaded}

The \texttt{startup.sh} script will put \texttt{clang++} and
\texttt{clang} from the oneAPI DPC++ Compiler into your path.

Please also install all Linux packages described in the previous
section.

Now, you can configure Embree using CMake by executing the following
command in the Embree root directory:

\begin{Shaded}
\begin{Highlighting}[]
\NormalTok{cmake }\OperatorTok{{-}}\NormalTok{B build \textbackslash{}}
      \OperatorTok{{-}}\NormalTok{DCMAKE\_CXX\_COMPILER}\OperatorTok{=}\NormalTok{clang}\OperatorTok{++}\NormalTok{ \textbackslash{}}
      \OperatorTok{{-}}\NormalTok{DCMAKE\_C\_COMPILER}\OperatorTok{=}\NormalTok{clang \textbackslash{}}
      \OperatorTok{{-}}\NormalTok{DEMBREE\_SYCL\_SUPPORT}\OperatorTok{=}\NormalTok{ON}
\end{Highlighting}
\end{Shaded}

This will create a directory \texttt{build} to use as the CMake build
directory, configure the usage of the oneAPI DPC++ Compiler, and turn on
SYCL support through \texttt{EMBREE\_SYCL\_SUPPORT=ON}.

Alternatively, you can download and run the installer of the

\begin{itemize}
\tightlist
\item
  \href{https://www.intel.com/content/www/us/en/developer/articles/tool/oneapi-standalone-components.html\#dpcpp-cpp}{Intel(R)
  oneAPI DPC++/C++ Compiler 2023.1.0}.
\end{itemize}

After installation, you can set up the compiler by sourcing the
\texttt{vars.sh} script in the \texttt{env} directory of the compiler
install directory, for example,

\begin{Shaded}
\begin{Highlighting}[]
\NormalTok{source }\OperatorTok{/}\NormalTok{opt}\OperatorTok{/}\NormalTok{intel}\OperatorTok{/}\NormalTok{oneAPI}\OperatorTok{/}\NormalTok{compiler}\OperatorTok{/}\FloatTok{2023.0}\ErrorTok{.0}\OperatorTok{/}\NormalTok{env}\OperatorTok{/}\NormalTok{vars}\OperatorTok{.}\NormalTok{sh}
\end{Highlighting}
\end{Shaded}

This script will put the \texttt{icpx} and \texttt{icx} compiler
executables from the Intel(R) oneAPI DPC++/C++ Compiler in your path.

Now, you can configure Embree using CMake by executing the following
command in the Embree root directory:

\begin{Shaded}
\begin{Highlighting}[]
\NormalTok{cmake }\OperatorTok{{-}}\NormalTok{B build \textbackslash{}}
      \OperatorTok{{-}}\NormalTok{DCMAKE\_CXX\_COMPILER}\OperatorTok{=}\NormalTok{icpx \textbackslash{}}
      \OperatorTok{{-}}\NormalTok{DCMAKE\_C\_COMPILER}\OperatorTok{=}\NormalTok{icx \textbackslash{}}
      \OperatorTok{{-}}\NormalTok{DEMBREE\_SYCL\_SUPPORT}\OperatorTok{=}\NormalTok{ON}
\end{Highlighting}
\end{Shaded}

More information about setting up the Intel(R) oneAPI DPC++/C++ compiler
can be found in the
\href{https://www.intel.com/content/www/us/en/develop/documentation/oneapi-dpcpp-cpp-compiler-dev-guide-and-reference/top/compiler-setup.html}{Development
Reference Guide}. Please note, that the Intel(R) oneAPI DPC++/C++
compiler requires
\href{https://www.intel.com/content/www/us/en/develop/documentation/oneapi-dpcpp-cpp-compiler-dev-guide-and-reference/top/compiler-setup/use-the-command-line/use-cmake-with-the-compiler.html}{at
least CMake version 3.20.5 on Linux}.

Independent of the DPC++ compiler choice, you can now build Embree using

\begin{Shaded}
\begin{Highlighting}[]
\NormalTok{cmake }\OperatorTok{{-}{-}}\NormalTok{build build }\OperatorTok{{-}}\NormalTok{j }\DecValTok{8}
\end{Highlighting}
\end{Shaded}

The executables will be generated inside the build folder. The
executable names of the SYCL versions of the tutorials end with
\texttt{\_sycl}.

\subsection{Linux Graphics Driver
Installation}\label{linux-graphics-driver-installation}

To run the SYCL code you need to install the latest GPGPU drivers for
your Intel Xe HPG/HPC GPUs from here \url{https://dgpu-docs.intel.com/}.
Follow the driver installation instructions for your graphics card and
operating system.

We tested Embree with the latest GPGPU driver Devel Release from
20220809. The Intel(R) Graphics Compute Runtime for oneAPI Level Zero
and OpenCL(TM) Driver from that release is too old for Embree to work
properly. Thus if no newer version of the GPGPU driver is available, you
need to additionally install the latest compute runtime from here
\href{https://github.com/intel/compute-runtime/releases/tag/23.09.25812.14}{23.09.25812.14}.

Unfortunately, these compute runtime packages are only available for
Ubuntu 22.04. You can also install a newer version of the compute
runtime if available.

\section{Windows}\label{windows}

Embree is tested using the following compilers under Windows:

\begin{itemize}
\tightlist
\item
  Intel® oneAPI DPC++/C++ Compiler 2023.1.0
\item
  Intel® oneAPI DPC++/C++ Compiler 2023.0.0
\item
  oneAPI DPC++/C++ Compiler 2023-04-17
\item
  Visual Studio 2019
\item
  Visual Studio 2017
\item
  Visual Studio 2015 (Update~1)
\item
  Intel® Implicit SPMD Program Compiler 1.19.0
\item
  Intel® Implicit SPMD Program Compiler 1.18.1
\item
  Intel® Implicit SPMD Program Compiler 1.17.0
\item
  Intel® Implicit SPMD Program Compiler 1.16.1
\item
  Intel® Implicit SPMD Program Compiler 1.15.0
\item
  Intel® Implicit SPMD Program Compiler 1.14.1
\item
  Intel® Implicit SPMD Program Compiler 1.13.0
\item
  Intel® Implicit SPMD Program Compiler 1.12.0
\end{itemize}

To compile Embree for AVX-512 you have to use the Intel® Compiler.

Embree supports using the Intel® Threading Building Blocks (TBB) as the
tasking system. For performance and flexibility reasons we recommend
using use Embree with the Intel® Threading Building Blocks (TBB) and
best also use TBB inside your application. Optionally you can disable
TBB in Embree through the \texttt{EMBREE\_TASKING\_SYSTEM} CMake
variable.

Embree will either find the Intel® Threading Building Blocks (TBB)
installation that comes with the Intel® Compiler, or you can install the
binary distribution of TBB directly from
\url{https://github.com/oneapi-src/oneTBB/releases} into a folder named
\texttt{tbb} into your Embree root directory. You also have to make sure
that the libraries \texttt{tbb.dll} and \texttt{tbb\_malloc.dll} can be
found when executing your Embree applications, e.g.~by putting the path
to these libraries into your \texttt{PATH} environment variable.

Embree supports the Intel® Implicit SPMD Program Compiler (Intel® ISPC),
which allows straightforward parallelization of an entire renderer. When
installing Intel® ISPC, make sure to download an Intel® ISPC version
from \href{https://ispc.github.io/downloads.html}{ispc.github.io} that
is compatible with your Visual Studio version. After installation, put
the path to \texttt{ispc.exe} permanently into your \texttt{PATH}
environment variable or you need to correctly set the
\texttt{EMBREE\_ISPC\_EXECUTABLE} variable during CMake configuration to
point to the ISPC executable. If you want to use Intel® ISPC, you have
to enable \texttt{EMBREE\_ISPC\_SUPPORT} in CMake.

You additionally have to install
\href{http://www.cmake.org/download/}{CMake} (version 3.1 or higher).
Note that you need a native Windows CMake installation because CMake
under Cygwin cannot generate solution files for Visual Studio.

\subsection{Using the IDE}\label{using-the-ide}

Run \texttt{cmake-gui}, browse to the Embree sources, set the build
directory and click Configure. Now you can select the Generator,
e.g.~``Visual Studio 12 2013'' for a 32-bit build or ``Visual Studio 12
2013 Win64'' for a 64-bit build.

To use a different compiler than the Microsoft Visual C++ compiler, you
additionally need to specify the proper compiler toolset through the
option ``Optional toolset to use (-T parameter)''. E.g. to use Clang for
compilation set the toolset to ``LLVM\_v142''.

Do not change the toolset manually in a solution file (neither through
the project properties dialog nor through the ``Use Intel Compiler''
project context menu), because then some compiler-specific command line
options cannot be set by CMake.

Most configuration parameters described in the
\hyperref[cmake-configuration]{CMake Configuration} can be set under
Windows as well. Finally, click ``Generate'' to create the Visual Studio
solution files.

The following CMake options are only available under Windows:

\begin{itemize}
\item
  \texttt{CMAKE\_CONFIGURATION\_TYPE}: List of generated configurations.
  The default value is Debug;Release;RelWithDebInfo.
\item
  \texttt{USE\_STATIC\_RUNTIME}: Use the static version of the C/C++
  runtime library. This option is turned OFF by default.
\end{itemize}

Use the generated Visual Studio solution file \texttt{embree4.sln} to
compile the project.

We recommend enabling syntax highlighting for the \texttt{.ispc} source
and \texttt{.isph} header files. To do so open Visual Studio, go to
Tools ⇒ Options ⇒ Text Editor ⇒ File Extension and add the \texttt{isph}
and \texttt{ispc} extensions for the ``Microsoft Visual C++'' editor.

\subsection{Using the Command Line}\label{using-the-command-line}

Embree can also be configured and built without the IDE using the Visual
Studio command prompt:

\begin{Shaded}
\begin{Highlighting}[]
\NormalTok{cd path\textbackslash{}to\textbackslash{}embree}
\NormalTok{mkdir build}
\NormalTok{cd build}
\NormalTok{cmake }\OperatorTok{{-}}\NormalTok{G }\StringTok{"Visual Studio 16 2019"} \OperatorTok{..}
\NormalTok{cmake }\OperatorTok{{-}{-}}\NormalTok{build }\OperatorTok{.} \OperatorTok{{-}{-}}\NormalTok{config Release}
\end{Highlighting}
\end{Shaded}

You can also build only some projects with the \texttt{-\/-target}
switch. Additional parameters after ``\texttt{-\/-}'' will be passed to
\texttt{msbuild}. For example, to build the Embree library in parallel
use

\begin{Shaded}
\begin{Highlighting}[]
\NormalTok{cmake }\OperatorTok{{-}{-}}\NormalTok{build }\OperatorTok{.} \OperatorTok{{-}{-}}\NormalTok{config Release }\OperatorTok{{-}{-}}\NormalTok{target embree }\OperatorTok{{-}{-}} \OperatorTok{/}\NormalTok{m}
\end{Highlighting}
\end{Shaded}

\subsection{Building Embree - Using
vcpkg}\label{building-embree---using-vcpkg}

You can download and install Embree using the
\href{https://github.com/Microsoft/vcpkg}{vcpkg} dependency manager:

\begin{Shaded}
\begin{Highlighting}[]
\NormalTok{git clone https}\OperatorTok{:}\CommentTok{//github.com/Microsoft/vcpkg.git}
\NormalTok{cd vcpkg}
\OperatorTok{./}\NormalTok{bootstrap}\OperatorTok{{-}}\NormalTok{vcpkg}\OperatorTok{.}\NormalTok{sh}
\OperatorTok{./}\NormalTok{vcpkg integrate install}
\OperatorTok{./}\NormalTok{vcpkg install embree3}
\end{Highlighting}
\end{Shaded}

The Embree port in vcpkg is kept up to date by Microsoft team members
and community contributors. If the version is out of date, please
\href{https://github.com/Microsoft/vcpkg}{create an issue or pull
request} on the vcpkg repository.

\section{Windows SYCL Compilation}\label{windows-sycl-compilation}

There are two options to compile Embree with SYCL support: The open
source \href{https://github.com/intel/llvm/}{``oneAPI DPC++ Compiler''}
or the
\href{https://www.intel.com/content/www/us/en/developer/articles/tool/oneapi-standalone-components.html\#dpcpp-cpp}{``Intel(R)
oneAPI DPC++/C++ Compiler''}. Other SYCL compilers are not supported.
You will also need an installed version of Visual Studio that supports
the C++17 standard, e.g.~Visual Studio 2019.

The ``oneAPI DPC++ Compiler'' is more up-to-date than the ``Intel(R)
oneAPI DPC++/C++ Compiler'' but less stable. The current tested version
of the oneAPI DPC++ compiler is

\begin{itemize}
\tightlist
\item
  \href{https://github.com/intel/llvm/releases/download/sycl-nightly\%2F20230417/dpcpp-compiler-win.tar.gz}{oneAPI
  DPC++ Compiler 2023-04-17}
\end{itemize}

Download and unpack the archive and open the ``x64 Native Tools Command
Prompt'' of Visual Studio and execute the following lines to properly
configure the environment to use the oneAPI DPC++ compiler:

\begin{Shaded}
\begin{Highlighting}[]
\NormalTok{set }\StringTok{"DPCPP\_DIR=path\_to\_dpcpp\_compiler"}
\NormalTok{set }\StringTok{"PATH=\%DPCPP\_DIR\%}\SpecialCharTok{\textbackslash{}b}\StringTok{in;\%PATH\%"}
\NormalTok{set }\StringTok{"PATH=\%DPCPP\_DIR\%}\ErrorTok{\textbackslash{}l}\StringTok{ib;\%PATH\%"}
\NormalTok{set }\StringTok{"CPATH=\%DPCPP\_DIR\%}\ErrorTok{\textbackslash{}i}\StringTok{nclude;\%CPATH\%"}
\NormalTok{set }\StringTok{"INCLUDE=\%DPCPP\_DIR\%}\ErrorTok{\textbackslash{}i}\StringTok{nclude;\%INCLUDE\%"}
\NormalTok{set }\StringTok{"LIB=\%DPCPP\_DIR\%}\ErrorTok{\textbackslash{}l}\StringTok{ib;\%LIB\%"}
\end{Highlighting}
\end{Shaded}

The \texttt{path\_to\_dpcpp\_compiler} should point to the unpacked
oneAPI DPC++ compiler.

Now, you can configure Embree using CMake by executing the following
command in the Embree root directory:

\begin{Shaded}
\begin{Highlighting}[]
\NormalTok{cmake }\OperatorTok{{-}}\NormalTok{B build}
      \OperatorTok{{-}}\NormalTok{G Ninja}
      \OperatorTok{{-}}\NormalTok{D CMAKE\_BUILD\_TYPE}\OperatorTok{=}\NormalTok{Release}
      \OperatorTok{{-}}\NormalTok{D CMAKE\_CXX\_COMPILER}\OperatorTok{=}\NormalTok{clang}\OperatorTok{++}
      \OperatorTok{{-}}\NormalTok{D CMAKE\_C\_COMPILER}\OperatorTok{=}\NormalTok{clang}
      \OperatorTok{{-}}\NormalTok{D EMBREE\_SYCL\_SUPPORT}\OperatorTok{=}\NormalTok{ON}
      \OperatorTok{{-}}\NormalTok{D TBB\_ROOT}\OperatorTok{=}\NormalTok{path\_to\_tbb\textbackslash{}lib\textbackslash{}cmake\textbackslash{}tbb}
\end{Highlighting}
\end{Shaded}

This will create a directory \texttt{build} to use as the CMake build
directory, and configure a release build that uses \texttt{clang++} and
\texttt{clang} from the oneAPI DPC++ compiler.

The \href{https://ninja-build.org/}{Ninja} generator is currently the
easiest way to use the oneAPI DPC++ compiler.

We also enable SYCL support in Embree using the
\texttt{EMBREE\_SYCL\_SUPPORT} CMake option.

Alternatively, you can download and run the installer of the

\begin{itemize}
\tightlist
\item
  \href{https://www.intel.com/content/www/us/en/developer/articles/tool/oneapi-standalone-components.html\#dpcpp-cpp}{Intel(R)
  oneAPI DPC++/C++ Compiler 2023.0.0}.
\end{itemize}

After installation, you can either open a regular
\texttt{Command\ Prompt} and execute the \texttt{vars.bat} script in the
\texttt{env} directory of the compiler install directory, for example

\begin{Shaded}
\begin{Highlighting}[]
\NormalTok{C}\OperatorTok{:}\NormalTok{\textbackslash{}Program Files }\OperatorTok{(}\NormalTok{x86}\OperatorTok{)}\NormalTok{\textbackslash{}Intel\textbackslash{}oneAPI\textbackslash{}compiler\textbackslash{}}\FloatTok{2023.0}\ErrorTok{.0}\NormalTok{\textbackslash{}env\textbackslash{}vars}\OperatorTok{.}\NormalTok{bat}
\end{Highlighting}
\end{Shaded}

or simply open the installed ``Intel oneAPI command prompt for Intel 64
for Visual Studio''.

Both ways will put the \texttt{icx} compiler executable from the
Intel(R) oneAPI DPC++/C++ compiler in your path.

Now, you can configure Embree using CMake by executing the following
command in the Embree root directory:

\begin{Shaded}
\begin{Highlighting}[]
\NormalTok{cmake }\OperatorTok{{-}}\NormalTok{B build}
      \OperatorTok{{-}}\NormalTok{G Ninja}
      \OperatorTok{{-}}\NormalTok{D CMAKE\_BUILD\_TYPE}\OperatorTok{=}\NormalTok{Release}
      \OperatorTok{{-}}\NormalTok{D CMAKE\_CXX\_COMPILER}\OperatorTok{=}\NormalTok{icx}
      \OperatorTok{{-}}\NormalTok{D CMAKE\_C\_COMPILER}\OperatorTok{=}\NormalTok{icx}
      \OperatorTok{{-}}\NormalTok{D EMBREE\_SYCL\_SUPPORT}\OperatorTok{=}\NormalTok{ON}
      \OperatorTok{{-}}\NormalTok{D TBB\_ROOT}\OperatorTok{=}\NormalTok{path\_to\_tbb\textbackslash{}lib\textbackslash{}cmake\textbackslash{}tbb}
\end{Highlighting}
\end{Shaded}

More information about setting up the Intel(R) oneAPI DPC++/C++ compiler
can be found in the
\href{https://www.intel.com/content/www/us/en/develop/documentation/oneapi-dpcpp-cpp-compiler-dev-guide-and-reference/top/compiler-setup.html}{Development
Reference Guide}. Please note, that the Intel(R) oneAPI DPC++/C++
compiler requires
\href{https://www.intel.com/content/www/us/en/develop/documentation/oneapi-dpcpp-cpp-compiler-dev-guide-and-reference/top/compiler-setup/use-the-command-line/use-cmake-with-the-compiler.html}{at
least CMake version 3.23 on Windows}.

Independent of the DPC++ compiler choice, you can now build Embree using

\begin{Shaded}
\begin{Highlighting}[]
\NormalTok{cmake }\OperatorTok{{-}{-}}\NormalTok{build build}
\end{Highlighting}
\end{Shaded}

If you have problems with Ninja re-running CMake in an infinite loop,
then first remove the ``Re-run CMake if any of its inputs changed.''
section from the \texttt{build.ninja} file and run the above command
again.

You can also create an Embree package using the following command:

\begin{Shaded}
\begin{Highlighting}[]
\NormalTok{cmake }\OperatorTok{{-}{-}}\NormalTok{build build }\OperatorTok{{-}{-}}\NormalTok{target package}
\end{Highlighting}
\end{Shaded}

Please see the \hyperref[building-embree-sycl-applications]{Building
Embree SYCL Applications} section on how to build your application with
such an Embree package.

\subsection{Windows Graphics Driver
Installation}\label{windows-graphics-driver-installation}

In order to run the SYCL tutorials on HPG hardware, you first need to
install the graphics drivers for your graphics card from
\url{https://www.intel.com}. Please make sure to have installed version
31.0.101.4314 or newer.

\section{CMake Configuration}\label{cmake-configuration}

The default CMake configuration in the configuration dialog should be
appropriate for most usages. The following list describes all parameters
that can be configured in CMake:

\begin{itemize}
\item
  \texttt{CMAKE\_BUILD\_TYPE}: Can be used to switch between Debug mode
  (Debug), Release mode (Release) (default), and Release mode with
  enabled assertions and debug symbols (RelWithDebInfo).
\item
  \texttt{EMBREE\_STACK\_PROTECTOR}: Enables protection of return
  address from buffer overwrites. This option is OFF by default.
\item
  \texttt{EMBREE\_ISPC\_SUPPORT}: Enables Intel® ISPC support of Embree.
  This option is OFF by default.
\item
  \texttt{EMBREE\_SYCL\_SUPPORT}: Enables GPU support using SYCL. When
  this option is enabled you have to use some DPC++ compiler. Please see
  the sections \hyperref[linux-sycl-compilation]{Linux SYCL Compilation}
  and \hyperref[windows-sycl-compilation]{Windows SYCL Compilation} on
  supported DPC++ compilers. This option is OFF by default.
\item
  \texttt{EMBREE\_SYCL\_AOT\_DEVICES}: Selects a list of GPU devices for
  ahead-of-time (AOT) compilation of device code. Possible values are
  either, ``none'' which enables only just in time (JIT) compilation, or
  a list of the Embree-supported Xe GPUs for AOT compilation:

  \begin{itemize}
  \tightlist
  \item
    XE\_HPG\_CORE : Xe HPG devices
  \item
    XE\_HPC\_CORE : Xe HPC devices
  \end{itemize}

  One can also specify multiple devices separated by comma to compile
  ahead of time for multiple devices,
  e.g.~``XE\_HPG\_CORE,XE\_HP\_CORE''. When enabling AOT compilation for
  one or multiple devices, JIT compilation will always additionally be
  enabled in case the code is executed on a device no code is
  precompiled for.

  Execute ``ocloc compile --help'' for more details of possible devices
  to pass. Embree is only supported on Xe HPG/HPC and newer devices.

  Per default, this option is set to ``none'' to enable JIT compilation.
  We recommend using JIT compilation as this enables the use of
  specialization constants to reduce code complexity.
\item
  \texttt{EMBREE\_STATIC\_LIB}: Builds Embree as a static library (OFF
  by default). Further multiple static libraries are generated for the
  different ISAs selected (e.g.~\texttt{embree4.a},
  \texttt{embree4\_sse42.a}, \texttt{embree4\_avx.a},
  \texttt{embree4\_avx2.a}, \texttt{embree4\_avx512.a}). You have to
  link these libraries in exactly this order of increasing ISA.
\item
  \texttt{EMBREE\_API\_NAMESPACE}: Specifies a namespace name to put all
  Embree API symbols inside. By default, no namespace is used and plain
  C symbols are exported.
\item
  \texttt{EMBREE\_LIBRARY\_NAME}: Specifies the name of the Embree
  library file created. By default, the name embree4 is used.
\item
  \texttt{EMBREE\_IGNORE\_CMAKE\_CXX\_FLAGS}: When enabled, Embree
  ignores default CMAKE\_CXX\_FLAGS. This option is turned ON by
  default.
\item
  \texttt{EMBREE\_TUTORIALS}: Enables build of Embree tutorials (default
  ON).
\item
  \texttt{EMBREE\_BACKFACE\_CULLING}: Enables backface culling,
  i.e.~only surfaces facing a ray can be hit. This option is turned OFF
  by default.
\item
  \texttt{EMBREE\_BACKFACE\_CULLING\_CURVES}: Enables backface culling
  for curves, i.e.~only surfaces facing a ray can be hit. This option is
  turned OFF by default.
\item
  \texttt{EMBREE\_BACKFACE\_CULLING\_SPHERES}: Enables backface culling
  for spheres, i.e.~only surfaces facing a ray can be hit. This option
  is turned OFF by default.
\item
  \texttt{EMBREE\_COMPACT\_POLYS}: Enables compact tris/quads, i.e.~only
  geomIDs and primIDs are stored inside the leaf nodes.
\item
  \texttt{EMBREE\_FILTER\_FUNCTION}: Enables the intersection filter
  function feature (ON by default).
\item
  \texttt{EMBREE\_RAY\_MASK}: Enables the ray masking feature (OFF by
  default).
\item
  \texttt{EMBREE\_RAY\_PACKETS}: Enables ray packet traversal kernels.
  This feature is turned ON by default. When turned on packet traversal
  is used internally and packets passed to rtcIntersect4/8/16 are kept
  intact in callbacks (when the ISA of appropriate width is enabled).
\item
  \texttt{EMBREE\_IGNORE\_INVALID\_RAYS}: Makes code robust against the
  risk of full-tree traversals caused by invalid rays (e.g.~rays
  containing INF/NaN as origins). This option is turned OFF by default.
\item
  \texttt{EMBREE\_TASKING\_SYSTEM}: Chooses between Intel® Threading TBB
  Building Blocks (TBB), Parallel Patterns Library (PPL) (Windows only),
  or an internal tasking system (INTERNAL). By default, TBB is used.
\item
  \texttt{EMBREE\_TBB\_ROOT}: If Intel® Threading Building Blocks (TBB)
  is used as a tasking system, search the library in this directory
  tree.
\item
  \texttt{EMBREE\_TBB\_COMPONENT}: The component/library name of Intel®
  Threading Building Blocks (TBB). Embree searches for this library name
  (default: tbb) when TBB is used as the tasking system.
\item
  \texttt{EMBREE\_TBB\_POSTFIX}: If Intel® Threading Building Blocks
  (TBB) is used as a tasking system, link to tbb.(so,dll,lib). Defaults
  to the empty string.
\item
  \texttt{EMBREE\_TBB\_DEBUG\_ROOT}: If Intel® Threading Building Blocks
  (TBB) is used as a tasking system, search the library in this
  directory tree in Debug mode. Defaults to \texttt{EMBREE\_TBB\_ROOT}.
\item
  \texttt{EMBREE\_TBB\_DEBUG\_POSTFIX}: If Intel® Threading Building
  Blocks (TBB) is used as a tasking system, link to tbb.(so,dll,lib) in
  Debug mode. Defaults to ``\_debug''.
\item
  \texttt{EMBREE\_MAX\_ISA}: Select highest supported ISA (SSE2, SSE4.2,
  AVX, AVX2, AVX512, or NONE). When set to NONE the EMBREE\_ISA\_*
  variables can be used to enable ISAs individually. By default, the
  option is set to AVX2.
\item
  \texttt{EMBREE\_ISA\_SSE2}: Enables SSE2 when EMBREE\_MAX\_ISA is set
  to NONE. By default, this option is turned OFF.
\item
  \texttt{EMBREE\_ISA\_SSE42}: Enables SSE4.2 when EMBREE\_MAX\_ISA is
  set to NONE. By default, this option is turned OFF.
\item
  \texttt{EMBREE\_ISA\_AVX}: Enables AVX when EMBREE\_MAX\_ISA is set to
  NONE. By default, this option is turned OFF.
\item
  \texttt{EMBREE\_ISA\_AVX2}: Enables AVX2 when EMBREE\_MAX\_ISA is set
  to NONE. By default, this option is turned OFF.
\item
  \texttt{EMBREE\_ISA\_AVX512}: Enables AVX-512 for Skylake when
  EMBREE\_MAX\_ISA is set to NONE. By default, this option is turned
  OFF.
\item
  \texttt{EMBREE\_GEOMETRY\_TRIANGLE}: Enables support for triangle
  geometries (ON by default).
\item
  \texttt{EMBREE\_GEOMETRY\_QUAD}: Enables support for quad geometries
  (ON by default).
\item
  \texttt{EMBREE\_GEOMETRY\_CURVE}: Enables support for curve geometries
  (ON by default).
\item
  \texttt{EMBREE\_GEOMETRY\_SUBDIVISION}: Enables support for
  subdivision geometries (ON by default).
\item
  \texttt{EMBREE\_GEOMETRY\_INSTANCE}: Enables support for instances (ON
  by default).
\item
  \texttt{EMBREE\_GEOMETRY\_USER}: Enables support for user-defined
  geometries (ON by default).
\item
  \texttt{EMBREE\_GEOMETRY\_POINT}: Enables support for point geometries
  (ON by default).
\item
  \texttt{EMBREE\_CURVE\_SELF\_INTERSECTION\_AVOIDANCE\_FACTOR}:
  Specifies a factor that controls the self-intersection avoidance
  feature for flat curves. Flat curve intersections which are closer
  than
  curve\_radius*\texttt{EMBREE\_CURVE\_SELF\_INTERSECTION\_AVOIDANCE\_FACTOR}
  to the ray origin are ignored. A value of 0.0f disables
  self-intersection avoidance while 2.0f is the default value.
\item
  \texttt{EMBREE\_DISC\_POINT\_SELF\_INTERSECTION\_AVOIDANCE}: Enables
  self-intersection avoidance for RTC\_GEOMETRY\_TYPE\_DISC\_POINT
  geometry type (ON by default). When enabled intersections are skipped
  if the ray origin lies inside the sphere defined by the point
  primitive.
\item
  \texttt{EMBREE\_MIN\_WIDTH}: Enabled the min-width feature, which
  allows increasing the radius of curves and points to match some amount
  of pixels. See
  \hyperref[rtcsetgeometrymaxradiusscale]{rtcSetGeometryMaxRadiusScale}
  for more details.
\item
  \texttt{EMBREE\_MAX\_INSTANCE\_LEVEL\_COUNT}: Specifies the maximum
  number of nested instance levels. Should be greater than 0; the
  default value is 1. Instances nested any deeper than this value will
  silently disappear in release mode, and cause assertions in debug
  mode.
\end{itemize}

\pagebreak

\chapter{Embree API}\label{embree-api}

The Embree API is a low-level C99 ray tracing API which can be used to
build spatial index structures for 3D scenes and perform ray queries of
different types.

The API can get used on the CPU using standard C, C++, and ISPC code and
Intel GPUs by using SYCL code.

The Intel® Implicit SPMD Program Compiler (Intel® ISPC) version of the
API, is almost identical to the standard C99 version, but contains
additional functions that operate on ray packets with a size of the
native SIMD width used by Intel® ISPC.

The SYCL version of the API is also mostly identical to the C99 version
of the API, with some exceptions listed in section
\hyperref[embree-sycl-api]{Embree SYCL API}.

For simplicity this document refers to the C99 version of the API
functions. For changes when upgrading from the Embree 3 to the current
Embree 4 API see Section
\hyperref[upgrading-from-embree-3-to-embree-4]{Upgrading from Embree 3
to Embree 4}.

All API calls carry the prefix \texttt{rtc} (or \texttt{RTC} for types)
which stands for \textbf{r}ay \textbf{t}racing \textbf{c}ore. The API
supports scenes consisting of different geometry types such as triangle
meshes, quad meshes (triangle pairs), grid meshes, flat curves, round
curves, oriented curves, subdivision meshes, instances, and user-defined
geometries. See Section \hyperref[scene-object]{Scene Object} for more
information.

Finding the closest hit of a ray segment with the scene
(\texttt{rtcIntersect}-type functions), and determining whether any hit
between a ray segment and the scene exists (\texttt{rtcOccluded}-type
functions) are both supported. The API supports queries for single rays
and ray packets. See Section \hyperref[ray-queries]{Ray Queries} for
more information.

The API is designed in an object-oriented manner, e.g.~it contains
device objects (\texttt{RTCDevice} type), scene objects
(\texttt{RTCScene} type), geometry objects (\texttt{RTCGeometry} type),
buffer objects (\texttt{RTCBuffer} type), and BVH objects
(\texttt{RTCBVH} type). All objects are reference counted, and handles
can be released by calling the appropriate release function
(e.g.~\texttt{rtcReleaseDevice}) or retained by incrementing the
reference count (e.g.~\texttt{rtcRetainDevice}). In general, API calls
that access the same object are not thread-safe, unless specified
otherwise. However, attaching geometries to the same scene and
performing ray queries in a scene is thread-safe.

\section{Device Object}\label{device-object}

Embree supports a device concept, which allows different components of
the application to use the Embree API without interfering with each
other. An application typically first creates a device using the
\hyperref[rtcnewdevice]{rtcNewDevice} function (or
\hyperref[rtcnewsycldevice]{rtcNewSYCLDevice} when using SYCL for the
GPU). This device can then be used to construct further objects, such as
scenes and geometries. Before the application exits, it should release
all devices by invoking \hyperref[rtcreleasedevice]{rtcReleaseDevice}.
An application typically creates only a single device. If required
differently, it should only use a small number of devices at any given
time.

Each user thread has its own error flag per device. If an error occurs
when invoking an API function, this flag is set to an error code (if it
isn't already set by a previous error). See Section
\hyperref[rtcgetdeviceerror]{rtcGetDeviceError} for information on how
to read the error code and Section
\hyperref[rtcsetdeviceerrorfunction]{rtcSetDeviceErrorFunction} on how
to register a callback that is invoked for each error encountered. It is
recommended to always set a error callback function, to detect all
errors.

\section{Scene Object}\label{scene-object}

A scene is a container for a set of geometries, and contains a spatial
acceleration structure which can be used to perform different types of
ray queries.

A scene is created using the \texttt{rtcNewScene} function call, and
released using the \texttt{rtcReleaseScene} function call. To populate a
scene with geometries use the \texttt{rtcAttachGeometry} call, and to
detach them use the \texttt{rtcDetachGeometry} call. Once all scene
geometries are attached, an \texttt{rtcCommitScene} call (or
\texttt{rtcJoinCommitScene} call) will finish the scene description and
trigger building of internal data structures. After the scene got
committed, it is safe to perform ray queries (see Section
\hyperref[ray-queries]{Ray Queries}) or to query the scene bounding box
(see \hyperref[rtcgetscenebounds]{rtcGetSceneBounds} and
\hyperref[rtcgetscenelinearbounds]{rtcGetSceneLinearBounds}).

If scene geometries get modified or attached or detached, the
\texttt{rtcCommitScene} call must be invoked before performing any
further ray queries for the scene; otherwise the effect of the ray query
is undefined. The modification of a geometry, committing the scene, and
tracing of rays must always happen sequentially, and never at the same
time. Any API call that sets a property of the scene or geometries
contained in the scene count as scene modification, e.g.~including
setting of intersection filter functions.

Scene flags can be used to configure a scene to use less memory
(\texttt{RTC\_SCENE\_FLAG\_COMPACT}), use more robust traversal
algorithms (\texttt{RTC\_SCENE\_FLAG\_ROBUST}), and to optimize for
dynamic content. See Section
\hyperref[rtcsetsceneflags]{rtcSetSceneFlags} for more details.

A build quality can be specified for a scene to balance between
acceleration structure build performance and ray query performance. See
Section \hyperref[rtcsetscenebuildquality]{rtcSetSceneBuildQuality} for
more details on build quality.

\section{Geometry Object}\label{geometry-object}

A new geometry is created using the \texttt{rtcNewGeometry} function.
Depending on the geometry type, different buffers must be bound (e.g.
using \texttt{rtcSetSharedGeometryBuffer}) to set up the geometry data.
In most cases, binding of a vertex and index buffer is required. The
number of primitives and vertices of that geometry is typically inferred
from the size of these bound buffers.

Changes to the geometry always must be committed using the
\texttt{rtcCommitGeometry} call before using the geometry. After
committing, a geometry is not included in any scene. A geometry can be
added to a scene by using the \texttt{rtcAttachGeometry} function (to
automatically assign a geometry ID) or using the
\texttt{rtcAttachGeometryById} function (to specify the geometry ID
manually). A geometry can get attached to multiple scenes.

All geometry types support multi-segment motion blur with an arbitrary
number of equidistant time steps (in the range of 2 to 129) inside a
user specified time range. Each geometry can have a different number of
time steps and a different time range. The motion blur geometry is
defined by linearly interpolating the geometries of neighboring time
steps. To construct a motion blur geometry, first the number of time
steps of the geometry must be specified using the
\texttt{rtcSetGeometryTimeStepCount} function, and then a vertex buffer
for each time step must be bound, e.g.~using the
\texttt{rtcSetSharedGeometryBuffer} function. Optionally, a time range
defining the start (and end time) of the first (and last) time step can
be set using the \texttt{rtcSetGeometryTimeRange} function. This feature
will also allow geometries to appear and disappear during the camera
shutter time if the time range is a sub range of {[}0,1{]}.

\section{Ray Queries}\label{ray-queries}

The API supports finding the closest hit of a ray segment with the scene
(\texttt{rtcIntersect}-type functions), and determining whether any hit
between a ray segment and the scene exists (\texttt{rtcOccluded}-type
functions).

Supported are single ray queries (\texttt{rtcIntersect1} and
\texttt{rtcOccluded1}) as well as ray packet queries for ray packets of
size 4 (\texttt{rtcIntersect4} and \texttt{rtcOccluded4}), ray packets
of size 8 (\texttt{rtcIntersect8} and \texttt{rtcOccluded8}), and ray
packets of size 16 (\texttt{rtcIntersect16} and \texttt{rtcOccluded16}).

See Sections \hyperref[rtcintersect1]{rtcIntersect1} and
\hyperref[rtcoccluded1]{rtcOccluded1} for a detailed description of how
to set up and trace a ray.

See tutorial \hyperref[triangle-geometry]{Triangle Geometry} for a
complete example of how to trace single rays and ray packets.

\section{Point Queries}\label{point-queries}

The API supports traversal of the BVH using a point query object that
specifies a location and a query radius. For all primitives intersecting
the according domain, a user defined callback function is called which
allows queries such as finding the closest point on the surface
geometries of the scene (see Tutorial \hyperref[closest-point]{Closest
Point}) or nearest neighbour queries (see Tutorial
\hyperref[voronoi]{Voronoi}).

See Section \hyperref[rtcpointquery]{rtcPointQuery} for a detailed
description of how to set up point queries.

\section{Collision Detection}\label{collision-detection}

The Embree API also supports collision detection queries between two
scenes consisting only of user geometries. Embree only performs
broadphase collision detection, the narrow phase detection can be
performed through a callback function.

See Section \hyperref[rtccollide]{rtcCollide} for a detailed description
of how to set up collision detection.

Seen tutorial \hyperref[collision-detection]{Collision Detection} for a
complete example of collision detection being used on a simple cloth
solver.

\section{Filter Functions}\label{filter-functions}

The API supports filter functions that are invoked for each intersection
found during the \texttt{rtcIntersect}-type or \texttt{rtcOccluded}-type
calls.

The filter functions can be set per-geometry using the
\texttt{rtcSetGeometryIntersectFilterFunction} and
\texttt{rtcSetGeometryOccludedFilterFunction} calls. The former ones are
called geometry intersection filter functions, the latter ones geometry
occlusion filter functions. These filter functions are designed to be
used to ignore intersections outside of a user-defined silhouette of a
primitive, e.g.~to model tree leaves using transparency textures.

The filter function can also get passed as arguments directly to the
traversal functions, see section
\hyperref[rtcinitintersectarguments]{rtcInitIntersectArguments} and
\hyperref[rtcinitoccludedarguments]{rtcInitOccludedArguments} for more
details. These argument filter functions are designed to change the
semantics of the ray query, e.g.~to accumulate opacity for transparent
shadows, count the number of surfaces along a ray, collect all hits
along a ray, etc. The argument filter function must be enabled to be
used for a scene using the
\texttt{RTC\_SCENE\_FLAG\_FILTER\_FUNCTION\_IN\_ARGUMENTS} scene flag.
The callback is only invoked for geometries that enable the callback
using the \texttt{rtcSetGeometryEnableFilterFunctionFromArguments} call,
or enabled for all geometries when the
\texttt{RTC\_RAY\_QUERY\_FLAG\_INVOKE\_ARGUMENT\_FILTER} ray query flag
is set.

\section{BVH Build API}\label{bvh-build-api}

The internal algorithms to build a BVH are exposed through the
\texttt{RTCBVH} object and \texttt{rtcBuildBVH} call. This call makes it
possible to build a BVH in a user-specified format over user-specified
primitives. See the documentation of the \texttt{rtcBuildBVH} call for
more details.

\chapter{Embree SYCL API}\label{embree-sycl-api}

Embree supports ray tracing on Intel GPUs by using the SYCL programming
language. SYCL is a Khronos standardized C++ based language for single
source heterogenous programming for acceleration offload, see the
\href{https://www.khronos.org/sycl/}{SYCL webpage} for details.

The Embree SYCL API is designed for photorealistic rendering use cases,
where scene setup is performed on the host, and rendering on the device.
The Embree SYCL API is very similar to the standard Embree C99 API, and
supports most of its features, such as all triangle-type geometries, all
curve types and basis functions, point geometry types, user geometries,
filter callbacks, multi-level instancing, and motion blur.

To enable SYCL support you have to include the \texttt{sycl.hpp} file
before the Embree API headers:

\begin{Shaded}
\begin{Highlighting}[]
\PreprocessorTok{\#include }\ImportTok{\textless{}sycl/sycl.hpp\textgreater{}}
\PreprocessorTok{\#include }\ImportTok{\textless{}embree4/rtcore.h\textgreater{}}
\end{Highlighting}
\end{Shaded}

Next you need to initializes an Embree SYCL device using the
\texttt{rtcNewSYCLDevice} API function by providing a SYCL context.

Embree provides the \texttt{rtcIsSYCLDeviceSupported} API function to
check if some SYCL device is supported by Embree. You can also use the
\texttt{rtcSYCLDeviceSelector} to conveniently select the first SYCL
device that is supported by Embree, e.g.:

\begin{Shaded}
\begin{Highlighting}[]
\NormalTok{sycl}\OperatorTok{::}\NormalTok{device device}\OperatorTok{(}\NormalTok{rtcSYCLDeviceSelector}\OperatorTok{);}
\NormalTok{sycl}\OperatorTok{::}\NormalTok{queue queue}\OperatorTok{(}\NormalTok{device}\OperatorTok{,}\NormalTok{ exception\_handler}\OperatorTok{);}
\NormalTok{sycl}\OperatorTok{::}\NormalTok{context context}\OperatorTok{(}\NormalTok{device}\OperatorTok{);}
\NormalTok{RTCDevice device }\OperatorTok{=}\NormalTok{ rtcNewSYCLDevice}\OperatorTok{(}\NormalTok{context}\OperatorTok{,}\StringTok{""}\OperatorTok{);}
\end{Highlighting}
\end{Shaded}

Scenes created with an Embree SYCL device can only get used to trace
rays using SYCL on the GPU, it is not possible to trace rays on the CPU
with such a device. To render on the CPU and GPU in parallel, the user
has to create a second Embree device and create a second scene to be
used on the CPU.

Files containing SYCL code, have to get compiled with the Intel® oneAPI
DPC++ compiler. Please see section
\hyperref[linux-sycl-compilation]{Linux SYCL Compilation} and
\hyperref[windows-sycl-compilation]{Windows SYCL Compilation} for
supported compilers. The DPC++ compiler performs a two-phase
compilation, where host code is compiled in a first phase, and device
code compiled in a second compilation phase.

Standard Embree API functions for scene construction can get used on the
host but not the device. Data buffers that are shared with Embree
(e.g.~for vertex of index buffers) have to get allocated as SYCL unified
shared memory (USM memory), using the \texttt{sycl::malloc} or
\texttt{sycl::aligned\_alloc} calls with
\texttt{sycl::usm::alloc::shared} property, or the
sycl::aligned\_alloc\_shared call, e.g:

\begin{Shaded}
\begin{Highlighting}[]
\DataTypeTok{void}\OperatorTok{*}\NormalTok{ ptr }\OperatorTok{=}\NormalTok{ sycl}\OperatorTok{::}\NormalTok{aligned\_alloc}\OperatorTok{(}\DecValTok{16}\OperatorTok{,}\NormalTok{ bytes}\OperatorTok{,}\NormalTok{ queue}\OperatorTok{,}\NormalTok{ sycl}\OperatorTok{::}\NormalTok{usm}\OperatorTok{::}\NormalTok{alloc}\OperatorTok{::}\NormalTok{shared}\OperatorTok{);}
\end{Highlighting}
\end{Shaded}

These shared allocations have to be valid during rendering, as Embree
may access contained data when tracing rays. Embree does not support
device-only memory allocations, as the BVH builder implemented on the
CPU relies on reading the data buffers.

Device side rendering can get invoked by submitting a SYCL
\texttt{parallel\_for} to the SYCL queue:

\begin{Shaded}
\begin{Highlighting}[]
\AttributeTok{const}\NormalTok{ sycl}\OperatorTok{::}\NormalTok{specialization\_id}\OperatorTok{\textless{}}\NormalTok{RTCFeatureFlags}\OperatorTok{\textgreater{}}\NormalTok{ feature\_mask}\OperatorTok{;}

\NormalTok{RTCFeatureFlags required\_features }\OperatorTok{=}\NormalTok{ RTC\_FEATURE\_FLAG\_TRIANGLE}\OperatorTok{;}

\NormalTok{queue}\OperatorTok{.}\NormalTok{submit}\OperatorTok{([=](}\NormalTok{sycl}\OperatorTok{::}\NormalTok{handler}\OperatorTok{\&}\NormalTok{ cgh}\OperatorTok{)}
\OperatorTok{\{}
\NormalTok{  cgh}\OperatorTok{.}\NormalTok{set\_specialization\_constant}\OperatorTok{\textless{}}\NormalTok{feature\_mask}\OperatorTok{\textgreater{}(}\NormalTok{required\_features}\OperatorTok{);}
  
\NormalTok{  cgh}\OperatorTok{.}\NormalTok{parallel\_for}\OperatorTok{(}\NormalTok{sycl}\OperatorTok{::}\NormalTok{range}\OperatorTok{\textless{}}\DecValTok{1}\OperatorTok{\textgreater{}(}\DecValTok{1}\OperatorTok{),[=](}\NormalTok{sycl}\OperatorTok{::}\NormalTok{id}\OperatorTok{\textless{}}\DecValTok{1}\OperatorTok{\textgreater{}}\NormalTok{ item}\OperatorTok{,}\NormalTok{ sycl}\OperatorTok{::}\NormalTok{kernel\_handler kh}\OperatorTok{)}
  \OperatorTok{\{}
\NormalTok{    RTCIntersectArguments args}\OperatorTok{;}
\NormalTok{    rtcInitIntersectArguments}\OperatorTok{(\&}\NormalTok{args}\OperatorTok{);}

    \AttributeTok{const}\NormalTok{ RTCFeatureFlags features }\OperatorTok{=}\NormalTok{ kh}\OperatorTok{.}\NormalTok{get\_specialization\_constant}\OperatorTok{\textless{}}\NormalTok{feature\_mask}\OperatorTok{\textgreater{}();}
\NormalTok{    args}\OperatorTok{.}\NormalTok{feature\_mask }\OperatorTok{=}\NormalTok{ features}\OperatorTok{;}

    \KeywordTok{struct}\NormalTok{ RTCRayHit rayhit}\OperatorTok{;}
\NormalTok{    rayhit}\OperatorTok{.}\NormalTok{ray}\OperatorTok{.}\NormalTok{org\_x }\OperatorTok{=}\NormalTok{ ox}\OperatorTok{;}
\NormalTok{    rayhit}\OperatorTok{.}\NormalTok{ray}\OperatorTok{.}\NormalTok{org\_y }\OperatorTok{=}\NormalTok{ oy}\OperatorTok{;}
\NormalTok{    rayhit}\OperatorTok{.}\NormalTok{ray}\OperatorTok{.}\NormalTok{org\_z }\OperatorTok{=}\NormalTok{ oz}\OperatorTok{;}
\NormalTok{    rayhit}\OperatorTok{.}\NormalTok{ray}\OperatorTok{.}\NormalTok{dir\_x }\OperatorTok{=}\NormalTok{ dx}\OperatorTok{;}
\NormalTok{    rayhit}\OperatorTok{.}\NormalTok{ray}\OperatorTok{.}\NormalTok{dir\_y }\OperatorTok{=}\NormalTok{ dy}\OperatorTok{;}
\NormalTok{    rayhit}\OperatorTok{.}\NormalTok{ray}\OperatorTok{.}\NormalTok{dir\_z }\OperatorTok{=}\NormalTok{ dz}\OperatorTok{;}
\NormalTok{    rayhit}\OperatorTok{.}\NormalTok{ray}\OperatorTok{.}\NormalTok{tnear }\OperatorTok{=} \DecValTok{0}\OperatorTok{;}
\NormalTok{    rayhit}\OperatorTok{.}\NormalTok{ray}\OperatorTok{.}\NormalTok{tfar }\OperatorTok{=} \BuiltInTok{std::}\NormalTok{numeric\_limits}\OperatorTok{\textless{}}\DataTypeTok{float}\OperatorTok{\textgreater{}::}\NormalTok{infinity}\OperatorTok{();}
\NormalTok{    rayhit}\OperatorTok{.}\NormalTok{ray}\OperatorTok{.}\NormalTok{mask }\OperatorTok{=} \OperatorTok{{-}}\DecValTok{1}\OperatorTok{;}
\NormalTok{    rayhit}\OperatorTok{.}\NormalTok{ray}\OperatorTok{.}\NormalTok{flags }\OperatorTok{=} \DecValTok{0}\OperatorTok{;}
\NormalTok{    rayhit}\OperatorTok{.}\NormalTok{hit}\OperatorTok{.}\NormalTok{geomID }\OperatorTok{=}\NormalTok{ RTC\_INVALID\_GEOMETRY\_ID}\OperatorTok{;}
\NormalTok{    rayhit}\OperatorTok{.}\NormalTok{hit}\OperatorTok{.}\NormalTok{instID}\OperatorTok{[}\DecValTok{0}\OperatorTok{]} \OperatorTok{=}\NormalTok{ RTC\_INVALID\_GEOMETRY\_ID}\OperatorTok{;}

\NormalTok{    rtcIntersect1}\OperatorTok{(}\NormalTok{scene}\OperatorTok{,} \OperatorTok{\&}\NormalTok{rayhit}\OperatorTok{,} \OperatorTok{\&}\NormalTok{args}\OperatorTok{);}

\NormalTok{    result}\OperatorTok{{-}\textgreater{}}\NormalTok{geomID }\OperatorTok{=}\NormalTok{ rayhit}\OperatorTok{.}\NormalTok{hit}\OperatorTok{.}\NormalTok{geomID}\OperatorTok{;}
\NormalTok{    result}\OperatorTok{{-}\textgreater{}}\NormalTok{primID }\OperatorTok{=}\NormalTok{ rayhit}\OperatorTok{.}\NormalTok{hit}\OperatorTok{.}\NormalTok{primID}\OperatorTok{;}
\NormalTok{    result}\OperatorTok{{-}\textgreater{}}\NormalTok{tfar }\OperatorTok{=}\NormalTok{ rayhit}\OperatorTok{.}\NormalTok{ray}\OperatorTok{.}\NormalTok{tfar}\OperatorTok{;}
  \OperatorTok{\});}
\OperatorTok{\});}
\NormalTok{queue}\OperatorTok{.}\NormalTok{wait\_and\_throw}\OperatorTok{();}
\end{Highlighting}
\end{Shaded}

This example passes a feature mask using a specialization contant to the
\texttt{rtcIntersect1} function, which is recommended for GPU rendering.
For best performance, this feature mask should get used to enable only
features required by the application to render the scene, e.g.~just
triangles in this example.

Inside the SYCL \texttt{parallel\_for} loop you can use rendering
related functions, such as the \texttt{rtcIntersect1} and
\texttt{rtcOccluded1} functions to trace rays,
\texttt{rtcForwardIntersect1} and \texttt{rtcForwardOccluded1} to
continue object traversal from inside a user geometry callback, and
\texttt{rtcGetGeometryUserDataFromScene} to get the user data pointer of
some geometry.

Have a look at the \hyperref[minimal]{Minimal} tutorial for a minimal
SYCL example.

\section{SYCL JIT caching}\label{sycl-jit-caching}

Compile times for just in time compilation (JIT compilation) can be
large. To resolve this issue we recommend enabling persistent JIT
compilation caching inside your application, by setting the
\texttt{SYCL\_CACHE\_PERSISTENT} environment variable to \texttt{1}, and
the \texttt{SYCL\_CACHE\_DIR} environment variable to some proper
directory where the JIT cache should get stored. These environment
variables have to get set before the SYCL device is created, e.g:

\begin{Shaded}
\begin{Highlighting}[]
\NormalTok{setenv}\OperatorTok{(}\StringTok{"SYCL\_CACHE\_PERSISTENT"}\OperatorTok{,}\StringTok{"1"}\OperatorTok{,}\DecValTok{1}\OperatorTok{);}
\NormalTok{setenv}\OperatorTok{(}\StringTok{"SYCL\_CACHE\_DIR"}\OperatorTok{,}\StringTok{"cache\_dir"}\OperatorTok{,}\DecValTok{1}\OperatorTok{);}

\NormalTok{sycl}\OperatorTok{::}\NormalTok{device device}\OperatorTok{(}\NormalTok{rtcSYCLDeviceSelector}\OperatorTok{);}
\OperatorTok{...}
\end{Highlighting}
\end{Shaded}

\section{SYCL Memory Pooling}\label{sycl-memory-pooling}

Memory Pooling is a mechanism where small USM memory allocations are
packed into larger allocation blocks. This mode is required when your
application performs many small USM allocations, as otherwise only a
small fraction of GPU memory is usable and data transfer performance
will be low.

Memory pooling is supported for USM allocations that are read-only by
the device. The following example allocated device read-only memory with
memory pooling support:

\begin{Shaded}
\begin{Highlighting}[]
\NormalTok{sycl}\OperatorTok{::}\NormalTok{aligned\_alloc\_shared}\OperatorTok{(}\NormalTok{align}\OperatorTok{,}\NormalTok{ bytes}\OperatorTok{,}\NormalTok{ queue}\OperatorTok{,}
\NormalTok{  sycl}\OperatorTok{::}\NormalTok{ext}\OperatorTok{::}\NormalTok{oneapi}\OperatorTok{::}\NormalTok{property}\OperatorTok{::}\NormalTok{usm}\OperatorTok{::}\NormalTok{device\_read\_only}\OperatorTok{());}
\end{Highlighting}
\end{Shaded}

\section{Embree SYCL Limitations}\label{embree-sycl-limitations}

Embree only supports Xe HPC and HPG GPUs as SYCL devices, thus in
particular the CPU and other GPUs cannot get used as a SYCL device. To
render on the CPU just use the standard C99 API without relying on SYCL.

The SYCL language spec puts some restrictions to device functions, such
as disallowing: global variable access, malloc, invokation of virtual
functions, function pointers, runtime type information, exceptions,
recursion, etc. See Section
\texttt{5.4.\ Language\ Restrictions\ for\ device\ functions} of the
\href{https://www.khronos.org/registry/SYCL/specs/sycl-2020/html/sycl-2020.html\#sec:language.restrictions.kernels}{SYCL
specification} for more details.

Using Intel's oneAPI DPC++ compiler invoking an indirectly called
function is allowed, but we do not recommend this for performance
reasons.

Some features are not supported by the Embree SYCL API thus cannot get
used on the GPU:

\begin{itemize}
\item
  The packet tracing functions \texttt{rtcIntersect4/8/16} and
  \texttt{rtcOccluded4/8/16} are not supported in SYCL device side code.
  Using these functions makes no sense for SYCL, as the programming
  model is implicitely executed in SIMT mode on the GPU anyway.
\item
  Filter and user geometry callbacks stored inside the geometry objects
  are not supported on SYCL. Please use the alternative approach of
  passing the function pointer through the
  \texttt{RTCIntersectArguments} (or \texttt{RTCOccludedArguments})
  structures to the tracing function, which enables inlining on the GPU.
\item
  The \texttt{rtcInterpolate} function cannot get used on the the
  device. For most primitive types the vertex data interpolation is
  anyway a trivial operation, and an API call just introduces overheads.
  On the CPU that overhead is acceptable, but on the GPU it is not. The
  \texttt{rtcInterpolate} function does not know the geometry type it is
  interpolating over, thus its implementation on the GPU would contain a
  large switch statement for all potential geometry types.
\item
  Tracing rays using \texttt{rtcIntersect1} and \texttt{rtcOccluded1}
  functions from user geometry callbacks is not supported in SYCL.
  Please use the tail recursive \texttt{rtcForwardIntersect1} and
  \texttt{rtcForwardOccluded1} calls instead.
\item
  Subdivision surfaces are not supported for Embree SYCL devices.
\item
  Collision detection (\texttt{rtcCollide} API call) is not supported in
  SYCL device side code.
\item
  Point queries (\texttt{rtcPointQuery} API call) are not supported in
  SYCL device side code.
\end{itemize}

\section{Embree SYCL Known Issues}\label{embree-sycl-known-issues}

\begin{itemize}
\item
  The SYCL support of Embree is in beta phase. Current functionality,
  quality, and GPU performance may not reflect that of the final
  product.
\item
  Compilation with build configuration ``debug'' is currently not
  working on Windows.
\end{itemize}

\chapter{Upgrading from Embree 3 to Embree
4}\label{upgrading-from-embree-3-to-embree-4}

This section summarizes API changes between Embree 3 and Embree4. Most
of these changes are motivated by GPU performance and having a
consistent API that works properly for the CPU and GPU.

\begin{itemize}
\item
  The API include folder got renamed from embree3 to embree4, to be able
  to install Embree 3 and Embree 4 side by side, without having
  conflicts in API folder.
\item
  The \texttt{RTCIntersectContext} is renamed to
  \texttt{RTCRayQueryContext} and the \texttt{RTCIntersectContextFlags}
  got renamed to \texttt{RTCRayQueryFlags}.
\item
  There are some changes to the \texttt{rtcIntersect} and
  \texttt{rtcOccluded} functions. Most members of the old intersect
  context have been moved to some optional
  \texttt{RTCIntersectArguments} (and \texttt{RTCOccludedArguments})
  structures, which also contains a pointer to the new ray query
  context. The argument structs fulfill the task of providing additional
  advanced arguments to the traversal functions. The ray query context
  can get used to pass additional data to callbacks, and to maintain an
  instID stack in case instancing is done manually inside user geometry
  callbacks. The arguments struct is not available inside callbacks.
  This change was in particular necessary for SYCL to allow inlining of
  function pointers provided to the traversal functions, and to reduce
  the amount of state passed to callbacks, which both improves GPU
  performance. Most applications can just drop passing the ray query
  context to port to Embree 4.
\item
  The \texttt{rtcFilterIntersection} and \texttt{rtcFilterOcclusion} API
  calls that invoke both, the geometry and argument version of the
  filter callback, from a user geometry callback are no longer
  supported. Instead applications should use the
  \texttt{rtcInvokeIntersectFilterFromGeometry} and
  \texttt{rtcInvokeOccludedFilterFromGeometry} API calls that invoke
  just the geometry version of the filter function, and invoke the
  argument filter function manually if required.
\item
  The filter function passed as arguments to \texttt{rtcIntersect} and
  \texttt{rtcOccluded} functions is only invoked for some geometry if
  enabled through
  \texttt{rtcSetGeometryEnableFilterFunctionFromArguments} for that
  geometry. Alternatively, argument filter functions can get enabled for
  all geometries using the
  \texttt{RTC\_RAY\_QUERY\_FLAG\_INVOKE\_ARGUMENT\_FILTER} ray query
  flag.
\item
  User geometry callbacks get a valid vector as input to identify valid
  and invalid rays. In Embree 3 the user geometry callback just had to
  update the ray hit members when an intersection was found and perform
  no operation otherwise. In Embree 4 the callback additionally has to
  return valid=-1 when a hit was found, and valid=0 when no hit was
  found. This allows Embree to properly pass the new hit distance to the
  ray tracing hardware only in the case a hit was found.
\item
  Further ray masking is enabled by default now as required by most
  applications and the default ray mask for geometries got changed from
  0xFFFFFFFF to 0x1.
\item
  The stream tracing functions \texttt{rtcIntersect1M},
  \texttt{rtcIntersect1Mp}, \texttt{rtcIntersectNM},
  \texttt{rtcIntersectNp}, \texttt{rtcOccluded1M},
  \texttt{rtcOccluded1Mp}, \texttt{rtcOccludedNM}, and
  \texttt{rtcOccludedNp} got removed as they were rarely used and did
  not provide relevant performance benefits. As alternative the
  application can just iterate over \texttt{rtcIntersect1} and
  potentially \texttt{rtcIntersect4/8/16} to get similar performance.
\end{itemize}

To use Embree through SYCL on the CPU and GPU additional changes are
required:

\begin{itemize}
\item
  Embree 3 allows to use \texttt{rtcIntersect} recursively from a user
  geometry or intersection filter callback to continue a ray inside an
  instantiated object. In Embree 4 using \texttt{rtcIntersect}
  recursively is disallowed on the GPU but still supported on the CPU.
  To properly continue a ray inside an instantiated object use the new
  \texttt{rtcForwardIntersect1} and \texttt{rtcForwardOccluded1}
  functions.
\item
  The geometry object of Embree 4 is a host side only object, thus
  accessing it during rendering from the GPU is not allowed. Thus all
  API functions that take an RTCGeometry object as argument cannot get
  used during rendering. Thus in particular the
  \texttt{rtcGetGeometryUserData(RTCGeometry)} call cannot get used, but
  there is an alternative function
  \texttt{rtcGetGeometryUserDataFromScene(RTCScene\ scene,uint\ geomID)}
  that should get used instead.
\item
  The user geometry callback and filter callback functions should get
  passed through the intersection and occlusion argument structures to
  the \texttt{rtcIntersect1} and \texttt{rtcOccluded1} functions
  directly to allow inlining. The experimental geometry version of the
  callbacks is disabled in SYCL and should not get used.
\item
  The feature flags should get used in SYCL to minimal GPU code for
  optimal performance.
\item
  The \texttt{rtcInterpolate} function cannot get used on the device,
  and vertex data interpolation should get implemented by the
  application.
\item
  Indirectly called functions must be declared with
  \texttt{RTC\_SYCL\_INDIRECTLY\_CALLABLE} when used as filter or user
  geometry callbacks.
\end{itemize}

\pagebreak

\chapter{Embree API Reference}\label{embree-api-reference}

\section{rtcNewDevice}\label{rtcnewdevice}

\subsubsection{NAME}\label{name}

\begin{Shaded}
\begin{Highlighting}[]
\NormalTok{rtcNewDevice }\OperatorTok{{-}}\NormalTok{ creates a }\KeywordTok{new}\NormalTok{ device}
\end{Highlighting}
\end{Shaded}

\subsubsection{SYNOPSIS}\label{synopsis}

\begin{Shaded}
\begin{Highlighting}[]
\PreprocessorTok{\#include }\ImportTok{\textless{}embree4/rtcore.h\textgreater{}}

\NormalTok{RTCDevice rtcNewDevice}\OperatorTok{(}\AttributeTok{const} \DataTypeTok{char}\OperatorTok{*}\NormalTok{ config}\OperatorTok{);}
\end{Highlighting}
\end{Shaded}

\subsubsection{DESCRIPTION}\label{description}

This function creates a new device to be used for CPU ray tracing and
returns a handle to this device. The device object is reference counted
with an initial reference count of 1. The handle can be released using
the \texttt{rtcReleaseDevice} API call.

The device object acts as a class factory for all other object types.
All objects created from the device (like scenes, geometries, etc.) hold
a reference to the device, thus the device will not be destroyed unless
these objects are destroyed first.

Objects are only compatible if they belong to the same device, e.g it is
not allowed to create a geometry in one device and attach it to a scene
created with a different device.

A configuration string (\texttt{config} argument) can be passed to the
device construction. This configuration string can be \texttt{NULL} to
use the default configuration.

The following configuration is supported:

\begin{itemize}
\item
  \texttt{threads={[}int{]}}: Specifies a number of build threads to
  use. A value of 0 enables all detected hardware threads. By default
  all hardware threads are used.
\item
  \texttt{user\_threads={[}int{]}}: Sets the number of user threads that
  can be used to join and participate in a scene commit using
  \texttt{rtcJoinCommitScene}. The tasking system will only use
  threads-user\_threads many worker threads, thus if the app wants to
  solely use its threads to commit scenes, just set threads equal to
  user\_threads. This option only has effect with the Intel(R) Threading
  Building Blocks (TBB) tasking system.
\item
  \texttt{set\_affinity={[}0/1{]}}: When enabled, build threads are
  affinitized to hardware threads. This option is disabled by default on
  standard CPUs, and enabled by default on Xeon Phi Processors.
\item
  \texttt{start\_threads={[}0/1{]}}: When enabled, the build threads are
  started upfront. This can be useful for benchmarking to exclude thread
  creation time. This option is disabled by default.
\item
  \texttt{isa={[}sse2,sse4.2,avx,avx2,avx512{]}}: Use specified ISA. By
  default the ISA is selected automatically.
\item
  \texttt{max\_isa={[}sse2,sse4.2,avx,avx2,avx512{]}}: Configures the
  automated ISA selection to use maximally the specified ISA.
\item
  \texttt{hugepages={[}0/1{]}}: Enables or disables usage of huge pages.
  Under Linux huge pages are used by default but under Windows and macOS
  they are disabled by default.
\item
  \texttt{enable\_selockmemoryprivilege={[}0/1{]}}: When set to 1, this
  enables the \texttt{SeLockMemoryPrivilege} privilege with is required
  to use huge pages on Windows. This option has an effect only under
  Windows and is ignored on other platforms. See Section
  \hyperref[huge-page-support]{Huge Page Support} for more details.
\item
  \texttt{verbose={[}0,1,2,3{]}}: Sets the verbosity of the output. When
  set to 0, no output is printed by Embree, when set to a higher level
  more output is printed. By default Embree does not print anything on
  the console.
\item
  \texttt{frequency\_level={[}simd128,simd256,simd512{]}}: Specifies the
  frequency level the application want to run on, which can be either:

  \begin{enumerate}
  \def\labelenumi{\alph{enumi})}
  \tightlist
  \item
    simd128 to run at highest frequency
  \item
    simd256 to run at AVX2-heavy frequency level
  \item
    simd512 to run at heavy AVX512 frequency level. When some frequency
    level is specified, Embree will avoid doing optimizations that may
    reduce the frequency level below the level specified. E.g. if your
    app does not use AVX instructions setting
    ``frequency\_level=simd128'' will cause some CPUs to run at highest
    frequency, which may result in higher application performance if you
    do much shading. If you application heavily uses AVX code, you
    should best set the frequency level to simd256. Per default Embree
    tries to avoid reducing the frequency of the CPU by setting the
    simd256 level only when the CPU has no significant down clocking.
  \end{enumerate}
\end{itemize}

Different configuration options should be separated by commas, e.g.:

\begin{Shaded}
\begin{Highlighting}[]
\NormalTok{rtcNewDevice}\OperatorTok{(}\StringTok{"threads=1,isa=avx"}\OperatorTok{);}
\end{Highlighting}
\end{Shaded}

\subsubsection{EXIT STATUS}\label{exit-status}

On success returns a handle of the created device. On failure returns
\texttt{NULL} as device and sets a per-thread error code that can be
queried using \texttt{rtcGetDeviceError(NULL)}.

\subsubsection{SEE ALSO}\label{see-also}

\hyperref[rtcretaindevice]{rtcRetainDevice},
\hyperref[rtcreleasedevice]{rtcReleaseDevice},
\hyperref[rtcnewsycldevice]{rtcNewSYCLDevice}

\pagebreak

\section{rtcNewSYCLDevice}\label{rtcnewsycldevice}

\subsubsection{NAME}\label{name}

\begin{Shaded}
\begin{Highlighting}[]
\NormalTok{rtcNewSYCLDevice }\OperatorTok{{-}}\NormalTok{ creates a }\KeywordTok{new}\NormalTok{ device to be used with SYCL}
\end{Highlighting}
\end{Shaded}

\subsubsection{SYNOPSIS}\label{synopsis}

\begin{Shaded}
\begin{Highlighting}[]
\PreprocessorTok{\#include }\ImportTok{\textless{}embree4/rtcore.h\textgreater{}}

\NormalTok{RTCDevice rtcNewSYCLDevice}\OperatorTok{(}\NormalTok{sycl}\OperatorTok{::}\NormalTok{context context}\OperatorTok{,} \AttributeTok{const} \DataTypeTok{char}\OperatorTok{*}\NormalTok{ config}\OperatorTok{);}
\end{Highlighting}
\end{Shaded}

\subsubsection{DESCRIPTION}\label{description}

This function creates a new device to be used with SYCL for GPU
rendering and returns a handle to this device. The device object is
reference counted with an initial reference count of 1. The handle can
get released using the \texttt{rtcReleaseDevice} API call.

The passed SYCL context (\texttt{context} argument) is used to allocate
GPU data, thus only devices contained inside this context can be used
for rendering. By default the GPU data is allocated on the first GPU
device of the context, but this behavior can get changed with the
\hyperref[rtcsetdevicesycldevice]{rtcSetDeviceSYCLDevice} function.

The device object acts as a class factory for all other object types.
All objects created from the device (like scenes, geometries, etc.) hold
a reference to the device, thus the device will not be destroyed unless
these objects are destroyed first.

Objects are only compatible if they belong to the same device, e.g it is
not allowed to create a geometry in one device and attach it to a scene
created with a different device.

For an overview of configurations that can get passed (\texttt{config}
argument) please see the \hyperref[rtcnewdevice]{rtcNewDevice} function
description.

\subsubsection{EXIT STATUS}\label{exit-status}

On success returns a handle of the created device. On failure returns
\texttt{NULL} as device and sets a per-thread error code that can be
queried using \texttt{rtcGetDeviceError(NULL)}.

\subsubsection{SEE ALSO}\label{see-also}

\hyperref[rtcretaindevice]{rtcRetainDevice},
\hyperref[rtcreleasedevice]{rtcReleaseDevice},
\hyperref[rtcnewdevice]{rtcNewDevice}

\pagebreak

\section{rtcIsSYCLDeviceSupported}\label{rtcissycldevicesupported}

\subsubsection{NAME}\label{name}

\begin{Shaded}
\begin{Highlighting}[]
\NormalTok{rtcIsSYCLDeviceSupported }\OperatorTok{{-}}\NormalTok{ checks }\ControlFlowTok{if}\NormalTok{ some SYCL device is supported by Embree}
\end{Highlighting}
\end{Shaded}

\subsubsection{SYNOPSIS}\label{synopsis}

\begin{Shaded}
\begin{Highlighting}[]
\PreprocessorTok{\#include }\ImportTok{\textless{}embree4/rtcore.h\textgreater{}}

\DataTypeTok{bool}\NormalTok{ rtcIsSYCLDeviceSupported}\OperatorTok{(}\AttributeTok{const}\NormalTok{ sycl}\OperatorTok{::}\NormalTok{device sycl\_device}\OperatorTok{);}
\end{Highlighting}
\end{Shaded}

\subsubsection{DESCRIPTION}\label{description}

This function can be used to check if some SYCL device
(\texttt{sycl\_device} argument) is supported by Embree.

\subsubsection{EXIT STATUS}\label{exit-status}

The function returns true if the SYCL device is supported by Embree and
false otherwise. On failure an error code is set that can get queried
using \texttt{rtcGetDeviceError}.

\subsubsection{SEE ALSO}\label{see-also}

\hyperref[rtcsycldeviceselector]{rtcSYCLDeviceSelector}

\pagebreak

\section{rtcSYCLDeviceSelector}\label{rtcsycldeviceselector}

\subsubsection{NAME}\label{name}

\begin{Shaded}
\begin{Highlighting}[]
\NormalTok{rtcSYCLDeviceSelector }\OperatorTok{{-}}\NormalTok{ SYCL device selector function to select}
\NormalTok{  devices supported by Embree}
\end{Highlighting}
\end{Shaded}

\subsubsection{SYNOPSIS}\label{synopsis}

\begin{Shaded}
\begin{Highlighting}[]
\PreprocessorTok{\#include }\ImportTok{\textless{}embree4/rtcore.h\textgreater{}}

\DataTypeTok{int}\NormalTok{ rtcSYCLDeviceSelector}\OperatorTok{(}\AttributeTok{const}\NormalTok{ sycl}\OperatorTok{::}\NormalTok{device sycl\_device}\OperatorTok{);}
\end{Highlighting}
\end{Shaded}

\subsubsection{DESCRIPTION}\label{description}

This function checks if the passed SYCL device (\texttt{sycl\_device}
arguments) is supported by Embree or not. This function can be used
directly to select some supported SYCL device by using it as SYCL device
selector function. For instance, the following code sequence selects an
Embree supported SYCL device and creates an Embree device from it:

\begin{Shaded}
\begin{Highlighting}[]
\NormalTok{sycl}\OperatorTok{::}\NormalTok{device sycl\_device}\OperatorTok{(}\NormalTok{rtcSYCLDeviceSelector}\OperatorTok{);}
\NormalTok{sycl}\OperatorTok{::}\NormalTok{queue sycl\_queue}\OperatorTok{(}\NormalTok{sycl\_device}\OperatorTok{);}
\NormalTok{sycl}\OperatorTok{::}\NormalTok{context}\OperatorTok{(}\NormalTok{sycl\_device}\OperatorTok{);}
\NormalTok{RTCDevice device }\OperatorTok{=}\NormalTok{ rtcNewSYCLDevice}\OperatorTok{(}\NormalTok{sycl\_context}\OperatorTok{,}\KeywordTok{nullptr}\OperatorTok{);}
\end{Highlighting}
\end{Shaded}

\subsubsection{EXIT STATUS}\label{exit-status}

The function returns -1 if the SYCL device is supported by Embree and 1
otherwise. On failure an error code is set that can get queried using
\texttt{rtcGetDeviceError}.

\subsubsection{SEE ALSO}\label{see-also}

\hyperref[rtcissycldevicesupported]{rtcIsSYCLDeviceSupported}

\pagebreak

\section{rtcSetDeviceSYCLDevice}\label{rtcsetdevicesycldevice}

\subsubsection{NAME}\label{name}

\begin{Shaded}
\begin{Highlighting}[]
\NormalTok{rtcSetDeviceSYCLDevice }\OperatorTok{{-}}\NormalTok{ sets the SYCL device to be used }\ControlFlowTok{for}\NormalTok{ memory allocations}
\end{Highlighting}
\end{Shaded}

\subsubsection{SYNOPSIS}\label{synopsis}

\begin{Shaded}
\begin{Highlighting}[]
\PreprocessorTok{\#include }\ImportTok{\textless{}embree4/rtcore.h\textgreater{}}

\DataTypeTok{void}\NormalTok{ rtcSetDeviceSYCLDevice}\OperatorTok{(}\NormalTok{RTCDevice device}\OperatorTok{,} \AttributeTok{const}\NormalTok{ sycl}\OperatorTok{::}\NormalTok{device sycl\_device}\OperatorTok{);}
\end{Highlighting}
\end{Shaded}

\subsubsection{DESCRIPTION}\label{description}

This function sets the SYCL device (\texttt{sycl\_device} argument) to
be used to allocate GPU memory when using the specified Embree device
(\texttt{device} argument). This SYCL device must be one of the SYCL
devices contained inside the SYCL context used to create the Embree
device.

\subsubsection{EXIT STATUS}\label{exit-status}

On failure an error code is set that can get queried using
\texttt{rtcGetDeviceError}.

\subsubsection{SEE ALSO}\label{see-also}

\hyperref[rtcnewsycldevice]{rtcNewSYCLDevice}

\pagebreak

\section{rtcRetainDevice}\label{rtcretaindevice}

\subsubsection{NAME}\label{name}

\begin{Shaded}
\begin{Highlighting}[]
\NormalTok{rtcRetainDevice }\OperatorTok{{-}}\NormalTok{ increments the device reference count}
\end{Highlighting}
\end{Shaded}

\subsubsection{SYNOPSIS}\label{synopsis}

\begin{Shaded}
\begin{Highlighting}[]
\PreprocessorTok{\#include }\ImportTok{\textless{}embree4/rtcore.h\textgreater{}}

\DataTypeTok{void}\NormalTok{ rtcRetainDevice}\OperatorTok{(}\NormalTok{RTCDevice device}\OperatorTok{);}
\end{Highlighting}
\end{Shaded}

\subsubsection{DESCRIPTION}\label{description}

Device objects are reference counted. The \texttt{rtcRetainDevice}
function increments the reference count of the passed device object
(\texttt{device} argument). This function together with
\texttt{rtcReleaseDevice} allows to use the internal reference counting
in a C++ wrapper class to manage the ownership of the object.

\subsubsection{EXIT STATUS}\label{exit-status}

On failure an error code is set that can be queried using
\texttt{rtcGetDeviceError}.

\subsubsection{SEE ALSO}\label{see-also}

\hyperref[rtcnewdevice]{rtcNewDevice},
\hyperref[rtcreleasedevice]{rtcReleaseDevice}

\pagebreak

\section{rtcReleaseDevice}\label{rtcreleasedevice}

\subsubsection{NAME}\label{name}

\begin{Shaded}
\begin{Highlighting}[]
\NormalTok{rtcReleaseDevice }\OperatorTok{{-}}\NormalTok{ decrements the device reference count}
\end{Highlighting}
\end{Shaded}

\subsubsection{SYNOPSIS}\label{synopsis}

\begin{Shaded}
\begin{Highlighting}[]
\PreprocessorTok{\#include }\ImportTok{\textless{}embree4/rtcore.h\textgreater{}}

\DataTypeTok{void}\NormalTok{ rtcReleaseDevice}\OperatorTok{(}\NormalTok{RTCDevice device}\OperatorTok{);}
\end{Highlighting}
\end{Shaded}

\subsubsection{DESCRIPTION}\label{description}

Device objects are reference counted. The \texttt{rtcReleaseDevice}
function decrements the reference count of the passed device object
(\texttt{device} argument). When the reference count falls to 0, the
device gets destroyed.

All objects created from the device (like scenes, geometries, etc.) hold
a reference to the device, thus the device will not get destroyed unless
these objects are destroyed first.

\subsubsection{EXIT STATUS}\label{exit-status}

On failure an error code is set that can be queried using
\texttt{rtcGetDeviceError}.

\subsubsection{SEE ALSO}\label{see-also}

\hyperref[rtcnewdevice]{rtcNewDevice},
\hyperref[rtcretaindevice]{rtcRetainDevice}

\pagebreak

\section{rtcGetDeviceProperty}\label{rtcgetdeviceproperty}

\subsubsection{NAME}\label{name}

\begin{Shaded}
\begin{Highlighting}[]
\NormalTok{rtcGetDeviceProperty }\OperatorTok{{-}}\NormalTok{ queries properties of the device}
\end{Highlighting}
\end{Shaded}

\subsubsection{SYNOPSIS}\label{synopsis}

\begin{Shaded}
\begin{Highlighting}[]
\PreprocessorTok{\#include }\ImportTok{\textless{}embree4/rtcore.h\textgreater{}}

\DataTypeTok{ssize\_t}\NormalTok{ rtcGetDeviceProperty}\OperatorTok{(}
\NormalTok{  RTCDevice device}\OperatorTok{,}
  \KeywordTok{enum}\NormalTok{ RTCDeviceProperty prop}
\OperatorTok{);}
\end{Highlighting}
\end{Shaded}

\subsubsection{DESCRIPTION}\label{description}

The \texttt{rtcGetDeviceProperty} function can be used to query
properties (\texttt{prop} argument) of a device object (\texttt{device}
argument). The returned property is an integer of type
\texttt{ssize\_t}.

Possible properties to query are:

\begin{itemize}
\item
  \texttt{RTC\_DEVICE\_PROPERTY\_VERSION}: Queries the combined version
  number (MAJOR.MINOR.PATCH) with two decimal digits per component. E.g.
  for Embree 2.8.3 the integer 208003 is returned.
\item
  \texttt{RTC\_DEVICE\_PROPERTY\_VERSION\_MAJOR}: Queries the major
  version number of Embree.
\item
  \texttt{RTC\_DEVICE\_PROPERTY\_VERSION\_MINOR}: Queries the minor
  version number of Embree.
\item
  \texttt{RTC\_DEVICE\_PROPERTY\_VERSION\_PATCH}: Queries the patch
  version number of Embree.
\item
  \texttt{RTC\_DEVICE\_PROPERTY\_NATIVE\_RAY4\_SUPPORTED}: Queries
  whether the \texttt{rtcIntersect4} and \texttt{rtcOccluded4} functions
  preserve packet size and ray order when invoking callback functions.
  This is only the case if Embree is compiled with
  \texttt{EMBREE\_RAY\_PACKETS} and \texttt{SSE2} (or \texttt{SSE4.2})
  enabled, and if the machine it is running on supports \texttt{SSE2}
  (or \texttt{SSE4.2}).
\item
  \texttt{RTC\_DEVICE\_PROPERTY\_NATIVE\_RAY8\_SUPPORTED}: Queries
  whether the \texttt{rtcIntersect8} and \texttt{rtcOccluded8} functions
  preserve packet size and ray order when invoking callback functions.
  This is only the case if Embree is compiled with
  \texttt{EMBREE\_RAY\_PACKETS} and \texttt{AVX} (or \texttt{AVX2})
  enabled, and if the machine it is running on supports \texttt{AVX} (or
  \texttt{AVX2}).
\item
  \texttt{RTC\_DEVICE\_PROPERTY\_NATIVE\_RAY16\_SUPPORTED}: Queries
  whether the \texttt{rtcIntersect16} and \texttt{rtcOccluded16}
  functions preserve packet size and ray order when invoking callback
  functions. This is only the case if Embree is compiled with
  \texttt{EMBREE\_RAY\_PACKETS} and \texttt{AVX512} enabled, and if the
  machine it is running on supports \texttt{AVX512}.
\item
  \texttt{RTC\_DEVICE\_PROPERTY\_RAY\_MASK\_SUPPORTED}: Queries whether
  ray masks are supported. This is only the case if Embree is compiled
  with \texttt{EMBREE\_RAY\_MASK} enabled.
\item
  \texttt{RTC\_DEVICE\_PROPERTY\_BACKFACE\_CULLING\_ENABLED}: Queries
  whether back face culling is enabled. This is only the case if Embree
  is compiled with \texttt{EMBREE\_BACKFACE\_CULLING} enabled.
\item
  \texttt{RTC\_DEVICE\_PROPERTY\_BACKFACE\_CULLING\_CURVES\_ENABLED}:
  Queries whether back face culling for curves is enabled. This is only
  the case if Embree is compiled with
  \texttt{EMBREE\_BACKFACE\_CULLING\_CURVES} enabled.
\item
  \texttt{RTC\_DEVICE\_PROPERTY\_BACKFACE\_CULLING\_SPHERES\_ENABLED}:
  Queries whether back face culling for spheres is enabled. This is only
  the case if Embree is compiled with
  \texttt{EMBREE\_BACKFACE\_CULLING\_SPHERES} enabled.
\item
  \texttt{RTC\_DEVICE\_PROPERTY\_COMPACT\_POLYS\_ENABLED}: Queries
  whether compact polys is enabled. This is only the case if Embree is
  compiled with \texttt{EMBREE\_COMPACT\_POLYS} enabled.
\item
  \texttt{RTC\_DEVICE\_PROPERTY\_FILTER\_FUNCTION\_SUPPORTED}: Queries
  whether filter functions are supported, which is the case if Embree is
  compiled with \texttt{EMBREE\_FILTER\_FUNCTION} enabled.
\item
  \texttt{RTC\_DEVICE\_PROPERTY\_IGNORE\_INVALID\_RAYS\_ENABLED}:
  Queries whether invalid rays are ignored, which is the case if Embree
  is compiled with \texttt{EMBREE\_IGNORE\_INVALID\_RAYS} enabled.
\item
  \texttt{RTC\_DEVICE\_PROPERTY\_TRIANGLE\_GEOMETRY\_SUPPORTED}: Queries
  whether triangles are supported, which is the case if Embree is
  compiled with \texttt{EMBREE\_GEOMETRY\_TRIANGLE} enabled.
\item
  \texttt{RTC\_DEVICE\_PROPERTY\_QUAD\_GEOMETRY\_SUPPORTED}: Queries
  whether quads are supported, which is the case if Embree is compiled
  with \texttt{EMBREE\_GEOMETRY\_QUAD} enabled.
\item
  \texttt{RTC\_DEVICE\_PROPERTY\_SUBDIVISION\_GEOMETRY\_SUPPORTED}:
  Queries whether subdivision meshes are supported, which is the case if
  Embree is compiled with \texttt{EMBREE\_GEOMETRY\_SUBDIVISION}
  enabled.
\item
  \texttt{RTC\_DEVICE\_PROPERTY\_CURVE\_GEOMETRY\_SUPPORTED}: Queries
  whether curves are supported, which is the case if Embree is compiled
  with \texttt{EMBREE\_GEOMETRY\_CURVE} enabled.
\item
  \texttt{RTC\_DEVICE\_PROPERTY\_POINT\_GEOMETRY\_SUPPORTED}: Queries
  whether points are supported, which is the case if Embree is compiled
  with \texttt{EMBREE\_GEOMETRY\_POINT} enabled.
\item
  \texttt{RTC\_DEVICE\_PROPERTY\_USER\_GEOMETRY\_SUPPORTED}: Queries
  whether user geometries are supported, which is the case if Embree is
  compiled with \texttt{EMBREE\_GEOMETRY\_USER} enabled.
\item
  \texttt{RTC\_DEVICE\_PROPERTY\_TASKING\_SYSTEM}: Queries the tasking
  system Embree is compiled with. Possible return values are:

  \begin{enumerate}
  \def\labelenumi{\arabic{enumi}.}
  \setcounter{enumi}{-1}
  \tightlist
  \item
    internal tasking system
  \item
    Intel Threading Building Blocks (TBB)
  \item
    Parallel Patterns Library (PPL)
  \end{enumerate}
\item
  \texttt{RTC\_DEVICE\_PROPERTY\_JOIN\_COMMIT\_SUPPORTED}: Queries
  whether \texttt{rtcJoinCommitScene} is supported. This is not the case
  when Embree is compiled with PPL or older versions of TBB.
\item
  \texttt{RTC\_DEVICE\_PROPERTY\_PARALLEL\_COMMIT\_SUPPORTED}: Queries
  whether \texttt{rtcCommitScene} can get invoked from multiple TBB
  worker threads concurrently. This feature is only supported starting
  with TBB 2019 Update 9.
\end{itemize}

\subsubsection{EXIT STATUS}\label{exit-status}

On success returns the value of the queried property. For properties
returning a boolean value, the return value 0 denotes \texttt{false} and
1 denotes \texttt{true}.

On failure zero is returned and an error code is set that can be queried
using \texttt{rtcGetDeviceError}.

\pagebreak

\section{rtcGetDeviceError}\label{rtcgetdeviceerror}

\subsubsection{NAME}\label{name}

\begin{Shaded}
\begin{Highlighting}[]
\NormalTok{rtcGetDeviceError }\OperatorTok{{-}}\NormalTok{ returns the error code of the device}
\end{Highlighting}
\end{Shaded}

\subsubsection{SYNOPSIS}\label{synopsis}

\begin{Shaded}
\begin{Highlighting}[]
\PreprocessorTok{\#include }\ImportTok{\textless{}embree4/rtcore.h\textgreater{}}

\NormalTok{RTCError rtcGetDeviceError}\OperatorTok{(}\NormalTok{RTCDevice device}\OperatorTok{);}
\end{Highlighting}
\end{Shaded}

\subsubsection{DESCRIPTION}\label{description}

Each thread has its own error code per device. If an error occurs when
calling an API function, this error code is set to the occurred error if
it stores no previous error. The \texttt{rtcGetDeviceError} function
reads and returns the currently stored error and clears the error code.
This assures that the returned error code is always the first error
occurred since the last invocation of \texttt{rtcGetDeviceError}.

Possible error codes returned by \texttt{rtcGetDeviceError} are:

\begin{itemize}
\item
  \texttt{RTC\_ERROR\_NONE}: No error occurred.
\item
  \texttt{RTC\_ERROR\_UNKNOWN}: An unknown error has occurred.
\item
  \texttt{RTC\_ERROR\_INVALID\_ARGUMENT}: An invalid argument was
  specified.
\item
  \texttt{RTC\_ERROR\_INVALID\_OPERATION}: The operation is not allowed
  for the specified object.
\item
  \texttt{RTC\_ERROR\_OUT\_OF\_MEMORY}: There is not enough memory left
  to complete the operation.
\item
  \texttt{RTC\_ERROR\_UNSUPPORTED\_CPU}: The CPU is not supported as it
  does not support the lowest ISA Embree is compiled for.
\item
  \texttt{RTC\_ERROR\_CANCELLED}: The operation got canceled by a memory
  monitor callback or progress monitor callback function.
\item
  \texttt{RTC\_ERROR\_UNSUPPORTED\_GPU}: The specified SYCL GPU device
  is not supported.
\end{itemize}

When the device construction fails, \texttt{rtcNewDevice} returns
\texttt{NULL} as device. To detect the error code of a such a failed
device construction, pass \texttt{NULL} as device to the
\texttt{rtcGetDeviceError} function. For all other invocations of
\texttt{rtcGetDeviceError}, a proper device pointer must be specified.

\subsubsection{EXIT STATUS}\label{exit-status}

Returns the error code for the device.

\subsubsection{SEE ALSO}\label{see-also}

\hyperref[rtcsetdeviceerrorfunction]{rtcSetDeviceErrorFunction}

\pagebreak

\section{rtcSetDeviceErrorFunction}\label{rtcsetdeviceerrorfunction}

\subsubsection{NAME}\label{name}

\begin{Shaded}
\begin{Highlighting}[]
\NormalTok{rtcSetDeviceErrorFunction }\OperatorTok{{-}}\NormalTok{ sets an error callback function }\ControlFlowTok{for}\NormalTok{ the device}
\end{Highlighting}
\end{Shaded}

\subsubsection{SYNOPSIS}\label{synopsis}

\begin{Shaded}
\begin{Highlighting}[]
\PreprocessorTok{\#include }\ImportTok{\textless{}embree4/rtcore.h\textgreater{}}

\KeywordTok{typedef} \DataTypeTok{void} \OperatorTok{(*}\NormalTok{RTCErrorFunction}\OperatorTok{)(}
  \DataTypeTok{void}\OperatorTok{*}\NormalTok{ userPtr}\OperatorTok{,}
\NormalTok{  RTCError code}\OperatorTok{,}
  \AttributeTok{const} \DataTypeTok{char}\OperatorTok{*}\NormalTok{ str}
\OperatorTok{);}

\DataTypeTok{void}\NormalTok{ rtcSetDeviceErrorFunction}\OperatorTok{(}
\NormalTok{  RTCDevice device}\OperatorTok{,}
\NormalTok{  RTCErrorFunction error}\OperatorTok{,}
  \DataTypeTok{void}\OperatorTok{*}\NormalTok{ userPtr}
\OperatorTok{);}
\end{Highlighting}
\end{Shaded}

\subsubsection{DESCRIPTION}\label{description}

Using the \texttt{rtcSetDeviceErrorFunction} call, it is possible to set
a callback function (\texttt{error} argument) with payload
(\texttt{userPtr} argument), which is called whenever an error occurs
for the specified device (\texttt{device} argument).

Only a single callback function can be registered per device, and
further invocations overwrite the previously set callback function.
Passing \texttt{NULL} as function pointer disables the registered
callback function.

When the registered callback function is invoked, it gets passed the
user-defined payload (\texttt{userPtr} argument as specified at
registration time), the error code (\texttt{code} argument) of the
occurred error, as well as a string (\texttt{str} argument) that further
describes the error.

The error code is also set if an error callback function is registered.

\subsubsection{EXIT STATUS}\label{exit-status}

On failure an error code is set that can be queried using
\texttt{rtcGetDeviceError}.

\subsubsection{SEE ALSO}\label{see-also}

\hyperref[rtcgetdeviceerror]{rtcGetDeviceError}

\pagebreak

\section{rtcSetDeviceMemoryMonitorFunction}\label{rtcsetdevicememorymonitorfunction}

\subsubsection{NAME}\label{name}

\begin{Shaded}
\begin{Highlighting}[]
\NormalTok{rtcSetDeviceMemoryMonitorFunction }\OperatorTok{{-}}\NormalTok{ registers a callback function}
\NormalTok{  to track memory consumption}
\end{Highlighting}
\end{Shaded}

\subsubsection{SYNOPSIS}\label{synopsis}

\begin{Shaded}
\begin{Highlighting}[]
\PreprocessorTok{\#include }\ImportTok{\textless{}embree4/rtcore.h\textgreater{}}

\KeywordTok{typedef} \DataTypeTok{bool} \OperatorTok{(*}\NormalTok{RTCMemoryMonitorFunction}\OperatorTok{)(}
  \DataTypeTok{void}\OperatorTok{*}\NormalTok{ userPtr}\OperatorTok{,}
  \DataTypeTok{ssize\_t}\NormalTok{ bytes}\OperatorTok{,}
  \DataTypeTok{bool}\NormalTok{ post}
\OperatorTok{);}

\DataTypeTok{void}\NormalTok{ rtcSetDeviceMemoryMonitorFunction}\OperatorTok{(}
\NormalTok{  RTCDevice device}\OperatorTok{,}
\NormalTok{  RTCMemoryMonitorFunction memoryMonitor}\OperatorTok{,}
  \DataTypeTok{void}\OperatorTok{*}\NormalTok{ userPtr}
\OperatorTok{);}
\end{Highlighting}
\end{Shaded}

\subsubsection{DESCRIPTION}\label{description}

Using the \texttt{rtcSetDeviceMemoryMonitorFunction} call, it is
possible to register a callback function (\texttt{memoryMonitor}
argument) with payload (\texttt{userPtr} argument) for a device
(\texttt{device} argument), which is called whenever internal memory is
allocated or deallocated by objects of that device. Using this memory
monitor callback mechanism, the application can track the memory
consumption of an Embree device, and optionally terminate API calls that
consume too much memory.

Only a single callback function can be registered per device, and
further invocations overwrite the previously set callback function.
Passing \texttt{NULL} as function pointer disables the registered
callback function.

Once registered, the Embree device will invoke the memory monitor
callback function before or after it allocates or frees important memory
blocks. The callback function gets passed the payload as specified at
registration time (\texttt{userPtr} argument), the number of bytes
allocated or deallocated (\texttt{bytes} argument), and whether the
callback is invoked after the allocation or deallocation took place
(\texttt{post} argument). The callback function might get called from
multiple threads concurrently.

The application can track the current memory usage of the Embree device
by atomically accumulating the \texttt{bytes} input parameter provided
to the callback function. This parameter will be \textgreater0 for
allocations and \textless0 for deallocations.

Embree will continue its operation normally when returning \texttt{true}
from the callback function. If \texttt{false} is returned, Embree will
cancel the current operation with the
\texttt{RTC\_ERROR\_OUT\_OF\_MEMORY} error code. Issuing multiple cancel
requests from different threads is allowed. Canceling will only happen
when the callback was called for allocations (bytes \textgreater{} 0),
otherwise the cancel request will be ignored.

If a callback to cancel was invoked before the allocation happens
(\texttt{post\ ==\ false}), then the \texttt{bytes} parameter should not
be accumulated, as the allocation will never happen. If the callback to
cancel was invoked after the allocation happened
(\texttt{post\ ==\ true}), then the \texttt{bytes} parameter should be
accumulated, as the allocation properly happened and a deallocation will
later free that data block.

\subsubsection{EXIT STATUS}\label{exit-status}

On failure an error code is set that can be queried using
\texttt{rtcGetDeviceError}.

\subsubsection{SEE ALSO}\label{see-also}

\hyperref[rtcnewdevice]{rtcNewDevice}

\pagebreak

\section{rtcNewScene}\label{rtcnewscene}

\subsubsection{NAME}\label{name}

\begin{Shaded}
\begin{Highlighting}[]
\NormalTok{rtcNewScene }\OperatorTok{{-}}\NormalTok{ creates a }\KeywordTok{new}\NormalTok{ scene}
\end{Highlighting}
\end{Shaded}

\subsubsection{SYNOPSIS}\label{synopsis}

\begin{Shaded}
\begin{Highlighting}[]
\PreprocessorTok{\#include }\ImportTok{\textless{}embree4/rtcore.h\textgreater{}}

\NormalTok{RTCScene rtcNewScene}\OperatorTok{(}\NormalTok{RTCDevice device}\OperatorTok{);}
\end{Highlighting}
\end{Shaded}

\subsubsection{DESCRIPTION}\label{description}

This function creates a new scene bound to the specified device
(\texttt{device} argument), and returns a handle to this scene. The
scene object is reference counted with an initial reference count of 1.
The scene handle can be released using the \texttt{rtcReleaseScene} API
call.

\subsubsection{EXIT STATUS}\label{exit-status}

On success a scene handle is returned. On failure \texttt{NULL} is
returned and an error code is set that can be queried using
\texttt{rtcGetDeviceError}.

\subsubsection{SEE ALSO}\label{see-also}

\hyperref[rtcretainscene]{rtcRetainScene},
\hyperref[rtcreleasescene]{rtcReleaseScene}

\pagebreak

\section{rtcGetSceneDevice}\label{rtcgetscenedevice}

\subsubsection{NAME}\label{name}

\begin{Shaded}
\begin{Highlighting}[]
\NormalTok{rtcGetSceneDevice }\OperatorTok{{-}}\NormalTok{ returns the device the scene got created in}
\end{Highlighting}
\end{Shaded}

\subsubsection{SYNOPSIS}\label{synopsis}

\begin{Shaded}
\begin{Highlighting}[]
\PreprocessorTok{\#include }\ImportTok{\textless{}embree4/rtcore.h\textgreater{}}

\NormalTok{RTCDevice rtcGetSceneDevice}\OperatorTok{(}\NormalTok{RTCScene scene}\OperatorTok{);}
\end{Highlighting}
\end{Shaded}

\subsubsection{DESCRIPTION}\label{description}

This function returns the device object the scene got created in. The
returned handle own one additional reference to the device object, thus
you should need to call \texttt{rtcReleaseDevice} when the returned
handle is no longer required.

\subsubsection{EXIT STATUS}\label{exit-status}

On failure an error code is set that can be queried using
\texttt{rtcGetDeviceError}.

\subsubsection{SEE ALSO}\label{see-also}

\hyperref[rtcreleasedevice]{rtcReleaseDevice}

\pagebreak

\section{rtcRetainScene}\label{rtcretainscene}

\subsubsection{NAME}\label{name}

\begin{Shaded}
\begin{Highlighting}[]
\NormalTok{rtcRetainScene }\OperatorTok{{-}}\NormalTok{ increments the scene reference count}
\end{Highlighting}
\end{Shaded}

\subsubsection{SYNOPSIS}\label{synopsis}

\begin{Shaded}
\begin{Highlighting}[]
\PreprocessorTok{\#include }\ImportTok{\textless{}embree4/rtcore.h\textgreater{}}

\DataTypeTok{void}\NormalTok{ rtcRetainScene}\OperatorTok{(}\NormalTok{RTCScene scene}\OperatorTok{);}
\end{Highlighting}
\end{Shaded}

\subsubsection{DESCRIPTION}\label{description}

Scene objects are reference counted. The \texttt{rtcRetainScene}
function increments the reference count of the passed scene object
(\texttt{scene} argument). This function together with
\texttt{rtcReleaseScene} allows to use the internal reference counting
in a C++ wrapper class to handle the ownership of the object.

\subsubsection{EXIT STATUS}\label{exit-status}

On failure an error code is set that can be queried using
\texttt{rtcGetDeviceError}.

\subsubsection{SEE ALSO}\label{see-also}

\hyperref[rtcnewscene]{rtcNewScene},
\hyperref[rtcreleasescene]{rtcReleaseScene}

\pagebreak

\section{rtcReleaseScene}\label{rtcreleasescene}

\subsubsection{NAME}\label{name}

\begin{Shaded}
\begin{Highlighting}[]
\NormalTok{rtcReleaseScene }\OperatorTok{{-}}\NormalTok{ decrements the scene reference count}
\end{Highlighting}
\end{Shaded}

\subsubsection{SYNOPSIS}\label{synopsis}

\begin{Shaded}
\begin{Highlighting}[]
\PreprocessorTok{\#include }\ImportTok{\textless{}embree4/rtcore.h\textgreater{}}

\DataTypeTok{void}\NormalTok{ rtcReleaseScene}\OperatorTok{(}\NormalTok{RTCScene scene}\OperatorTok{);}
\end{Highlighting}
\end{Shaded}

\subsubsection{DESCRIPTION}\label{description}

Scene objects are reference counted. The \texttt{rtcReleaseScene}
function decrements the reference count of the passed scene object
(\texttt{scene} argument). When the reference count falls to 0, the
scene gets destroyed.

The scene holds a reference to all attached geometries, thus if the
scene gets destroyed, all geometries get detached and their reference
count decremented.

\subsubsection{EXIT STATUS}\label{exit-status}

On failure an error code is set that can be queried using
\texttt{rtcGetDeviceError}.

\subsubsection{SEE ALSO}\label{see-also}

\hyperref[rtcnewscene]{rtcNewScene},
\hyperref[rtcretainscene]{rtcRetainScene}

\pagebreak

\section{rtcAttachGeometry}\label{rtcattachgeometry}

\subsubsection{NAME}\label{name}

\begin{Shaded}
\begin{Highlighting}[]
\NormalTok{rtcAttachGeometry }\OperatorTok{{-}}\NormalTok{ attaches a geometry to the scene}
\end{Highlighting}
\end{Shaded}

\subsubsection{SYNOPSIS}\label{synopsis}

\begin{Shaded}
\begin{Highlighting}[]
\PreprocessorTok{\#include }\ImportTok{\textless{}embree4/rtcore.h\textgreater{}}

\DataTypeTok{unsigned} \DataTypeTok{int}\NormalTok{ rtcAttachGeometry}\OperatorTok{(}
\NormalTok{  RTCScene scene}\OperatorTok{,}
\NormalTok{  RTCGeometry geometry}
\OperatorTok{);}
\end{Highlighting}
\end{Shaded}

\subsubsection{DESCRIPTION}\label{description}

The \texttt{rtcAttachGeometry} function attaches a geometry
(\texttt{geometry} argument) to a scene (\texttt{scene} argument) and
assigns a geometry ID to that geometry. All geometries attached to a
scene are defined to be included inside the scene. A geometry can get
attached to multiple scenes. The geometry ID is unique for the scene,
and is used to identify the geometry when hit by a ray during ray
queries.

This function is thread-safe, thus multiple threads can attach
geometries to a scene in parallel.

The geometry IDs are assigned sequentially, starting from 0, as long as
no geometry got detached. If geometries got detached, the implementation
will reuse IDs in an implementation dependent way. Consequently
sequential assignment is no longer guaranteed, but a compact range of
IDs.

These rules allow the application to manage a dynamic array to
efficiently map from geometry IDs to its own geometry representation.
Alternatively, the application can also use per-geometry user data to
map to its geometry representation. See \texttt{rtcSetGeometryUserData}
and \texttt{rtcGetGeometryUserData} for more information.

\subsubsection{EXIT STATUS}\label{exit-status}

On failure an error code is set that can be queried using
\texttt{rtcGetDeviceError}.

\subsubsection{SEE ALSO}\label{see-also}

\hyperref[rtcsetgeometryuserdata]{rtcSetGeometryUserData},
\hyperref[rtcgetgeometryuserdata]{rtcGetGeometryUserData}

\pagebreak

\section{rtcAttachGeometryByID}\label{rtcattachgeometrybyid}

\subsubsection{NAME}\label{name}

\begin{Shaded}
\begin{Highlighting}[]
\NormalTok{rtcAttachGeometryByID }\OperatorTok{{-}}\NormalTok{ attaches a geometry to the scene}
  \KeywordTok{using}\NormalTok{ a specified geometry ID}
\end{Highlighting}
\end{Shaded}

\subsubsection{SYNOPSIS}\label{synopsis}

\begin{Shaded}
\begin{Highlighting}[]
\PreprocessorTok{\#include }\ImportTok{\textless{}embree4/rtcore.h\textgreater{}}

\DataTypeTok{void}\NormalTok{ rtcAttachGeometryByID}\OperatorTok{(}
\NormalTok{  RTCScene scene}\OperatorTok{,}
\NormalTok{  RTCGeometry geometry}\OperatorTok{,}
  \DataTypeTok{unsigned} \DataTypeTok{int}\NormalTok{ geomID}
\OperatorTok{);}
\end{Highlighting}
\end{Shaded}

\subsubsection{DESCRIPTION}\label{description}

The \texttt{rtcAttachGeometryByID} function attaches a geometry
(\texttt{geometry} argument) to a scene (\texttt{scene} argument) and
assigns a user provided geometry ID (\texttt{geomID} argument) to that
geometry. All geometries attached to a scene are defined to be included
inside the scene. A geometry can get attached to multiple scenes. The
passed user-defined geometry ID is used to identify the geometry when
hit by a ray during ray queries. Using this function, it is possible to
share the same IDs to refer to geometries inside the application and
Embree.

This function is thread-safe, thus multiple threads can attach
geometries to a scene in parallel.

The user-provided geometry ID must be unused in the scene, otherwise the
creation of the geometry will fail. Further, the user-provided geometry
IDs should be compact, as Embree internally creates a vector which size
is equal to the largest geometry ID used. Creating very large geometry
IDs for small scenes would thus cause a memory consumption and
performance overhead.

\subsubsection{EXIT STATUS}\label{exit-status}

On failure an error code is set that can be queried using
\texttt{rtcGetDeviceError}.

\subsubsection{SEE ALSO}\label{see-also}

\hyperref[rtcattachgeometry]{rtcAttachGeometry}

\pagebreak

\section{rtcDetachGeometry}\label{rtcdetachgeometry}

\subsubsection{NAME}\label{name}

\begin{Shaded}
\begin{Highlighting}[]
\NormalTok{rtcDetachGeometry }\OperatorTok{{-}}\NormalTok{ detaches a geometry from the scene}
\end{Highlighting}
\end{Shaded}

\subsubsection{SYNOPSIS}\label{synopsis}

\begin{Shaded}
\begin{Highlighting}[]
\PreprocessorTok{\#include }\ImportTok{\textless{}embree4/rtcore.h\textgreater{}}

\DataTypeTok{void}\NormalTok{ rtcDetachGeometry}\OperatorTok{(}\NormalTok{RTCScene scene}\OperatorTok{,} \DataTypeTok{unsigned} \DataTypeTok{int}\NormalTok{ geomID}\OperatorTok{);}
\end{Highlighting}
\end{Shaded}

\subsubsection{DESCRIPTION}\label{description}

This function detaches a geometry identified by its geometry ID
(\texttt{geomID} argument) from a scene (\texttt{scene} argument). When
detached, the geometry is no longer contained in the scene.

This function is thread-safe, thus multiple threads can detach
geometries from a scene at the same time.

\subsubsection{EXIT STATUS}\label{exit-status}

On failure an error code is set that can be queried using
\texttt{rtcGetDeviceError}.

\subsubsection{SEE ALSO}\label{see-also}

\hyperref[rtcattachgeometry]{rtcAttachGeometry},
\hyperref[rtcattachgeometrybyid]{rtcAttachGeometryByID}

\pagebreak

\section{rtcGetGeometry}\label{rtcgetgeometry}

\subsubsection{NAME}\label{name}

\begin{Shaded}
\begin{Highlighting}[]
\NormalTok{rtcGetGeometry }\OperatorTok{{-}}\NormalTok{ returns the geometry bound to}
\NormalTok{  the specified geometry ID}
\end{Highlighting}
\end{Shaded}

\subsubsection{SYNOPSIS}\label{synopsis}

\begin{Shaded}
\begin{Highlighting}[]
\PreprocessorTok{\#include }\ImportTok{\textless{}embree4/rtcore.h\textgreater{}}

\NormalTok{RTCGeometry rtcGetGeometry}\OperatorTok{(}\NormalTok{RTCScene scene}\OperatorTok{,} \DataTypeTok{unsigned} \DataTypeTok{int}\NormalTok{ geomID}\OperatorTok{);}
\end{Highlighting}
\end{Shaded}

\subsubsection{DESCRIPTION}\label{description}

The \texttt{rtcGetGeometry} function returns the geometry that is bound
to the specified geometry ID (\texttt{geomID} argument) for the
specified scene (\texttt{scene} argument). This function just looks up
the handle and does \emph{not} increment the reference count. If you
want to get ownership of the handle, you need to additionally call
\texttt{rtcRetainGeometry}.

This function is not thread safe and thus can be used during rendering.
However, it is generally recommended to store the geometry handle inside
the application's geometry representation and look up the geometry
handle from that representation directly.

If you need a thread safe version of this function please use
\hyperref[rtcgetgeometrythreadsafe]{rtcGetGeometryThreadSafe}.

\subsubsection{EXIT STATUS}\label{exit-status}

On failure \texttt{NULL} is returned and an error code is set that can
be queried using \texttt{rtcGetDeviceError}.

\subsubsection{SEE ALSO}\label{see-also}

\hyperref[rtcattachgeometry]{rtcAttachGeometry},
\hyperref[rtcattachgeometrybyid]{rtcAttachGeometryByID},
\hyperref[rtcgetgeometrythreadsafe]{rtcGetGeometryThreadSafe}

\pagebreak

\section{rtcGetGeometryThreadSafe}\label{rtcgetgeometrythreadsafe}

\subsubsection{NAME}\label{name}

\begin{Shaded}
\begin{Highlighting}[]
\NormalTok{rtcGetGeometryThreadSafe }\OperatorTok{{-}}\NormalTok{ returns the geometry bound to}
\NormalTok{  the specified geometry ID}
\end{Highlighting}
\end{Shaded}

\subsubsection{SYNOPSIS}\label{synopsis}

\begin{Shaded}
\begin{Highlighting}[]
\PreprocessorTok{\#include }\ImportTok{\textless{}embree4/rtcore.h\textgreater{}}

\NormalTok{RTCGeometry rtcGetGeometryThreadSafe}\OperatorTok{(}\NormalTok{RTCScene scene}\OperatorTok{,} \DataTypeTok{unsigned} \DataTypeTok{int}\NormalTok{ geomID}\OperatorTok{);}
\end{Highlighting}
\end{Shaded}

\subsubsection{DESCRIPTION}\label{description}

The \texttt{rtcGetGeometryThreadSafe} function returns the geometry that
is bound to the specified geometry ID (\texttt{geomID} argument) for the
specified scene (\texttt{scene} argument). This function just looks up
the handle and does \emph{not} increment the reference count. If you
want to get ownership of the handle, you need to additionally call
\texttt{rtcRetainGeometry}.

This function is thread safe and should NOT get used during rendering.
If you need a fast non-thread safe version during rendering please use
the \hyperref[rtcgetgeometry]{rtcGetGeometry} function.

\subsubsection{EXIT STATUS}\label{exit-status}

On failure \texttt{NULL} is returned and an error code is set that can
be queried using \texttt{rtcGetDeviceError}.

\subsubsection{SEE ALSO}\label{see-also}

\hyperref[rtcattachgeometry]{rtcAttachGeometry},
\hyperref[rtcattachgeometrybyid]{rtcAttachGeometryByID},
\hyperref[rtcgetgeometry]{rtcGetGeometry}

\pagebreak

\section{rtcCommitScene}\label{rtccommitscene}

\subsubsection{NAME}\label{name}

\begin{Shaded}
\begin{Highlighting}[]
\NormalTok{rtcCommitScene }\OperatorTok{{-}}\NormalTok{ commits scene changes}
\end{Highlighting}
\end{Shaded}

\subsubsection{SYNOPSIS}\label{synopsis}

\begin{Shaded}
\begin{Highlighting}[]
\PreprocessorTok{\#include }\ImportTok{\textless{}embree4/rtcore.h\textgreater{}}

\DataTypeTok{void}\NormalTok{ rtcCommitScene}\OperatorTok{(}\NormalTok{RTCScene scene}\OperatorTok{);}
\end{Highlighting}
\end{Shaded}

\subsubsection{DESCRIPTION}\label{description}

The \texttt{rtcCommitScene} function commits all changes for the
specified scene (\texttt{scene} argument). This internally triggers
building of a spatial acceleration structure for the scene using all
available worker threads. Ray queries can be performed only after
committing all scene changes.

If the application uses TBB 2019 Update 9 or later for parallelization
of rendering, lazy scene construction during rendering is supported by
\texttt{rtcCommitScene}. Therefore \texttt{rtcCommitScene} can get
called from multiple TBB worker threads concurrently for the same scene.
The \texttt{rtcCommitScene} function will then internally isolate the
scene construction using a tbb::isolated\_task\_group. The alternative
approach of using \texttt{rtcJoinCommitScene} which uses an
tbb:task\_arena internally, is not recommended due to it's high runtime
overhead.

If scene geometries get modified or attached or detached, the
\texttt{rtcCommitScene} call must be invoked before performing any
further ray queries for the scene; otherwise the effect of the ray query
is undefined. The modification of a geometry, committing the scene, and
tracing of rays must always happen sequentially, and never at the same
time. Any API call that sets a property of the scene or geometries
contained in the scene count as scene modification, e.g.~including
setting of intersection filter functions.

The kind of acceleration structure built can be influenced using scene
flags (see \texttt{rtcSetSceneFlags}), and the quality can be specified
using the \texttt{rtcSetSceneBuildQuality} function.

Embree silently ignores primitives during spatial acceleration structure
construction that would cause numerical issues, e.g.~primitives
containing NaNs, INFs, or values greater than 1.844E18f (as no
reasonable calculations can be performed with such values without
causing overflows).

\subsubsection{EXIT STATUS}\label{exit-status}

On failure an error code is set that can be queried using
\texttt{rtcGetDeviceError}.

\subsubsection{SEE ALSO}\label{see-also}

\hyperref[rtcjoincommitscene]{rtcJoinCommitScene}

\pagebreak

\section{rtcJoinCommitScene}\label{rtcjoincommitscene}

\subsubsection{NAME}\label{name}

\begin{Shaded}
\begin{Highlighting}[]
\NormalTok{rtcJoinCommitScene }\OperatorTok{{-}}\NormalTok{ commits the scene from multiple threads}
\end{Highlighting}
\end{Shaded}

\subsubsection{SYNOPSIS}\label{synopsis}

\begin{Shaded}
\begin{Highlighting}[]
\PreprocessorTok{\#include }\ImportTok{\textless{}embree4/rtcore.h\textgreater{}}

\DataTypeTok{void}\NormalTok{ rtcJoinCommitScene}\OperatorTok{(}\NormalTok{RTCScene scene}\OperatorTok{);}
\end{Highlighting}
\end{Shaded}

\subsubsection{DESCRIPTION}\label{description}

The \texttt{rtcJoinCommitScene} function commits all changes for the
specified scene (\texttt{scene} argument). The scene commit internally
triggers building of a spatial acceleration structure for the scene. Ray
queries can be performed after scene changes got properly committed.

The \texttt{rtcJoinCommitScene} function can get called from multiple
user threads which will all cooperate in the build operation. All
threads calling into this function will return from
\texttt{rtcJoinCommitScene} after the scene commit is finished. All
threads must consistently call \texttt{rtcJoinCommitScene} and not
\texttt{rtcCommitScene}.

In contrast to the \texttt{rtcCommitScene} function, the
\texttt{rtcJoinCommitScene} function can be called from multiple user
threads, while the \texttt{rtcCommitScene} can only get called from
multiple TBB worker threads when used concurrently. For optimal
performance we strongly recommend using TBB inside the application
together with the \texttt{rtcCommitScene} function and to avoid using
the \texttt{rtcJoinCommitScene} function.

The \texttt{rtcJoinCommitScene} feature allows a flexible way to lazily
create hierarchies during rendering. A thread reaching a
not-yet-constructed sub-scene of a two-level scene can generate the
sub-scene geometry and call \texttt{rtcJoinCommitScene} on that just
generated scene. During construction, further threads reaching the
not-yet-built scene can join the build operation by also invoking
\texttt{rtcJoinCommitScene}. A thread that calls
\texttt{rtcJoinCommitScene} after the build finishes will directly
return from the \texttt{rtcJoinCommitScene} call.

Multiple scene commit operations on different scenes can be running at
the same time, hence it is possible to commit many small scenes in
parallel, distributing the commits to many threads.

When using Embree with the Intel® Threading Building Blocks (which is
the default), threads that call \texttt{rtcJoinCommitScene} will join
the build operation, but other TBB worker threads might also participate
in the build. To avoid thread oversubscription, we recommend using TBB
also inside the application. Further, the join mode only works properly
starting with TBB v4.4 Update 1. For earlier TBB versions, threads that
call \texttt{rtcJoinCommitScene} to join a running build will just
trigger the build and wait for the build to finish. Further, old TBB
versions with \texttt{TBB\_INTERFACE\_VERSION\_MAJOR\ \textless{}\ 8} do
not support \texttt{rtcJoinCommitScene}, and invoking this function will
result in an error.

When using Embree with the internal tasking system, only threads that
call \texttt{rtcJoinCommitScene} will perform the build operation, and
no additional worker threads will be scheduled.

When using Embree with the Parallel Patterns Library (PPL),
\texttt{rtcJoinCommitScene} is not supported and calling that function
will result in an error.

To detect whether \texttt{rtcJoinCommitScene} is supported, use the
\texttt{rtcGetDeviceProperty} function.

\subsubsection{EXIT STATUS}\label{exit-status}

On failure an error code is set that can be queried using
\texttt{rtcGetDeviceError}.

\subsubsection{SEE ALSO}\label{see-also}

\hyperref[rtccommitscene]{rtcCommitScene},
\hyperref[rtcgetdeviceproperty]{rtcGetDeviceProperty}

\pagebreak

\section{rtcSetSceneProgressMonitorFunction}\label{rtcsetsceneprogressmonitorfunction}

\subsubsection{NAME}\label{name}

\begin{Shaded}
\begin{Highlighting}[]
\NormalTok{rtcSetSceneProgressMonitorFunction }\OperatorTok{{-}}\NormalTok{ registers a callback}
\NormalTok{  to track build progress}
\end{Highlighting}
\end{Shaded}

\subsubsection{SYNOPSIS}\label{synopsis}

\begin{Shaded}
\begin{Highlighting}[]
\PreprocessorTok{\#include }\ImportTok{\textless{}embree4/rtcore.h\textgreater{}}

\KeywordTok{typedef} \DataTypeTok{bool} \OperatorTok{(*}\NormalTok{RTCProgressMonitorFunction}\OperatorTok{)(}
  \DataTypeTok{void}\OperatorTok{*}\NormalTok{ ptr}\OperatorTok{,}
  \DataTypeTok{double}\NormalTok{ n}
\OperatorTok{);}

\DataTypeTok{void}\NormalTok{ rtcSetSceneProgressMonitorFunction}\OperatorTok{(}
\NormalTok{  RTCScene scene}\OperatorTok{,}
\NormalTok{  RTCProgressMonitorFunction progress}\OperatorTok{,}
  \DataTypeTok{void}\OperatorTok{*}\NormalTok{ userPtr}
\OperatorTok{);}
\end{Highlighting}
\end{Shaded}

\subsubsection{DESCRIPTION}\label{description}

Embree supports a progress monitor callback mechanism that can be used
to report progress of hierarchy build operations and to cancel build
operations.

The \texttt{rtcSetSceneProgressMonitorFunction} registers a progress
monitor callback function (\texttt{progress} argument) with payload
(\texttt{userPtr} argument) for the specified scene (\texttt{scene}
argument).

Only a single callback function can be registered per scene, and further
invocations overwrite the previously set callback function. Passing
\texttt{NULL} as function pointer disables the registered callback
function.

Once registered, Embree will invoke the callback function multiple times
during hierarchy build operations of the scene, by passing the payload
as set at registration time (\texttt{userPtr} argument), and a double in
the range \([0, 1]\) which estimates the progress of the operation
(\texttt{n} argument). The callback function might be called from
multiple threads concurrently.

When returning \texttt{true} from the callback function, Embree will
continue the build operation normally. When returning \texttt{false},
Embree will cancel the build operation with the
\texttt{RTC\_ERROR\_CANCELLED} error code. Issuing multiple cancel
requests for the same build operation is allowed.

\subsubsection{EXIT STATUS}\label{exit-status}

On failure an error code is set that can be queried using
\texttt{rtcGetDeviceError}.

\subsubsection{SEE ALSO}\label{see-also}

\hyperref[rtcnewscene]{rtcNewScene}

\pagebreak

\section{rtcSetSceneBuildQuality}\label{rtcsetscenebuildquality}

\subsubsection{NAME}\label{name}

\begin{Shaded}
\begin{Highlighting}[]
\NormalTok{rtcSetSceneBuildQuality }\OperatorTok{{-}}\NormalTok{ sets the build quality }\ControlFlowTok{for}
\NormalTok{  the scene}
\end{Highlighting}
\end{Shaded}

\subsubsection{SYNOPSIS}\label{synopsis}

\begin{Shaded}
\begin{Highlighting}[]
\PreprocessorTok{\#include }\ImportTok{\textless{}embree4/rtcore.h\textgreater{}}

\DataTypeTok{void}\NormalTok{ rtcSetSceneBuildQuality}\OperatorTok{(}
\NormalTok{  RTCScene scene}\OperatorTok{,}
  \KeywordTok{enum}\NormalTok{ RTCBuildQuality quality}
\OperatorTok{);}
\end{Highlighting}
\end{Shaded}

\subsubsection{DESCRIPTION}\label{description}

The \texttt{rtcSetSceneBuildQuality} function sets the build quality
(\texttt{quality} argument) for the specified scene (\texttt{scene}
argument). Possible values for the build quality are:

\begin{itemize}
\item
  \texttt{RTC\_BUILD\_QUALITY\_LOW}: Create lower quality data
  structures, e.g.~for dynamic scenes. A two-level spatial index
  structure is built when enabling this mode, which supports fast
  partial scene updates, and allows for setting a per-geometry build
  quality through the \texttt{rtcSetGeometryBuildQuality} function.
\item
  \texttt{RTC\_BUILD\_QUALITY\_MEDIUM}: Default build quality for most
  usages. Gives a good compromise between build and render performance.
\item
  \texttt{RTC\_BUILD\_QUALITY\_HIGH}: Create higher quality data
  structures for final-frame rendering. For certain geometry types this
  enables a spatial split BVH. When high quality mode is enabled, filter
  callbacks may be invoked multiple times for the same geometry.
\end{itemize}

Selecting a higher build quality results in better rendering performance
but slower scene commit times. The default build quality for a scene is
\texttt{RTC\_BUILD\_QUALITY\_MEDIUM}.

\subsubsection{EXIT STATUS}\label{exit-status}

On failure an error code is set that can be queried using
\texttt{rtcGetDeviceError}.

\subsubsection{SEE ALSO}\label{see-also}

\hyperref[rtcsetgeometrybuildquality]{rtcSetGeometryBuildQuality}

\pagebreak

\section{rtcSetSceneFlags}\label{rtcsetsceneflags}

\subsubsection{NAME}\label{name}

\begin{Shaded}
\begin{Highlighting}[]
\NormalTok{rtcSetSceneFlags }\OperatorTok{{-}}\NormalTok{ sets the flags }\ControlFlowTok{for}\NormalTok{ the scene}
\end{Highlighting}
\end{Shaded}

\subsubsection{SYNOPSIS}\label{synopsis}

\begin{Shaded}
\begin{Highlighting}[]
\PreprocessorTok{\#include }\ImportTok{\textless{}embree4/rtcore.h\textgreater{}}

\KeywordTok{enum}\NormalTok{ RTCSceneFlags}
\OperatorTok{\{}
\NormalTok{  RTC\_SCENE\_FLAG\_NONE                    }\OperatorTok{=} \DecValTok{0}\OperatorTok{,}
\NormalTok{  RTC\_SCENE\_FLAG\_DYNAMIC                 }\OperatorTok{=} \OperatorTok{(}\DecValTok{1} \OperatorTok{\textless{}\textless{}} \DecValTok{0}\OperatorTok{),}
\NormalTok{  RTC\_SCENE\_FLAG\_COMPACT                 }\OperatorTok{=} \OperatorTok{(}\DecValTok{1} \OperatorTok{\textless{}\textless{}} \DecValTok{1}\OperatorTok{),}
\NormalTok{  RTC\_SCENE\_FLAG\_ROBUST                  }\OperatorTok{=} \OperatorTok{(}\DecValTok{1} \OperatorTok{\textless{}\textless{}} \DecValTok{2}\OperatorTok{),}
\NormalTok{  RTC\_SCENE\_FLAG\_FILTER\_FUNCTION\_IN\_ARGUMENTS }\OperatorTok{=} \OperatorTok{(}\DecValTok{1} \OperatorTok{\textless{}\textless{}} \DecValTok{3}\OperatorTok{)}
\OperatorTok{\};}

\DataTypeTok{void}\NormalTok{ rtcSetSceneFlags}\OperatorTok{(}\NormalTok{RTCScene scene}\OperatorTok{,} \KeywordTok{enum}\NormalTok{ RTCSceneFlags flags}\OperatorTok{);}
\end{Highlighting}
\end{Shaded}

\subsubsection{DESCRIPTION}\label{description}

The \texttt{rtcSetSceneFlags} function sets the scene flags
(\texttt{flags} argument) for the specified scene (\texttt{scene}
argument). Possible scene flags are:

\begin{itemize}
\item
  \texttt{RTC\_SCENE\_FLAG\_NONE}: No flags set.
\item
  \texttt{RTC\_SCENE\_FLAG\_DYNAMIC}: Provides better build performance
  for dynamic scenes (but also higher memory consumption).
\item
  \texttt{RTC\_SCENE\_FLAG\_COMPACT}: Uses compact acceleration
  structures and avoids algorithms that consume much memory.
\item
  \texttt{RTC\_SCENE\_FLAG\_ROBUST}: Uses acceleration structures that
  allow for robust traversal, and avoids optimizations that reduce
  arithmetic accuracy. This mode is typically used for avoiding
  artifacts caused by rays shooting through edges of neighboring
  primitives.
\item
  \texttt{RTC\_SCENE\_FLAG\_FILTER\_FUNCTION\_IN\_ARGUMENTS}: Enables
  scene support for filter functions passed as argument to the traversal
  functions. See Section
  \hyperref[rtcinitintersectarguments]{rtcInitIntersectArguments} and
  \hyperref[rtcinitoccludedarguments]{rtcInitOccludedArguments} for more
  details.
\end{itemize}

Multiple flags can be enabled using an \texttt{or} operation,
e.g.~\texttt{RTC\_SCENE\_FLAG\_COMPACT\ \textbar{}\ RTC\_SCENE\_FLAG\_ROBUST}.

\subsubsection{EXIT STATUS}\label{exit-status}

On failure an error code is set that can be queried using
\texttt{rtcGetDeviceError}.

\subsubsection{SEE ALSO}\label{see-also}

\hyperref[rtcgetsceneflags]{rtcGetSceneFlags}

\pagebreak

\section{rtcGetSceneFlags}\label{rtcgetsceneflags}

\subsubsection{NAME}\label{name}

\begin{Shaded}
\begin{Highlighting}[]
\NormalTok{rtcGetSceneFlags }\OperatorTok{{-}}\NormalTok{ returns the flags of the scene}
\end{Highlighting}
\end{Shaded}

\subsubsection{SYNOPSIS}\label{synopsis}

\begin{Shaded}
\begin{Highlighting}[]
\PreprocessorTok{\#include }\ImportTok{\textless{}embree4/rtcore.h\textgreater{}}

\KeywordTok{enum}\NormalTok{ RTCSceneFlags rtcGetSceneFlags}\OperatorTok{(}\NormalTok{RTCScene scene}\OperatorTok{);}
\end{Highlighting}
\end{Shaded}

\subsubsection{DESCRIPTION}\label{description}

Queries the flags of a scene. This function can be useful when setting
individual flags, e.g.~to just set the robust mode without changing
other flags the following way:

\begin{Shaded}
\begin{Highlighting}[]
\NormalTok{RTCSceneFlags flags }\OperatorTok{=}\NormalTok{ rtcGetSceneFlags}\OperatorTok{(}\NormalTok{scene}\OperatorTok{);}
\NormalTok{rtcSetSceneFlags}\OperatorTok{(}\NormalTok{scene}\OperatorTok{,}\NormalTok{ RTC\_SCENE\_FLAG\_ROBUST }\OperatorTok{|}\NormalTok{ flags}\OperatorTok{);}
\end{Highlighting}
\end{Shaded}

\subsubsection{EXIT STATUS}\label{exit-status}

On failure \texttt{RTC\_SCENE\_FLAG\_NONE} is returned and an error code
is set that can be queried using \texttt{rtcGetDeviceError}.

\subsubsection{SEE ALSO}\label{see-also}

\hyperref[rtcsetsceneflags]{rtcSetSceneFlags}

\pagebreak

\section{rtcGetSceneBounds}\label{rtcgetscenebounds}

\subsubsection{NAME}\label{name}

\begin{Shaded}
\begin{Highlighting}[]
\NormalTok{rtcGetSceneBounds }\OperatorTok{{-}}\NormalTok{ returns the axis}\OperatorTok{{-}}\NormalTok{aligned bounding box of the scene}
\end{Highlighting}
\end{Shaded}

\subsubsection{SYNOPSIS}\label{synopsis}

\begin{Shaded}
\begin{Highlighting}[]
\PreprocessorTok{\#include }\ImportTok{\textless{}embree4/rtcore.h\textgreater{}}

\KeywordTok{struct}\NormalTok{ RTCORE\_ALIGN}\OperatorTok{(}\DecValTok{16}\OperatorTok{)}\NormalTok{ RTCBounds}
\OperatorTok{\{}
  \DataTypeTok{float}\NormalTok{ lower\_x}\OperatorTok{,}\NormalTok{ lower\_y}\OperatorTok{,}\NormalTok{ lower\_z}\OperatorTok{,}\NormalTok{ align0}\OperatorTok{;}
  \DataTypeTok{float}\NormalTok{ upper\_x}\OperatorTok{,}\NormalTok{ upper\_y}\OperatorTok{,}\NormalTok{ upper\_z}\OperatorTok{,}\NormalTok{ align1}\OperatorTok{;}
\OperatorTok{\};}

\DataTypeTok{void}\NormalTok{ rtcGetSceneBounds}\OperatorTok{(}
\NormalTok{  RTCScene scene}\OperatorTok{,}
  \KeywordTok{struct}\NormalTok{ RTCBounds}\OperatorTok{*}\NormalTok{ bounds\_o}
\OperatorTok{);}
\end{Highlighting}
\end{Shaded}

\subsubsection{DESCRIPTION}\label{description}

The \texttt{rtcGetSceneBounds} function queries the axis-aligned
bounding box of the specified scene (\texttt{scene} argument) and stores
that bounding box to the provided destination pointer
(\texttt{bounds\_o} argument). The stored bounding box consists of lower
and upper bounds for the x, y, and z dimensions as specified by the
\texttt{RTCBounds} structure.

The provided destination pointer must be aligned to 16 bytes. The
function may be invoked only after committing the scene; otherwise the
result is undefined.

\subsubsection{EXIT STATUS}\label{exit-status}

On failure an error code is set that can be queried using
\texttt{rtcGetDeviceError}.

\subsubsection{SEE ALSO}\label{see-also}

\hyperref[rtcgetscenelinearbounds]{rtcGetSceneLinearBounds},
\hyperref[rtccommitscene]{rtcCommitScene},
\hyperref[rtcjoincommitscene]{rtcJoinCommitScene}

\pagebreak

\section{rtcGetSceneLinearBounds}\label{rtcgetscenelinearbounds}

\subsubsection{NAME}\label{name}

\begin{Shaded}
\begin{Highlighting}[]
\NormalTok{rtcGetSceneLinearBounds }\OperatorTok{{-}}\NormalTok{ returns the linear bounds of the scene}
\end{Highlighting}
\end{Shaded}

\subsubsection{SYNOPSIS}\label{synopsis}

\begin{Shaded}
\begin{Highlighting}[]
\PreprocessorTok{\#include }\ImportTok{\textless{}embree4/rtcore.h\textgreater{}}

\KeywordTok{struct}\NormalTok{ RTCORE\_ALIGN}\OperatorTok{(}\DecValTok{16}\OperatorTok{)}\NormalTok{ RTCLinearBounds}
\OperatorTok{\{}
\NormalTok{  RTCBounds bounds0}\OperatorTok{;}
\NormalTok{  RTCBounds bounds1}\OperatorTok{;}
\OperatorTok{\};}

\DataTypeTok{void}\NormalTok{ rtcGetSceneLinearBounds}\OperatorTok{(}
\NormalTok{  RTCScene scene}\OperatorTok{,}
  \KeywordTok{struct}\NormalTok{ RTCLinearBounds}\OperatorTok{*}\NormalTok{ bounds\_o}
\OperatorTok{);}
\end{Highlighting}
\end{Shaded}

\subsubsection{DESCRIPTION}\label{description}

The \texttt{rtcGetSceneLinearBounds} function queries the linear bounds
of the specified scene (\texttt{scene} argument) and stores them to the
provided destination pointer (\texttt{bounds\_o} argument). The stored
linear bounds consist of bounding boxes for time 0 (\texttt{bounds0}
member) and time 1 (\texttt{bounds1} member) as specified by the
\texttt{RTCLinearBounds} structure. Linearly interpolating these bounds
to a specific time \texttt{t} yields bounds for the geometry at that
time.

The provided destination pointer must be aligned to 16 bytes. The
function may be called only after committing the scene, otherwise the
result is undefined.

\subsubsection{EXIT STATUS}\label{exit-status}

On failure an error code is set that can be queried using
\texttt{rtcGetDeviceError}.

\subsubsection{SEE ALSO}\label{see-also}

\hyperref[rtcgetscenebounds]{rtcGetSceneBounds},
\hyperref[rtccommitscene]{rtcCommitScene},
\hyperref[rtcjoincommitscene]{rtcJoinCommitScene}

\pagebreak

\section{rtcNewGeometry}\label{rtcnewgeometry}

\subsubsection{NAME}\label{name}

\begin{Shaded}
\begin{Highlighting}[]
\NormalTok{rtcNewGeometry }\OperatorTok{{-}}\NormalTok{ creates a }\KeywordTok{new}\NormalTok{ geometry object}
\end{Highlighting}
\end{Shaded}

\subsubsection{SYNOPSIS}\label{synopsis}

\begin{Shaded}
\begin{Highlighting}[]
\PreprocessorTok{\#include }\ImportTok{\textless{}embree4/rtcore.h\textgreater{}}

\KeywordTok{enum}\NormalTok{ RTCGeometryType}
\OperatorTok{\{}
\NormalTok{ RTC\_GEOMETRY\_TYPE\_TRIANGLE}\OperatorTok{,}
\NormalTok{ RTC\_GEOMETRY\_TYPE\_QUAD}\OperatorTok{,}
\NormalTok{ RTC\_GEOMETRY\_TYPE\_SUBDIVISION}\OperatorTok{,}
\NormalTok{ RTC\_GEOMETRY\_TYPE\_FLAT\_LINEAR\_CURVE}\OperatorTok{,}
\NormalTok{ RTC\_GEOMETRY\_TYPE\_FLAT\_BEZIER\_CURVE}\OperatorTok{,}
\NormalTok{ RTC\_GEOMETRY\_TYPE\_FLAT\_BSPLINE\_CURVE}\OperatorTok{,}
\NormalTok{ RTC\_GEOMETRY\_TYPE\_FLAT\_HERMITE\_CURVE}\OperatorTok{,}
\NormalTok{ RTC\_GEOMETRY\_TYPE\_FLAT\_CATMULL\_ROM\_CURVE}\OperatorTok{,}
\NormalTok{ RTC\_GEOMETRY\_TYPE\_NORMAL\_ORIENTED\_BEZIER\_CURVE}\OperatorTok{,}
\NormalTok{ RTC\_GEOMETRY\_TYPE\_NORMAL\_ORIENTED\_BSPLINE\_CURVE}\OperatorTok{,}
\NormalTok{ RTC\_GEOMETRY\_TYPE\_NORMAL\_ORIENTED\_HERMITE\_CURVE}\OperatorTok{,}
\NormalTok{ RTC\_GEOMETRY\_TYPE\_NORMAL\_ORIENTED\_CATMULL\_ROM\_CURVE}\OperatorTok{,}
\NormalTok{ RTC\_GEOMETRY\_TYPE\_CONE\_LINEAR\_CURVE}\OperatorTok{,}
\NormalTok{ RTC\_GEOMETRY\_TYPE\_ROUND\_LINEAR\_CURVE}\OperatorTok{,}
\NormalTok{ RTC\_GEOMETRY\_TYPE\_ROUND\_BEZIER\_CURVE}\OperatorTok{,}
\NormalTok{ RTC\_GEOMETRY\_TYPE\_ROUND\_BSPLINE\_CURVE}\OperatorTok{,}
\NormalTok{ RTC\_GEOMETRY\_TYPE\_ROUND\_HERMITE\_CURVE}\OperatorTok{,}
\NormalTok{ RTC\_GEOMETRY\_TYPE\_ROUND\_CATMULL\_ROM\_CURVE}\OperatorTok{,}
\NormalTok{ RTC\_GEOMETRY\_TYPE\_GRID}\OperatorTok{,}
\NormalTok{ RTC\_GEOMETRY\_TYPE\_SPHERE\_POINT}\OperatorTok{,}
\NormalTok{ RTC\_GEOMETRY\_TYPE\_DISC\_POINT}\OperatorTok{,}
\NormalTok{ RTC\_GEOMETRY\_TYPE\_ORIENTED\_DISC\_POINT}\OperatorTok{,}
\NormalTok{ RTC\_GEOMETRY\_TYPE\_USER}\OperatorTok{,}
\NormalTok{ RTC\_GEOMETRY\_TYPE\_INSTANCE}
\OperatorTok{\};}

\NormalTok{RTCGeometry rtcNewGeometry}\OperatorTok{(}
\NormalTok{  RTCDevice device}\OperatorTok{,}
  \KeywordTok{enum}\NormalTok{ RTCGeometryType type}
\OperatorTok{);}
\end{Highlighting}
\end{Shaded}

\subsubsection{DESCRIPTION}\label{description}

Geometries are objects that represent an array of primitives of the same
type. The \texttt{rtcNewGeometry} function creates a new geometry of
specified type (\texttt{type} argument) bound to the specified device
(\texttt{device} argument) and returns a handle to this geometry. The
geometry object is reference counted with an initial reference count of
1. The geometry handle can be released using the
\texttt{rtcReleaseGeometry} API call.

Supported geometry types are triangle meshes
(\texttt{RTC\_GEOMETRY\_TYPE\_TRIANGLE} type), quad meshes (triangle
pairs) (\texttt{RTC\_GEOMETRY\_TYPE\_QUAD} type), Catmull-Clark
subdivision surfaces (\texttt{RTC\_GEOMETRY\_TYPE\_SUBDIVISION} type),
curve geometries with different bases
(\texttt{RTC\_GEOMETRY\_TYPE\_FLAT\_LINEAR\_CURVE},
\texttt{RTC\_GEOMETRY\_TYPE\_FLAT\_BEZIER\_CURVE},\\
\texttt{RTC\_GEOMETRY\_TYPE\_FLAT\_BSPLINE\_CURVE},
\texttt{RTC\_GEOMETRY\_TYPE\_FLAT\_HERMITE\_CURVE},\\
\texttt{RTC\_GEOMETRY\_TYPE\_FLAT\_CATMULL\_ROM\_CURVE},
\texttt{RTC\_GEOMETRY\_TYPE\_NORMAL\_ORIENTED\_BEZIER\_CURVE},
\texttt{RTC\_GEOMETRY\_TYPE\_NORMAL\_ORIENTED\_BSPLINE\_CURVE},
\texttt{RTC\_GEOMETRY\_TYPE\_NORMAL\_ORIENTED\_HERMITE\_CURVE},
\texttt{RTC\_GEOMETRY\_TYPE\_NORMAL\_ORIENTED\_CATMULL\_ROM\_CURVE},
\texttt{RTC\_GEOMETRY\_TYPE\_CONE\_LINEAR\_CURVE},
\texttt{RTC\_GEOMETRY\_TYPE\_ROUND\_LINEAR\_CURVE},
\texttt{RTC\_GEOMETRY\_TYPE\_ROUND\_BEZIER\_CURVE},
\texttt{RTC\_GEOMETRY\_TYPE\_ROUND\_BSPLINE\_CURVE},
\texttt{RTC\_GEOMETRY\_TYPE\_ROUND\_HERMITE\_CURVE},
\texttt{RTC\_GEOMETRY\_TYPE\_ROUND\_CATMULL\_ROM\_CURVE} types) grid
meshes (\texttt{RTC\_GEOMETRY\_TYPE\_GRID}), point geometries
(\texttt{RTC\_GEOMETRY\_TYPE\_SPHERE\_POINT},
\texttt{RTC\_GEOMETRY\_TYPE\_DISC\_POINT},
\texttt{RTC\_TYPE\_ORIENTED\_DISC\_POINT}), user-defined geometries
(\texttt{RTC\_GEOMETRY\_TYPE\_USER}), and instances
(\texttt{RTC\_GEOMETRY\_TYPE\_INSTANCE}).

The types \texttt{RTC\_GEOMETRY\_TYPE\_ROUND\_BEZIER\_CURVE},
\texttt{RTC\_GEOMETRY\_TYPE\_ROUND\_BSPLINE\_CURVE}, and
\texttt{RTC\_GEOMETRY\_TYPE\_ROUND\_CATMULL\_ROM\_CURVE} will treat the
curve as a sweep surface of a varying-radius circle swept tangentially
along the curve. The types
\texttt{RTC\_GEOMETRY\_TYPE\_FLAT\_BEZIER\_CURVE},
\texttt{RTC\_GEOMETRY\_TYPE\_FLAT\_BSPLINE\_CURVE}, and
\texttt{RTC\_GEOMETRY\_TYPE\_FLAT\_CATMULL\_ROM\_CURVE} use ray-facing
ribbons as a faster-to-intersect approximation.

After construction, geometries are enabled by default and not attached
to any scene. Geometries can be disabled (\texttt{rtcDisableGeometry}
call), and enabled again (\texttt{rtcEnableGeometry} call). A geometry
can be attached to multiple scenes using the \texttt{rtcAttachGeometry}
call (or \texttt{rtcAttachGeometryByID} call), and detached using the
\texttt{rtcDetachGeometry} call. During attachment, a geometry ID is
assigned to the geometry (or assigned by the user when using the
\texttt{rtcAttachGeometryByID} call), which uniquely identifies the
geometry inside that scene. This identifier is returned when primitives
of the geometry are hit in later ray queries for the scene.

Geometries can also be modified, including their vertex and index
buffers. After modifying a buffer, \texttt{rtcUpdateGeometryBuffer} must
be called to notify that the buffer got modified.

The application can use the \texttt{rtcSetGeometryUserData} function to
set a user data pointer to its own geometry representation, and later
read out this pointer using the \texttt{rtcGetGeometryUserData}
function.

After setting up the geometry or modifying it,
\texttt{rtcCommitGeometry} must be called to finish the geometry setup.
After committing the geometry, vertex data interpolation can be
performed using the \texttt{rtcInterpolate} and \texttt{rtcInterpolateN}
functions.

A build quality can be specified for a geometry using the
\texttt{rtcSetGeometryBuildQuality} function, to balance between
acceleration structure build performance and ray query performance. The
build quality per geometry will be used if a two-level acceleration
structure is built internally, which is the case if the
\texttt{RTC\_BUILD\_QUALITY\_LOW} is set as the scene build quality. See
Section \hyperref[rtcsetscenebuildquality]{rtcSetSceneBuildQuality} for
more details.

\subsubsection{EXIT STATUS}\label{exit-status}

On failure \texttt{NULL} is returned and an error code is set that can
be queried using \texttt{rtcGetDeviceError}.

\subsubsection{SEE ALSO}\label{see-also}

\hyperref[rtcenablegeometry]{rtcEnableGeometry},
\hyperref[rtcdisablegeometry]{rtcDisableGeometry},
\hyperref[rtcattachgeometry]{rtcAttachGeometry},
\hyperref[rtcattachgeometrybyid]{rtcAttachGeometryByID},
\hyperref[rtcupdategeometrybuffer]{rtcUpdateGeometryBuffer},
\hyperref[rtcsetgeometryuserdata]{rtcSetGeometryUserData},
\hyperref[rtcgetgeometryuserdata]{rtcGetGeometryUserData},
\hyperref[rtccommitgeometry]{rtcCommitGeometry},
\hyperref[rtcinterpolate]{rtcInterpolate},
\hyperref[rtcinterpolaten]{rtcInterpolateN},
\hyperref[rtcsetgeometrybuildquality]{rtcSetGeometryBuildQuality},
\hyperref[rtcsetscenebuildquality]{rtcSetSceneBuildQuality},
\hyperref[rtc_geometry_type_triangle]{RTC\_GEOMETRY\_TYPE\_TRIANGLE},
\hyperref[rtc_geometry_type_quad]{RTC\_GEOMETRY\_TYPE\_QUAD},
\hyperref[rtc_geometry_type_subdivision]{RTC\_GEOMETRY\_TYPE\_SUBDIVISION},
\hyperref[rtc_geometry_type_curve]{RTC\_GEOMETRY\_TYPE\_CURVE},
\hyperref[rtc_geometry_type_grid]{RTC\_GEOMETRY\_TYPE\_GRID},
\hyperref[rtc_geometry_type_point]{RTC\_GEOMETRY\_TYPE\_POINT},
\hyperref[rtc_geometry_type_user]{RTC\_GEOMETRY\_TYPE\_USER},
\hyperref[rtc_geometry_type_instance]{RTC\_GEOMETRY\_TYPE\_INSTANCE}

\pagebreak

\section{RTC\_GEOMETRY\_TYPE\_TRIANGLE}\label{rtc_geometry_type_triangle}

\subsubsection{NAME}\label{name}

\begin{Shaded}
\begin{Highlighting}[]
\NormalTok{RTC\_GEOMETRY\_TYPE\_TRIANGLE }\OperatorTok{{-}}\NormalTok{ triangle geometry type}
\end{Highlighting}
\end{Shaded}

\subsubsection{SYNOPSIS}\label{synopsis}

\begin{Shaded}
\begin{Highlighting}[]
\PreprocessorTok{\#include }\ImportTok{\textless{}embree4/rtcore.h\textgreater{}}

\NormalTok{RTCGeometry geometry }\OperatorTok{=}
\NormalTok{  rtcNewGeometry}\OperatorTok{(}\NormalTok{device}\OperatorTok{,}\NormalTok{ RTC\_GEOMETRY\_TYPE\_TRIANGLE}\OperatorTok{);}
\end{Highlighting}
\end{Shaded}

\subsubsection{DESCRIPTION}\label{description}

Triangle meshes are created by passing
\texttt{RTC\_GEOMETRY\_TYPE\_TRIANGLE} to the \texttt{rtcNewGeometry}
function call. The triangle indices can be specified by setting an index
buffer (\texttt{RTC\_BUFFER\_TYPE\_INDEX} type) and the triangle
vertices by setting a vertex buffer (\texttt{RTC\_BUFFER\_TYPE\_VERTEX}
type). See \texttt{rtcSetGeometryBuffer} and
\texttt{rtcSetSharedGeometryBuffer} for more details on how to set
buffers. The index buffer must contain an array of three 32-bit indices
per triangle (\texttt{RTC\_FORMAT\_UINT3} format) and the number of
primitives is inferred from the size of that buffer. The vertex buffer
must contain an array of single precision \texttt{x}, \texttt{y},
\texttt{z} floating point coordinates (\texttt{RTC\_FORMAT\_FLOAT3}
format), and the number of vertices are inferred from the size of that
buffer. The vertex buffer can be at most 16 GB large.

The parametrization of a triangle uses the first vertex \texttt{p0} as
base point, the vector \texttt{p1\ -\ p0} as u-direction and the vector
\texttt{p2\ -\ p0} as v-direction. Thus vertex attributes
\texttt{t0,t1,t2} can be linearly interpolated over the triangle the
following way:

\begin{Shaded}
\begin{Highlighting}[]
\NormalTok{t\_uv }\OperatorTok{=} \OperatorTok{(}\DecValTok{1}\OperatorTok{{-}}\NormalTok{u}\OperatorTok{{-}}\NormalTok{v}\OperatorTok{)*}\NormalTok{t0 }\OperatorTok{+}\NormalTok{ u}\OperatorTok{*}\NormalTok{t1 }\OperatorTok{+}\NormalTok{ v}\OperatorTok{*}\NormalTok{t2}
     \OperatorTok{=}\NormalTok{ t0 }\OperatorTok{+}\NormalTok{ u}\OperatorTok{*(}\NormalTok{t1}\OperatorTok{{-}}\NormalTok{t0}\OperatorTok{)} \OperatorTok{+}\NormalTok{ v}\OperatorTok{*(}\NormalTok{t2}\OperatorTok{{-}}\NormalTok{t0}\OperatorTok{)}
\end{Highlighting}
\end{Shaded}

A triangle whose vertices are laid out counter-clockwise has its
geometry normal pointing upwards outside the front face, like
illustrated in the following picture:

!\href{images/triangle_uv.pdf}{}

For multi-segment motion blur, the number of time steps must be first
specified using the \texttt{rtcSetGeometryTimeStepCount} call. Then a
vertex buffer for each time step can be set using different buffer
slots, and all these buffers have to have the same stride and size.

Also see tutorial \hyperref[triangle-geometry]{Triangle Geometry} for an
example of how to create triangle meshes.

\subsubsection{EXIT STATUS}\label{exit-status}

On failure \texttt{NULL} is returned and an error code is set that be
get queried using \texttt{rtcGetDeviceError}.

\subsubsection{SEE ALSO}\label{see-also}

\hyperref[rtcnewgeometry]{rtcNewGeometry}

\pagebreak

\section{RTC\_GEOMETRY\_TYPE\_QUAD}\label{rtc_geometry_type_quad}

\subsubsection{NAME}\label{name}

\begin{Shaded}
\begin{Highlighting}[]
\NormalTok{RTC\_GEOMETRY\_TYPE\_QUAD }\OperatorTok{{-}}\NormalTok{ quad geometry type}
\end{Highlighting}
\end{Shaded}

\subsubsection{SYNOPSIS}\label{synopsis}

\begin{Shaded}
\begin{Highlighting}[]
\PreprocessorTok{\#include }\ImportTok{\textless{}embree4/rtcore.h\textgreater{}}

\NormalTok{RTCGeometry geometry }\OperatorTok{=}
\NormalTok{  rtcNewGeometry}\OperatorTok{(}\NormalTok{device}\OperatorTok{,}\NormalTok{ RTC\_GEOMETRY\_TYPE\_QUAD}\OperatorTok{);}
\end{Highlighting}
\end{Shaded}

\subsubsection{DESCRIPTION}\label{description}

Quad meshes are created by passing \texttt{RTC\_GEOMETRY\_TYPE\_QUAD} to
the \texttt{rtcNewGeometry} function call. The quad indices can be
specified by setting an index buffer (\texttt{RTC\_BUFFER\_TYPE\_INDEX}
type) and the quad vertices by setting a vertex buffer
(\texttt{RTC\_BUFFER\_TYPE\_VERTEX} type). See
\texttt{rtcSetGeometryBuffer} and \texttt{rtcSetSharedGeometryBuffer}
for more details on how to set buffers. The index buffer contains an
array of four 32-bit indices per quad (\texttt{RTC\_FORMAT\_UINT4}
format), and the number of primitives is inferred from the size of that
buffer. The vertex buffer contains an array of single precision
\texttt{x}, \texttt{y}, \texttt{z} floating point coordinates
(\texttt{RTC\_FORMAT\_FLOAT3} format), and the number of vertices is
inferred from the size of that buffer. The vertex buffer can be at most
16 GB large.

A quad is internally handled as a pair of two triangles
\texttt{v0,v1,v3} and \texttt{v2,v3,v1}, with the
\texttt{u\textquotesingle{}}/\texttt{v\textquotesingle{}} coordinates of
the second triangle corrected by \texttt{u\ =\ 1-u\textquotesingle{}}
and \texttt{v\ =\ 1-v\textquotesingle{}} to produce a quad
parametrization where \texttt{u} and \texttt{v} are in the range 0 to 1.
Thus the parametrization of a quad uses the first vertex \texttt{p0} as
base point, and the vector \texttt{p1\ -\ p0} as \texttt{u}-direction,
and \texttt{p3\ -\ p0} as v-direction. Thus vertex attributes
\texttt{t0,t1,t2,t3} can be bilinearly interpolated over the
quadrilateral the following way:

\begin{Shaded}
\begin{Highlighting}[]
\NormalTok{t\_uv }\OperatorTok{=} \OperatorTok{(}\DecValTok{1}\OperatorTok{{-}}\NormalTok{v}\OperatorTok{)((}\DecValTok{1}\OperatorTok{{-}}\NormalTok{u}\OperatorTok{)*}\NormalTok{t0 }\OperatorTok{+}\NormalTok{ u}\OperatorTok{*}\NormalTok{t1}\OperatorTok{)} \OperatorTok{+}\NormalTok{ v}\OperatorTok{*((}\DecValTok{1}\OperatorTok{{-}}\NormalTok{u}\OperatorTok{)*}\NormalTok{t3 }\OperatorTok{+}\NormalTok{ u}\OperatorTok{*}\NormalTok{t2}\OperatorTok{)}
\end{Highlighting}
\end{Shaded}

Mixed triangle/quad meshes are supported by encoding a triangle as a
quad, which can be achieved by replicating the last triangle vertex
(\texttt{v0,v1,v2} -\textgreater{} \texttt{v0,v1,v2,v2}). This way the
second triangle is a line (which can never get hit), and the
parametrization of the first triangle is compatible with the standard
triangle parametrization.

A quad whose vertices are laid out counter-clockwise has its geometry
normal pointing upwards outside the front face, like illustrated in the
following picture.

!\href{images/quad_uv.pdf}{}

For multi-segment motion blur, the number of time steps must be first
specified using the \texttt{rtcSetGeometryTimeStepCount} call. Then a
vertex buffer for each time step can be set using different buffer
slots, and all these buffers must have the same stride and size.

\subsubsection{EXIT STATUS}\label{exit-status}

On failure \texttt{NULL} is returned and an error code is set that can
be queried using \texttt{rtcGetDeviceError}.

\subsubsection{SEE ALSO}\label{see-also}

\hyperref[rtcnewgeometry]{rtcNewGeometry}

\pagebreak

\section{RTC\_GEOMETRY\_TYPE\_GRID}\label{rtc_geometry_type_grid}

\subsubsection{NAME}\label{name}

\begin{Shaded}
\begin{Highlighting}[]
\NormalTok{RTC\_GEOMETRY\_TYPE\_GRID }\OperatorTok{{-}}\NormalTok{ grid geometry type}
\end{Highlighting}
\end{Shaded}

\subsubsection{SYNOPSIS}\label{synopsis}

\begin{Shaded}
\begin{Highlighting}[]
\PreprocessorTok{\#include }\ImportTok{\textless{}embree4/rtcore.h\textgreater{}}

\NormalTok{RTCGeometry geometry }\OperatorTok{=}
\NormalTok{  rtcNewGeometry}\OperatorTok{(}\NormalTok{device}\OperatorTok{,}\NormalTok{ RTC\_GEOMETRY\_TYPE\_GRID}\OperatorTok{);}
\end{Highlighting}
\end{Shaded}

\subsubsection{DESCRIPTION}\label{description}

Grid meshes are created by passing \texttt{RTC\_GEOMETRY\_TYPE\_GRID} to
the \texttt{rtcNewGeometry} function call, and contain an array of grid
primitives. This array of grids can be specified by setting up a grid
buffer (with \texttt{RTC\_BUFFER\_TYPE\_GRID} type and
\texttt{RTC\_FORMAT\_GRID} format) and the grid mesh vertices by setting
a vertex buffer (\texttt{RTC\_BUFFER\_TYPE\_VERTEX} type). See
\texttt{rtcSetGeometryBuffer} and \texttt{rtcSetSharedGeometryBuffer}
for more details on how to set buffers. The number of grid primitives in
the grid mesh is inferred from the size of the grid buffer.

The vertex buffer contains an array of single precision \texttt{x},
\texttt{y}, \texttt{z} floating point coordinates
(\texttt{RTC\_FORMAT\_FLOAT3} format), and the number of vertices is
inferred from the size of that buffer.

Each grid in the grid buffer is of the type \texttt{RTCGrid}:

\begin{Shaded}
\begin{Highlighting}[]
\KeywordTok{struct}\NormalTok{ RTCGrid}
\OperatorTok{\{}
  \DataTypeTok{unsigned} \DataTypeTok{int}\NormalTok{ startVertexID}\OperatorTok{;}
  \DataTypeTok{unsigned} \DataTypeTok{int}\NormalTok{ stride}\OperatorTok{;}
  \DataTypeTok{unsigned} \DataTypeTok{short}\NormalTok{ width}\OperatorTok{,}\NormalTok{height}\OperatorTok{;} 
\OperatorTok{\};}
\end{Highlighting}
\end{Shaded}

The \texttt{RTCGrid} structure describes a 2D grid of vertices (with
respect to the vertex buffer of the grid mesh). The \texttt{width} and
\texttt{height} members specify the number of vertices in u and v
direction, e.g.~setting both \texttt{width} and \texttt{height} to 3
sets up a 3×3 vertex grid. The maximum allowed \texttt{width} and
\texttt{height} is 32767. The \texttt{startVertexID} specifies the ID of
the top-left vertex in the vertex grid, while the \texttt{stride}
parameter specifies a stride (in number of vertices) used to step to the
next row.

A vertex grid of dimensions \texttt{width} and \texttt{height} is
treated as a \texttt{(width-1)} x \texttt{(height-1)} grid of
\texttt{quads} (triangle-pairs), with the same shared edge handling as
for regular quad meshes. However, the \texttt{u}/\texttt{v} coordinates
have the uniform range \texttt{{[}0..1{]}} for an entire vertex grid.
The \texttt{u} direction follows the \texttt{width} of the grid while
the \texttt{v} direction the \texttt{height}.

For multi-segment motion blur, the number of time steps must be first
specified using the \texttt{rtcSetGeometryTimeStepCount} call. Then a
vertex buffer for each time step can be set using different buffer
slots, and all these buffers must have the same stride and size.

\subsubsection{EXIT STATUS}\label{exit-status}

On failure \texttt{NULL} is returned and an error code is set that can
be queried using \texttt{rtcGetDeviceError}.

\subsubsection{SEE ALSO}\label{see-also}

\hyperref[rtcnewgeometry]{rtcNewGeometry}

\pagebreak

\section{RTC\_GEOMETRY\_TYPE\_SUBDIVISION}\label{rtc_geometry_type_subdivision}

\subsubsection{NAME}\label{name}

\begin{Shaded}
\begin{Highlighting}[]
\NormalTok{RTC\_GEOMETRY\_TYPE\_SUBDIVISION }\OperatorTok{{-}}\NormalTok{ subdivision geometry type}
\end{Highlighting}
\end{Shaded}

\subsubsection{SYNOPSIS}\label{synopsis}

\begin{Shaded}
\begin{Highlighting}[]
\PreprocessorTok{\#include }\ImportTok{\textless{}embree4/rtcore.h\textgreater{}}

\NormalTok{RTCGeometry geometry }\OperatorTok{=}
\NormalTok{  rtcNewGeometry}\OperatorTok{(}\NormalTok{device}\OperatorTok{,}\NormalTok{ RTC\_GEOMETRY\_TYPE\_SUBDIVISION}\OperatorTok{);}
\end{Highlighting}
\end{Shaded}

\subsubsection{DESCRIPTION}\label{description}

Catmull-Clark subdivision meshes are supported, including support for
edge creases, vertex creases, holes, non-manifold geometry, and
face-varying interpolation. The number of vertices per face can be in
the range of 3 to 15 vertices (triangles, quadrilateral, pentagons,
etc).

Subdivision meshes are created by passing
\texttt{RTC\_GEOMETRY\_TYPE\_SUBDIVISION} to the \texttt{rtcNewGeometry}
function. Various buffers need to be set by the application to set up
the subdivision mesh. See \texttt{rtcSetGeometryBuffer} and
\texttt{rtcSetSharedGeometryBuffer} for more details on how to set
buffers. The face buffer (\texttt{RTC\_BUFFER\_TYPE\_FACE} type and
\texttt{RTC\_FORMAT\_UINT} format) contains the number of edges/indices
of each face (3 to 15), and the number of faces is inferred from the
size of this buffer. The index buffer (\texttt{RTC\_BUFFER\_TYPE\_INDEX}
type) contains multiple (3 to 15) 32-bit vertex indices
(\texttt{RTC\_FORMAT\_UINT} format) for each face, and the number of
edges is inferred from the size of this buffer. The vertex buffer
(\texttt{RTC\_BUFFER\_TYPE\_VERTEX} type) stores an array of single
precision \texttt{x}, \texttt{y}, \texttt{z} floating point coordinates
(\texttt{RTC\_FORMAT\_FLOAT3} format), and the number of vertices is
inferred from the size of this buffer.

Optionally, the application may set additional index buffers using
different buffer slots if multiple topologies are required for
face-varying interpolation. The standard vertex buffers
(\texttt{RTC\_BUFFER\_TYPE\_VERTEX}) are always bound to the geometry
topology (topology 0) thus use \texttt{RTC\_BUFFER\_TYPE\_INDEX} with
buffer slot 0. User vertex data interpolation may use different
topologies as described later.

Optionally, the application can set up the hole buffer
(\texttt{RTC\_BUFFER\_TYPE\_HOLE}) which contains an array of 32-bit
indices (\texttt{RTC\_FORMAT\_UINT} format) of faces that should be
considered non-existing in all topologies. The number of holes is
inferred from the size of this buffer.

Optionally, the application can fill the level buffer
(\texttt{RTC\_BUFFER\_TYPE\_LEVEL}) with a tessellation rate for each of
the edges of each face. This buffer must have the same size as the index
buffer. The tessellation level is a positive floating point value
(\texttt{RTC\_FORMAT\_FLOAT} format) that specifies how many quads along
the edge should be generated during tessellation. If no level buffer is
specified, a level of 1 is used. The maximally supported edge level is
4096, and larger levels are clamped to that value. Note that edges may
be shared between (typically 2) faces. To guarantee a watertight
tessellation, the level of these shared edges should be identical. A
uniform tessellation rate for an entire subdivision mesh can be set by
using the \texttt{rtcSetGeometryTessellationRate} function. The
existence of a level buffer has precedence over the uniform tessellation
rate.

Optionally, the application can fill the sparse edge crease buffers to
make edges appear sharper. The edge crease index buffer
(\texttt{RTC\_BUFFER\_TYPE\_EDGE\_CREASE\_INDEX}) contains an array of
pairs of 32-bit vertex indices (\texttt{RTC\_FORMAT\_UINT2} format) that
specify unoriented edges in the geometry topology. The edge crease
weight buffer (\texttt{RTC\_BUFFER\_TYPE\_EDGE\_CREASE\_WEIGHT}) stores
for each of these crease edges a positive floating point weight
(\texttt{RTC\_FORMAT\_FLOAT} format). The number of edge creases is
inferred from the size of these buffers, which has to be identical. The
larger a weight, the sharper the edge. Specifying a weight of infinity
is supported and marks an edge as infinitely sharp. Storing an edge
multiple times with the same crease weight is allowed, but has lower
performance. Storing an edge multiple times with different crease
weights results in undefined behavior. For a stored edge (i,j), the
reverse direction edges (j,i) do not have to be stored, as both are
considered the same unoriented edge. Edge crease features are shared
between all topologies.

Optionally, the application can fill the sparse vertex crease buffers to
make vertices appear sharper. The vertex crease index buffer
(\texttt{RTC\_BUFFER\_TYPE\_VERTEX\_CREASE\_INDEX}), contains an array
of 32-bit vertex indices (\texttt{RTC\_FORMAT\_UINT} format) to specify
a set of vertices from the geometry topology. The vertex crease weight
buffer (\texttt{RTC\_BUFFER\_TYPE\_VERTEX\_CREASE\_WEIGHT}) specifies
for each of these vertices a positive floating point weight
(\texttt{RTC\_FORMAT\_FLOAT} format). The number of vertex creases is
inferred from the size of these buffers, and has to be identical. The
larger a weight, the sharper the vertex. Specifying a weight of infinity
is supported and makes the vertex infinitely sharp. Storing a vertex
multiple times with the same crease weight is allowed, but has lower
performance. Storing a vertex multiple times with different crease
weights results in undefined behavior. Vertex crease features are shared
between all topologies.

Subdivision modes can be used to force linear interpolation for parts of
the subdivision mesh; see \texttt{rtcSetGeometrySubdivisionMode} for
more details.

For multi-segment motion blur, the number of time steps must be first
specified using the \texttt{rtcSetGeometryTimeStepCount} call. Then a
vertex buffer for each time step can be set using different buffer
slots, and all these buffers have to have the same stride and size.

Also see tutorial \hyperref[subdivision-geometry]{Subdivision Geometry}
for an example of how to create subdivision surfaces.

\subsubsection{Parametrization}\label{parametrization}

The parametrization for subdivision faces is different for
quadrilaterals and non-quadrilateral faces.

The parametrization of a quadrilateral face uses the first vertex
\texttt{p0} as base point, and the vector \texttt{p1\ -\ p0} as
u-direction and \texttt{p3\ -\ p0} as v-direction.

The parametrization for all other face types (with number of vertices
not equal 4), have a special parametrization where the subpatch ID
\texttt{n} (of the \texttt{n}-th quadrilateral that would be obtained by
a single subdivision step) and the local hit location inside this
quadrilateral are encoded in the UV coordinates. The following code
extracts the sub-patch ID \texttt{i} and local UVs of this subpatch:

\begin{Shaded}
\begin{Highlighting}[]
\DataTypeTok{unsigned} \DataTypeTok{int}\NormalTok{ l }\OperatorTok{=}\NormalTok{ floorf}\OperatorTok{(}\FloatTok{0.5}\BuiltInTok{f}\OperatorTok{*}\NormalTok{U}\OperatorTok{);}
\DataTypeTok{unsigned} \DataTypeTok{int}\NormalTok{ h }\OperatorTok{=}\NormalTok{ floorf}\OperatorTok{(}\FloatTok{0.5}\BuiltInTok{f}\OperatorTok{*}\NormalTok{V}\OperatorTok{);}
\DataTypeTok{unsigned} \DataTypeTok{int}\NormalTok{ i }\OperatorTok{=} \DecValTok{4}\OperatorTok{*}\NormalTok{h}\OperatorTok{+}\NormalTok{l}\OperatorTok{;}
\DataTypeTok{float}\NormalTok{ u }\OperatorTok{=} \FloatTok{2.0}\BuiltInTok{f}\OperatorTok{*}\NormalTok{fracf}\OperatorTok{(}\FloatTok{0.5}\BuiltInTok{f}\OperatorTok{*}\NormalTok{U}\OperatorTok{){-}}\FloatTok{0.5}\BuiltInTok{f}\OperatorTok{;}
\DataTypeTok{float}\NormalTok{ v }\OperatorTok{=} \FloatTok{2.0}\BuiltInTok{f}\OperatorTok{*}\NormalTok{fracf}\OperatorTok{(}\FloatTok{0.5}\BuiltInTok{f}\OperatorTok{*}\NormalTok{V}\OperatorTok{){-}}\FloatTok{0.5}\BuiltInTok{f}\OperatorTok{;}
\end{Highlighting}
\end{Shaded}

This encoding allows local subpatch UVs to be in the range
\texttt{{[}-0.5,1.5{[}} thus negative subpatch UVs can be passed to
\texttt{rtcInterpolate} to sample subpatches slightly out of bounds.
This can be useful to calculate derivatives using finite differences if
required. The encoding further has the property that one can just move
the value \texttt{u} (or \texttt{v}) on a subpatch by adding \texttt{du}
(or \texttt{dv}) to the special UV encoding as long as it does not fall
out of the \texttt{{[}-0.5,1.5{[}} range.

To smoothly interpolate vertex attributes over the subdivision surface
we recommend using the \texttt{rtcInterpolate} function, which will
apply the standard subdivision rules for interpolation and automatically
takes care of the special UV encoding for non-quadrilaterals.

\subsubsection{Face-Varying Data}\label{face-varying-data}

Face-varying interpolation is supported through multiple topologies per
subdivision mesh and binding such topologies to vertex attribute buffers
to interpolate. This way, texture coordinates may use a different
topology with additional boundaries to construct separate UV regions
inside one subdivision mesh.

Each such topology \texttt{i} has a separate index buffer (specified
using \texttt{RTC\_BUFFER\_TYPE\_INDEX} with buffer slot \texttt{i}) and
separate subdivision mode that can be set using
\texttt{rtcSetGeometrySubdivisionMode}. A vertex attribute buffer
\texttt{RTC\_BUFFER\_TYPE\_VERTEX\_ATTRIBUTE} bound to a buffer slot
\texttt{j} can be assigned to use a topology for interpolation using the
\texttt{rtcSetGeometryVertexAttributeTopology} call.

The face buffer (\texttt{RTC\_BUFFER\_TYPE\_FACE} type) is shared
between all topologies, which means that the \texttt{n}-th primitive
always has the same number of vertices (e.g.~being a triangle or a quad)
for each topology. However, the indices of the topologies themselves may
be different.

\subsubsection{EXIT STATUS}\label{exit-status}

On failure \texttt{NULL} is returned and an error code is set that can
be queried using \texttt{rtcGetDeviceError}.

\subsubsection{SEE ALSO}\label{see-also}

\hyperref[rtcnewgeometry]{rtcNewGeometry}

\pagebreak

\section{RTC\_GEOMETRY\_TYPE\_CURVE}\label{rtc_geometry_type_curve}

\subsubsection{NAME}\label{name}

\begin{Shaded}
\begin{Highlighting}[]
\NormalTok{RTC\_GEOMETRY\_TYPE\_FLAT\_LINEAR\_CURVE }\OperatorTok{{-}}
\NormalTok{  flat curve geometry with linear basis}

\NormalTok{RTC\_GEOMETRY\_TYPE\_FLAT\_BEZIER\_CURVE }\OperatorTok{{-}}
\NormalTok{  flat curve geometry with cubic Bézier basis}

\NormalTok{RTC\_GEOMETRY\_TYPE\_FLAT\_BSPLINE\_CURVE }\OperatorTok{{-}} 
\NormalTok{  flat curve geometry with cubic B}\OperatorTok{{-}}\NormalTok{spline basis}

\NormalTok{RTC\_GEOMETRY\_TYPE\_FLAT\_HERMITE\_CURVE }\OperatorTok{{-}} 
\NormalTok{  flat curve geometry with cubic Hermite basis}

\NormalTok{RTC\_GEOMETRY\_TYPE\_FLAT\_CATMULL\_ROM\_CURVE }\OperatorTok{{-}} 
\NormalTok{  flat curve geometry with Catmull}\OperatorTok{{-}}\NormalTok{Rom basis}

\NormalTok{RTC\_GEOMETRY\_TYPE\_NORMAL\_ORIENTED\_BEZIER\_CURVE }\OperatorTok{{-}}
\NormalTok{  flat normal oriented curve geometry with cubic Bézier basis}

\NormalTok{RTC\_GEOMETRY\_TYPE\_NORMAL\_ORIENTED\_BSPLINE\_CURVE }\OperatorTok{{-}} 
\NormalTok{  flat normal oriented curve geometry with cubic B}\OperatorTok{{-}}\NormalTok{spline basis}

\NormalTok{RTC\_GEOMETRY\_TYPE\_NORMAL\_ORIENTED\_HERMITE\_CURVE }\OperatorTok{{-}} 
\NormalTok{  flat normal oriented curve geometry with cubic Hermite basis}

\NormalTok{RTC\_GEOMETRY\_TYPE\_NORMAL\_ORIENTED\_CATMULL\_ROM\_CURVE }\OperatorTok{{-}} 
\NormalTok{  flat normal oriented curve geometry with Catmull}\OperatorTok{{-}}\NormalTok{Rom basis}

\NormalTok{RTC\_GEOMETRY\_TYPE\_CONE\_LINEAR\_CURVE }\OperatorTok{{-}}
\NormalTok{  capped cone curve geometry with linear basis }\OperatorTok{{-}}\NormalTok{ discontinuous at edge boundaries}

\NormalTok{RTC\_GEOMETRY\_TYPE\_ROUND\_LINEAR\_CURVE }\OperatorTok{{-}}
\NormalTok{  capped cone curve geometry with linear basis }\KeywordTok{and}\NormalTok{ spherical ending}

\NormalTok{RTC\_GEOMETRY\_TYPE\_ROUND\_BEZIER\_CURVE }\OperatorTok{{-}}
\NormalTok{  swept surface curve geometry with cubic Bézier basis}

\NormalTok{RTC\_GEOMETRY\_TYPE\_ROUND\_BSPLINE\_CURVE }\OperatorTok{{-}}
\NormalTok{  swept surface curve geometry with cubic B}\OperatorTok{{-}}\NormalTok{spline basis}

\NormalTok{RTC\_GEOMETRY\_TYPE\_ROUND\_HERMITE\_CURVE }\OperatorTok{{-}}
\NormalTok{  swept surface curve geometry with cubic Hermite basis}

\NormalTok{RTC\_GEOMETRY\_TYPE\_ROUND\_CATMULL\_ROM\_CURVE }\OperatorTok{{-}}
\NormalTok{  swept surface curve geometry with Catmull}\OperatorTok{{-}}\NormalTok{Rom basis}
\end{Highlighting}
\end{Shaded}

\subsubsection{SYNOPSIS}\label{synopsis}

\begin{Shaded}
\begin{Highlighting}[]
\PreprocessorTok{\#include }\ImportTok{\textless{}embree4/rtcore.h\textgreater{}}

\NormalTok{rtcNewGeometry}\OperatorTok{(}\NormalTok{device}\OperatorTok{,}\NormalTok{ RTC\_GEOMETRY\_TYPE\_FLAT\_LINEAR\_CURVE}\OperatorTok{);}
\NormalTok{rtcNewGeometry}\OperatorTok{(}\NormalTok{device}\OperatorTok{,}\NormalTok{ RTC\_GEOMETRY\_TYPE\_FLAT\_BEZIER\_CURVE}\OperatorTok{);}
\NormalTok{rtcNewGeometry}\OperatorTok{(}\NormalTok{device}\OperatorTok{,}\NormalTok{ RTC\_GEOMETRY\_TYPE\_FLAT\_BSPLINE\_CURVE}\OperatorTok{);}
\NormalTok{rtcNewGeometry}\OperatorTok{(}\NormalTok{device}\OperatorTok{,}\NormalTok{ RTC\_GEOMETRY\_TYPE\_FLAT\_HERMITE\_CURVE}\OperatorTok{);}
\NormalTok{rtcNewGeometry}\OperatorTok{(}\NormalTok{device}\OperatorTok{,}\NormalTok{ RTC\_GEOMETRY\_TYPE\_FLAT\_CATMULL\_ROM\_CURVE}\OperatorTok{);}
\NormalTok{rtcNewGeometry}\OperatorTok{(}\NormalTok{device}\OperatorTok{,}\NormalTok{ RTC\_GEOMETRY\_TYPE\_NORMAL\_ORIENTED\_BEZIER\_CURVE}\OperatorTok{);}
\NormalTok{rtcNewGeometry}\OperatorTok{(}\NormalTok{device}\OperatorTok{,}\NormalTok{ RTC\_GEOMETRY\_TYPE\_NORMAL\_ORIENTED\_BSPLINE\_CURVE}\OperatorTok{);}
\NormalTok{rtcNewGeometry}\OperatorTok{(}\NormalTok{device}\OperatorTok{,}\NormalTok{ RTC\_GEOMETRY\_TYPE\_NORMAL\_ORIENTED\_HERMITE\_CURVE}\OperatorTok{);}
\NormalTok{rtcNewGeometry}\OperatorTok{(}\NormalTok{device}\OperatorTok{,}\NormalTok{ RTC\_GEOMETRY\_TYPE\_NORMAL\_ORIENTED\_CATMULL\_ROM\_CURVE}\OperatorTok{);}
\NormalTok{rtcNewGeometry}\OperatorTok{(}\NormalTok{device}\OperatorTok{,}\NormalTok{ RTC\_GEOMETRY\_TYPE\_CONE\_LINEAR\_CURVE}\OperatorTok{);}
\NormalTok{rtcNewGeometry}\OperatorTok{(}\NormalTok{device}\OperatorTok{,}\NormalTok{ RTC\_GEOMETRY\_TYPE\_ROUND\_LINEAR\_CURVE}\OperatorTok{);}
\NormalTok{rtcNewGeometry}\OperatorTok{(}\NormalTok{device}\OperatorTok{,}\NormalTok{ RTC\_GEOMETRY\_TYPE\_ROUND\_BEZIER\_CURVE}\OperatorTok{);}
\NormalTok{rtcNewGeometry}\OperatorTok{(}\NormalTok{device}\OperatorTok{,}\NormalTok{ RTC\_GEOMETRY\_TYPE\_ROUND\_BSPLINE\_CURVE}\OperatorTok{);}
\NormalTok{rtcNewGeometry}\OperatorTok{(}\NormalTok{device}\OperatorTok{,}\NormalTok{ RTC\_GEOMETRY\_TYPE\_ROUND\_HERMITE\_CURVE}\OperatorTok{);}
\NormalTok{rtcNewGeometry}\OperatorTok{(}\NormalTok{device}\OperatorTok{,}\NormalTok{ RTC\_GEOMETRY\_TYPE\_ROUND\_CATMULL\_ROM\_CURVE}\OperatorTok{);}
\end{Highlighting}
\end{Shaded}

\subsubsection{DESCRIPTION}\label{description}

Curves with per vertex radii are supported with linear, cubic Bézier,
cubic B-spline, and cubic Hermite bases. Such curve geometries are
created by passing \texttt{RTC\_GEOMETRY\_TYPE\_FLAT\_LINEAR\_CURVE},
\texttt{RTC\_GEOMETRY\_TYPE\_FLAT\_BEZIER\_CURVE},
\texttt{RTC\_GEOMETRY\_TYPE\_FLAT\_BSPLINE\_CURVE},
\texttt{RTC\_GEOMETRY\_TYPE\_FLAT\_HERMITE\_CURVE},
\texttt{RTC\_GEOMETRY\_TYPE\_FLAT\_CATMULL\_ROM\_CURVE},
\texttt{RTC\_GEOMETRY\_TYPE\_NORMAL\_ORIENTED\_FLAT\_BEZIER\_CURVE},
\texttt{RTC\_GEOMETRY\_TYPE\_NORMAL\_ORIENTED\_FLAT\_BSPLINE\_CURVE},
\texttt{RTC\_GEOMETRY\_TYPE\_NORMAL\_ORIENTED\_FLAT\_HERMITE\_CURVE},
\texttt{RTC\_GEOMETRY\_TYPE\_NORMAL\_ORIENTED\_FLAT\_CATMULL\_ROM\_CURVE},
\texttt{RTC\_GEOMETRY\_TYPE\_CONE\_LINEAR\_CURVE},
\texttt{RTC\_GEOMETRY\_TYPE\_ROUND\_LINEAR\_CURVE},
\texttt{RTC\_GEOMETRY\_TYPE\_ROUND\_BEZIER\_CURVE},
\texttt{RTC\_GEOMETRY\_TYPE\_ROUND\_BSPLINE\_CURVE},
\texttt{RTC\_GEOMETRY\_TYPE\_ROUND\_HERMITE\_CURVE}, or
\texttt{RTC\_GEOMETRY\_TYPE\_ROUND\_CATMULL\_ROM\_CURVE} to the
\texttt{rtcNewGeometry} function. The curve indices can be specified
through an index buffer (\texttt{RTC\_BUFFER\_TYPE\_INDEX}) and the
curve vertices through a vertex buffer
(\texttt{RTC\_BUFFER\_TYPE\_VERTEX}). For the Hermite basis a tangent
buffer (\texttt{RTC\_BUFFER\_TYPE\_TANGENT}), normal oriented curves a
normal buffer (\texttt{RTC\_BUFFER\_TYPE\_NORMAL}), and for normal
oriented Hermite curves a normal derivative buffer
(\texttt{RTC\_BUFFER\_TYPE\_NORMAL\_DERIVATIVE}) has to get specified
additionally. See \texttt{rtcSetGeometryBuffer} and
\texttt{rtcSetSharedGeometryBuffer} for more details on how to set
buffers.

The index buffer contains an array of 32-bit indices
(\texttt{RTC\_FORMAT\_UINT} format), each pointing to the first control
vertex in the vertex buffer, but also to the first tangent in the
tangent buffer, and first normal in the normal buffer if these buffers
are present.

The vertex buffer stores each control vertex in the form of a single
precision position and radius stored in (\texttt{x}, \texttt{y},
\texttt{z}, \texttt{r}) order in memory (\texttt{RTC\_FORMAT\_FLOAT4}
format). The number of vertices is inferred from the size of this
buffer. The radii may be smaller than zero but the interpolated radii
should always be greater or equal to zero. Similarly, the tangent buffer
stores the derivative of each control vertex (\texttt{x}, \texttt{y},
\texttt{z}, \texttt{r} order and \texttt{RTC\_FORMAT\_FLOAT4} format)
and the normal buffer stores a single precision normal per control
vertex (\texttt{x}, \texttt{y}, \texttt{z} order and
\texttt{RTC\_FORMAT\_FLOAT3} format).

\paragraph{Linear Basis}\label{linear-basis}

For the linear basis the indices point to the first of 2 consecutive
control points in the vertex buffer. The first control point is the
start and the second control point the end of the line segment. When
constructing hair strands in this basis, the end-point can be shared
with the start of the next line segment.

For the linear basis the user optionally can provide a flags buffer of
type \texttt{RTC\_BUFFER\_TYPE\_FLAGS} which contains bytes that encode
if the left neighbor segment (\texttt{RTC\_CURVE\_FLAG\_NEIGHBOR\_LEFT}
flag) and/or right neighbor segment
(\texttt{RTC\_CURVE\_FLAG\_NEIGHBOR\_RIGHT} flags) exist (see
\hyperref[rtccurveflags]{RTCCurveFlags}). If this buffer is not set,
than the left/right neighbor bits are automatically calculated base on
the index buffer (left segment exists if segment(id-1)+1 == segment(id)
and right segment exists if segment(id+1)-1 == segment(id)).

A left neighbor segment is assumed to end at the start vertex of the
current segment, and to start at the previous vertex in the vertex
buffer. Similarly, the right neighbor segment is assumed to start at the
end vertex of the current segment, and to end at the next vertex in the
vertex buffer.

Only when the left and right bits are properly specified the current
segment can properly attach to the left and/or right neighbor, otherwise
the touching area may not get rendered properly.

\paragraph{Bézier Basis}\label{buxe9zier-basis}

For the cubic Bézier basis the indices point to the first of 4
consecutive control points in the vertex buffer. These control points
use the cubic Bézier basis, where the first control point represents the
start point of the curve, and the 4th control point the end point of the
curve. The Bézier basis is interpolating, thus the curve does go exactly
through the first and fourth control vertex.

\paragraph{B-spline Basis}\label{b-spline-basis}

For the cubic B-spline basis the indices point to the first of 4
consecutive control points in the vertex buffer. These control points
make up a cardinal cubic B-spline (implicit equidistant knot vector).
This basis is not interpolating, thus the curve does in general not go
through any of the control points directly. A big advantage of this
basis is that 3 control points can be shared for two continuous
neighboring curve segments, e.g.~the curves (p0,p1,p2,p3) and
(p1,p2,p3,p4) are C1 continuous. This feature makes this basis a good
choice to construct continuous multi-segment curves, as memory
consumption can be kept minimal.

\paragraph{Hermite Basis}\label{hermite-basis}

For the cubic Hermite basis the indices point to the first of 2
consecutive points in the vertex buffer, and the first of 2 consecutive
tangents in the tangent buffer. These two points and two tangents make
up a cubic Hermite curve. This basis is interpolating, thus does exactly
go through the first and second control point, and the first order
derivative at the begin and end matches exactly the value specified in
the tangent buffer. When connecting two segments continuously, the end
point and tangent of the previous segment can be shared. Different
versions of Catmull-Rom splines can be easily constructed using the
Hermite basis, by calculating a proper tangent buffer from the control
points.

\paragraph{Catmull-Rom Basis}\label{catmull-rom-basis}

For the Catmull-Rom basis the indices point to the first of 4
consecutive control points in the vertex buffer. This basis goes through
p1 and p2, with tangents (p2-p0)/2 and (p3-p1)/2.

\paragraph{Flat Curves}\label{flat-curves}

The \texttt{RTC\_GEOMETRY\_TYPE\_FLAT\_*} flat mode is a fast mode
designed to render distant hair. In this mode the curve is rendered as a
connected sequence of ray facing quads. Individual quads are considered
to have subpixel size, and zooming onto the curve might show geometric
artifacts. The number of quads to subdivide into can be specified
through the \texttt{rtcSetGeometryTessellationRate} function. By default
the tessellation rate is 4.

\paragraph{Normal Oriented Curves}\label{normal-oriented-curves}

The \texttt{RTC\_GEOMETRY\_TYPE\_NORMAL\_ORIENTED\_*} mode is a mode
designed to render blades of grass. In this mode a vertex spline has to
get specified as for the previous modes, but additionally a normal
spline is required. If the Hermite basis is used, the
\texttt{RTC\_BUFFER\_TYPE\_NORMAL} and
\texttt{RTC\_BUFFER\_TYPE\_NORMAL\_DERIVATIVE} buffers have both to be
set.

The curve is rendered as a flat band whose center approximately follows
the provided vertex spline, whose half width approximately follows the
provided radius spline, and whose normal orientation approximately
follows the provided normal spline.

To intersect the normal oriented curve, we perform a newton-raphson
style intersection of a ray with a tensor product surface of a linear
basis (perpendicular to the curve) and cubic Bézier basis (along the
curve). We use a guide curve and its derivatives to construct the
control points of that surface. The guide curve is defined by a sweep
surface defined by sweeping a line centered at the vertex spline
location along the curve. At each parameter value the half width of the
line matches the radius spline, and the direction matches the cross
product of the normal from the normal spline and tangent of the vertex
spline. Note that this construction does not work when the provided
normals are parallel to the curve direction. For this reason the
provided normals should best be kept as perpendicular to the curve
direction as possible. We further assume second order derivatives of the
center curve to be zero for this construction, as otherwise very large
curvatures occurring in corner cases, can thicken the constructed curve
significantly.

\paragraph{Round Curves}\label{round-curves}

In the \texttt{RTC\_GEOMETRY\_TYPE\_ROUND\_*} round mode, a real
geometric surface is rendered for the curve, which is more expensive but
allows closeup views.

For the linear basis the round mode renders a cone that tangentially
touches a start-sphere and end-sphere. The start sphere is rendered when
no previous segments is indicated by the neighbor bits. The end sphere
is always rendered but parts that lie inside the next segment are
clipped away (if that next segment exists). This way a curve is closed
on both ends and the interior will render properly as long as only
neighboring segments penetrate into a segment. For this to work properly
it is important that the flags buffer is properly populated with
neighbor information.

For the cubic polynomial bases, the round mode renders a sweep surface
by sweeping a varying radius circle tangential along the curve. As a
limitation, the radius of the curve has to be smaller than the curvature
radius of the curve at each location on the curve.

The intersection with the curve segment stores the parametric hit
location along the curve segment as u-coordinate (range 0 to +1).

For flat curves, the v-coordinate is set to the normalized distance in
the range -1 to +1. For normal oriented curves the v-coordinate is in
the range 0 to 1. For the linear basis and in round mode the
v-coordinate is set to zero.

In flat mode, the geometry normal \texttt{Ng} is set to the tangent of
the curve at the hit location. In round mode and for normal oriented
curves, the geometry normal \texttt{Ng} is set to the non-normalized
geometric normal of the surface.

For multi-segment motion blur, the number of time steps must be first
specified using the \texttt{rtcSetGeometryTimeStepCount} call. Then a
vertex buffer for each time step can be set using different buffer
slots, and all these buffers must have the same stride and size. For the
Hermite basis also a tangent buffer has to be set for each time step and
for normal oriented curves a normal buffer has to get specified for each
time step.

Also see tutorials \hyperref[hair]{Hair} and
\hyperref[buxe9zier-curves]{Curves} for examples of how to create and
use curve geometries.

\subsubsection{EXIT STATUS}\label{exit-status}

On failure \texttt{NULL} is returned and an error code is set that can
be queried using \texttt{rtcGetDeviceError}.

\subsubsection{SEE ALSO}\label{see-also}

\hyperref[rtcnewgeometry]{rtcNewGeometry},
\hyperref[rtccurveflags]{RTCCurveFlags}

\pagebreak

\section{RTC\_GEOMETRY\_TYPE\_POINT}\label{rtc_geometry_type_point}

\subsubsection{NAME}\label{name}

\begin{Shaded}
\begin{Highlighting}[]
\NormalTok{RTC\_GEOMETRY\_TYPE\_SPHERE\_POINT }\OperatorTok{{-}}
\NormalTok{  point geometry spheres}

\NormalTok{RTC\_GEOMETRY\_TYPE\_DISC\_POINT }\OperatorTok{{-}}
\NormalTok{  point geometry with ray}\OperatorTok{{-}}\NormalTok{oriented discs}

\NormalTok{RTC\_GEOMETRY\_TYPE\_ORIENTED\_DISC\_POINT }\OperatorTok{{-}}
\NormalTok{  point geometry with normal}\OperatorTok{{-}}\NormalTok{oriented discs}
\end{Highlighting}
\end{Shaded}

\subsubsection{SYNOPSIS}\label{synopsis}

\begin{Shaded}
\begin{Highlighting}[]
\PreprocessorTok{\#include }\ImportTok{\textless{}embree4/rtcore.h\textgreater{}}

\NormalTok{rtcNewGeometry}\OperatorTok{(}\NormalTok{device}\OperatorTok{,}\NormalTok{ RTC\_GEOMETRY\_TYPE\_SPHERE\_POINT}\OperatorTok{);}
\NormalTok{rtcNewGeometry}\OperatorTok{(}\NormalTok{device}\OperatorTok{,}\NormalTok{ RTC\_GEOMETRY\_TYPE\_DISC\_POINT}\OperatorTok{);}
\NormalTok{rtcNewGeometry}\OperatorTok{(}\NormalTok{device}\OperatorTok{,}\NormalTok{ RTC\_GEOMETRY\_TYPE\_ORIENTED\_DISC\_POINT}\OperatorTok{);}
\end{Highlighting}
\end{Shaded}

\subsubsection{DESCRIPTION}\label{description}

Points with per vertex radii are supported with sphere, ray-oriented
discs, and normal-oriented discs geometric representations. Such point
geometries are created by passing
\texttt{RTC\_GEOMETRY\_TYPE\_SPHERE\_POINT},
\texttt{RTC\_GEOMETRY\_TYPE\_DISC\_POINT}, or
\texttt{RTC\_GEOMETRY\_TYPE\_ORIENTED\_DISC\_POINT} to the
\texttt{rtcNewGeometry} function. The point vertices can be specified t
through a vertex buffer (\texttt{RTC\_BUFFER\_TYPE\_VERTEX}). For the
normal oriented discs a normal buffer
(\texttt{RTC\_BUFFER\_TYPE\_NORMAL}) has to get specified additionally.
See \texttt{rtcSetGeometryBuffer} and
\texttt{rtcSetSharedGeometryBuffer} for more details on how to set
buffers.

The vertex buffer stores each control vertex in the form of a single
precision position and radius stored in (\texttt{x}, \texttt{y},
\texttt{z}, \texttt{r}) order in memory (\texttt{RTC\_FORMAT\_FLOAT4}
format). The number of vertices is inferred from the size of this
buffer. Similarly, the normal buffer stores a single precision normal
per control vertex (\texttt{x}, \texttt{y}, \texttt{z} order and
\texttt{RTC\_FORMAT\_FLOAT3} format).

In the \texttt{RTC\_GEOMETRY\_TYPE\_SPHERE\_POINT} mode, a real
geometric surface is rendered for the curve, which is more expensive but
allows closeup views.

The \texttt{RTC\_GEOMETRY\_TYPE\_DISC\_POINT} flat mode is a fast mode
designed to render distant points. In this mode the point is rendered as
a ray facing disc.

The \texttt{RTC\_GEOMETRY\_TYPE\_ORIENTED\_DISC\_POINT} mode is a mode
designed as a midpoint geometrically between ray facing discs and
spheres. In this mode the point is rendered as a normal oriented disc.

For all point types, only the hit distance and geometry normal is
returned as hit information, u and v are set to zero.

For multi-segment motion blur, the number of time steps must be first
specified using the \texttt{rtcSetGeometryTimeStepCount} call. Then a
vertex buffer for each time step can be set using different buffer
slots, and all these buffers must have the same stride and size.

Also see tutorial {[}Points{]} for an example of how to create and use
point geometries.

\subsubsection{EXIT STATUS}\label{exit-status}

On failure \texttt{NULL} is returned and an error code is set that can
be queried using \texttt{rtcGetDeviceError}.

\subsubsection{SEE ALSO}\label{see-also}

\hyperref[rtcnewgeometry]{rtcNewGeometry}

\pagebreak

\section{RTC\_GEOMETRY\_TYPE\_USER}\label{rtc_geometry_type_user}

\subsubsection{NAME}\label{name}

\begin{Shaded}
\begin{Highlighting}[]
\NormalTok{RTC\_GEOMETRY\_TYPE\_USER }\OperatorTok{{-}}\NormalTok{ user geometry type}
\end{Highlighting}
\end{Shaded}

\subsubsection{SYNOPSIS}\label{synopsis}

\begin{Shaded}
\begin{Highlighting}[]
\PreprocessorTok{\#include }\ImportTok{\textless{}embree4/rtcore.h\textgreater{}}

\NormalTok{RTCGeometry geometry }\OperatorTok{=}
\NormalTok{  rtcNewGeometry}\OperatorTok{(}\NormalTok{device}\OperatorTok{,}\NormalTok{ RTC\_GEOMETRY\_TYPE\_USER}\OperatorTok{);}
\end{Highlighting}
\end{Shaded}

\subsubsection{DESCRIPTION}\label{description}

User-defined geometries contain a number of user-defined primitives,
just like triangle meshes contain multiple triangles. The shape of the
user-defined primitives is specified through registered callback
functions, which enable extending Embree with arbitrary types of
primitives.

User-defined geometries are created by passing
\texttt{RTC\_GEOMETRY\_TYPE\_USER} to the \texttt{rtcNewGeometry}
function call. One has to set the number of primitives (see
\texttt{rtcSetGeometryUserPrimitiveCount}), a user data pointer (see
\texttt{rtcSetGeometryUserData}), a bounding function closure (see
\texttt{rtcSetGeometryBoundsFunction}), as well as user-defined
intersect (see \texttt{rtcSetGeometryIntersectFunction}) and occluded
(see \texttt{rtcSetGeometryOccludedFunction}) callback functions. The
bounding function is used to query the bounds of all time steps of a
user primitive, while the intersect and occluded callback functions are
called to intersect the primitive with a ray. The user data pointer is
passed to each callback invocation and can be used to point to the
application's representation of the user geometry.

The creation of a user geometry typically looks the following:

\begin{Shaded}
\begin{Highlighting}[]
\NormalTok{RTCGeometry geometry }\OperatorTok{=}\NormalTok{ rtcNewGeometry}\OperatorTok{(}\NormalTok{device}\OperatorTok{,}\NormalTok{ RTC\_GEOMETRY\_TYPE\_USER}\OperatorTok{);}
\NormalTok{rtcSetGeometryUserPrimitiveCount}\OperatorTok{(}\NormalTok{geometry}\OperatorTok{,}\NormalTok{ numPrimitives}\OperatorTok{);}
\NormalTok{rtcSetGeometryUserData}\OperatorTok{(}\NormalTok{geometry}\OperatorTok{,}\NormalTok{ userGeometryRepresentation}\OperatorTok{);}
\NormalTok{rtcSetGeometryBoundsFunction}\OperatorTok{(}\NormalTok{geometry}\OperatorTok{,}\NormalTok{ boundsFunction}\OperatorTok{);}
\NormalTok{rtcSetGeometryIntersectFunction}\OperatorTok{(}\NormalTok{geometry}\OperatorTok{,}\NormalTok{ intersectFunction}\OperatorTok{);}
\NormalTok{rtcSetGeometryOccludedFunction}\OperatorTok{(}\NormalTok{geometry}\OperatorTok{,}\NormalTok{ occludedFunction}\OperatorTok{);}
\end{Highlighting}
\end{Shaded}

Please have a look at the \texttt{rtcSetGeometryBoundsFunction},
\texttt{rtcSetGeometryIntersectFunction}, and
\texttt{rtcSetGeometryOccludedFunction} functions on the implementation
of the callback functions.

Primitives of a user geometry are ignored during rendering when their
bounds are empty, thus bounds have lower\textgreater upper in at least
one dimension.

See tutorial \hyperref[user-geometry]{User Geometry} for an example of
how to use the user-defined geometries.

\subsubsection{EXIT STATUS}\label{exit-status}

On failure \texttt{NULL} is returned and an error code is set that can
be queried using \texttt{rtcGetDeviceError}.

\subsubsection{SEE ALSO}\label{see-also}

\hyperref[rtcnewgeometry]{rtcNewGeometry},
\hyperref[rtcsetgeometryuserprimitivecount]{rtcSetGeometryUserPrimitiveCount},
\hyperref[rtcsetgeometryuserdata]{rtcSetGeometryUserData},
\hyperref[rtcsetgeometryboundsfunction]{rtcSetGeometryBoundsFunction},
\hyperref[rtcsetgeometryintersectfunction]{rtcSetGeometryIntersectFunction},
\hyperref[rtcsetgeometryoccludedfunction]{rtcSetGeometryOccludedFunction}

\pagebreak

\section{RTC\_GEOMETRY\_TYPE\_INSTANCE}\label{rtc_geometry_type_instance}

\subsubsection{NAME}\label{name}

\begin{Shaded}
\begin{Highlighting}[]
\NormalTok{RTC\_GEOMETRY\_TYPE\_INSTANCE }\OperatorTok{{-}}\NormalTok{ instance geometry type}
\end{Highlighting}
\end{Shaded}

\subsubsection{SYNOPSIS}\label{synopsis}

\begin{Shaded}
\begin{Highlighting}[]
\PreprocessorTok{\#include }\ImportTok{\textless{}embree4/rtcore.h\textgreater{}}

\NormalTok{RTCGeometry geometry }\OperatorTok{=}
\NormalTok{   rtcNewGeometry}\OperatorTok{(}\NormalTok{device}\OperatorTok{,}\NormalTok{ RTC\_GEOMETRY\_TYPE\_INSTANCE}\OperatorTok{);}
\end{Highlighting}
\end{Shaded}

\subsubsection{DESCRIPTION}\label{description}

Embree supports instancing of scenes using affine transformations (3×3
matrix plus translation). As the instanced scene is stored only a single
time, even if instanced to multiple locations, this feature can be used
to create very complex scenes with small memory footprint.

Embree supports both single-level instancing and multi-level instancing.
The maximum instance nesting depth is
\texttt{RTC\_MAX\_INSTANCE\_LEVEL\_COUNT}; it can be configured at
compile-time using the constant
\texttt{EMBREE\_MAX\_INSTANCE\_LEVEL\_COUNT}. Users should adapt this
constant to their needs: instances nested any deeper are silently
ignored in release mode, and cause assertions in debug mode.

Instances are created by passing \texttt{RTC\_GEOMETRY\_TYPE\_INSTANCE}
to the \texttt{rtcNewGeometry} function call. The instanced scene can be
set using the \texttt{rtcSetGeometryInstancedScene} call, and the affine
transformation can be set using the \texttt{rtcSetGeometryTransform}
function.

Please note that \texttt{rtcCommitScene} on the instanced scene should
be called first, followed by \texttt{rtcCommitGeometry} on the instance,
followed by \texttt{rtcCommitScene} for the top-level scene containing
the instance.

If a ray hits the instance, the \texttt{geomID} and \texttt{primID}
members of the hit are set to the geometry ID and primitive ID of the
hit primitive in the instanced scene, and the \texttt{instID} member of
the hit is set to the geometry ID of the instance in the top-level
scene.

The instancing scheme can also be implemented using user geometries. To
achieve this, the user geometry code should set the \texttt{instID}
member of the ray query context to the geometry ID of the instance, then
trace the transformed ray, and finally set the \texttt{instID} field of
the ray query context again to -1. The \texttt{instID} field is copied
automatically by each primitive intersector into the \texttt{instID}
field of the hit structure when the primitive is hit. See the
\hyperref[user-geometry]{User Geometry} tutorial for an example.

For multi-segment motion blur, the number of time steps must be first
specified using the \texttt{rtcSetGeometryTimeStepCount} function. Then
a transformation for each time step can be specified using the
\texttt{rtcSetGeometryTransform} function.

See tutorials \hyperref[instanced-geometry]{Instanced Geometry} and
\hyperref[multi-level-instancing]{Multi Level Instancing} for examples
of how to use instances.

\subsubsection{EXIT STATUS}\label{exit-status}

On failure \texttt{NULL} is returned and an error code is set that can
be queried using \texttt{rtcGetDeviceError}.

\subsubsection{SEE ALSO}\label{see-also}

\hyperref[rtcnewgeometry]{rtcNewGeometry},
\hyperref[rtcsetgeometryinstancedscene]{rtcSetGeometryInstancedScene},
\hyperref[rtcsetgeometrytransform]{rtcSetGeometryTransform}

\pagebreak

\section{RTCCurveFlags}\label{rtccurveflags}

\subsubsection{NAME}\label{name}

\begin{Shaded}
\begin{Highlighting}[]
\NormalTok{RTCCurveFlags }\OperatorTok{{-}}\NormalTok{ per segment flags }\ControlFlowTok{for}\NormalTok{ curve geometry}
\end{Highlighting}
\end{Shaded}

\subsubsection{SYNOPSIS}\label{synopsis}

\begin{Shaded}
\begin{Highlighting}[]
\PreprocessorTok{\#include }\ImportTok{\textless{}embree4/rtcore.h\textgreater{}}

\KeywordTok{enum}\NormalTok{ RTCCurveFlags}
\OperatorTok{\{}
\NormalTok{  RTC\_CURVE\_FLAG\_NEIGHBOR\_LEFT  }\OperatorTok{=} \OperatorTok{(}\DecValTok{1} \OperatorTok{\textless{}\textless{}} \DecValTok{0}\OperatorTok{),} 
\NormalTok{  RTC\_CURVE\_FLAG\_NEIGHBOR\_RIGHT }\OperatorTok{=} \OperatorTok{(}\DecValTok{1} \OperatorTok{\textless{}\textless{}} \DecValTok{1}\OperatorTok{)} 
\OperatorTok{\};}
\end{Highlighting}
\end{Shaded}

\subsubsection{DESCRIPTION}\label{description}

The RTCCurveFlags type is used for linear curves to determine if the
left and/or right neighbor segment exist. Therefore one attaches a
buffer of type RTC\_BUFFER\_TYPE\_FLAGS to the curve geometry which
stores an individual byte per curve segment.

If the RTC\_CURVE\_FLAG\_NEIGHBOR\_LEFT flag in that byte is enabled for
a curve segment, then the left segment exists (which starts one vertex
before the start vertex of the current curve) and the current segment is
rendered to properly attach to that segment.

If the RTC\_CURVE\_FLAG\_NEIGHBOR\_RIGHT flag in that byte is enabled
for a curve segment, then the right segment exists (which ends one
vertex after the end vertex of the current curve) and the current
segment is rendered to properly attach to that segment.

When not properly specifying left and right flags for linear curves, the
rendering at the ending of these curves may not look correct, in
particular when round linear curves are viewed from the inside.

\subsubsection{EXIT STATUS}\label{exit-status}

\subsubsection{SEE ALSO}\label{see-also}

\hyperref[rtc_geometry_type_curve]{RTC\_GEOMETRY\_TYPE\_CURVE}

\pagebreak

\section{rtcRetainGeometry}\label{rtcretaingeometry}

\subsubsection{NAME}\label{name}

\begin{Shaded}
\begin{Highlighting}[]
\NormalTok{rtcRetainGeometry }\OperatorTok{{-}}\NormalTok{ increments the geometry reference count}
\end{Highlighting}
\end{Shaded}

\subsubsection{SYNOPSIS}\label{synopsis}

\begin{Shaded}
\begin{Highlighting}[]
\PreprocessorTok{\#include }\ImportTok{\textless{}embree4/rtcore.h\textgreater{}}

\DataTypeTok{void}\NormalTok{ rtcRetainGeometry}\OperatorTok{(}\NormalTok{RTCGeometry geometry}\OperatorTok{);}
\end{Highlighting}
\end{Shaded}

\subsubsection{DESCRIPTION}\label{description}

Geometry objects are reference counted. The \texttt{rtcRetainGeometry}
function increments the reference count of the passed geometry object
(\texttt{geometry} argument). This function together with
\texttt{rtcReleaseGeometry} allows to use the internal reference
counting in a C++ wrapper class to handle the ownership of the object.

\subsubsection{EXIT STATUS}\label{exit-status}

On failure an error code is set that can be queried using
\texttt{rtcGetDeviceError}.

\subsubsection{SEE ALSO}\label{see-also}

\hyperref[rtcnewgeometry]{rtcNewGeometry},
\hyperref[rtcreleasegeometry]{rtcReleaseGeometry}

\pagebreak

\section{rtcReleaseGeometry}\label{rtcreleasegeometry}

\subsubsection{NAME}\label{name}

\begin{Shaded}
\begin{Highlighting}[]
\NormalTok{rtcReleaseGeometry }\OperatorTok{{-}}\NormalTok{ decrements the geometry reference count}
\end{Highlighting}
\end{Shaded}

\subsubsection{SYNOPSIS}\label{synopsis}

\begin{Shaded}
\begin{Highlighting}[]
\PreprocessorTok{\#include }\ImportTok{\textless{}embree4/rtcore.h\textgreater{}}

\DataTypeTok{void}\NormalTok{ rtcReleaseGeometry}\OperatorTok{(}\NormalTok{RTCGeometry geometry}\OperatorTok{);}
\end{Highlighting}
\end{Shaded}

\subsubsection{DESCRIPTION}\label{description}

Geometry objects are reference counted. The \texttt{rtcReleaseGeometry}
function decrements the reference count of the passed geometry object
(\texttt{geometry} argument). When the reference count falls to 0, the
geometry gets destroyed.

\subsubsection{EXIT STATUS}\label{exit-status}

On failure an error code is set that can be queried using
\texttt{rtcGetDeviceError}.

\subsubsection{SEE ALSO}\label{see-also}

\hyperref[rtcnewgeometry]{rtcNewGeometry},
\hyperref[rtcretaingeometry]{rtcRetainGeometry}

\pagebreak

\section{rtcCommitGeometry}\label{rtccommitgeometry}

\subsubsection{NAME}\label{name}

\begin{Shaded}
\begin{Highlighting}[]
\NormalTok{rtcCommitGeometry }\OperatorTok{{-}}\NormalTok{ commits geometry changes}
\end{Highlighting}
\end{Shaded}

\subsubsection{SYNOPSIS}\label{synopsis}

\begin{Shaded}
\begin{Highlighting}[]
\PreprocessorTok{\#include }\ImportTok{\textless{}embree4/rtcore.h\textgreater{}}

\DataTypeTok{void}\NormalTok{ rtcCommitGeometry}\OperatorTok{(}\NormalTok{RTCGeometry geometry}\OperatorTok{);}
\end{Highlighting}
\end{Shaded}

\subsubsection{DESCRIPTION}\label{description}

The \texttt{rtcCommitGeometry} function is used to commit all geometry
changes performed to a geometry (\texttt{geometry} parameter). After a
geometry gets modified, this function must be called to properly update
the internal state of the geometry to perform interpolations using
\texttt{rtcInterpolate} or to commit a scene containing the geometry
using \texttt{rtcCommitScene}.

\subsubsection{EXIT STATUS}\label{exit-status}

On failure an error code is set that can be queried using
\texttt{rtcGetDeviceError}.

\subsubsection{SEE ALSO}\label{see-also}

\hyperref[rtcinterpolate]{rtcInterpolate},
\hyperref[rtccommitscene]{rtcCommitScene}

\pagebreak

\section{rtcEnableGeometry}\label{rtcenablegeometry}

\subsubsection{NAME}\label{name}

\begin{Shaded}
\begin{Highlighting}[]
\NormalTok{rtcEnableGeometry }\OperatorTok{{-}}\NormalTok{ enables the geometry}
\end{Highlighting}
\end{Shaded}

\subsubsection{SYNOPSIS}\label{synopsis}

\begin{Shaded}
\begin{Highlighting}[]
\PreprocessorTok{\#include }\ImportTok{\textless{}embree4/rtcore.h\textgreater{}}

\DataTypeTok{void}\NormalTok{ rtcEnableGeometry}\OperatorTok{(}\NormalTok{RTCGeometry geometry}\OperatorTok{);}
\end{Highlighting}
\end{Shaded}

\subsubsection{DESCRIPTION}\label{description}

The \texttt{rtcEnableGeometry} function enables the specified geometry
(\texttt{geometry} argument). Only enabled geometries are rendered. Each
geometry is enabled by default at construction time.

After enabling a geometry, the scene containing that geometry must be
committed using \texttt{rtcCommitScene} for the change to have effect.

\subsubsection{EXIT STATUS}\label{exit-status}

On failure an error code is set that can be queried using
\texttt{rtcGetDeviceError}.

\subsubsection{SEE ALSO}\label{see-also}

\hyperref[rtcnewgeometry]{rtcNewGeometry},
\hyperref[rtcdisablegeometry]{rtcDisableGeometry},
\hyperref[rtccommitscene]{rtcCommitScene}

\pagebreak

\section{rtcDisableGeometry}\label{rtcdisablegeometry}

\subsubsection{NAME}\label{name}

\begin{Shaded}
\begin{Highlighting}[]
\NormalTok{rtcDisableGeometry }\OperatorTok{{-}}\NormalTok{ disables the geometry}
\end{Highlighting}
\end{Shaded}

\subsubsection{SYNOPSIS}\label{synopsis}

\begin{Shaded}
\begin{Highlighting}[]
\PreprocessorTok{\#include }\ImportTok{\textless{}embree4/rtcore.h\textgreater{}}

\DataTypeTok{void}\NormalTok{ rtcDisableGeometry}\OperatorTok{(}\NormalTok{RTCGeometry geometry}\OperatorTok{);}
\end{Highlighting}
\end{Shaded}

\subsubsection{DESCRIPTION}\label{description}

The \texttt{rtcDisableGeometry} function disables the specified geometry
(\texttt{geometry} argument). A disabled geometry is not rendered. Each
geometry is enabled by default at construction time.

After disabling a geometry, the scene containing that geometry must be
committed using \texttt{rtcCommitScene} for the change to have effect.

\subsubsection{EXIT STATUS}\label{exit-status}

On failure an error code is set that can be queried using
\texttt{rtcGetDeviceError}.

\subsubsection{SEE ALSO}\label{see-also}

\hyperref[rtcnewgeometry]{rtcNewGeometry},
\hyperref[rtcenablegeometry]{rtcEnableGeometry},
\hyperref[rtccommitscene]{rtcCommitScene}

\pagebreak

\section{rtcSetGeometryTimeStepCount}\label{rtcsetgeometrytimestepcount}

\subsubsection{NAME}\label{name}

\begin{Shaded}
\begin{Highlighting}[]
\NormalTok{rtcSetGeometryTimeStepCount }\OperatorTok{{-}}\NormalTok{ sets the number of time steps of the}
\NormalTok{  geometry}
\end{Highlighting}
\end{Shaded}

\subsubsection{SYNOPSIS}\label{synopsis}

\begin{Shaded}
\begin{Highlighting}[]
\PreprocessorTok{\#include }\ImportTok{\textless{}embree4/rtcore.h\textgreater{}}

\DataTypeTok{void}\NormalTok{ rtcSetGeometryTimeStepCount}\OperatorTok{(}
\NormalTok{  RTCGeometry geometry}\OperatorTok{,}
  \DataTypeTok{unsigned} \DataTypeTok{int}\NormalTok{ timeStepCount}
\OperatorTok{);}
\end{Highlighting}
\end{Shaded}

\subsubsection{DESCRIPTION}\label{description}

The \texttt{rtcSetGeometryTimeStepCount} function sets the number of
time steps for multi-segment motion blur (\texttt{timeStepCount}
parameter) of the specified geometry (\texttt{geometry} parameter).

For triangle meshes (\texttt{RTC\_GEOMETRY\_TYPE\_TRIANGLE}), quad
meshes (\texttt{RTC\_GEOMETRY\_TYPE\_QUAD}), curves
(\texttt{RTC\_GEOMETRY\_TYPE\_CURVE}), points
(\texttt{RTC\_GEOMETRY\_TYPE\_POINT}), and subdivision geometries
(\texttt{RTC\_GEOMETRY\_TYPE\_SUBDIVISION}), the number of time steps
directly corresponds to the number of vertex buffer slots available
(\texttt{RTC\_BUFFER\_TYPE\_VERTEX} buffer type). For these geometries,
one vertex buffer per time step must be specified when creating
multi-segment motion blur geometries.

For instance geometries (\texttt{RTC\_GEOMETRY\_TYPE\_INSTANCE}), a
transformation must be specified for each time step (see
\texttt{rtcSetGeometryTransform}).

For user geometries, the registered bounding callback function must
provide a bounding box per primitive and time step, and the intersection
and occlusion callback functions should properly intersect the
motion-blurred geometry at the ray time.

\subsubsection{EXIT STATUS}\label{exit-status}

On failure an error code is set that can be queried using
\texttt{rtcGetDeviceError}.

\subsubsection{SEE ALSO}\label{see-also}

\hyperref[rtcnewgeometry]{rtcNewGeometry},
\hyperref[rtcsetgeometrytimerange]{rtcSetGeometryTimeRange}

\pagebreak

\section{rtcSetGeometryTimeRange}\label{rtcsetgeometrytimerange}

\subsubsection{NAME}\label{name}

\begin{Shaded}
\begin{Highlighting}[]
\NormalTok{rtcSetGeometryTimeRange }\OperatorTok{{-}}\NormalTok{ sets the time range }\ControlFlowTok{for}\NormalTok{ a motion blur geometry}
\end{Highlighting}
\end{Shaded}

\subsubsection{SYNOPSIS}\label{synopsis}

\begin{Shaded}
\begin{Highlighting}[]
\PreprocessorTok{\#include }\ImportTok{\textless{}embree4/rtcore.h\textgreater{}}

\DataTypeTok{void}\NormalTok{ rtcSetGeometryTimeRange}\OperatorTok{(}
\NormalTok{  RTCGeometry geometry}\OperatorTok{,}
  \DataTypeTok{float}\NormalTok{ startTime}\OperatorTok{,}
  \DataTypeTok{float}\NormalTok{ endTime}
\OperatorTok{);}
\end{Highlighting}
\end{Shaded}

\subsubsection{DESCRIPTION}\label{description}

The \texttt{rtcSetGeometryTimeRange} function sets a time range which
defines the start (and end time) of the first (and last) time step of a
motion blur geometry. The time range is defined relative to the camera
shutter interval {[}0,1{]} but it can be arbitrary. Thus the startTime
can be smaller, equal, or larger 0, indicating a geometry whose
animation definition start before, at, or after the camera shutter
opens. Similar the endTime can be smaller, equal, or larger than 1,
indicating a geometry whose animation definition ends after, at, or
before the camera shutter closes. The startTime has to be smaller or
equal to the endTime.

The default time range when this function is not called is the entire
camera shutter {[}0,1{]}. For best performance at most one time segment
of the piece wise linear definition of the motion should fall outside
the shutter window to the left and to the right. Thus do not set the
startTime or endTime too far outside the {[}0,1{]} interval for best
performance.

This time range feature will also allow geometries to appear and
disappear during the camera shutter time if the specified time range is
a sub range of {[}0,1{]}.

Please also have a look at the \texttt{rtcSetGeometryTimeStepCount}
function to see how to define the time steps for the specified time
range.

\subsubsection{EXIT STATUS}\label{exit-status}

On failure an error code is set that can be queried using
\texttt{rtcGetDeviceError}.

\subsubsection{SEE ALSO}\label{see-also}

\hyperref[rtcsetgeometrytimestepcount]{rtcSetGeometryTimeStepCount}

\pagebreak

\section{rtcSetGeometryVertexAttributeCount}\label{rtcsetgeometryvertexattributecount}

\subsubsection{NAME}\label{name}

\begin{Shaded}
\begin{Highlighting}[]
\NormalTok{rtcSetGeometryVertexAttributeCount }\OperatorTok{{-}}\NormalTok{ sets the number of vertex}
\NormalTok{  attributes of the geometry}
\end{Highlighting}
\end{Shaded}

\subsubsection{SYNOPSIS}\label{synopsis}

\begin{Shaded}
\begin{Highlighting}[]
\PreprocessorTok{\#include }\ImportTok{\textless{}embree4/rtcore.h\textgreater{}}

\DataTypeTok{void}\NormalTok{ rtcSetGeometryVertexAttributeCount}\OperatorTok{(}
\NormalTok{  RTCGeometry geometry}\OperatorTok{,}
  \DataTypeTok{unsigned} \DataTypeTok{int}\NormalTok{ vertexAttributeCount}
\OperatorTok{);}
\end{Highlighting}
\end{Shaded}

\subsubsection{DESCRIPTION}\label{description}

The \texttt{rtcSetGeometryVertexAttributeCount} function sets the number
of slots (\texttt{vertexAttributeCount} parameter) for vertex attribute
buffers (\texttt{RTC\_BUFFER\_TYPE\_VERTEX\_ATTRIBUTE}) that can be used
for the specified geometry (\texttt{geometry} parameter).

This function is supported only for triangle meshes
(\texttt{RTC\_GEOMETRY\_TYPE\_TRIANGLE}), quad meshes
(\texttt{RTC\_GEOMETRY\_TYPE\_QUAD}), curves
(\texttt{RTC\_GEOMETRY\_TYPE\_CURVE}), points
(\texttt{RTC\_GEOMETRY\_TYPE\_POINT}), and subdivision geometries
(\texttt{RTC\_GEOMETRY\_TYPE\_SUBDIVISION}).

\subsubsection{EXIT STATUS}\label{exit-status}

On failure an error code is set that can be queried using
\texttt{rtcGetDeviceError}.

\subsubsection{SEE ALSO}\label{see-also}

\hyperref[rtcnewgeometry]{rtcNewGeometry},
\hyperref[rtcbuffertype]{RTCBufferType}

\pagebreak

\section{rtcSetGeometryMask}\label{rtcsetgeometrymask}

\subsubsection{NAME}\label{name}

\begin{Shaded}
\begin{Highlighting}[]
\NormalTok{rtcSetGeometryMask }\OperatorTok{{-}}\NormalTok{ sets the geometry mask}
\end{Highlighting}
\end{Shaded}

\subsubsection{SYNOPSIS}\label{synopsis}

\begin{Shaded}
\begin{Highlighting}[]
\PreprocessorTok{\#include }\ImportTok{\textless{}embree4/rtcore.h\textgreater{}}

\DataTypeTok{void}\NormalTok{ rtcSetGeometryMask}\OperatorTok{(}
\NormalTok{  RTCGeometry geometry}\OperatorTok{,}
  \DataTypeTok{unsigned} \DataTypeTok{int}\NormalTok{ mask}
\OperatorTok{);}
\end{Highlighting}
\end{Shaded}

\subsubsection{DESCRIPTION}\label{description}

The \texttt{rtcSetGeometryMask} function sets a 32-bit geometry mask
(\texttt{mask} argument) for the specified geometry (\texttt{geometry}
argument).

This geometry mask is used together with the ray mask stored inside the
\texttt{mask} field of the ray. The primitives of the geometry are hit
by the ray only if the bitwise \texttt{and} operation of the geometry
mask with the ray mask is not 0. This feature can be used to disable
selected geometries for specifically tagged rays, e.g.~to disable shadow
casting for certain geometries.

\subsubsection{EXIT STATUS}\label{exit-status}

On failure an error code is set that can be queried using
\texttt{rtcGetDeviceError}.

\subsubsection{SEE ALSO}\label{see-also}

\hyperref[rtcray]{RTCRay},
\hyperref[rtcgetdeviceproperty]{rtcGetDeviceProperty}

\pagebreak

\section{rtcSetGeometryBuildQuality}\label{rtcsetgeometrybuildquality}

\subsubsection{NAME}\label{name}

\begin{Shaded}
\begin{Highlighting}[]
\NormalTok{rtcSetGeometryBuildQuality }\OperatorTok{{-}}\NormalTok{ sets the build quality }\ControlFlowTok{for}\NormalTok{ the geometry}
\end{Highlighting}
\end{Shaded}

\subsubsection{SYNOPSIS}\label{synopsis}

\begin{Shaded}
\begin{Highlighting}[]
\PreprocessorTok{\#include }\ImportTok{\textless{}embree4/rtcore.h\textgreater{}}

\DataTypeTok{void}\NormalTok{ rtcSetGeometryBuildQuality}\OperatorTok{(}
\NormalTok{  RTCGeometry geometry}\OperatorTok{,}
  \KeywordTok{enum}\NormalTok{ RTCBuildQuality quality}
\OperatorTok{);}
\end{Highlighting}
\end{Shaded}

\subsubsection{DESCRIPTION}\label{description}

The \texttt{rtcSetGeometryBuildQuality} function sets the build quality
(\texttt{quality} argument) for the specified geometry
(\texttt{geometry} argument). The per-geometry build quality is only a
hint and may be ignored. Embree currently uses the per-geometry build
quality when the scene build quality is set to
\texttt{RTC\_BUILD\_QUALITY\_LOW}. In this mode a two-level acceleration
structure is build, and geometries build a separate acceleration
structure using the geometry build quality. The per-geometry build
quality can be one of:

\begin{itemize}
\item
  \texttt{RTC\_BUILD\_QUALITY\_LOW}: Creates lower quality data
  structures, e.g.~for dynamic scenes.
\item
  \texttt{RTC\_BUILD\_QUALITY\_MEDIUM}: Default build quality for most
  usages. Gives a good compromise between build and render performance.
\item
  \texttt{RTC\_BUILD\_QUALITY\_HIGH}: Creates higher quality data
  structures for final-frame rendering. Enables a spatial split builder
  for certain primitive types.
\item
  \texttt{RTC\_BUILD\_QUALITY\_REFIT}: Uses a BVH refitting approach
  when changing only the vertex buffer.
\end{itemize}

\subsubsection{EXIT STATUS}\label{exit-status}

On failure an error code is set that can be queried using
\texttt{rtcGetDeviceError}.

\subsubsection{SEE ALSO}\label{see-also}

\hyperref[rtcsetscenebuildquality]{rtcSetSceneBuildQuality}

\pagebreak

\section{rtcSetGeometryMaxRadiusScale}\label{rtcsetgeometrymaxradiusscale}

\subsubsection{NAME}\label{name}

\begin{Shaded}
\begin{Highlighting}[]
\NormalTok{rtcSetGeometryMaxRadiusScale }\OperatorTok{{-}}\NormalTok{ assigns a maximal curve radius scale factor }\ControlFlowTok{for}\NormalTok{ min}\OperatorTok{{-}}\NormalTok{width feature}
\end{Highlighting}
\end{Shaded}

\subsubsection{SYNOPSIS}\label{synopsis}

\begin{Shaded}
\begin{Highlighting}[]
\PreprocessorTok{\#include }\ImportTok{\textless{}embree4/rtcore.h\textgreater{}}

\DataTypeTok{void}\NormalTok{ rtcSetGeometryMaxRadiusScale}\OperatorTok{(}\NormalTok{RTCGeometry geometry}\OperatorTok{,} \DataTypeTok{float}\NormalTok{ maxRadiusScale}\OperatorTok{);}
\end{Highlighting}
\end{Shaded}

\subsubsection{DESCRIPTION}\label{description}

The \texttt{rtcSetMaxGeometryScale} function specifies a maximal scaling
factor for curve radii used by the min-width feature.

The min-width feature can increase the radius of curves and points, in
order to reduce aliasing and improve render times. The feature is
disabled by default and has to get enabled using the EMBREE\_MIN\_WIDTH
cmake option.

To use the feature, one has to specify a maximal curve radius scaling
factor using the
\hyperref[rtcsetgeometrymaxradiusscale]{rtcSetGeometryMaxRadiusScale}
function. This factor should be a small number (e.g.~4) as the
constructed BVH bounds get increased in order to bound the curve in the
worst case of maximal radii.

One also has to set the minWidthDistanceFactor in the RTCRayQueryContext
when tracing a ray. This factor controls the target radius size of a
curve or point at some distance away of the ray origin.

For each control point p with radius r of a curve or point primitive,
the primitive intersectors first calculate a target radius r' as:

\begin{Shaded}
\begin{Highlighting}[]
\NormalTok{r}\CharTok{\textquotesingle{} }\ErrorTok{= length(p{-}ray\_org) * minWidthDistanceFactor}
\end{Highlighting}
\end{Shaded}

Typically the minWidthDistanceFactor is set by the application such that
the target radius projects to the width of half a pixel (thus primitive
diameter is pixel sized).

The target radius r' is then clamped against the minimal bound r and
maximal bound maxRadiusScale*r to obtain the final radius r'\,':

\begin{Shaded}
\begin{Highlighting}[]
\NormalTok{r}\CharTok{\textquotesingle{}}\ErrorTok{\textquotesingle{}} \OperatorTok{=}\NormalTok{ max}\OperatorTok{(}\NormalTok{r}\OperatorTok{,}\NormalTok{ min}\OperatorTok{(}\NormalTok{r}\CharTok{\textquotesingle{},}\ErrorTok{ maxRadiusScale*r))}
\end{Highlighting}
\end{Shaded}

Thus curves or points close to the camera are rendered with a normal
radii r, and curves or points far from the camera are not enlarged too
much, as this would be very expensive to render.

When \texttt{rtcSetGeometryMaxRadiusScale} function is not invoked for a
curve or point geometry (or if the maximal scaling factor is set to
1.0), then the curve or point geometry renders normally, with radii not
modified by the min-width feature.

\subsubsection{EXIT STATUS}\label{exit-status}

On failure an error code is set that can be queried using
\texttt{rtcGetDeviceError}.

\subsubsection{SEE ALSO}\label{see-also}

\hyperref[rtcinitrayquerycontext]{rtcInitRayQueryContext}

\pagebreak

\section{rtcSetGeometryBuffer}\label{rtcsetgeometrybuffer}

\subsubsection{NAME}\label{name}

\begin{Shaded}
\begin{Highlighting}[]
\NormalTok{rtcSetGeometryBuffer }\OperatorTok{{-}}\NormalTok{ assigns a view of a buffer to the geometry}
\end{Highlighting}
\end{Shaded}

\subsubsection{SYNOPSIS}\label{synopsis}

\begin{Shaded}
\begin{Highlighting}[]
\PreprocessorTok{\#include }\ImportTok{\textless{}embree4/rtcore.h\textgreater{}}

\DataTypeTok{void}\NormalTok{ rtcSetGeometryBuffer}\OperatorTok{(}
\NormalTok{  RTCGeometry geometry}\OperatorTok{,}
  \KeywordTok{enum}\NormalTok{ RTCBufferType type}\OperatorTok{,}
  \DataTypeTok{unsigned} \DataTypeTok{int}\NormalTok{ slot}\OperatorTok{,}
  \KeywordTok{enum}\NormalTok{ RTCFormat format}\OperatorTok{,}
\NormalTok{  RTCBuffer buffer}\OperatorTok{,}
  \DataTypeTok{size\_t}\NormalTok{ byteOffset}\OperatorTok{,}
  \DataTypeTok{size\_t}\NormalTok{ byteStride}\OperatorTok{,}
  \DataTypeTok{size\_t}\NormalTok{ itemCount}
\OperatorTok{);}
\end{Highlighting}
\end{Shaded}

\subsubsection{DESCRIPTION}\label{description}

The \texttt{rtcSetGeometryBuffer} function binds a view of a buffer
object (\texttt{buffer} argument) to a geometry buffer type and slot
(\texttt{type} and \texttt{slot} argument) of the specified geometry
(\texttt{geometry} argument).

One can specify the start of the first buffer element in bytes
(\texttt{byteOffset} argument), the byte stride between individual
buffer elements (\texttt{byteStride} argument), the format of the buffer
elements (\texttt{format} argument), and the number of elements to bind
(\texttt{itemCount}).

The start address (\texttt{byteOffset} argument) and stride
(\texttt{byteStride} argument) must be both aligned to 4 bytes,
otherwise the \texttt{rtcSetGeometryBuffer} function will fail.

After successful completion of this function, the geometry will hold a
reference to the buffer object.

\subsubsection{EXIT STATUS}\label{exit-status}

On failure an error code is set that can be queried using
\texttt{rtcGetDeviceError}.

\subsubsection{SEE ALSO}\label{see-also}

\hyperref[rtcsetsharedgeometrybuffer]{rtcSetSharedGeometryBuffer},
\hyperref[rtcsetnewgeometrybuffer]{rtcSetNewGeometryBuffer}

\pagebreak

\section{rtcSetSharedGeometryBuffer}\label{rtcsetsharedgeometrybuffer}

\subsubsection{NAME}\label{name}

\begin{Shaded}
\begin{Highlighting}[]
\NormalTok{rtcSetSharedGeometryBuffer }\OperatorTok{{-}}\NormalTok{ assigns a view of a shared data buffer}
\NormalTok{  to a geometry}
\end{Highlighting}
\end{Shaded}

\subsubsection{SYNOPSIS}\label{synopsis}

\begin{Shaded}
\begin{Highlighting}[]
\PreprocessorTok{\#include }\ImportTok{\textless{}embree4/rtcore.h\textgreater{}}

\DataTypeTok{void}\NormalTok{ rtcSetSharedGeometryBuffer}\OperatorTok{(}
\NormalTok{  RTCGeometry geometry}\OperatorTok{,}
  \KeywordTok{enum}\NormalTok{ RTCBufferType type}\OperatorTok{,}
  \DataTypeTok{unsigned} \DataTypeTok{int}\NormalTok{ slot}\OperatorTok{,}
  \KeywordTok{enum}\NormalTok{ RTCFormat format}\OperatorTok{,}
  \AttributeTok{const} \DataTypeTok{void}\OperatorTok{*}\NormalTok{ ptr}\OperatorTok{,}
  \DataTypeTok{size\_t}\NormalTok{ byteOffset}\OperatorTok{,}
  \DataTypeTok{size\_t}\NormalTok{ byteStride}\OperatorTok{,}
  \DataTypeTok{size\_t}\NormalTok{ itemCount}
\OperatorTok{);}
\end{Highlighting}
\end{Shaded}

\subsubsection{DESCRIPTION}\label{description}

The \texttt{rtcSetSharedGeometryBuffer} function binds a view of a
shared user-managed data buffer (\texttt{ptr} argument) to a geometry
buffer type and slot (\texttt{type} and \texttt{slot} argument) of the
specified geometry (\texttt{geometry} argument).

One can specify the start of the first buffer element in bytes
(\texttt{byteOffset} argument), the byte stride between individual
buffer elements (\texttt{byteStride} argument), the format of the buffer
elements (\texttt{format} argument), and the number of elements to bind
(\texttt{itemCount}).

The start address (\texttt{byteOffset} argument) and stride
(\texttt{byteStride} argument) must be both aligned to 4 bytes;
otherwise the \texttt{rtcSetSharedGeometryBuffer} function will fail.

When the buffer will be used as a vertex buffer
(\texttt{RTC\_BUFFER\_TYPE\_VERTEX} and
\texttt{RTC\_BUFFER\_TYPE\_VERTEX\_ATTRIBUTE}), the last buffer element
must be readable using 16-byte SSE load instructions, thus padding the
last element is required for certain layouts. E.g. a standard
\texttt{float3} vertex buffer layout should add storage for at least one
more float to the end of the buffer.

The buffer data must remain valid for as long as the buffer may be used,
and the user is responsible for freeing the buffer data when no longer
required.

Sharing buffers can significantly reduce the memory required by the
application, thus we recommend using this feature. When enabling the
\texttt{RTC\_SCENE\_FLAG\_COMPACT} scene flag, the spatial index
structures index into the vertex buffer, resulting in even higher memory
savings.

\subsubsection{EXIT STATUS}\label{exit-status}

On failure an error code is set that can be queried using
\texttt{rtcGetDeviceError}.

\subsubsection{SEE ALSO}\label{see-also}

\hyperref[rtcsetgeometrybuffer]{rtcSetGeometryBuffer},
\hyperref[rtcsetnewgeometrybuffer]{rtcSetNewGeometryBuffer}

\pagebreak

\section{rtcSetNewGeometryBuffer}\label{rtcsetnewgeometrybuffer}

\subsubsection{NAME}\label{name}

\begin{Shaded}
\begin{Highlighting}[]
\NormalTok{rtcSetNewGeometryBuffer }\OperatorTok{{-}}\NormalTok{ creates }\KeywordTok{and}\NormalTok{ assigns a }\KeywordTok{new}\NormalTok{ data buffer to}
\NormalTok{  the geometry}
\end{Highlighting}
\end{Shaded}

\subsubsection{SYNOPSIS}\label{synopsis}

\begin{Shaded}
\begin{Highlighting}[]
\PreprocessorTok{\#include }\ImportTok{\textless{}embree4/rtcore.h\textgreater{}}

\DataTypeTok{void}\OperatorTok{*}\NormalTok{ rtcSetNewGeometryBuffer}\OperatorTok{(}
\NormalTok{  RTCGeometry geometry}\OperatorTok{,}
  \KeywordTok{enum}\NormalTok{ RTCBufferType type}\OperatorTok{,}
  \DataTypeTok{unsigned} \DataTypeTok{int}\NormalTok{ slot}\OperatorTok{,}
  \KeywordTok{enum}\NormalTok{ RTCFormat format}\OperatorTok{,}
  \DataTypeTok{size\_t}\NormalTok{ byteStride}\OperatorTok{,}
  \DataTypeTok{size\_t}\NormalTok{ itemCount}
\OperatorTok{);}
\end{Highlighting}
\end{Shaded}

\subsubsection{DESCRIPTION}\label{description}

The \texttt{rtcSetNewGeometryBuffer} function creates a new data buffer
of specified format (\texttt{format} argument), byte stride
(\texttt{byteStride} argument), and number of items (\texttt{itemCount}
argument), and assigns it to a geometry buffer slot (\texttt{type} and
\texttt{slot} argument) of the specified geometry (\texttt{geometry}
argument). The buffer data is managed internally and automatically freed
when the geometry is destroyed.

The byte stride (\texttt{byteStride} argument) must be aligned to 4
bytes; otherwise the \texttt{rtcSetNewGeometryBuffer} function will
fail.

The allocated buffer will be automatically over-allocated slightly when
used as a vertex buffer, where a requirement is that each buffer element
should be readable using 16-byte SSE load instructions.

\subsubsection{EXIT STATUS}\label{exit-status}

On failure an error code is set that can be queried using
\texttt{rtcGetDeviceError}.

\subsubsection{SEE ALSO}\label{see-also}

\hyperref[rtcsetgeometrybuffer]{rtcSetGeometryBuffer},
\hyperref[rtcsetsharedgeometrybuffer]{rtcSetSharedGeometryBuffer}

\pagebreak

\section{RTCFormat}\label{rtcformat}

\subsubsection{NAME}\label{name}

\begin{Shaded}
\begin{Highlighting}[]
\NormalTok{RTCFormat }\OperatorTok{{-}}\NormalTok{ specifies format of data in buffers}
\end{Highlighting}
\end{Shaded}

\subsubsection{SYNOPSIS}\label{synopsis}

\begin{Shaded}
\begin{Highlighting}[]
\PreprocessorTok{\#include }\ImportTok{\textless{}embree4/rtcore\_ray.h\textgreater{}}

\KeywordTok{enum}\NormalTok{ RTCFormat}
\OperatorTok{\{}
\NormalTok{  RTC\_FORMAT\_UINT}\OperatorTok{,}
\NormalTok{  RTC\_FORMAT\_UINT2}\OperatorTok{,}
\NormalTok{  RTC\_FORMAT\_UINT3}\OperatorTok{,}
\NormalTok{  RTC\_FORMAT\_UINT4}\OperatorTok{,}

\NormalTok{  RTC\_FORMAT\_FLOAT}\OperatorTok{,}
\NormalTok{  RTC\_FORMAT\_FLOAT2}\OperatorTok{,}
\NormalTok{  RTC\_FORMAT\_FLOAT3}\OperatorTok{,}
\NormalTok{  RTC\_FORMAT\_FLOAT4}\OperatorTok{,}
\NormalTok{  RTC\_FORMAT\_FLOAT5}\OperatorTok{,}
\NormalTok{  RTC\_FORMAT\_FLOAT6}\OperatorTok{,}
\NormalTok{  RTC\_FORMAT\_FLOAT7}\OperatorTok{,}
\NormalTok{  RTC\_FORMAT\_FLOAT8}\OperatorTok{,}
\NormalTok{  RTC\_FORMAT\_FLOAT9}\OperatorTok{,}
\NormalTok{  RTC\_FORMAT\_FLOAT10}\OperatorTok{,}
\NormalTok{  RTC\_FORMAT\_FLOAT11}\OperatorTok{,}
\NormalTok{  RTC\_FORMAT\_FLOAT12}\OperatorTok{,}
\NormalTok{  RTC\_FORMAT\_FLOAT13}\OperatorTok{,}
\NormalTok{  RTC\_FORMAT\_FLOAT14}\OperatorTok{,}
\NormalTok{  RTC\_FORMAT\_FLOAT15}\OperatorTok{,}
\NormalTok{  RTC\_FORMAT\_FLOAT16}\OperatorTok{,}

\NormalTok{  RTC\_FORMAT\_FLOAT3X4\_ROW\_MAJOR}\OperatorTok{,}
\NormalTok{  RTC\_FORMAT\_FLOAT4X4\_ROW\_MAJOR}\OperatorTok{,}

\NormalTok{  RTC\_FORMAT\_FLOAT3X4\_COLUMN\_MAJOR}\OperatorTok{,}
\NormalTok{  RTC\_FORMAT\_FLOAT4X4\_COLUMN\_MAJOR}\OperatorTok{,}

\NormalTok{  RTC\_FORMAT\_GRID}\OperatorTok{,}
\OperatorTok{\};}
\end{Highlighting}
\end{Shaded}

\subsubsection{DESCRIPTION}\label{description}

The \texttt{RTFormat} structure defines the data format stored in data
buffers provided to Embree using the
\hyperref[rtcsetgeometrybuffer]{rtcSetGeometryBuffer},
\hyperref[rtcsetsharedgeometrybuffer]{rtcSetSharedGeometryBuffer}, and
\hyperref[rtcsetnewgeometrybuffer]{rtcSetNewGeometryBuffer} API calls.

The \texttt{RTC\_FORMAT\_UINT/2/3/4} format are used to specify that
data buffers store unsigned integers, or unsigned integer vectors of
size 2,3 or 4. This format has typically to get used when specifying
index buffers, e.g.~\texttt{RTC\_FORMAT\_UINT3} for triangle meshes.

The \texttt{RTC\_FORMAT\_FLOAT/2/3/4...} format are used to specify that
data buffers store single precision floating point values, or vectors
there of (size 2,3,4, etc.). This format is typcally used to specify to
format of vertex buffers, e.g.~the \texttt{RTC\_FORMAT\_FLOAT3} type for
vertex buffers of triangle meshes.

The \texttt{RTC\_FORMAT\_FLOAT3X4\_ROW\_MAJOR} and
\texttt{RTC\_FORMAT\_FLOAT3X4\_COLUMN\_MAJOR} formats, specify a 3x4
floating point matrix layed out either row major or column major. The
\texttt{RTC\_FORMAT\_FLOAT4X4\_ROW\_MAJOR} and
\texttt{RTC\_FORMAT\_FLOAT4X4\_COLUMN\_MAJOR} formats, specify a 4x4
floating point matrix layed out either row major or column major. These
matrix formats are used in the
\hyperref[rtcsetgeometrytransform]{rtcSetGeometryTransform} function in
order to set a transformation matrix for geometries.

The \texttt{RTC\_FORMAT\_GRID} is a special data format used to specify
grid primitives of layout RTCGrid when creating grid geometries (see
\hyperref[rtc_geometry_type_grid]{RTC\_GEOMETRY\_TYPE\_GRID}).

\subsubsection{EXIT STATUS}\label{exit-status}

\subsubsection{SEE ALSO}\label{see-also}

\hyperref[rtcsetgeometrybuffer]{rtcSetGeometryBuffer},
\hyperref[rtcsetsharedgeometrybuffer]{rtcSetSharedGeometryBuffer},
\hyperref[rtcsetnewgeometrybuffer]{rtcSetNewGeometryBuffer},
\hyperref[rtcsetgeometrytransform]{rtcSetGeometryTransform}

\pagebreak

\section{RTCBufferType}\label{rtcbuffertype}

\subsubsection{NAME}\label{name}

\begin{Shaded}
\begin{Highlighting}[]
\NormalTok{RTCFormat }\OperatorTok{{-}}\NormalTok{ specifies format of data in buffers}
\end{Highlighting}
\end{Shaded}

\subsubsection{SYNOPSIS}\label{synopsis}

\begin{Shaded}
\begin{Highlighting}[]
\PreprocessorTok{\#include }\ImportTok{\textless{}embree4/rtcore\_ray.h\textgreater{}}

\KeywordTok{enum}\NormalTok{ RTCBufferType}
\OperatorTok{\{}
\NormalTok{  RTC\_BUFFER\_TYPE\_INDEX            }\OperatorTok{=} \DecValTok{0}\OperatorTok{,}
\NormalTok{  RTC\_BUFFER\_TYPE\_VERTEX           }\OperatorTok{=} \DecValTok{1}\OperatorTok{,}
\NormalTok{  RTC\_BUFFER\_TYPE\_VERTEX\_ATTRIBUTE }\OperatorTok{=} \DecValTok{2}\OperatorTok{,}
\NormalTok{  RTC\_BUFFER\_TYPE\_NORMAL           }\OperatorTok{=} \DecValTok{3}\OperatorTok{,}
\NormalTok{  RTC\_BUFFER\_TYPE\_TANGENT          }\OperatorTok{=} \DecValTok{4}\OperatorTok{,}
\NormalTok{  RTC\_BUFFER\_TYPE\_NORMAL\_DERIVATIVE }\OperatorTok{=} \DecValTok{5}\OperatorTok{,}

\NormalTok{  RTC\_BUFFER\_TYPE\_GRID                 }\OperatorTok{=} \DecValTok{8}\OperatorTok{,}

\NormalTok{  RTC\_BUFFER\_TYPE\_FACE                 }\OperatorTok{=} \DecValTok{16}\OperatorTok{,}
\NormalTok{  RTC\_BUFFER\_TYPE\_LEVEL                }\OperatorTok{=} \DecValTok{17}\OperatorTok{,}
\NormalTok{  RTC\_BUFFER\_TYPE\_EDGE\_CREASE\_INDEX    }\OperatorTok{=} \DecValTok{18}\OperatorTok{,}
\NormalTok{  RTC\_BUFFER\_TYPE\_EDGE\_CREASE\_WEIGHT   }\OperatorTok{=} \DecValTok{19}\OperatorTok{,}
\NormalTok{  RTC\_BUFFER\_TYPE\_VERTEX\_CREASE\_INDEX  }\OperatorTok{=} \DecValTok{20}\OperatorTok{,}
\NormalTok{  RTC\_BUFFER\_TYPE\_VERTEX\_CREASE\_WEIGHT }\OperatorTok{=} \DecValTok{21}\OperatorTok{,}
\NormalTok{  RTC\_BUFFER\_TYPE\_HOLE                 }\OperatorTok{=} \DecValTok{22}\OperatorTok{,}

\NormalTok{  RTC\_BUFFER\_TYPE\_FLAGS }\OperatorTok{=} \DecValTok{32}
\OperatorTok{\};}
\end{Highlighting}
\end{Shaded}

\subsubsection{DESCRIPTION}\label{description}

The \texttt{RTBufferType} structure defines slots to assign data buffers
to using the \hyperref[rtcsetgeometrybuffer]{rtcSetGeometryBuffer},
\hyperref[rtcsetsharedgeometrybuffer]{rtcSetSharedGeometryBuffer}, and
\hyperref[rtcsetnewgeometrybuffer]{rtcSetNewGeometryBuffer} API calls.

For most geometry types the \texttt{RTC\_BUFFER\_TYPE\_INDEX} slot is
used to assign an index buffer, while the
\texttt{RTC\_BUFFER\_TYPE\_VERTEX} is used to assign the corresponding
vertex buffer.

The \texttt{RTC\_BUFFER\_TYPE\_VERTEX\_ATTRIBUTE} slot can get used to
assign arbitrary additional vertex data which can get interpolated using
the \hyperref[rtcinterpolate]{rtcInterpolate} API call.

The \texttt{RTC\_BUFFER\_TYPE\_NORMAL},
\texttt{RTC\_BUFFER\_TYPE\_TANGENT}, and
\texttt{RTC\_BUFFER\_TYPE\_NORMAL\_DERIVATIVE} are special buffers
required to assign per vertex normals, tangents, and normal derivatives
for some curve types.

The \texttt{RTC\_BUFFER\_TYPE\_GRID} buffer is used to assign the grid
primitive buffer for grid geometries (see
\hyperref[rtc_geometry_type_grid]{RTC\_GEOMETRY\_TYPE\_GRID}).

The \texttt{RTC\_BUFFER\_TYPE\_FACE}, \texttt{RTC\_BUFFER\_TYPE\_LEVEL},
\texttt{RTC\_BUFFER\_TYPE\_EDGE\_CREASE\_INDEX},
\texttt{RTC\_BUFFER\_TYPE\_EDGE\_CREASE\_WEIGHT},
\texttt{RTC\_BUFFER\_TYPE\_VERTEX\_CREASE\_INDEX},
\texttt{RTC\_BUFFER\_TYPE\_VERTEX\_CREASE\_WEIGHT}, and
\texttt{RTC\_BUFFER\_TYPE\_HOLE} are special buffers required to create
subdivision meshes (see
\hyperref[rtc_geometry_type_subdivision]{RTC\_GEOMETRY\_TYPE\_SUBDIVISION}).

The \texttt{RTC\_BUFFER\_TYPE\_FLAGS} can get used to add additional
flag per primitive of a geometry, and is currently only used for linear
curves.

\subsubsection{EXIT STATUS}\label{exit-status}

\subsubsection{SEE ALSO}\label{see-also}

\hyperref[rtcsetgeometrybuffer]{rtcSetGeometryBuffer},
\hyperref[rtcsetsharedgeometrybuffer]{rtcSetSharedGeometryBuffer},
\hyperref[rtcsetnewgeometrybuffer]{rtcSetNewGeometryBuffer}

\pagebreak

\section{rtcGetGeometryBufferData}\label{rtcgetgeometrybufferdata}

\subsubsection{NAME}\label{name}

\begin{Shaded}
\begin{Highlighting}[]
\NormalTok{rtcGetGeometryBufferData }\OperatorTok{{-}}\NormalTok{ gets pointer to}
\NormalTok{  the first buffer view element}
\end{Highlighting}
\end{Shaded}

\subsubsection{SYNOPSIS}\label{synopsis}

\begin{Shaded}
\begin{Highlighting}[]
\PreprocessorTok{\#include }\ImportTok{\textless{}embree4/rtcore.h\textgreater{}}

\DataTypeTok{void}\OperatorTok{*}\NormalTok{ rtcGetGeometryBufferData}\OperatorTok{(}
\NormalTok{  RTCGeometry geometry}\OperatorTok{,}
  \KeywordTok{enum}\NormalTok{ RTCBufferType type}\OperatorTok{,}
  \DataTypeTok{unsigned} \DataTypeTok{int}\NormalTok{ slot}
\OperatorTok{);}
\end{Highlighting}
\end{Shaded}

\subsubsection{DESCRIPTION}\label{description}

The \texttt{rtcGetGeometryBufferData} function returns a pointer to the
first element of the buffer view attached to the specified buffer type
and slot (\texttt{type} and \texttt{slot} argument) of the geometry
(\texttt{geometry} argument).

\subsubsection{EXIT STATUS}\label{exit-status}

On failure an error code is set that can be queried using
\texttt{rtcGetDeviceError}.

\subsubsection{SEE ALSO}\label{see-also}

\hyperref[rtcsetgeometrybuffer]{rtcSetGeometryBuffer},
\hyperref[rtcsetsharedgeometrybuffer]{rtcSetSharedGeometryBuffer},
\hyperref[rtcsetnewgeometrybuffer]{rtcSetNewGeometryBuffer}

\pagebreak

\section{rtcUpdateGeometryBuffer}\label{rtcupdategeometrybuffer}

\subsubsection{NAME}\label{name}

\begin{Shaded}
\begin{Highlighting}[]
\NormalTok{rtcUpdateGeometryBuffer }\OperatorTok{{-}}\NormalTok{ marks a buffer view bound to the geometry}
\NormalTok{  as modified}
\end{Highlighting}
\end{Shaded}

\subsubsection{SYNOPSIS}\label{synopsis}

\begin{Shaded}
\begin{Highlighting}[]
\PreprocessorTok{\#include }\ImportTok{\textless{}embree4/rtcore.h\textgreater{}}

\DataTypeTok{void}\NormalTok{ rtcUpdateGeometryBuffer}\OperatorTok{(}
\NormalTok{  RTCGeometry geometry}\OperatorTok{,}
  \KeywordTok{enum}\NormalTok{ RTCBufferType type}\OperatorTok{,}
  \DataTypeTok{unsigned} \DataTypeTok{int}\NormalTok{ slot}
\OperatorTok{);}
\end{Highlighting}
\end{Shaded}

\subsubsection{DESCRIPTION}\label{description}

The \texttt{rtcUpdateGeometryBuffer} function marks the buffer view
bound to the specified buffer type and slot (\texttt{type} and
\texttt{slot} argument) of a geometry (\texttt{geometry} argument) as
modified.

If a data buffer is changed by the application, the
\texttt{rtcUpdateGeometryBuffer} call must be invoked for that buffer.
Each buffer view assigned to a buffer slot is initially marked as
modified, thus this function needs to be called only when doing buffer
modifications after the first \texttt{rtcCommitScene}.

\subsubsection{EXIT STATUS}\label{exit-status}

On failure an error code is set that can be queried using
\texttt{rtcGetDeviceError}.

\subsubsection{SEE ALSO}\label{see-also}

\hyperref[rtcnewgeometry]{rtcNewGeometry},
\hyperref[rtccommitscene]{rtcCommitScene}

\pagebreak

\section{rtcSetGeometryIntersectFilterFunction}\label{rtcsetgeometryintersectfilterfunction}

\subsubsection{NAME}\label{name}

\begin{Shaded}
\begin{Highlighting}[]
\NormalTok{rtcSetGeometryIntersectFilterFunction }\OperatorTok{{-}}\NormalTok{ sets the intersection filter}
  \ControlFlowTok{for}\NormalTok{ the geometry}
\end{Highlighting}
\end{Shaded}

\subsubsection{SYNOPSIS}\label{synopsis}

\begin{Shaded}
\begin{Highlighting}[]
\PreprocessorTok{\#include }\ImportTok{\textless{}embree4/rtcore.h\textgreater{}}

\KeywordTok{struct}\NormalTok{ RTCFilterFunctionNArguments}
\OperatorTok{\{}
  \DataTypeTok{int}\OperatorTok{*}\NormalTok{ valid}\OperatorTok{;}
  \DataTypeTok{void}\OperatorTok{*}\NormalTok{ geometryUserPtr}\OperatorTok{;}
  \AttributeTok{const} \KeywordTok{struct}\NormalTok{ RTCRayQueryContext}\OperatorTok{*}\NormalTok{ context}\OperatorTok{;}
  \KeywordTok{struct}\NormalTok{ RTCRayN}\OperatorTok{*}\NormalTok{ ray}\OperatorTok{;}
  \KeywordTok{struct}\NormalTok{ RTCHitN}\OperatorTok{*}\NormalTok{ hit}\OperatorTok{;}
  \DataTypeTok{unsigned} \DataTypeTok{int}\NormalTok{ N}\OperatorTok{;}
\OperatorTok{\};}

\KeywordTok{typedef} \DataTypeTok{void} \OperatorTok{(*}\NormalTok{RTCFilterFunctionN}\OperatorTok{)(}
  \AttributeTok{const} \KeywordTok{struct}\NormalTok{ RTCFilterFunctionNArguments}\OperatorTok{*}\NormalTok{ args}
\OperatorTok{);}

\DataTypeTok{void}\NormalTok{ rtcSetGeometryIntersectFilterFunction}\OperatorTok{(}
\NormalTok{  RTCGeometry geometry}\OperatorTok{,}
\NormalTok{  RTCFilterFunctionN filter}
\OperatorTok{);}
\end{Highlighting}
\end{Shaded}

\subsubsection{DESCRIPTION}\label{description}

The \texttt{rtcSetGeometryIntersectFilterFunction} function registers an
intersection filter callback function (\texttt{filter} argument) for the
specified geometry (\texttt{geometry} argument).

Only a single callback function can be registered per geometry, and
further invocations overwrite the previously set callback function.
Passing \texttt{NULL} as function pointer disables the registered
callback function.

The registered intersection filter function is invoked for every hit
encountered during the \texttt{rtcIntersect}-type ray queries and can
accept or reject that hit. The feature can be used to define a
silhouette for a primitive and reject hits that are outside the
silhouette. E.g. a tree leaf could be modeled with an alpha texture that
decides whether hit points lie inside or outside the leaf.

If the \texttt{RTC\_BUILD\_QUALITY\_HIGH} mode is set, the filter
functions may be called multiple times for the same primitive hit.
Further, rays hitting exactly the edge might also report two hits for
the same surface. For certain use cases, the application may have to
work around this limitation by collecting already reported hits
(\texttt{geomID}/\texttt{primID} pairs) and ignoring duplicates.

The filter function callback of type \texttt{RTCFilterFunctionN} gets
passed a number of arguments through the
\texttt{RTCFilterFunctionNArguments} structure. The \texttt{valid}
parameter of that structure points to an integer valid mask (0 means
invalid and -1 means valid). The \texttt{geometryUserPtr} member is a
user pointer optionally set per geometry through the
\texttt{rtcSetGeometryUserData} function. The \texttt{context} member
points to the ray query context passed to the ray query function. The
\texttt{ray} parameter points to \texttt{N} rays in SOA layout. The
\texttt{hit} parameter points to \texttt{N} hits in SOA layout to test.
The \texttt{N} parameter is the number of rays and hits in \texttt{ray}
and \texttt{hit}. The hit distance is provided as the \texttt{tfar}
value of the ray. If the hit geometry is instanced, the \texttt{instID}
member of the ray is valid, and the ray and the potential hit are in
object space.

The filter callback function has the task to check for each valid ray
whether it wants to accept or reject the corresponding hit. To reject a
hit, the filter callback function just has to write \texttt{0} to the
integer valid mask of the corresponding ray. To accept the hit, it just
has to leave the valid mask set to \texttt{-1}. The filter function is
further allowed to change the hit and decrease the \texttt{tfar} value
of the ray but it should not modify other ray data nor any inactive
components of the ray or hit.

When performing ray queries using \texttt{rtcIntersect1}, it is
guaranteed that the packet size is 1 when the callback is invoked. When
performing ray queries using the \texttt{rtcIntersect4/8/16} functions,
it is not generally guaranteed that the ray packet size (and order of
rays inside the packet) passed to the callback matches the initial ray
packet. However, under some circumstances these properties are
guaranteed, and whether this is the case can be queried using
\texttt{rtcGetDeviceProperty}.

For many usage scenarios, repacking and re-ordering of rays does not
cause difficulties in implementing the callback function. However,
algorithms that need to extend the ray with additional data must use the
\texttt{rayID} component of the ray to identify the original ray to
access the per-ray data.

The implementation of the filter function can choose to implement a
single code path that uses the ray access helper functions
\texttt{RTCRay\_XXX} and hit access helper functions
\texttt{RTCHit\_XXX} to access ray and hit data. Alternatively the code
can branch to optimized implementations for specific sizes of \texttt{N}
and cast the \texttt{ray} and \texttt{hit} inputs to the proper packet
types.

\subsubsection{EXIT STATUS}\label{exit-status}

On failure an error code is set that can be queried using
\texttt{rtcGetDeviceError}.

\subsubsection{SEE ALSO}\label{see-also}

\hyperref[rtcsetgeometryoccludedfilterfunction]{rtcSetGeometryOccludedFilterFunction}

\pagebreak

\section{rtcSetGeometryOccludedFilterFunction}\label{rtcsetgeometryoccludedfilterfunction}

\subsubsection{NAME}\label{name}

\begin{Shaded}
\begin{Highlighting}[]
\NormalTok{rtcSetGeometryOccludedFilterFunction }\OperatorTok{{-}}\NormalTok{ sets the occlusion filter}
  \ControlFlowTok{for}\NormalTok{ the geometry}
\end{Highlighting}
\end{Shaded}

\subsubsection{SYNOPSIS}\label{synopsis}

\begin{Shaded}
\begin{Highlighting}[]
\PreprocessorTok{\#include }\ImportTok{\textless{}embree4/rtcore.h\textgreater{}}

\DataTypeTok{void}\NormalTok{ rtcSetGeometryOccludedFilterFunction}\OperatorTok{(}
\NormalTok{  RTCGeometry geometry}\OperatorTok{,}
\NormalTok{  RTCFilterFunctionN filter}
\OperatorTok{);}
\end{Highlighting}
\end{Shaded}

\subsubsection{DESCRIPTION}\label{description}

The \texttt{rtcSetGeometryOccludedFilterFunction} function registers an
occlusion filter callback function (\texttt{filter} argument) for the
specified geometry (\texttt{geometry} argument).

Only a single callback function can be registered per geometry, and
further invocations overwrite the previously set callback function.
Passing \texttt{NULL} as function pointer disables the registered
callback function.

The registered intersection filter function is invoked for every hit
encountered during the \texttt{rtcOccluded}-type ray queries and can
accept or reject that hit. The feature can be used to define a
silhouette for a primitive and reject hits that are outside the
silhouette. E.g. a tree leaf could be modeled with an alpha texture that
decides whether hit points lie inside or outside the leaf.

Please see the description of the
\texttt{rtcSetGeometryIntersectFilterFunction} for a description of the
filter callback function.

\subsubsection{EXIT STATUS}\label{exit-status}

On failure an error code is set that can be queried using
\texttt{rtcGetDeviceError}.

\subsubsection{SEE ALSO}\label{see-also}

\hyperref[rtcsetgeometryintersectfilterfunction]{rtcSetGeometryIntersectFilterFunction}

\pagebreak

\section{rtcSetGeometryEnableFilterFunctionFromArguments}\label{rtcsetgeometryenablefilterfunctionfromarguments}

\subsubsection{NAME}\label{name}

\begin{Shaded}
\begin{Highlighting}[]
\NormalTok{rtcSetGeometryEnableFilterFunctionFromArguments }\OperatorTok{{-}}\NormalTok{ enables}
\NormalTok{  argument filter functions }\ControlFlowTok{for}\NormalTok{ the geometry}
\end{Highlighting}
\end{Shaded}

\subsubsection{SYNOPSIS}\label{synopsis}

\begin{Shaded}
\begin{Highlighting}[]
\PreprocessorTok{\#include }\ImportTok{\textless{}embree4/rtcore.h\textgreater{}}

\DataTypeTok{void}\NormalTok{ rtcSetGeometryEnableFilterFunctionFromArguments}\OperatorTok{(}
\NormalTok{   RTCGeometry geometry}\OperatorTok{,} \DataTypeTok{bool}\NormalTok{ enable}\OperatorTok{);}
\end{Highlighting}
\end{Shaded}

\subsubsection{DESCRIPTION}\label{description}

This function enables invokation the filter function passed through
\texttt{RTCIntersectArguments} or \texttt{RTCOccludedArguments} to the
intersect and occluded queries. If enable is true the argument filter
function invokation is enabled for the geometry or disabled otherwise.
By default the invokation of the argument filter function is disabled
for some geometry.

The argument filter function invokation can also get enforced for each
geometry by using the
\texttt{RTC\_RAY\_QUERY\_FLAG\_INVOKE\_ARGUMENT\_FILTER} ray query flag
that can get passed to \texttt{rtcIntersect} and \texttt{rtcOccluded}
functions. See Section
\hyperref[rtcinitintersectarguments]{rtcInitIntersectArguments} and
\hyperref[rtcinitoccludedarguments]{rtcInitOccludedArguments} for more
details.

In order to use the argument filter function for some scene, that
feature additionally has to get enabled using the
\texttt{RTC\_SCENE\_FLAG\_FILTER\_FUNCTION\_IN\_ARGUMENTS} scene flag.
See Section \hyperref[rtcsetsceneflags]{rtcSetSceneFlags} for more
details.

\subsubsection{EXIT STATUS}\label{exit-status}

On failure an error code is set that can get queried using
\texttt{rtcGetDeviceError}.

\subsubsection{SEE ALSO}\label{see-also}

\hyperref[rtcinitintersectarguments]{rtcInitIntersectArguments},
\hyperref[rtcinitoccludedarguments]{rtcInitOccludedArguments},
\hyperref[rtcsetsceneflags]{rtcSetSceneFlags}

\pagebreak

\section{rtcInvokeIntersectFilterFromGeometry}\label{rtcinvokeintersectfilterfromgeometry}

\subsubsection{NAME}\label{name}

\begin{Shaded}
\begin{Highlighting}[]
\NormalTok{rtcInvokeIntersectFilterFromGeometry }\OperatorTok{{-}}\NormalTok{ invokes the}
\NormalTok{  intersection filter function from the geometry}
\end{Highlighting}
\end{Shaded}

\subsubsection{SYNOPSIS}\label{synopsis}

\begin{Shaded}
\begin{Highlighting}[]
\PreprocessorTok{\#include }\ImportTok{\textless{}embree4/rtcore.h\textgreater{}}

\DataTypeTok{void}\NormalTok{ rtcInvokeIntersectFilterFromGeometry}\OperatorTok{(}
  \AttributeTok{const} \KeywordTok{struct}\NormalTok{ RTCIntersectFunctionNArguments}\OperatorTok{*}\NormalTok{ args}\OperatorTok{,}
  \AttributeTok{const} \KeywordTok{struct}\NormalTok{ RTCFilterFunctionNArguments}\OperatorTok{*}\NormalTok{ filterArgs}
\OperatorTok{);}
\end{Highlighting}
\end{Shaded}

\subsubsection{DESCRIPTION}\label{description}

The \texttt{rtcInvokeIntersectFilterFromGeometry} function can be called
inside an \texttt{RTCIntersectFunctionN} user geometry callback function
to invoke the intersection filter registered to the geometry. For this
an \texttt{RTCFilterFunctionNArguments} structure must be created (see
\texttt{rtcSetGeometryIntersectFilterFunction}) which basically consists
of a valid mask, a hit packet to filter, the corresponding ray packet,
and the packet size. After the invocation of
\texttt{rtcInvokeIntersectFilterFromGeometry}, only rays that are still
valid (valid mask set to -1) should update a hit.

\subsubsection{EXIT STATUS}\label{exit-status}

For performance reasons this function does not do any error checks, thus
will not set any error flags on failure.

\subsubsection{SEE ALSO}\label{see-also}

\hyperref[rtcinvokeoccludedfilterfromgeometry]{rtcInvokeOccludedFilterFromGeometry},
\hyperref[rtcsetgeometryintersectfunction]{rtcSetGeometryIntersectFunction}

\pagebreak

\section{rtcInvokeOccludedFilterFromGeometry}\label{rtcinvokeoccludedfilterfromgeometry}

\subsubsection{NAME}\label{name}

\begin{Shaded}
\begin{Highlighting}[]
\NormalTok{rtcInvokeOccludedFilterFromGeometry }\OperatorTok{{-}}\NormalTok{ invokes the occlusion}
\NormalTok{  filter function from the geometry}
\end{Highlighting}
\end{Shaded}

\subsubsection{SYNOPSIS}\label{synopsis}

\begin{Shaded}
\begin{Highlighting}[]
\PreprocessorTok{\#include }\ImportTok{\textless{}embree4/rtcore.h\textgreater{}}

\DataTypeTok{void}\NormalTok{ rtcInvokeOccludedFilterFromGeometry}\OperatorTok{(}
  \AttributeTok{const} \KeywordTok{struct}\NormalTok{ RTCOccludedFunctionNArguments}\OperatorTok{*}\NormalTok{ args}\OperatorTok{,}
  \AttributeTok{const} \KeywordTok{struct}\NormalTok{ RTCFilterFunctionNArguments}\OperatorTok{*}\NormalTok{ filterArgs}
\OperatorTok{);}
\end{Highlighting}
\end{Shaded}

\subsubsection{DESCRIPTION}\label{description}

The \texttt{rtcInvokeOccludedFilterFromGeometry} function can be called
inside an \texttt{RTCOccludedFunctionN} user geometry callback function
to invoke the occlusion filter registered to the geometry. For this an
\texttt{RTCFilterFunctionNArguments} structure must be created (see
\texttt{rtcSetGeometryIntersectFilterFunction}) which basically consists
of a valid mask, a hit packet to filter, the corresponding ray packet,
and the packet size. After the invocation of
\texttt{rtcInvokeOccludedFilterFromGeometry} only rays that are still
valid (valid mask set to -1) should signal an occlusion.

\subsubsection{EXIT STATUS}\label{exit-status}

For performance reasons this function does not do any error checks, thus
will not set any error flags on failure.

\subsubsection{SEE ALSO}\label{see-also}

\hyperref[rtcinvokeintersectfilterfromgeometry]{rtcInvokeIntersectFilterFromGeometry},
\hyperref[rtcsetgeometryoccludedfunction]{rtcSetGeometryOccludedFunction}

\pagebreak

\section{rtcSetGeometryUserData}\label{rtcsetgeometryuserdata}

\subsubsection{NAME}\label{name}

\begin{Shaded}
\begin{Highlighting}[]
\NormalTok{rtcSetGeometryUserData }\OperatorTok{{-}}\NormalTok{ sets the user}\OperatorTok{{-}}\NormalTok{defined data pointer of the}
\NormalTok{  geometry}
\end{Highlighting}
\end{Shaded}

\subsubsection{SYNOPSIS}\label{synopsis}

\begin{Shaded}
\begin{Highlighting}[]
\PreprocessorTok{\#include }\ImportTok{\textless{}embree4/rtcore.h\textgreater{}}

\DataTypeTok{void}\NormalTok{ rtcSetGeometryUserData}\OperatorTok{(}\NormalTok{RTCGeometry geometry}\OperatorTok{,} \DataTypeTok{void}\OperatorTok{*}\NormalTok{ userPtr}\OperatorTok{);}
\end{Highlighting}
\end{Shaded}

\subsubsection{DESCRIPTION}\label{description}

The \texttt{rtcSetGeometryUserData} function sets the user-defined data
pointer (\texttt{userPtr} argument) for a geometry (\texttt{geometry}
argument). This user data pointer is intended to be pointing to the
application's representation of the geometry, and is passed to various
callback functions. The application can use this pointer inside the
callback functions to access its geometry representation.

The \texttt{rtcGetGeometryUserData} function can be used to query an
already set user data pointer of a geometry.

\subsubsection{EXIT STATUS}\label{exit-status}

On failure an error code is set that can be queried using
\texttt{rtcGetDeviceError}.

\subsubsection{SEE ALSO}\label{see-also}

\hyperref[rtcgetgeometryuserdata]{rtcGetGeometryUserData}

\pagebreak

\section{rtcGetGeometryUserData}\label{rtcgetgeometryuserdata}

\subsubsection{NAME}\label{name}

\begin{Shaded}
\begin{Highlighting}[]
\NormalTok{rtcGetGeometryUserData }\OperatorTok{{-}}\NormalTok{ returns the user data pointer}
\NormalTok{  of the geometry}
\end{Highlighting}
\end{Shaded}

\subsubsection{SYNOPSIS}\label{synopsis}

\begin{Shaded}
\begin{Highlighting}[]
\PreprocessorTok{\#include }\ImportTok{\textless{}embree4/rtcore.h\textgreater{}}

\DataTypeTok{void}\OperatorTok{*}\NormalTok{ rtcGetGeometryUserData}\OperatorTok{(}\NormalTok{RTCGeometry geometry}\OperatorTok{);}
\end{Highlighting}
\end{Shaded}

\subsubsection{DESCRIPTION}\label{description}

The \texttt{rtcGetGeometryUserData} function queries the user data
pointer previously set with \texttt{rtcSetGeometryUserData}. When
\texttt{rtcSetGeometryUserData} was not called yet, \texttt{NULL} is
returned.

This function is supposed to be used during rendering, but only
supported on the CPU and in SYCL on the GPU.

\subsubsection{EXIT STATUS}\label{exit-status}

On failure an error code is set that can be queried using
\texttt{rtcGetDeviceError}.

\subsubsection{SEE ALSO}\label{see-also}

\hyperref[rtcsetgeometryuserdata]{rtcSetGeometryUserData}

\pagebreak

\section{rtcGetGeometryUserDataFromScene}\label{rtcgetgeometryuserdatafromscene}

\subsubsection{NAME}\label{name}

\begin{Shaded}
\begin{Highlighting}[]
\NormalTok{rtcGetGeometryUserDataFromScene }\OperatorTok{{-}}\NormalTok{ returns the user data pointer}
\NormalTok{  of the geometry through the scene object}
\end{Highlighting}
\end{Shaded}

\subsubsection{SYNOPSIS}\label{synopsis}

\begin{Shaded}
\begin{Highlighting}[]
\PreprocessorTok{\#include }\ImportTok{\textless{}embree4/rtcore.h\textgreater{}}

\DataTypeTok{void}\OperatorTok{*}\NormalTok{ rtcGetGeometryUserDataFromScene}\OperatorTok{(}\NormalTok{RTCScene scene}\OperatorTok{,} \DataTypeTok{unsigned} \DataTypeTok{int}\NormalTok{ geomID}\OperatorTok{);}
\end{Highlighting}
\end{Shaded}

\subsubsection{DESCRIPTION}\label{description}

The \texttt{rtcGetGeometryUserDataFromScene} function queries the user
data pointer previously set with \texttt{rtcSetGeometryUserData} from
the geometry with index \texttt{geomID} from the specified scene
\texttt{scene}. When \texttt{rtcSetGeometryUserData} was not called yet,
\texttt{NULL} is returned.

This function is supposed to be used during rendering.

\subsubsection{EXIT STATUS}\label{exit-status}

On failure an error code is set that can be queried using
\texttt{rtcGetDeviceError}.

\subsubsection{SEE ALSO}\label{see-also}

\hyperref[rtcsetgeometryuserdata]{rtcSetGeometryUserData},
\hyperref[rtcgetgeometryuserdata]{rtcGetGeometryUserData}

\pagebreak

\section{rtcSetGeometryUserPrimitiveCount}\label{rtcsetgeometryuserprimitivecount}

\subsubsection{NAME}\label{name}

\begin{Shaded}
\begin{Highlighting}[]
\NormalTok{rtcSetGeometryUserPrimitiveCount }\OperatorTok{{-}}\NormalTok{ sets the number of primitives}
\NormalTok{  of a user}\OperatorTok{{-}}\NormalTok{defined geometry}
\end{Highlighting}
\end{Shaded}

\subsubsection{SYNOPSIS}\label{synopsis}

\begin{Shaded}
\begin{Highlighting}[]
\PreprocessorTok{\#include }\ImportTok{\textless{}embree4/rtcore.h\textgreater{}}

\DataTypeTok{void}\NormalTok{ rtcSetGeometryUserPrimitiveCount}\OperatorTok{(}
\NormalTok{  RTCGeometry geometry}\OperatorTok{,}
  \DataTypeTok{unsigned} \DataTypeTok{int}\NormalTok{ userPrimitiveCount}
\OperatorTok{);}
\end{Highlighting}
\end{Shaded}

\subsubsection{DESCRIPTION}\label{description}

The \texttt{rtcSetGeometryUserPrimitiveCount} function sets the number
of user-defined primitives (\texttt{userPrimitiveCount} parameter) of
the specified user-defined geometry (\texttt{geometry} parameter).

\subsubsection{EXIT STATUS}\label{exit-status}

On failure an error code is set that can be queried using
\texttt{rtcGetDeviceError}.

\subsubsection{SEE ALSO}\label{see-also}

\hyperref[rtc_geometry_type_user]{RTC\_GEOMETRY\_TYPE\_USER}

\pagebreak

\section{rtcSetGeometryBoundsFunction}\label{rtcsetgeometryboundsfunction}

\subsubsection{NAME}\label{name}

\begin{Shaded}
\begin{Highlighting}[]
\NormalTok{rtcSetGeometryBoundsFunction }\OperatorTok{{-}}\NormalTok{ sets a callback to query the}
\NormalTok{  bounding box of user}\OperatorTok{{-}}\NormalTok{defined primitives}
\end{Highlighting}
\end{Shaded}

\subsubsection{SYNOPSIS}\label{synopsis}

\begin{Shaded}
\begin{Highlighting}[]
\PreprocessorTok{\#include }\ImportTok{\textless{}embree4/rtcore.h\textgreater{}}

\KeywordTok{struct}\NormalTok{ RTCBoundsFunctionArguments}
\OperatorTok{\{}
  \DataTypeTok{void}\OperatorTok{*}\NormalTok{ geometryUserPtr}\OperatorTok{;}
  \DataTypeTok{unsigned} \DataTypeTok{int}\NormalTok{ primID}\OperatorTok{;}
  \DataTypeTok{unsigned} \DataTypeTok{int}\NormalTok{ timeStep}\OperatorTok{;}
  \KeywordTok{struct}\NormalTok{ RTCBounds}\OperatorTok{*}\NormalTok{ bounds\_o}\OperatorTok{;}
\OperatorTok{\};}

\KeywordTok{typedef} \DataTypeTok{void} \OperatorTok{(*}\NormalTok{RTCBoundsFunction}\OperatorTok{)(}
  \AttributeTok{const} \KeywordTok{struct}\NormalTok{ RTCBoundsFunctionArguments}\OperatorTok{*}\NormalTok{ args}
\OperatorTok{);}

\DataTypeTok{void}\NormalTok{ rtcSetGeometryBoundsFunction}\OperatorTok{(}
\NormalTok{  RTCGeometry geometry}\OperatorTok{,}
\NormalTok{  RTCBoundsFunction bounds}\OperatorTok{,}
  \DataTypeTok{void}\OperatorTok{*}\NormalTok{ userPtr}
\OperatorTok{);}
\end{Highlighting}
\end{Shaded}

\subsubsection{DESCRIPTION}\label{description}

The \texttt{rtcSetGeometryBoundsFunction} function registers a bounding
box callback function (\texttt{bounds} argument) with payload
(\texttt{userPtr} argument) for the specified user geometry
(\texttt{geometry} argument).

Only a single callback function can be registered per geometry, and
further invocations overwrite the previously set callback function.
Passing \texttt{NULL} as function pointer disables the registered
callback function.

In SYCL mode the BVH construction is done on the host and the passed
function pointer must be a host-side function pointer.

The registered bounding box callback function is invoked to calculate
axis-aligned bounding boxes of the primitives of the user-defined
geometry during spatial acceleration structure construction. The
bounding box callback of \texttt{RTCBoundsFunction} type is invoked with
a pointer to a structure of type \texttt{RTCBoundsFunctionArguments}
which contains various arguments, such as: the user data of the geometry
(\texttt{geometryUserPtr} member), the ID of the primitive to calculate
the bounds for (\texttt{primID} member), the time step at which to
calculate the bounds (\texttt{timeStep} member), and a memory location
to write the calculated bound to (\texttt{bounds\_o} member).

In a typical usage scenario one would store a pointer to the internal
representation of the user geometry object using
\texttt{rtcSetGeometryUserData}. The callback function can then read
that pointer from the \texttt{geometryUserPtr} field and calculate the
proper bounding box for the requested primitive and time, and store that
bounding box to the destination structure (\texttt{bounds\_o} member).

\subsubsection{EXIT STATUS}\label{exit-status}

On failure an error code is set that can be queried using
\texttt{rtcGetDeviceError}.

\subsubsection{SEE ALSO}\label{see-also}

\hyperref[rtc_geometry_type_user]{RTC\_GEOMETRY\_TYPE\_USER}

\pagebreak

\section{rtcSetGeometryIntersectFunction}\label{rtcsetgeometryintersectfunction}

\subsubsection{NAME}\label{name}

\begin{Shaded}
\begin{Highlighting}[]
\NormalTok{rtcSetGeometryIntersectFunction }\OperatorTok{{-}}\NormalTok{ sets the callback function to}
\NormalTok{  intersect a user geometry}
\end{Highlighting}
\end{Shaded}

\subsubsection{SYNOPSIS}\label{synopsis}

\begin{Shaded}
\begin{Highlighting}[]
\PreprocessorTok{\#include }\ImportTok{\textless{}embree4/rtcore.h\textgreater{}}

\KeywordTok{struct}\NormalTok{ RTCIntersectFunctionNArguments}
\OperatorTok{\{}
  \DataTypeTok{int}\OperatorTok{*}\NormalTok{ valid}\OperatorTok{;}
  \DataTypeTok{void}\OperatorTok{*}\NormalTok{ geometryUserPtr}\OperatorTok{;}
  \DataTypeTok{unsigned} \DataTypeTok{int}\NormalTok{ primID}\OperatorTok{;}
  \KeywordTok{struct}\NormalTok{ RTCRayQueryContext}\OperatorTok{*}\NormalTok{ context}\OperatorTok{;}
  \KeywordTok{struct}\NormalTok{ RTCRayHitN}\OperatorTok{*}\NormalTok{ rayhit}\OperatorTok{;}
  \DataTypeTok{unsigned} \DataTypeTok{int}\NormalTok{ N}\OperatorTok{;}
  \DataTypeTok{unsigned} \DataTypeTok{int}\NormalTok{ geomID}\OperatorTok{;}
\OperatorTok{\};}

\KeywordTok{typedef} \DataTypeTok{void} \OperatorTok{(*}\NormalTok{RTCIntersectFunctionN}\OperatorTok{)(}
  \AttributeTok{const} \KeywordTok{struct}\NormalTok{ RTCIntersectFunctionNArguments}\OperatorTok{*}\NormalTok{ args}
\OperatorTok{);}

\DataTypeTok{void}\NormalTok{ rtcSetGeometryIntersectFunction}\OperatorTok{(}
\NormalTok{  RTCGeometry geometry}\OperatorTok{,}
\NormalTok{  RTCIntersectFunctionN intersect}
\OperatorTok{);}
\end{Highlighting}
\end{Shaded}

\subsubsection{DESCRIPTION}\label{description}

The \texttt{rtcSetGeometryIntersectFunction} function registers a
ray/primitive intersection callback function (\texttt{intersect}
argument) for the specified user geometry (\texttt{geometry} argument).

Only a single callback function can be registered per geometry and
further invocations overwrite the previously set callback function.
Passing \texttt{NULL} as function pointer disables the registered
callback function.

The registered callback function is invoked by
\texttt{rtcIntersect}-type ray queries to calculate the intersection of
a ray packet of variable size with one user-defined primitive. The
callback function of type \texttt{RTCIntersectFunctionN} gets passed a
number of arguments through the \texttt{RTCIntersectFunctionNArguments}
structure. The value \texttt{N} specifies the ray packet size,
\texttt{valid} points to an array of integers that specify whether the
corresponding ray is valid (-1) or invalid (0), the
\texttt{geometryUserPtr} member points to the geometry user data
previously set through \texttt{rtcSetGeometryUserData}, the
\texttt{context} member points to the ray query context passed to the
ray query, the \texttt{rayhit} member points to a ray and hit packet of
variable size \texttt{N}, and the \texttt{geomID} and \texttt{primID}
member identifies the geometry ID and primitive ID of the primitive to
intersect.

The \texttt{ray} component of the \texttt{rayhit} structure contains
valid data, in particular the \texttt{tfar} value is the current closest
hit distance found. All data inside the \texttt{hit} component of the
\texttt{rayhit} structure are undefined and should not be read by the
function.

The task of the callback function is to intersect each active ray from
the ray packet with the specified user primitive. If the user-defined
primitive is missed by a ray of the ray packet, the function should
return without modifying the ray or hit. If an intersection of the
user-defined primitive with the ray was found in the valid range (from
\texttt{tnear} to \texttt{tfar}), it should update the hit distance of
the ray (\texttt{tfar} member) and the hit (\texttt{u}, \texttt{v},
\texttt{Ng}, \texttt{instID}, \texttt{geomID}, \texttt{primID} members).
In particular, the currently intersected instance is stored in the
\texttt{instID} field of the ray query context, which must be deep
copied into the \texttt{instID} member of the hit.

As a primitive might have multiple intersections with a ray, the
intersection filter function needs to be invoked by the user geometry
intersection callback for each encountered intersection, if filtering of
intersections is desired. This can be achieved through the
\texttt{rtcInvokeIntersectFilterFromGeometry} call.

Within the user geometry intersect function, it is safe to trace new
rays and create new scenes and geometries.

When performing ray queries using \texttt{rtcIntersect1}, it is
guaranteed that the packet size is 1 when the callback is invoked. When
performing ray queries using the \texttt{rtcIntersect4/8/16} functions,
it is not generally guaranteed that the ray packet size (and order of
rays inside the packet) passed to the callback matches the initial ray
packet. However, under some circumstances these properties are
guaranteed, and whether this is the case can be queried using
\texttt{rtcGetDeviceProperty}.

For many usage scenarios, repacking and re-ordering of rays does not
cause difficulties in implementing the callback function. However,
algorithms that need to extend the ray with additional data must use the
\texttt{rayID} component of the ray to identify the original ray to
access the per-ray data.

\subsubsection{EXIT STATUS}\label{exit-status}

On failure an error code is set that can be queried using
\texttt{rtcGetDeviceError}.

\subsubsection{SEE ALSO}\label{see-also}

\hyperref[rtcsetgeometryoccludedfunction]{rtcSetGeometryOccludedFunction},
\hyperref[rtcsetgeometryuserdata]{rtcSetGeometryUserData},
\hyperref[rtcinvokeintersectfilterfromgeometry]{rtcInvokeIntersectFilterFromGeometry}

\pagebreak

\section{rtcSetGeometryOccludedFunction}\label{rtcsetgeometryoccludedfunction}

\subsubsection{NAME}\label{name}

\begin{Shaded}
\begin{Highlighting}[]
\NormalTok{rtcSetGeometryOccludedFunction }\OperatorTok{{-}}\NormalTok{ sets the callback function to}
\NormalTok{  test a user geometry }\ControlFlowTok{for}\NormalTok{ occlusion}
\end{Highlighting}
\end{Shaded}

\subsubsection{SYNOPSIS}\label{synopsis}

\begin{Shaded}
\begin{Highlighting}[]
\PreprocessorTok{\#include }\ImportTok{\textless{}embree4/rtcore.h\textgreater{}}

\KeywordTok{struct}\NormalTok{ RTCOccludedFunctionNArguments}
\OperatorTok{\{}
  \DataTypeTok{int}\OperatorTok{*}\NormalTok{ valid}\OperatorTok{;}
  \DataTypeTok{void}\OperatorTok{*}\NormalTok{ geometryUserPtr}\OperatorTok{;}
  \DataTypeTok{unsigned} \DataTypeTok{int}\NormalTok{ primID}\OperatorTok{;}
  \KeywordTok{struct}\NormalTok{ RTCRayQueryContext}\OperatorTok{*}\NormalTok{ context}\OperatorTok{;}
  \KeywordTok{struct}\NormalTok{ RTCRayN}\OperatorTok{*}\NormalTok{ ray}\OperatorTok{;}
  \DataTypeTok{unsigned} \DataTypeTok{int}\NormalTok{ N}\OperatorTok{;}
  \DataTypeTok{unsigned} \DataTypeTok{int}\NormalTok{ geomID}\OperatorTok{;}
\OperatorTok{\};}

\KeywordTok{typedef} \DataTypeTok{void} \OperatorTok{(*}\NormalTok{RTCOccludedFunctionN}\OperatorTok{)(}
  \AttributeTok{const} \KeywordTok{struct}\NormalTok{ RTCOccludedFunctionNArguments}\OperatorTok{*}\NormalTok{ args}
\OperatorTok{);}

\DataTypeTok{void}\NormalTok{ rtcSetGeometryOccludedFunction}\OperatorTok{(}
\NormalTok{  RTCGeometry geometry}\OperatorTok{,}
\NormalTok{  RTCOccludedFunctionN filter}
\OperatorTok{);}
\end{Highlighting}
\end{Shaded}

\subsubsection{DESCRIPTION}\label{description}

The \texttt{rtcSetGeometryOccludedFunction} function registers a
ray/primitive occlusion callback function (\texttt{filter} argument) for
the specified user geometry (\texttt{geometry} argument).

Only a single callback function can be registered per geometry, and
further invocations overwrite the previously set callback function.
Passing \texttt{NULL} as function pointer disables the registered
callback function.

The registered callback function is invoked by \texttt{rtcOccluded}-type
ray queries to test whether the rays of a packet of variable size are
occluded by a user-defined primitive. The callback function of type
\texttt{RTCOccludedFunctionN} gets passed a number of arguments through
the \texttt{RTCOccludedFunctionNArguments} structure. The value
\texttt{N} specifies the ray packet size, \texttt{valid} points to an
array of integers which specify whether the corresponding ray is valid
(-1) or invalid (0), the \texttt{geometryUserPtr} member points to the
geometry user data previously set through
\texttt{rtcSetGeometryUserData}, the \texttt{context} member points to
the ray query context passed to the ray query, the \texttt{ray} member
points to a ray packet of variable size \texttt{N}, and the
\texttt{geomID} and \texttt{primID} member identifies the geometry ID
and primitive ID of the primitive to intersect.

The task of the callback function is to intersect each active ray from
the ray packet with the specified user primitive. If the user-defined
primitive is missed by a ray of the ray packet, the function should
return without modifying the ray. If an intersection of the user-defined
primitive with the ray was found in the valid range (from \texttt{tnear}
to \texttt{tfar}), it should set the \texttt{tfar} member of the ray to
\texttt{-inf}.

As a primitive might have multiple intersections with a ray, the
occlusion filter function needs to be invoked by the user geometry
occlusion callback for each encountered intersection, if filtering of
intersections is desired. This can be achieved through the
\texttt{rtcInvokeOccludedFilterFromGeometry} call.

Within the user geometry occlusion function, it is safe to trace new
rays and create new scenes and geometries.

When performing ray queries using \texttt{rtcOccluded1}, it is
guaranteed that the packet size is 1 when the callback is invoked. When
performing ray queries using the \texttt{rtcOccluded4/8/16} functions,
it is not generally guaranteed that the ray packet size (and order of
rays inside the packet) passed to the callback matches the initial ray
packet. However, under some circumstances these properties are
guaranteed, and whether this is the case can be queried using
\texttt{rtcGetDeviceProperty}.

For many usage scenarios, repacking and re-ordering of rays does not
cause difficulties in implementing the callback function. However,
algorithms that need to extend the ray with additional data must use the
\texttt{rayID} component of the ray to identify the original ray to
access the per-ray data.

\subsubsection{EXIT STATUS}\label{exit-status}

On failure an error code is set that can be queried using
\texttt{rtcGetDeviceError}.

\subsubsection{SEE ALSO}\label{see-also}

\hyperref[rtcsetgeometryintersectfunction]{rtcSetGeometryIntersectFunction},
\hyperref[rtcsetgeometryuserdata]{rtcSetGeometryUserData},
\hyperref[rtcinvokeoccludedfilterfromgeometry]{rtcInvokeOccludedFilterFromGeometry}

\pagebreak

\section{rtcSetGeometryPointQueryFunction}\label{rtcsetgeometrypointqueryfunction}

\subsubsection{NAME}\label{name}

\begin{Shaded}
\begin{Highlighting}[]
\NormalTok{rtcSetGeometryPointQueryFunction }\OperatorTok{{-}}\NormalTok{ sets the point query callback function}
  \ControlFlowTok{for}\NormalTok{ a geometry}
\end{Highlighting}
\end{Shaded}

\subsubsection{SYNOPSIS}\label{synopsis}

\begin{Shaded}
\begin{Highlighting}[]
\PreprocessorTok{\#include }\ImportTok{\textless{}embree4/rtcore.h\textgreater{}}

\KeywordTok{struct}\NormalTok{ RTCPointQueryFunctionArguments}
\OperatorTok{\{}
  \CommentTok{// the (world space) query object that was passed as an argument of rtcPointQuery.}
  \KeywordTok{struct}\NormalTok{ RTCPointQuery}\OperatorTok{*}\NormalTok{ query}\OperatorTok{;}

  \CommentTok{// used for user input/output data. Will not be read or modified internally.}
  \DataTypeTok{void}\OperatorTok{*}\NormalTok{ userPtr}\OperatorTok{;}

  \CommentTok{// primitive and geometry ID of primitive}
  \DataTypeTok{unsigned} \DataTypeTok{int}\NormalTok{  primID}\OperatorTok{;}        
  \DataTypeTok{unsigned} \DataTypeTok{int}\NormalTok{  geomID}\OperatorTok{;}    

  \CommentTok{// the context with transformation and instance ID stack}
  \KeywordTok{struct}\NormalTok{ RTCPointQueryContext}\OperatorTok{*}\NormalTok{ context}\OperatorTok{;}

  \CommentTok{// scaling factor indicating whether the current instance transformation}
  \CommentTok{// is a similarity transformation.}
  \DataTypeTok{float}\NormalTok{ similarityScale}\OperatorTok{;}
\OperatorTok{\};}

\KeywordTok{typedef} \DataTypeTok{bool} \OperatorTok{(*}\NormalTok{RTCPointQueryFunction}\OperatorTok{)(}
  \KeywordTok{struct}\NormalTok{ RTCPointQueryFunctionArguments}\OperatorTok{*}\NormalTok{ args}
\OperatorTok{);}

\DataTypeTok{void}\NormalTok{ rtcSetGeometryPointQueryFunction}\OperatorTok{(}
\NormalTok{  RTCGeometry geometry}\OperatorTok{,}
\NormalTok{  RTCPointQueryFunction queryFunc}
\OperatorTok{);}
\end{Highlighting}
\end{Shaded}

\subsubsection{DESCRIPTION}\label{description}

The \texttt{rtcSetGeometryPointQueryFunction} function registers a point
query callback function (\texttt{queryFunc} argument) for the specified
geometry (\texttt{geometry} argument).

Only a single callback function can be registered per geometry and
further invocations overwrite the previously set callback function.
Passing \texttt{NULL} as function pointer disables the registered
callback function.

The registered callback function is invoked by
\hyperref[rtcpointquery]{rtcPointQuery} for every primitive of the
geometry that intersects the corresponding point query domain. The
callback function of type \texttt{RTCPointQueryFunction} gets passed a
number of arguments through the \texttt{RTCPointQueryFunctionArguments}
structure. The \texttt{query} object is the original point query object
passed into \hyperref[rtcpointquery]{rtcPointQuery}, \texttt{usrPtr} is
an arbitrary pointer to pass input into and store results of the
callback function. The \texttt{primID}, \texttt{geomID} and
\texttt{context} (see
\hyperref[rtcinitpointquerycontext]{rtcInitPointQueryContext} for
details) can be used to identify the geometry data of the primitive.

A \texttt{RTCPointQueryFunction} can also be passed directly as an
argument to \hyperref[rtcpointquery]{rtcPointQuery}. In this case the
callback is invoked for all primitives in the scene that intersect the
query domain. If a callback function is passed as an argument to
\hyperref[rtcpointquery]{rtcPointQuery} and (a potentially different)
callback function is set for a geometry with
\hyperref[rtcsetgeometrypointqueryfunction]{rtcSetGeometryPointQueryFunction}
both callback functions are invoked and the callback function passed to
\hyperref[rtcpointquery]{rtcPointQuery} will be called before the
geometry specific callback function.

If instancing is used, the parameter \texttt{simliarityScale} indicates
whether the current instance transform (top element of the stack in
\texttt{context}) is a similarity transformation or not. Similarity
transformations are composed of translation, rotation and uniform
scaling and if a matrix M defines a similarity transformation, there is
a scaling factor D such that for all x,y: dist(Mx, My) = D * dist(x, y).
In this case the parameter \texttt{scalingFactor} is this scaling factor
D and otherwise it is 0. A valid similarity scale
(\texttt{similarityScale} \textgreater{} 0) allows to compute distance
information in instance space and scale the distances into world space
(for example, to update the query radius, see below) by dividing the
instance space distance with the similarity scale. If the current
instance transform is not a similarity transform
(\texttt{similarityScale} is 0), the distance computation has to be
performed in world space to ensure correctness. In this case the
instance to world transformations given with the \texttt{context} should
be used to transform the primitive data into world space. Otherwise, the
query location can be transformed into instance space which can be more
efficient. If there is no instance transform, the similarity scale is 1.

The callback function will potentially be called for primitives outside
the query domain for two reasons: First, the callback is invoked for all
primitives inside a BVH leaf node since no geometry data of primitives
is determined internally and therefore individual primitives are not
culled (only their (aggregated) bounding boxes). Second, in case non
similarity transformations are used, the resulting ellipsoidal query
domain (in instance space) is approximated by its axis aligned bounding
box internally and therefore inner nodes that do not intersect the
original domain might intersect the approximative bounding box which
results in unnecessary callbacks. In any case, the callbacks are
conservative, i.e.~if a primitive is inside the query domain a callback
will be invoked but the reverse is not necessarily true.

For efficiency, the radius of the \texttt{query} object can be decreased
(in world space) inside the callback function to improve culling of
geometry during BVH traversal. If the query radius was updated, the
callback function should return \texttt{true} to issue an update of
internal traversal information. Increasing the radius or modifying the
time or position of the query results in undefined behaviour.

Within the callback function, it is safe to call
\hyperref[rtcpointquery]{rtcPointQuery} again, for example when
implementing instancing manually. In this case the instance
transformation should be pushed onto the stack in \texttt{context}.
Embree will internally compute the point query information in instance
space using the top element of the stack in \texttt{context} when
\hyperref[rtcpointquery]{rtcPointQuery} is called.

For a reference implementation of a closest point traversal of triangle
meshes using instancing and user defined instancing see the tutorial
{[}ClosestPoint{]}.

\subsubsection{SEE ALSO}\label{see-also}

\hyperref[rtcpointquery]{rtcPointQuery},
\hyperref[rtcinitpointquerycontext]{rtcInitPointQueryContext}

\pagebreak

\section{rtcGetSYCLDeviceFunctionPointer}\label{rtcgetsycldevicefunctionpointer}

\subsubsection{NAME}\label{name}

\begin{Shaded}
\begin{Highlighting}[]
\NormalTok{rtcGetSYCLDeviceFunctionPointer }\OperatorTok{{-}}\NormalTok{ obtains a device side}
\NormalTok{  function pointer }\ControlFlowTok{for}\NormalTok{ some SYCL function}
\end{Highlighting}
\end{Shaded}

\subsubsection{SYNOPSIS}\label{synopsis}

\begin{Shaded}
\begin{Highlighting}[]
\PreprocessorTok{\#include }\ImportTok{\textless{}embree4/rtcore.h\textgreater{}}

\KeywordTok{template}\OperatorTok{\textless{}}\KeywordTok{auto}\NormalTok{ F}\OperatorTok{\textgreater{}}
\KeywordTok{inline} \KeywordTok{decltype}\OperatorTok{(}\NormalTok{F}\OperatorTok{)}\NormalTok{ rtcGetSYCLDeviceFunctionPointer}\OperatorTok{(}\NormalTok{sycl}\OperatorTok{::}\NormalTok{queue}\OperatorTok{\&}\NormalTok{ queue}\OperatorTok{);}
\end{Highlighting}
\end{Shaded}

\subsubsection{DESCRIPTION}\label{description}

This function returns a device side function pointer for some function
F. This function F must be defined using the
\texttt{RTC\_SYCL\_INDIRECTLY\_CALLABLE} attribute, e.g.:

\begin{Shaded}
\begin{Highlighting}[]
\NormalTok{RTC\_SYCL\_INDIRECTLY\_CALLABLE }\DataTypeTok{void}\NormalTok{ filter}\OperatorTok{(}
  \AttributeTok{const}\NormalTok{ RTCFilterFunctionNArguments}\OperatorTok{*}\NormalTok{ args}\OperatorTok{)} \OperatorTok{\{} \OperatorTok{...} \OperatorTok{\}}

\NormalTok{RTCFilterFunctionN fptr }\OperatorTok{=}\NormalTok{ rtcGetSYCLDeviceFunctionPointer}\OperatorTok{\textless{}}\NormalTok{filter}\OperatorTok{\textgreater{}(}\NormalTok{queue}\OperatorTok{);}
\end{Highlighting}
\end{Shaded}

Such a device side function pointers of some filter callbacks can get
assigned to a geometry using the
\texttt{rtcSetGeometryIntersectFilterFunction} and
\texttt{rtcSetGeometryOccludedFilterFunction} API functions.

Further, device side function pointers for user geometry callbacks can
be assigned to geometries using the
\texttt{rtcSetGeometryIntersectFunction} and
\texttt{rtcSetGeometryOccludedFunction} API calls.

These geometry versions of the callback functions are disabled in SYCL
by default, and we recommend not using them for performance reasons.

\subsubsection{EXIT STATUS}\label{exit-status}

On failure an error code is set that can be queried using
\texttt{rtcGetDeviceError}.

\subsubsection{SEE ALSO}\label{see-also}

\hyperref[rtcsetgeometryintersectfunction]{rtcSetGeometryIntersectFunction},
\hyperref[rtcsetgeometryoccludedfunction]{rtcSetGeometryOccludedFunction},
\hyperref[rtcsetgeometryintersectfilterfunction]{rtcSetGeometryIntersectFilterFunction},
\hyperref[rtcsetgeometryoccludedfilterfunction]{rtcSetGeometryOccludedFilterFunction}

\pagebreak

\section{rtcSetGeometryInstancedScene}\label{rtcsetgeometryinstancedscene}

\subsubsection{NAME}\label{name}

\begin{Shaded}
\begin{Highlighting}[]
\NormalTok{rtcSetGeometryInstancedScene }\OperatorTok{{-}}\NormalTok{ sets the instanced scene of}
\NormalTok{  an instance geometry}
\end{Highlighting}
\end{Shaded}

\subsubsection{SYNOPSIS}\label{synopsis}

\begin{Shaded}
\begin{Highlighting}[]
\PreprocessorTok{\#include }\ImportTok{\textless{}embree4/rtcore.h\textgreater{}}

\DataTypeTok{void}\NormalTok{ rtcSetGeometryInstancedScene}\OperatorTok{(}
\NormalTok{  RTCGeometry geometry}\OperatorTok{,}
\NormalTok{  RTCScene scene}
\OperatorTok{);}
\end{Highlighting}
\end{Shaded}

\subsubsection{DESCRIPTION}\label{description}

The \texttt{rtcSetGeometryInstancedScene} function sets the instanced
scene (\texttt{scene} argument) of the specified instance geometry
(\texttt{geometry} argument).

\subsubsection{EXIT STATUS}\label{exit-status}

On failure an error code is set that can be queried using
\texttt{rtcGetDeviceError}.

\subsubsection{SEE ALSO}\label{see-also}

\hyperref[rtc_geometry_type_instance]{RTC\_GEOMETRY\_TYPE\_INSTANCE},
\hyperref[rtcsetgeometrytransform]{rtcSetGeometryTransform}

\pagebreak

\section{rtcSetGeometryTransform}\label{rtcsetgeometrytransform}

\subsubsection{NAME}\label{name}

\begin{Shaded}
\begin{Highlighting}[]
\NormalTok{rtcSetGeometryTransform }\OperatorTok{{-}}\NormalTok{ sets the transformation }\ControlFlowTok{for}\NormalTok{ a particular}
\NormalTok{  time step of an instance geometry}
\end{Highlighting}
\end{Shaded}

\subsubsection{SYNOPSIS}\label{synopsis}

\begin{Shaded}
\begin{Highlighting}[]
\PreprocessorTok{\#include }\ImportTok{\textless{}embree4/rtcore.h\textgreater{}}

\DataTypeTok{void}\NormalTok{ rtcSetGeometryTransform}\OperatorTok{(}
\NormalTok{  RTCGeometry geometry}\OperatorTok{,}
  \DataTypeTok{unsigned} \DataTypeTok{int}\NormalTok{ timeStep}\OperatorTok{,}
  \KeywordTok{enum}\NormalTok{ RTCFormat format}\OperatorTok{,}
  \AttributeTok{const} \DataTypeTok{float}\OperatorTok{*}\NormalTok{ xfm}
\OperatorTok{);}
\end{Highlighting}
\end{Shaded}

\subsubsection{DESCRIPTION}\label{description}

The \texttt{rtcSetGeometryTransform} function sets the local-to-world
affine transformation (\texttt{xfm} parameter) of an instance geometry
(\texttt{geometry} parameter) for a particular time step
(\texttt{timeStep} parameter). The transformation is specified as a 3×4
matrix (3×3 linear transformation plus translation), for which the
following formats (\texttt{format} parameter) are supported:

\begin{itemize}
\item
  \texttt{RTC\_FORMAT\_FLOAT3X4\_ROW\_MAJOR}: The 3×4 float matrix is
  laid out in row-major form.
\item
  \texttt{RTC\_FORMAT\_FLOAT3X4\_COLUMN\_MAJOR}: The 3×4 float matrix is
  laid out in column-major form.
\item
  \texttt{RTC\_FORMAT\_FLOAT4X4\_COLUMN\_MAJOR}: The 3×4 float matrix is
  laid out in column-major form as a 4×4 homogeneous matrix with the
  last row being equal to (0, 0, 0, 1).
\end{itemize}

\subsubsection{EXIT STATUS}\label{exit-status}

On failure an error code is set that can be queried using
\texttt{rtcGetDeviceError}.

\subsubsection{SEE ALSO}\label{see-also}

\hyperref[rtc_geometry_type_instance]{RTC\_GEOMETRY\_TYPE\_INSTANCE}

\pagebreak

\section{rtcSetGeometryTransformQuaternion}\label{rtcsetgeometrytransformquaternion}

\subsubsection{NAME}\label{name}

\begin{Shaded}
\begin{Highlighting}[]
\NormalTok{rtcSetGeometryTransformQuaternion }\OperatorTok{{-}}\NormalTok{ sets the transformation }\ControlFlowTok{for}\NormalTok{ a particular}
\NormalTok{  time step of an instance geometry as a decomposition of the}
\NormalTok{  transformation matrix }\KeywordTok{using}\NormalTok{ quaternions to represent the rotation}\OperatorTok{.}
\end{Highlighting}
\end{Shaded}

\subsubsection{SYNOPSIS}\label{synopsis}

\begin{Shaded}
\begin{Highlighting}[]
\PreprocessorTok{\#include }\ImportTok{\textless{}embree4/rtcore.h\textgreater{}}

\DataTypeTok{void}\NormalTok{ rtcSetGeometryTransformQuaternion}\OperatorTok{(}
\NormalTok{  RTCGeometry geometry}\OperatorTok{,}
  \DataTypeTok{unsigned} \DataTypeTok{int}\NormalTok{ timeStep}\OperatorTok{,}
  \AttributeTok{const} \KeywordTok{struct}\NormalTok{ RTCQuaternionDecomposition}\OperatorTok{*}\NormalTok{ qd}
\OperatorTok{);}
\end{Highlighting}
\end{Shaded}

\subsubsection{DESCRIPTION}\label{description}

The \texttt{rtcSetGeometryTransformQuaternion} function sets the
local-to-world affine transformation (\texttt{qd} parameter) of an
instance geometry (\texttt{geometry} parameter) for a particular time
step (\texttt{timeStep} parameter). The transformation is specified as a
\hyperref[rtcquaterniondecomposition]{RTCQuaternionDecomposition}, which
is a decomposition of an affine transformation that represents the
rotational component of an affine transformation as a quaternion. This
allows interpolating rotational transformations exactly using spherical
linear interpolation (such as a turning wheel).

For more information about the decomposition see
\hyperref[rtcquaterniondecomposition]{RTCQuaternionDecomposition}. The
quaternion given in the \texttt{RTCQuaternionDecomposition} struct will
be normalized internally.

For correct results, the transformation matrices for all time steps must
be set either using \texttt{rtcSetGeometryTransform} or
\texttt{rtcSetGeometryTransformQuaternion}. Mixing both representations
is not allowed. Spherical linear interpolation will be used, iff the
transformation matizes are set with
\texttt{rtcSetGeometryTransformQuaternion}.

For an example of this feature see the tutorial
\hyperref[quaternion-motion-blur]{Quaternion Motion Blur}.

\subsubsection{EXIT STATUS}\label{exit-status}

On failure an error code is set that can be queried using
\texttt{rtcGetDeviceError}.

\subsubsection{SEE ALSO}\label{see-also}

\hyperref[rtcinitquaterniondecomposition]{rtcInitQuaternionDecomposition},
\hyperref[rtcsetgeometrytransform]{rtcSetGeometryTransform}

\pagebreak

\section{rtcGetGeometryTransform}\label{rtcgetgeometrytransform}

\subsubsection{NAME}\label{name}

\begin{Shaded}
\begin{Highlighting}[]
\NormalTok{rtcGetGeometryTransform }\OperatorTok{{-}}\NormalTok{ returns the interpolated instance}
\NormalTok{  transformation }\ControlFlowTok{for}\NormalTok{ the specified time}
\end{Highlighting}
\end{Shaded}

\subsubsection{SYNOPSIS}\label{synopsis}

\begin{Shaded}
\begin{Highlighting}[]
\PreprocessorTok{\#include }\ImportTok{\textless{}embree4/rtcore.h\textgreater{}}

\DataTypeTok{void}\NormalTok{ rtcGetGeometryTransform}\OperatorTok{(}
\NormalTok{  RTCGeometry geometry}\OperatorTok{,}
  \DataTypeTok{float}\NormalTok{ time}\OperatorTok{,}
  \KeywordTok{enum}\NormalTok{ RTCFormat format}\OperatorTok{,}
  \DataTypeTok{void}\OperatorTok{*}\NormalTok{ xfm}
\OperatorTok{);}
\end{Highlighting}
\end{Shaded}

\subsubsection{DESCRIPTION}\label{description}

The \texttt{rtcGetGeometryTransform} function returns the interpolated
local to world transformation (\texttt{xfm} parameter) of an instance
geometry (\texttt{geometry} parameter) for a particular time
(\texttt{time} parameter in range \([0,1]\)) in the specified format
(\texttt{format} parameter).

Possible formats for the returned matrix are:

\begin{itemize}
\item
  \texttt{RTC\_FORMAT\_FLOAT3X4\_ROW\_MAJOR}: The 3×4 float matrix is
  laid out in row-major form.
\item
  \texttt{RTC\_FORMAT\_FLOAT3X4\_COLUMN\_MAJOR}: The 3×4 float matrix is
  laid out in column-major form.
\item
  \texttt{RTC\_FORMAT\_FLOAT4X4\_COLUMN\_MAJOR}: The 3×4 float matrix is
  laid out in column-major form as a 4×4 homogeneous matrix with last
  row equal to (0, 0, 0, 1).
\end{itemize}

\subsubsection{EXIT STATUS}\label{exit-status}

On failure an error code is set that can be queried using
\texttt{rtcGetDeviceError}.

\subsubsection{SEE ALSO}\label{see-also}

\hyperref[rtc_geometry_type_instance]{RTC\_GEOMETRY\_TYPE\_INSTANCE},
\hyperref[rtcsetgeometrytransform]{rtcSetGeometryTransform}

\pagebreak

\section{rtcSetGeometryTessellationRate}\label{rtcsetgeometrytessellationrate}

\subsubsection{NAME}\label{name}

\begin{Shaded}
\begin{Highlighting}[]
\NormalTok{rtcSetGeometryTessellationRate }\OperatorTok{{-}}\NormalTok{ sets the tessellation rate of the}
\NormalTok{  geometry}
\end{Highlighting}
\end{Shaded}

\subsubsection{SYNOPSIS}\label{synopsis}

\begin{Shaded}
\begin{Highlighting}[]
\PreprocessorTok{\#include }\ImportTok{\textless{}embree4/rtcore.h\textgreater{}}

\DataTypeTok{void}\NormalTok{ rtcSetGeometryTessellationRate}\OperatorTok{(}
\NormalTok{  RTCGeometry geometry}\OperatorTok{,}
  \DataTypeTok{float}\NormalTok{ tessellationRate}
\OperatorTok{);}
\end{Highlighting}
\end{Shaded}

\subsubsection{DESCRIPTION}\label{description}

The \texttt{rtcSetGeometryTessellationRate} function sets the
tessellation rate (\texttt{tessellationRate} argument) for the specified
geometry (\texttt{geometry} argument). The tessellation rate can only be
set for flat curves and subdivision geometries. For curves, the
tessellation rate specifies the number of ray-facing quads per curve
segment. For subdivision surfaces, the tessellation rate specifies the
number of quads along each edge.

\subsubsection{EXIT STATUS}\label{exit-status}

On failure an error code is set that can be queried using
\texttt{rtcGetDeviceError}.

\subsubsection{SEE ALSO}\label{see-also}

\hyperref[rtc_geometry_type_curve]{RTC\_GEOMETRY\_TYPE\_CURVE},
\hyperref[rtc_geometry_type_subdivision]{RTC\_GEOMETRY\_TYPE\_SUBDIVISION}

\pagebreak

\section{rtcSetGeometryTopologyCount}\label{rtcsetgeometrytopologycount}

\subsubsection{NAME}\label{name}

\begin{Shaded}
\begin{Highlighting}[]
\NormalTok{rtcSetGeometryTopologyCount }\OperatorTok{{-}}\NormalTok{ sets the number of topologies of}
\NormalTok{  a subdivision geometry}
\end{Highlighting}
\end{Shaded}

\subsubsection{SYNOPSIS}\label{synopsis}

\begin{Shaded}
\begin{Highlighting}[]
\PreprocessorTok{\#include }\ImportTok{\textless{}embree4/rtcore.h\textgreater{}}

\DataTypeTok{void}\NormalTok{ rtcSetGeometryTopologyCount}\OperatorTok{(}
\NormalTok{  RTCGeometry geometry}\OperatorTok{,}
  \DataTypeTok{unsigned} \DataTypeTok{int}\NormalTok{ topologyCount}
\OperatorTok{);}
\end{Highlighting}
\end{Shaded}

\subsubsection{DESCRIPTION}\label{description}

The \texttt{rtcSetGeometryTopologyCount} function sets the number of
topologies (\texttt{topologyCount} parameter) for the specified
subdivision geometry (\texttt{geometry} parameter). The number of
topologies of a subdivision geometry must be greater or equal to 1.

To use multiple topologies, first the number of topologies must be
specified, then the individual topologies can be configured using
\texttt{rtcSetGeometrySubdivisionMode} and by setting an index buffer
(\texttt{RTC\_BUFFER\_TYPE\_INDEX}) using the topology ID as the buffer
slot.

\subsubsection{EXIT STATUS}\label{exit-status}

On failure an error code is set that can be queried using
\texttt{rtcGetDeviceError}.

\subsubsection{SEE ALSO}\label{see-also}

\hyperref[rtc_geometry_type_subdivision]{RTC\_GEOMETRY\_TYPE\_SUBDIVISION},
\hyperref[rtcsetgeometrysubdivisionmode]{rtcSetGeometrySubdivisionMode}

\pagebreak

\section{rtcSetGeometrySubdivisionMode}\label{rtcsetgeometrysubdivisionmode}

\subsubsection{NAME}\label{name}

\begin{Shaded}
\begin{Highlighting}[]
\NormalTok{rtcSetGeometrySubdivisionMode }\OperatorTok{{-}}\NormalTok{ sets the subdivision mode}
\NormalTok{  of a subdivision geometry}
\end{Highlighting}
\end{Shaded}

\subsubsection{SYNOPSIS}\label{synopsis}

\begin{Shaded}
\begin{Highlighting}[]
\PreprocessorTok{\#include }\ImportTok{\textless{}embree4/rtcore.h\textgreater{}}

\DataTypeTok{void}\NormalTok{ rtcSetGeometrySubdivisionMode}\OperatorTok{(}
\NormalTok{  RTCGeometry geometry}\OperatorTok{,}
  \DataTypeTok{unsigned} \DataTypeTok{int}\NormalTok{ topologyID}\OperatorTok{,}
  \KeywordTok{enum}\NormalTok{ RTCSubdivisionMode mode}
\OperatorTok{);}
\end{Highlighting}
\end{Shaded}

\subsubsection{DESCRIPTION}\label{description}

The \texttt{rtcSetGeometrySubdivisionMode} function sets the subdivision
mode (\texttt{mode} parameter) for the topology (\texttt{topologyID}
parameter) of the specified subdivision geometry (\texttt{geometry}
parameter).

The subdivision modes can be used to force linear interpolation for
certain parts of the subdivision mesh:

\begin{itemize}
\item
  \texttt{RTC\_SUBDIVISION\_MODE\_NO\_BOUNDARY}: Boundary patches are
  ignored. This way each rendered patch has a full set of control
  vertices.
\item
  \texttt{RTC\_SUBDIVISION\_MODE\_SMOOTH\_BOUNDARY}: The sequence of
  boundary control points are used to generate a smooth B-spline
  boundary curve (default mode).
\item
  \texttt{RTC\_SUBDIVISION\_MODE\_PIN\_CORNERS}: Corner vertices are
  pinned to their location during subdivision.
\item
  \texttt{RTC\_SUBDIVISION\_MODE\_PIN\_BOUNDARY}: All vertices at the
  border are pinned to their location during subdivision. This way the
  boundary is interpolated linearly. This mode is typically used for
  texturing to also map texels at the border of the texture to the mesh.
\item
  \texttt{RTC\_SUBDIVISION\_MODE\_PIN\_ALL}: All vertices at the border
  are pinned to their location during subdivision. This way all patches
  are linearly interpolated.
\end{itemize}

\subsubsection{EXIT STATUS}\label{exit-status}

On failure an error code is set that can be queried using
\texttt{rtcGetDeviceError}.

\subsubsection{SEE ALSO}\label{see-also}

\hyperref[rtc_geometry_type_subdivision]{RTC\_GEOMETRY\_TYPE\_SUBDIVISION}

\pagebreak

\section{rtcSetGeometryVertexAttributeTopology}\label{rtcsetgeometryvertexattributetopology}

\subsubsection{NAME}\label{name}

\begin{Shaded}
\begin{Highlighting}[]
\NormalTok{rtcSetGeometryVertexAttributeTopology }\OperatorTok{{-}}\NormalTok{ binds a vertex}
\NormalTok{  attribute to a topology of the geometry}
\end{Highlighting}
\end{Shaded}

\subsubsection{SYNOPSIS}\label{synopsis}

\begin{Shaded}
\begin{Highlighting}[]
\PreprocessorTok{\#include }\ImportTok{\textless{}embree4/rtcore.h\textgreater{}}

\DataTypeTok{void}\NormalTok{ rtcSetGeometryVertexAttributeTopology}\OperatorTok{(}
\NormalTok{  RTCGeometry geometry}\OperatorTok{,}
  \DataTypeTok{unsigned} \DataTypeTok{int}\NormalTok{ vertexAttributeID}\OperatorTok{,}
  \DataTypeTok{unsigned} \DataTypeTok{int}\NormalTok{ topologyID}
\OperatorTok{);}
\end{Highlighting}
\end{Shaded}

\subsubsection{DESCRIPTION}\label{description}

The \texttt{rtcSetGeometryVertexAttributeTopology} function binds a
vertex attribute buffer slot (\texttt{vertexAttributeID} argument) to a
topology (\texttt{topologyID} argument) for the specified subdivision
geometry (\texttt{geometry} argument). Standard vertex buffers are
always bound to the default topology (topology 0) and cannot be bound
differently. A vertex attribute buffer always uses the topology it is
bound to when used in the \texttt{rtcInterpolate} and
\texttt{rtcInterpolateN} calls.

A topology with ID \texttt{i} consists of a subdivision mode set through
\texttt{rtcSetGeometrySubdivisionMode} and the index buffer bound to the
index buffer slot \texttt{i}. This index buffer can assign indices for
each face of the subdivision geometry that are different to the indices
of the default topology. These new indices can for example be used to
introduce additional borders into the subdivision mesh to map multiple
textures onto one subdivision geometry.

\subsubsection{EXIT STATUS}\label{exit-status}

On failure an error code is set that can be queried using
\texttt{rtcGetDeviceError}.

\subsubsection{SEE ALSO}\label{see-also}

\hyperref[rtcsetgeometrysubdivisionmode]{rtcSetGeometrySubdivisionMode},
\hyperref[rtcinterpolate]{rtcInterpolate},
\hyperref[rtcinterpolaten]{rtcInterpolateN}

\pagebreak

\section{rtcSetGeometryDisplacementFunction}\label{rtcsetgeometrydisplacementfunction}

\subsubsection{NAME}\label{name}

\begin{Shaded}
\begin{Highlighting}[]
\NormalTok{rtcSetGeometryDisplacementFunction }\OperatorTok{{-}}\NormalTok{ sets the displacement function}
  \ControlFlowTok{for}\NormalTok{ a subdivision geometry}
\end{Highlighting}
\end{Shaded}

\subsubsection{SYNOPSIS}\label{synopsis}

\begin{Shaded}
\begin{Highlighting}[]
\PreprocessorTok{\#include }\ImportTok{\textless{}embree4/rtcore.h\textgreater{}}

\KeywordTok{struct}\NormalTok{ RTCDisplacementFunctionNArguments}
\OperatorTok{\{}
  \DataTypeTok{void}\OperatorTok{*}\NormalTok{ geometryUserPtr}\OperatorTok{;}
\NormalTok{  RTCGeometry geometry}\OperatorTok{;}
  \DataTypeTok{unsigned} \DataTypeTok{int}\NormalTok{ primID}\OperatorTok{;}
  \DataTypeTok{unsigned} \DataTypeTok{int}\NormalTok{ timeStep}\OperatorTok{;}
  \AttributeTok{const} \DataTypeTok{float}\OperatorTok{*}\NormalTok{ u}\OperatorTok{;}
  \AttributeTok{const} \DataTypeTok{float}\OperatorTok{*}\NormalTok{ v}\OperatorTok{;}
  \AttributeTok{const} \DataTypeTok{float}\OperatorTok{*}\NormalTok{ Ng\_x}\OperatorTok{;}
  \AttributeTok{const} \DataTypeTok{float}\OperatorTok{*}\NormalTok{ Ng\_y}\OperatorTok{;}
  \AttributeTok{const} \DataTypeTok{float}\OperatorTok{*}\NormalTok{ Ng\_z}\OperatorTok{;}
  \DataTypeTok{float}\OperatorTok{*}\NormalTok{ P\_x}\OperatorTok{;}
  \DataTypeTok{float}\OperatorTok{*}\NormalTok{ P\_y}\OperatorTok{;}
  \DataTypeTok{float}\OperatorTok{*}\NormalTok{ P\_z}\OperatorTok{;}
  \DataTypeTok{unsigned} \DataTypeTok{int}\NormalTok{ N}\OperatorTok{;}
\OperatorTok{\};}

\KeywordTok{typedef} \DataTypeTok{void} \OperatorTok{(*}\NormalTok{RTCDisplacementFunctionN}\OperatorTok{)(}
   \AttributeTok{const} \KeywordTok{struct}\NormalTok{ RTCDisplacementFunctionNArguments}\OperatorTok{*}\NormalTok{ args}
\OperatorTok{);}

\DataTypeTok{void}\NormalTok{ rtcSetGeometryDisplacementFunction}\OperatorTok{(}
\NormalTok{  RTCGeometry geometry}\OperatorTok{,}
\NormalTok{  RTCDisplacementFunctionN displacement}
\OperatorTok{);}
\end{Highlighting}
\end{Shaded}

\subsubsection{DESCRIPTION}\label{description}

The \texttt{rtcSetGeometryDisplacementFunction} function registers a
displacement callback function (\texttt{displacement} argument) for the
specified subdivision geometry (\texttt{geometry} argument).

Only a single callback function can be registered per geometry, and
further invocations overwrite the previously set callback function.
Passing \texttt{NULL} as function pointer disables the registered
callback function.

The registered displacement callback function is invoked to displace
points on the subdivision geometry during spatial acceleration structure
construction, during the \texttt{rtcCommitScene} call.

The callback function of type \texttt{RTCDisplacementFunctionN} is
invoked with a number of arguments stored inside the
\texttt{RTCDisplacementFunctionNArguments} structure. The provided user
data pointer of the geometry (\texttt{geometryUserPtr} member) can be
used to point to the application's representation of the subdivision
mesh. A number \texttt{N} of points to displace are specified in a
structure of array layout. For each point to displace, the local patch
UV coordinates (\texttt{u} and \texttt{v} arrays), the normalized
geometry normal (\texttt{Ng\_x}, \texttt{Ng\_y}, and \texttt{Ng\_z}
arrays), and the position (\texttt{P\_x}, \texttt{P\_y}, and
\texttt{P\_z} arrays) are provided. The task of the displacement
function is to use this information and change the position data.

The geometry handle (\texttt{geometry} member) and primitive ID
(\texttt{primID} member) of the patch to displace are additionally
provided as well as the time step \texttt{timeStep}, which can be
important if the displacement is time-dependent and motion blur is used.

All passed arrays must be aligned to 64 bytes and properly padded to
make wide vector processing inside the displacement function easily
possible.

Also see tutorial \hyperref[displacement-geometry]{Displacement
Geometry} for an example of how to use the displacement mapping
functions.

\subsubsection{EXIT STATUS}\label{exit-status}

On failure an error code is set that can be queried using
\texttt{rtcGetDeviceError}.

\subsubsection{SEE ALSO}\label{see-also}

\hyperref[rtc_geometry_type_subdivision]{RTC\_GEOMETRY\_TYPE\_SUBDIVISION}

\pagebreak

\section{rtcGetGeometryFirstHalfEdge}\label{rtcgetgeometryfirsthalfedge}

\subsubsection{NAME}\label{name}

\begin{Shaded}
\begin{Highlighting}[]
\NormalTok{rtcGetGeometryFirstHalfEdge }\OperatorTok{{-}}\NormalTok{ returns the first half edge of a face}
\end{Highlighting}
\end{Shaded}

\subsubsection{SYNOPSIS}\label{synopsis}

\begin{Shaded}
\begin{Highlighting}[]
\PreprocessorTok{\#include }\ImportTok{\textless{}embree4/rtcore.h\textgreater{}}

\DataTypeTok{unsigned} \DataTypeTok{int}\NormalTok{ rtcGetGeometryFirstHalfEdge}\OperatorTok{(}
\NormalTok{  RTCGeometry geometry}\OperatorTok{,}
  \DataTypeTok{unsigned} \DataTypeTok{int}\NormalTok{ faceID}
\OperatorTok{);}
\end{Highlighting}
\end{Shaded}

\subsubsection{DESCRIPTION}\label{description}

The \texttt{rtcGetGeometryFirstHalfEdge} function returns the ID of the
first half edge belonging to the specified face (\texttt{faceID}
argument). For instance in the following example the first half edge of
face \texttt{f1} is \texttt{e4}.

!\href{images/half_edges.pdf}{}

This function can only be used for subdivision geometries. As all
topologies of a subdivision geometry share the same face buffer the
function does not depend on the topology ID.

Here f0 to f7 are 8 quadrilateral faces with 4 vertices each. The edges
e0 to e23 of these faces are shown with their orientation. For each face
the ID of the edges corresponds to the slots the face occupies in the
index array of the geometry. E.g. as the indices of face f1 start at
location 4 of the index array, the first edge is edge e4, the next edge
e5, etc.

\subsubsection{EXIT STATUS}\label{exit-status}

On failure an error code is set that can be queried using
\texttt{rtcGetDeviceError}.

\subsubsection{SEE ALSO}\label{see-also}

\hyperref[rtcgetgeometryfirsthalfedge]{rtcGetGeometryFirstHalfEdge},
\hyperref[rtcgetgeometryface]{rtcGetGeometryFace},
\hyperref[rtcgetgeometryoppositehalfedge]{rtcGetGeometryOppositeHalfEdge},
\hyperref[rtcgetgeometrynexthalfedge]{rtcGetGeometryNextHalfEdge},
\hyperref[rtcgetgeometryprevioushalfedge]{rtcGetGeometryPreviousHalfEdge}

\pagebreak

\section{rtcGetGeometryFace}\label{rtcgetgeometryface}

\subsubsection{NAME}\label{name}

\begin{Shaded}
\begin{Highlighting}[]
\NormalTok{rtcGetGeometryFace }\OperatorTok{{-}}\NormalTok{ returns the face of some half edge}
\end{Highlighting}
\end{Shaded}

\subsubsection{SYNOPSIS}\label{synopsis}

\begin{Shaded}
\begin{Highlighting}[]
\PreprocessorTok{\#include }\ImportTok{\textless{}embree4/rtcore.h\textgreater{}}

\DataTypeTok{unsigned} \DataTypeTok{int}\NormalTok{ rtcGetGeometryFace}\OperatorTok{(}
\NormalTok{  RTCGeometry geometry}\OperatorTok{,}
  \DataTypeTok{unsigned} \DataTypeTok{int}\NormalTok{ edgeID}
\OperatorTok{);}
\end{Highlighting}
\end{Shaded}

\subsubsection{DESCRIPTION}\label{description}

The \texttt{rtcGetGeometryFace} function returns the ID of the face the
specified half edge (\texttt{edgeID} argument) belongs to. For instance
in the following example the face \texttt{f1} is returned for edges
\texttt{e4}, \texttt{e5}, \texttt{e6}, and \texttt{e7}.

!\href{images/half_edges.pdf}{}

This function can only be used for subdivision geometries. As all
topologies of a subdivision geometry share the same face buffer the
function does not depend on the topology ID.

\subsubsection{EXIT STATUS}\label{exit-status}

On failure an error code is set that can be queried using
\texttt{rtcGetDeviceError}.

\subsubsection{SEE ALSO}\label{see-also}

\hyperref[rtcgetgeometryfirsthalfedge]{rtcGetGeometryFirstHalfEdge},
\hyperref[rtcgetgeometryface]{rtcGetGeometryFace},
\hyperref[rtcgetgeometryoppositehalfedge]{rtcGetGeometryOppositeHalfEdge},
\hyperref[rtcgetgeometrynexthalfedge]{rtcGetGeometryNextHalfEdge},
\hyperref[rtcgetgeometryprevioushalfedge]{rtcGetGeometryPreviousHalfEdge}

\pagebreak

\section{rtcGetGeometryNextHalfEdge}\label{rtcgetgeometrynexthalfedge}

\subsubsection{NAME}\label{name}

\begin{Shaded}
\begin{Highlighting}[]
\NormalTok{rtcGetGeometryNextHalfEdge }\OperatorTok{{-}}\NormalTok{ returns the next half edge}
\end{Highlighting}
\end{Shaded}

\subsubsection{SYNOPSIS}\label{synopsis}

\begin{Shaded}
\begin{Highlighting}[]
\PreprocessorTok{\#include }\ImportTok{\textless{}embree4/rtcore.h\textgreater{}}

\DataTypeTok{unsigned} \DataTypeTok{int}\NormalTok{ rtcGetGeometryNextHalfEdge}\OperatorTok{(}
\NormalTok{  RTCGeometry geometry}\OperatorTok{,}
  \DataTypeTok{unsigned} \DataTypeTok{int}\NormalTok{ edgeID}
\OperatorTok{);}
\end{Highlighting}
\end{Shaded}

\subsubsection{DESCRIPTION}\label{description}

The \texttt{rtcGetGeometryNextHalfEdge} function returns the ID of the
next half edge of the specified half edge (\texttt{edgeID} argument).
For instance in the following example the next half edge of \texttt{e10}
is \texttt{e11}.

!\href{images/half_edges.pdf}{}

This function can only be used for subdivision geometries. As all
topologies of a subdivision geometry share the same face buffer the
function does not depend on the topology ID.

\subsubsection{EXIT STATUS}\label{exit-status}

On failure an error code is set that can be queried using
\texttt{rtcGetDeviceError}.

\subsubsection{SEE ALSO}\label{see-also}

\hyperref[rtcgetgeometryfirsthalfedge]{rtcGetGeometryFirstHalfEdge},
\hyperref[rtcgetgeometryface]{rtcGetGeometryFace},
\hyperref[rtcgetgeometryoppositehalfedge]{rtcGetGeometryOppositeHalfEdge},
\hyperref[rtcgetgeometrynexthalfedge]{rtcGetGeometryNextHalfEdge},
\hyperref[rtcgetgeometryprevioushalfedge]{rtcGetGeometryPreviousHalfEdge}

\pagebreak

\section{rtcGetGeometryPreviousHalfEdge}\label{rtcgetgeometryprevioushalfedge}

\subsubsection{NAME}\label{name}

\begin{Shaded}
\begin{Highlighting}[]
\NormalTok{rtcGetGeometryPreviousHalfEdge }\OperatorTok{{-}}\NormalTok{ returns the previous half edge}
\end{Highlighting}
\end{Shaded}

\subsubsection{SYNOPSIS}\label{synopsis}

\begin{Shaded}
\begin{Highlighting}[]
\PreprocessorTok{\#include }\ImportTok{\textless{}embree4/rtcore.h\textgreater{}}

\DataTypeTok{unsigned} \DataTypeTok{int}\NormalTok{ rtcGetGeometryPreviousHalfEdge}\OperatorTok{(}
\NormalTok{  RTCGeometry geometry}\OperatorTok{,}
  \DataTypeTok{unsigned} \DataTypeTok{int}\NormalTok{ edgeID}
\OperatorTok{);}
\end{Highlighting}
\end{Shaded}

\subsubsection{DESCRIPTION}\label{description}

The \texttt{rtcGetGeometryPreviousHalfEdge} function returns the ID of
the previous half edge of the specified half edge (\texttt{edgeID}
argument). For instance in the following example the previous half edge
of \texttt{e6} is \texttt{e5}.

!\href{images/half_edges.pdf}{}

This function can only be used for subdivision geometries. As all
topologies of a subdivision geometry share the same face buffer the
function does not depend on the topology ID.

\subsubsection{EXIT STATUS}\label{exit-status}

On failure an error code is set that can be queried using
\texttt{rtcGetDeviceError}.

\subsubsection{SEE ALSO}\label{see-also}

\hyperref[rtcgetgeometryfirsthalfedge]{rtcGetGeometryFirstHalfEdge},
\hyperref[rtcgetgeometryface]{rtcGetGeometryFace},
\hyperref[rtcgetgeometryoppositehalfedge]{rtcGetGeometryOppositeHalfEdge},
\hyperref[rtcgetgeometrynexthalfedge]{rtcGetGeometryNextHalfEdge},
\hyperref[rtcgetgeometryprevioushalfedge]{rtcGetGeometryPreviousHalfEdge}

\pagebreak

\section{rtcGetGeometryOppositeHalfEdge}\label{rtcgetgeometryoppositehalfedge}

\subsubsection{NAME}\label{name}

\begin{Shaded}
\begin{Highlighting}[]
\NormalTok{rtcGetGeometryOppositeHalfEdge }\OperatorTok{{-}}\NormalTok{ returns the opposite half edge}
\end{Highlighting}
\end{Shaded}

\subsubsection{SYNOPSIS}\label{synopsis}

\begin{Shaded}
\begin{Highlighting}[]
\PreprocessorTok{\#include }\ImportTok{\textless{}embree4/rtcore.h\textgreater{}}

\DataTypeTok{unsigned} \DataTypeTok{int}\NormalTok{ rtcGetGeometryOppositeHalfEdge}\OperatorTok{(}
\NormalTok{  RTCGeometry geometry}\OperatorTok{,}
  \DataTypeTok{unsigned} \DataTypeTok{int}\NormalTok{ topologyID}\OperatorTok{,}
  \DataTypeTok{unsigned} \DataTypeTok{int}\NormalTok{ edgeID}
\OperatorTok{);}
\end{Highlighting}
\end{Shaded}

\subsubsection{DESCRIPTION}\label{description}

The \texttt{rtcGetGeometryOppositeHalfEdge} function returns the ID of
the opposite half edge of the specified half edge (\texttt{edgeID}
argument) in the specified topology (\texttt{topologyID} argument). For
instance in the following example the opposite half edge of \texttt{e6}
is \texttt{e16}.

!\href{images/half_edges.pdf}{}

An opposite half edge does not exist if the specified half edge has
either no neighboring face, or more than 2 neighboring faces. In these
cases the function just returns the same edge \texttt{edgeID} again.

This function can only be used for subdivision geometries. The function
depends on the topology as the topologies of a subdivision geometry have
different index buffers assigned.

\subsubsection{EXIT STATUS}\label{exit-status}

On failure an error code is set that can be queried using
\texttt{rtcGetDeviceError}.

\subsubsection{SEE ALSO}\label{see-also}

\hyperref[rtcgetgeometryfirsthalfedge]{rtcGetGeometryFirstHalfEdge},
\hyperref[rtcgetgeometryface]{rtcGetGeometryFace},
\hyperref[rtcgetgeometryoppositehalfedge]{rtcGetGeometryOppositeHalfEdge},
\hyperref[rtcgetgeometrynexthalfedge]{rtcGetGeometryNextHalfEdge},
\hyperref[rtcgetgeometryprevioushalfedge]{rtcGetGeometryPreviousHalfEdge}

\pagebreak

\section{rtcInterpolate}\label{rtcinterpolate}

\subsubsection{NAME}\label{name}

\begin{Shaded}
\begin{Highlighting}[]
\NormalTok{rtcInterpolate }\OperatorTok{{-}}\NormalTok{ interpolates vertex attributes}
\end{Highlighting}
\end{Shaded}

\subsubsection{SYNOPSIS}\label{synopsis}

\begin{Shaded}
\begin{Highlighting}[]
\PreprocessorTok{\#include }\ImportTok{\textless{}embree4/rtcore.h\textgreater{}}

\KeywordTok{struct}\NormalTok{ RTCInterpolateArguments}
\OperatorTok{\{}
\NormalTok{  RTCGeometry geometry}\OperatorTok{;}
  \DataTypeTok{unsigned} \DataTypeTok{int}\NormalTok{ primID}\OperatorTok{;}
  \DataTypeTok{float}\NormalTok{ u}\OperatorTok{;}
  \DataTypeTok{float}\NormalTok{ v}\OperatorTok{;}
  \KeywordTok{enum}\NormalTok{ RTCBufferType bufferType}\OperatorTok{;}
  \DataTypeTok{unsigned} \DataTypeTok{int}\NormalTok{ bufferSlot}\OperatorTok{;}
  \DataTypeTok{float}\OperatorTok{*}\NormalTok{ P}\OperatorTok{;}
  \DataTypeTok{float}\OperatorTok{*}\NormalTok{ dPdu}\OperatorTok{;}
  \DataTypeTok{float}\OperatorTok{*}\NormalTok{ dPdv}\OperatorTok{;}
  \DataTypeTok{float}\OperatorTok{*}\NormalTok{ ddPdudu}\OperatorTok{;}
  \DataTypeTok{float}\OperatorTok{*}\NormalTok{ ddPdvdv}\OperatorTok{;}
  \DataTypeTok{float}\OperatorTok{*}\NormalTok{ ddPdudv}\OperatorTok{;}
  \DataTypeTok{unsigned} \DataTypeTok{int}\NormalTok{ valueCount}\OperatorTok{;}
\OperatorTok{\};}

\DataTypeTok{void}\NormalTok{ rtcInterpolate}\OperatorTok{(}
  \AttributeTok{const} \KeywordTok{struct}\NormalTok{ RTCInterpolateArguments}\OperatorTok{*}\NormalTok{ args}
\OperatorTok{);}
\end{Highlighting}
\end{Shaded}

\subsubsection{DESCRIPTION}\label{description}

The \texttt{rtcInterpolate} function smoothly interpolates per-vertex
data over the geometry. This interpolation is supported for triangle
meshes, quad meshes, curve geometries, and subdivision geometries. Apart
from interpolating the vertex attribute itself, it is also possible to
get the first and second order derivatives of that value. This
interpolation ignores displacements of subdivision surfaces and always
interpolates the underlying base surface.

The \texttt{rtcInterpolate} call gets passed a number of arguments
inside a structure of type \texttt{RTCInterpolateArguments}. For some
geometry (\texttt{geometry} parameter) this function smoothly
interpolates the per-vertex data stored inside the specified geometry
buffer (\texttt{bufferType} and \texttt{bufferSlot} parameters) to the
u/v location (\texttt{u} and \texttt{v} parameters) of the primitive
(\texttt{primID} parameter). The number of floating point values to
interpolate and store to the destination arrays can be specified using
the \texttt{valueCount} parameter. As interpolation buffer, one can
specify vertex buffers (\texttt{RTC\_BUFFER\_TYPE\_VERTEX}) and vertex
attribute buffers (\texttt{RTC\_BUFFER\_TYPE\_VERTEX\_ATTRIBUTE}) as
well.

The \texttt{rtcInterpolate} call stores \texttt{valueCount} number of
interpolated floating point values to the memory location pointed to by
\texttt{P}. One can avoid storing the interpolated value by setting
\texttt{P} to \texttt{NULL}.

The first order derivative of the interpolation by u and v are stored at
the \texttt{dPdu} and \texttt{dPdv} memory locations. One can avoid
storing first order derivatives by setting both \texttt{dPdu} and
\texttt{dPdv} to \texttt{NULL}.

The second order derivatives are stored at the \texttt{ddPdudu},
\texttt{ddPdvdv}, and \texttt{ddPdudv} memory locations. One can avoid
storing second order derivatives by setting these three pointers to
\texttt{NULL}.

To use \texttt{rtcInterpolate} for a geometry, all changes to that
geometry must be properly committed using \texttt{rtcCommitGeometry}.

All input buffers and output arrays must be padded to 16 bytes, as the
implementation uses 16-byte SSE instructions to read and write into
these buffers.

See tutorial \hyperref[interpolation]{Interpolation} for an example of
using the \texttt{rtcInterpolate} function.

\subsubsection{EXIT STATUS}\label{exit-status}

For performance reasons this function does not do any error checks, thus
will not set any error flags on failure.

\subsubsection{SEE ALSO}\label{see-also}

\hyperref[rtcinterpolaten]{rtcInterpolateN}

\pagebreak

\section{rtcInterpolateN}\label{rtcinterpolaten}

\subsubsection{NAME}\label{name}

\begin{Shaded}
\begin{Highlighting}[]
\NormalTok{rtcInterpolateN }\OperatorTok{{-}}\NormalTok{ performs N interpolations of vertex attribute data}
\end{Highlighting}
\end{Shaded}

\subsubsection{SYNOPSIS}\label{synopsis}

\begin{Shaded}
\begin{Highlighting}[]
\PreprocessorTok{\#include }\ImportTok{\textless{}embree4/rtcore.h\textgreater{}}

\KeywordTok{struct}\NormalTok{ RTCInterpolateNArguments}
\OperatorTok{\{}
\NormalTok{  RTCGeometry geometry}\OperatorTok{;}
  \AttributeTok{const} \DataTypeTok{void}\OperatorTok{*}\NormalTok{ valid}\OperatorTok{;}
  \AttributeTok{const} \DataTypeTok{unsigned} \DataTypeTok{int}\OperatorTok{*}\NormalTok{ primIDs}\OperatorTok{;}
  \AttributeTok{const} \DataTypeTok{float}\OperatorTok{*}\NormalTok{ u}\OperatorTok{;}
  \AttributeTok{const} \DataTypeTok{float}\OperatorTok{*}\NormalTok{ v}\OperatorTok{;}
  \DataTypeTok{unsigned} \DataTypeTok{int}\NormalTok{ N}\OperatorTok{;}
  \KeywordTok{enum}\NormalTok{ RTCBufferType bufferType}\OperatorTok{;}
  \DataTypeTok{unsigned} \DataTypeTok{int}\NormalTok{ bufferSlot}\OperatorTok{;}
  \DataTypeTok{float}\OperatorTok{*}\NormalTok{ P}\OperatorTok{;}
  \DataTypeTok{float}\OperatorTok{*}\NormalTok{ dPdu}\OperatorTok{;}
  \DataTypeTok{float}\OperatorTok{*}\NormalTok{ dPdv}\OperatorTok{;}
  \DataTypeTok{float}\OperatorTok{*}\NormalTok{ ddPdudu}\OperatorTok{;}
  \DataTypeTok{float}\OperatorTok{*}\NormalTok{ ddPdvdv}\OperatorTok{;}
  \DataTypeTok{float}\OperatorTok{*}\NormalTok{ ddPdudv}\OperatorTok{;}
  \DataTypeTok{unsigned} \DataTypeTok{int}\NormalTok{ valueCount}\OperatorTok{;}
\OperatorTok{\};}

\DataTypeTok{void}\NormalTok{ rtcInterpolateN}\OperatorTok{(}
  \AttributeTok{const} \KeywordTok{struct}\NormalTok{ RTCInterpolateNArguments}\OperatorTok{*}\NormalTok{ args}
\OperatorTok{);}
\end{Highlighting}
\end{Shaded}

\subsubsection{DESCRIPTION}\label{description}

The \texttt{rtcInterpolateN} is similar to \texttt{rtcInterpolate}, but
performs \texttt{N} many interpolations at once. It additionally gets an
array of u/v coordinates and a valid mask (\texttt{valid} parameter)
that specifies which of these coordinates are valid. The valid mask
points to \texttt{N} integers, and a value of -1 denotes valid and 0
invalid. If the valid pointer is \texttt{NULL} all elements are
considers valid. The destination arrays are filled in structure of array
(SOA) layout. The value \texttt{N} must be divisible by 4.

To use \texttt{rtcInterpolateN} for a geometry, all changes to that
geometry must be properly committed using \texttt{rtcCommitGeometry}.

\subsubsection{EXIT STATUS}\label{exit-status}

For performance reasons this function does not do any error checks, thus
will not set any error flags on failure.

\subsubsection{SEE ALSO}\label{see-also}

\hyperref[rtcinterpolate]{rtcInterpolate}

\pagebreak

\section{rtcNewBuffer}\label{rtcnewbuffer}

\subsubsection{NAME}\label{name}

\begin{Shaded}
\begin{Highlighting}[]
\NormalTok{rtcNewBuffer }\OperatorTok{{-}}\NormalTok{ creates a }\KeywordTok{new}\NormalTok{ data buffer}
\end{Highlighting}
\end{Shaded}

\subsubsection{SYNOPSIS}\label{synopsis}

\begin{Shaded}
\begin{Highlighting}[]
\PreprocessorTok{\#include }\ImportTok{\textless{}embree4/rtcore.h\textgreater{}}

\NormalTok{RTCBuffer rtcNewBuffer}\OperatorTok{(}
\NormalTok{  RTCDevice device}\OperatorTok{,}
  \DataTypeTok{size\_t}\NormalTok{ byteSize}
\OperatorTok{);}
\end{Highlighting}
\end{Shaded}

\subsubsection{DESCRIPTION}\label{description}

The \texttt{rtcNewBuffer} function creates a new data buffer object of
specified size in bytes (\texttt{byteSize} argument) that is bound to
the specified device (\texttt{device} argument). The buffer object is
reference counted with an initial reference count of 1. The returned
buffer object can be released using the \texttt{rtcReleaseBuffer} API
call. The specified number of bytes are allocated at buffer construction
time and deallocated when the buffer is destroyed.

When the buffer will be used as a vertex buffer
(\texttt{RTC\_BUFFER\_TYPE\_VERTEX} and
\texttt{RTC\_BUFFER\_TYPE\_VERTEX\_ATTRIBUTE}), the last buffer element
must be readable using 16-byte SSE load instructions, thus padding the
last element is required for certain layouts. E.g. a standard
\texttt{float3} vertex buffer layout should add storage for at least one
more float to the end of the buffer.

\subsubsection{EXIT STATUS}\label{exit-status}

On failure \texttt{NULL} is returned and an error code is set that can
be queried using \texttt{rtcGetDeviceError}.

\subsubsection{SEE ALSO}\label{see-also}

\hyperref[rtcretainbuffer]{rtcRetainBuffer},
\hyperref[rtcreleasebuffer]{rtcReleaseBuffer}

\pagebreak

\section{rtcNewSharedBuffer}\label{rtcnewsharedbuffer}

\subsubsection{NAME}\label{name}

\begin{Shaded}
\begin{Highlighting}[]
\NormalTok{rtcNewSharedBuffer }\OperatorTok{{-}}\NormalTok{ creates a }\KeywordTok{new}\NormalTok{ shared data buffer}
\end{Highlighting}
\end{Shaded}

\subsubsection{SYNOPSIS}\label{synopsis}

\begin{Shaded}
\begin{Highlighting}[]
\PreprocessorTok{\#include }\ImportTok{\textless{}embree4/rtcore.h\textgreater{}}

\NormalTok{RTCBuffer rtcNewSharedBuffer}\OperatorTok{(}
\NormalTok{  RTCDevice device}\OperatorTok{,}
  \DataTypeTok{void}\OperatorTok{*}\NormalTok{ ptr}\OperatorTok{,}
  \DataTypeTok{size\_t}\NormalTok{ byteSize}
\OperatorTok{);}
\end{Highlighting}
\end{Shaded}

\subsubsection{DESCRIPTION}\label{description}

The \texttt{rtcNewSharedBuffer} function creates a new shared data
buffer object bound to the specified device (\texttt{device} argument).
The buffer object is reference counted with an initial reference count
of 1. The buffer can be released using the \texttt{rtcReleaseBuffer}
function.

At construction time, the pointer to the user-managed buffer data
(\texttt{ptr} argument) including its size in bytes (\texttt{byteSize}
argument) is provided to create the buffer. At buffer construction time
no buffer data is allocated, but the buffer data provided by the
application is used. The buffer data must remain valid for as long as
the buffer may be used, and the user is responsible to free the buffer
data when no longer required.

When the buffer will be used as a vertex buffer
(\texttt{RTC\_BUFFER\_TYPE\_VERTEX} and
\texttt{RTC\_BUFFER\_TYPE\_VERTEX\_ATTRIBUTE}), the last buffer element
must be readable using 16-byte SSE load instructions, thus padding the
last element is required for certain layouts. E.g. a standard
\texttt{float3} vertex buffer layout should add storage for at least one
more float to the end of the buffer.

The data pointer (\texttt{ptr} argument) must be aligned to 4 bytes;
otherwise the \texttt{rtcNewSharedBuffer} function will fail.

\subsubsection{EXIT STATUS}\label{exit-status}

On failure \texttt{NULL} is returned and an error code is set that can
be queried using \texttt{rtcGetDeviceError}.

\subsubsection{SEE ALSO}\label{see-also}

\hyperref[rtcretainbuffer]{rtcRetainBuffer},
\hyperref[rtcreleasebuffer]{rtcReleaseBuffer}

\pagebreak

\section{rtcRetainBuffer}\label{rtcretainbuffer}

\subsubsection{NAME}\label{name}

\begin{Shaded}
\begin{Highlighting}[]
\NormalTok{rtcRetainBuffer }\OperatorTok{{-}}\NormalTok{ increments the buffer reference count}
\end{Highlighting}
\end{Shaded}

\subsubsection{SYNOPSIS}\label{synopsis}

\begin{Shaded}
\begin{Highlighting}[]
\PreprocessorTok{\#include }\ImportTok{\textless{}embree4/rtcore.h\textgreater{}}

\DataTypeTok{void}\NormalTok{ rtcRetainBuffer}\OperatorTok{(}\NormalTok{RTCBuffer buffer}\OperatorTok{);}
\end{Highlighting}
\end{Shaded}

\subsubsection{DESCRIPTION}\label{description}

Buffer objects are reference counted. The \texttt{rtcRetainBuffer}
function increments the reference count of the passed buffer object
(\texttt{buffer} argument). This function together with
\texttt{rtcReleaseBuffer} allows to use the internal reference counting
in a C++ wrapper class to handle the ownership of the object.

\subsubsection{EXIT STATUS}\label{exit-status}

On failure an error code is set that can be queried using
\texttt{rtcGetDeviceError}.

\subsubsection{SEE ALSO}\label{see-also}

\hyperref[rtcnewbuffer]{rtcNewBuffer},
\hyperref[rtcreleasebuffer]{rtcReleaseBuffer}

\pagebreak

\section{rtcReleaseBuffer}\label{rtcreleasebuffer}

\subsubsection{NAME}\label{name}

\begin{Shaded}
\begin{Highlighting}[]
\NormalTok{rtcReleaseBuffer }\OperatorTok{{-}}\NormalTok{ decrements the buffer reference count}
\end{Highlighting}
\end{Shaded}

\subsubsection{SYNOPSIS}\label{synopsis}

\begin{Shaded}
\begin{Highlighting}[]
\PreprocessorTok{\#include }\ImportTok{\textless{}embree4/rtcore.h\textgreater{}}

\DataTypeTok{void}\NormalTok{ rtcReleaseBuffer}\OperatorTok{(}\NormalTok{RTCBuffer buffer}\OperatorTok{);}
\end{Highlighting}
\end{Shaded}

\subsubsection{DESCRIPTION}\label{description}

Buffer objects are reference counted. The \texttt{rtcReleaseBuffer}
function decrements the reference count of the passed buffer object
(\texttt{buffer} argument). When the reference count falls to 0, the
buffer gets destroyed.

\subsubsection{EXIT STATUS}\label{exit-status}

On failure an error code is set that can be queried using
\texttt{rtcGetDeviceError}.

\subsubsection{SEE ALSO}\label{see-also}

\hyperref[rtcnewbuffer]{rtcNewBuffer},
\hyperref[rtcretainbuffer]{rtcRetainBuffer}

\pagebreak

\section{rtcGetBufferData}\label{rtcgetbufferdata}

\subsubsection{NAME}\label{name}

\begin{Shaded}
\begin{Highlighting}[]
\NormalTok{rtcGetBufferData }\OperatorTok{{-}}\NormalTok{ gets a pointer to the buffer data}
\end{Highlighting}
\end{Shaded}

\subsubsection{SYNOPSIS}\label{synopsis}

\begin{Shaded}
\begin{Highlighting}[]
\PreprocessorTok{\#include }\ImportTok{\textless{}embree4/rtcore.h\textgreater{}}

\DataTypeTok{void}\OperatorTok{*}\NormalTok{ rtcGetBufferData}\OperatorTok{(}\NormalTok{RTCBuffer buffer}\OperatorTok{);}
\end{Highlighting}
\end{Shaded}

\subsubsection{DESCRIPTION}\label{description}

The \texttt{rtcGetBufferData} function returns a pointer to the buffer
data of the specified buffer object (\texttt{buffer} argument).

\subsubsection{EXIT STATUS}\label{exit-status}

On failure an error code is set that can be queried using
\texttt{rtcGetDeviceError}.

\subsubsection{SEE ALSO}\label{see-also}

\hyperref[rtcnewbuffer]{rtcNewBuffer}

\pagebreak

\section{RTCRay}\label{rtcray}

\subsubsection{NAME}\label{name}

\begin{Shaded}
\begin{Highlighting}[]
\NormalTok{RTCRay }\OperatorTok{{-}}\NormalTok{ single ray structure}
\end{Highlighting}
\end{Shaded}

\subsubsection{SYNOPSIS}\label{synopsis}

\begin{Shaded}
\begin{Highlighting}[]
\PreprocessorTok{\#include }\ImportTok{\textless{}embree4/rtcore\_ray.h\textgreater{}}

\KeywordTok{struct}\NormalTok{ RTC\_ALIGN}\OperatorTok{(}\DecValTok{16}\OperatorTok{)}\NormalTok{ RTCRay}
\OperatorTok{\{}
  \DataTypeTok{float}\NormalTok{ org\_x}\OperatorTok{;}        \CommentTok{// x coordinate of ray origin}
  \DataTypeTok{float}\NormalTok{ org\_y}\OperatorTok{;}        \CommentTok{// y coordinate of ray origin}
  \DataTypeTok{float}\NormalTok{ org\_z}\OperatorTok{;}        \CommentTok{// z coordinate of ray origin}
  \DataTypeTok{float}\NormalTok{ tnear}\OperatorTok{;}        \CommentTok{// start of ray segment}

  \DataTypeTok{float}\NormalTok{ dir\_x}\OperatorTok{;}        \CommentTok{// x coordinate of ray direction}
  \DataTypeTok{float}\NormalTok{ dir\_y}\OperatorTok{;}        \CommentTok{// y coordinate of ray direction}
  \DataTypeTok{float}\NormalTok{ dir\_z}\OperatorTok{;}        \CommentTok{// z coordinate of ray direction}
  \DataTypeTok{float}\NormalTok{ time}\OperatorTok{;}         \CommentTok{// time of this ray for motion blur}

  \DataTypeTok{float}\NormalTok{ tfar}\OperatorTok{;}         \CommentTok{// end of ray segment (set to hit distance)}
  \DataTypeTok{unsigned} \DataTypeTok{int}\NormalTok{ mask}\OperatorTok{;}  \CommentTok{// ray mask}
  \DataTypeTok{unsigned} \DataTypeTok{int}\NormalTok{ id}\OperatorTok{;}    \CommentTok{// ray ID}
  \DataTypeTok{unsigned} \DataTypeTok{int}\NormalTok{ flags}\OperatorTok{;} \CommentTok{// ray flags}
\OperatorTok{\};}
\end{Highlighting}
\end{Shaded}

\subsubsection{DESCRIPTION}\label{description}

The \texttt{RTCRay} structure defines the ray layout for a single ray.
The ray contains the origin (\texttt{org\_x}, \texttt{org\_y},
\texttt{org\_z} members), direction vector (\texttt{dir\_x},
\texttt{dir\_y}, \texttt{dir\_z} members), and ray segment
(\texttt{tnear} and \texttt{tfar} members). The ray direction does not
have to be normalized, and only the parameter range specified by the
\texttt{tnear}/\texttt{tfar} interval is considered valid.

The ray segment must be in the range \([0, \infty]\), thus ranges that
start behind the ray origin are not allowed, but ranges can reach to
infinity.

The ray further contains a motion blur time in the range \([0, 1]\)
(\texttt{time} member), a ray mask (\texttt{mask} member), a ray ID
(\texttt{id} member), and ray flags (\texttt{flags} member). The ray
mask can be used to mask out some geometries for some rays (see
\texttt{rtcSetGeometryMask} for more details). The ray ID can be used to
identify a ray inside a callback function, even if the order of rays
inside a ray packet has changed.

The \texttt{embree4/rtcore\_ray.h} header additionally defines the same
ray structure in structure of array (SOA) layout for API functions
accepting ray packets of size 4 (\texttt{RTCRay4} type), size 8
(\texttt{RTCRay8} type), and size 16 (\texttt{RTCRay16} type). The
header additionally defines an \texttt{RTCRayNt} template for ray
packets of an arbitrary compile-time size.

\subsubsection{EXIT STATUS}\label{exit-status}

\subsubsection{SEE ALSO}\label{see-also}

\hyperref[rtchit]{RTCHit}

\pagebreak

\section{RTCHit}\label{rtchit}

\subsubsection{NAME}\label{name}

\begin{Shaded}
\begin{Highlighting}[]
\NormalTok{RTCHit }\OperatorTok{{-}}\NormalTok{ single hit structure}
\end{Highlighting}
\end{Shaded}

\subsubsection{SYNOPSIS}\label{synopsis}

\begin{Shaded}
\begin{Highlighting}[]
\PreprocessorTok{\#include }\ImportTok{\textless{}embree4/rtcore.h\textgreater{}}

\KeywordTok{struct}\NormalTok{ RTCHit}
\OperatorTok{\{}
  \DataTypeTok{float}\NormalTok{ Ng\_x}\OperatorTok{;}                                        \CommentTok{// x coordinate of geometry normal}
  \DataTypeTok{float}\NormalTok{ Ng\_y}\OperatorTok{;}                                        \CommentTok{// y coordinate of geometry normal}
  \DataTypeTok{float}\NormalTok{ Ng\_z}\OperatorTok{;}                                        \CommentTok{// z coordinate of geometry normal}

  \DataTypeTok{float}\NormalTok{ u}\OperatorTok{;}                                           \CommentTok{// barycentric u coordinate of hit}
  \DataTypeTok{float}\NormalTok{ v}\OperatorTok{;}                                           \CommentTok{// barycentric v coordinate of hit}

  \DataTypeTok{unsigned} \DataTypeTok{int}\NormalTok{ primID}\OperatorTok{;}                               \CommentTok{// geometry ID}
  \DataTypeTok{unsigned} \DataTypeTok{int}\NormalTok{ geomID}\OperatorTok{;}                               \CommentTok{// primitive ID}
  \DataTypeTok{unsigned} \DataTypeTok{int}\NormalTok{ instID}\OperatorTok{[}\NormalTok{RTC\_MAX\_INSTANCE\_LEVEL\_COUNT}\OperatorTok{];} \CommentTok{// instance ID}
\OperatorTok{\};}
\end{Highlighting}
\end{Shaded}

\subsubsection{DESCRIPTION}\label{description}

The \texttt{RTCHit} type defines the type of a ray/primitive
intersection result. The hit contains the unnormalized geometric normal
in object space at the hit location (\texttt{Ng\_x}, \texttt{Ng\_y},
\texttt{Ng\_z} members), the barycentric u/v coordinates of the hit
(\texttt{u} and \texttt{v} members), as well as the primitive ID
(\texttt{primID} member), geometry ID (\texttt{geomID} member), and
instance ID stack (\texttt{instID} member) of the hit. The parametric
intersection distance is not stored inside the hit, but stored inside
the \texttt{tfar} member of the ray.

The \texttt{embree4/rtcore\_ray.h} header additionally defines the same
hit structure in structure of array (SOA) layout for hit packets of size
4 (\texttt{RTCHit4} type), size 8 (\texttt{RTCHit8} type), and size 16
(\texttt{RTCHit16} type). The header additionally defines an
\texttt{RTCHitNt} template for hit packets of an arbitrary compile-time
size.

\subsubsection{EXIT STATUS}\label{exit-status}

\subsubsection{SEE ALSO}\label{see-also}

\hyperref[rtcray]{RTCRay}, {[}Multi-Level Instancing{]}

\pagebreak

\section{RTCRayHit}\label{rtcrayhit}

\subsubsection{NAME}\label{name}

\begin{Shaded}
\begin{Highlighting}[]
\NormalTok{RTCRayHit }\OperatorTok{{-}}\NormalTok{ combined single ray}\OperatorTok{/}\NormalTok{hit structure}
\end{Highlighting}
\end{Shaded}

\subsubsection{SYNOPSIS}\label{synopsis}

\begin{Shaded}
\begin{Highlighting}[]
\PreprocessorTok{\#include }\ImportTok{\textless{}embree4/rtcore\_ray.h\textgreater{}}

\KeywordTok{struct}\NormalTok{ RTCORE\_ALIGN}\OperatorTok{(}\DecValTok{16}\OperatorTok{)}\NormalTok{ RTCRayHit}
\OperatorTok{\{}
  \KeywordTok{struct}\NormalTok{ RTCRay ray}\OperatorTok{;}
  \KeywordTok{struct}\NormalTok{ RTCHit hit}\OperatorTok{;}
\OperatorTok{\};}
\end{Highlighting}
\end{Shaded}

\subsubsection{DESCRIPTION}\label{description}

The \texttt{RTCRayHit} structure is used as input for the
\texttt{rtcIntersect}-type functions and stores the ray to intersect and
some hit fields that hold the intersection result afterwards.

The \texttt{embree4/rtcore\_ray.h} header additionally defines the same
ray/hit structure in structure of array (SOA) layout for API functions
accepting ray packets of size 4 (\texttt{RTCRayHit4} type), size 8
(\texttt{RTCRayHit8} type), and size 16 (\texttt{RTCRayHit16} type). The
header additionally defines an \texttt{RTCRayHitNt} template to generate
ray/hit packets of an arbitrary compile-time size.

\subsubsection{EXIT STATUS}\label{exit-status}

\subsubsection{SEE ALSO}\label{see-also}

\hyperref[rtcray]{RTCRay}, \hyperref[rtchit]{RTCHit}

\pagebreak

\section{RTCRayN}\label{rtcrayn}

\subsubsection{NAME}\label{name}

\begin{Shaded}
\begin{Highlighting}[]
\NormalTok{RTCRayN }\OperatorTok{{-}}\NormalTok{ ray packet of runtime size}
\end{Highlighting}
\end{Shaded}

\subsubsection{SYNOPSIS}\label{synopsis}

\begin{Shaded}
\begin{Highlighting}[]
\PreprocessorTok{\#include }\ImportTok{\textless{}embree4/rtcore\_ray.h\textgreater{}}

\KeywordTok{struct}\NormalTok{ RTCRayN}\OperatorTok{;}

\DataTypeTok{float}\OperatorTok{\&}\NormalTok{ RTCRayN\_org\_x}\OperatorTok{(}\NormalTok{RTCRayN}\OperatorTok{*}\NormalTok{ ray}\OperatorTok{,} \DataTypeTok{unsigned} \DataTypeTok{int}\NormalTok{ N}\OperatorTok{,} \DataTypeTok{unsigned} \DataTypeTok{int}\NormalTok{ i}\OperatorTok{);}
\DataTypeTok{float}\OperatorTok{\&}\NormalTok{ RTCRayN\_org\_y}\OperatorTok{(}\NormalTok{RTCRayN}\OperatorTok{*}\NormalTok{ ray}\OperatorTok{,} \DataTypeTok{unsigned} \DataTypeTok{int}\NormalTok{ N}\OperatorTok{,} \DataTypeTok{unsigned} \DataTypeTok{int}\NormalTok{ i}\OperatorTok{);}
\DataTypeTok{float}\OperatorTok{\&}\NormalTok{ RTCRayN\_org\_z}\OperatorTok{(}\NormalTok{RTCRayN}\OperatorTok{*}\NormalTok{ ray}\OperatorTok{,} \DataTypeTok{unsigned} \DataTypeTok{int}\NormalTok{ N}\OperatorTok{,} \DataTypeTok{unsigned} \DataTypeTok{int}\NormalTok{ i}\OperatorTok{);}
\DataTypeTok{float}\OperatorTok{\&}\NormalTok{ RTCRayN\_tnear}\OperatorTok{(}\NormalTok{RTCRayN}\OperatorTok{*}\NormalTok{ ray}\OperatorTok{,} \DataTypeTok{unsigned} \DataTypeTok{int}\NormalTok{ N}\OperatorTok{,} \DataTypeTok{unsigned} \DataTypeTok{int}\NormalTok{ i}\OperatorTok{);}

\DataTypeTok{float}\OperatorTok{\&}\NormalTok{ RTCRayN\_dir\_x}\OperatorTok{(}\NormalTok{RTCRayN}\OperatorTok{*}\NormalTok{ ray}\OperatorTok{,} \DataTypeTok{unsigned} \DataTypeTok{int}\NormalTok{ N}\OperatorTok{,} \DataTypeTok{unsigned} \DataTypeTok{int}\NormalTok{ i}\OperatorTok{);}
\DataTypeTok{float}\OperatorTok{\&}\NormalTok{ RTCRayN\_dir\_y}\OperatorTok{(}\NormalTok{RTCRayN}\OperatorTok{*}\NormalTok{ ray}\OperatorTok{,} \DataTypeTok{unsigned} \DataTypeTok{int}\NormalTok{ N}\OperatorTok{,} \DataTypeTok{unsigned} \DataTypeTok{int}\NormalTok{ i}\OperatorTok{);}
\DataTypeTok{float}\OperatorTok{\&}\NormalTok{ RTCRayN\_dir\_z}\OperatorTok{(}\NormalTok{RTCRayN}\OperatorTok{*}\NormalTok{ ray}\OperatorTok{,} \DataTypeTok{unsigned} \DataTypeTok{int}\NormalTok{ N}\OperatorTok{,} \DataTypeTok{unsigned} \DataTypeTok{int}\NormalTok{ i}\OperatorTok{);}
\DataTypeTok{float}\OperatorTok{\&}\NormalTok{ RTCRayN\_time }\OperatorTok{(}\NormalTok{RTCRayN}\OperatorTok{*}\NormalTok{ ray}\OperatorTok{,} \DataTypeTok{unsigned} \DataTypeTok{int}\NormalTok{ N}\OperatorTok{,} \DataTypeTok{unsigned} \DataTypeTok{int}\NormalTok{ i}\OperatorTok{);}

\DataTypeTok{float}\OperatorTok{\&}\NormalTok{        RTCRayN\_tfar }\OperatorTok{(}\NormalTok{RTCRayN}\OperatorTok{*}\NormalTok{ ray}\OperatorTok{,} \DataTypeTok{unsigned} \DataTypeTok{int}\NormalTok{ N}\OperatorTok{,} \DataTypeTok{unsigned} \DataTypeTok{int}\NormalTok{ i}\OperatorTok{);}
\DataTypeTok{unsigned} \DataTypeTok{int}\OperatorTok{\&}\NormalTok{ RTCRayN\_mask }\OperatorTok{(}\NormalTok{RTCRayN}\OperatorTok{*}\NormalTok{ ray}\OperatorTok{,} \DataTypeTok{unsigned} \DataTypeTok{int}\NormalTok{ N}\OperatorTok{,} \DataTypeTok{unsigned} \DataTypeTok{int}\NormalTok{ i}\OperatorTok{);}
\DataTypeTok{unsigned} \DataTypeTok{int}\OperatorTok{\&}\NormalTok{ RTCRayN\_id   }\OperatorTok{(}\NormalTok{RTCRayN}\OperatorTok{*}\NormalTok{ ray}\OperatorTok{,} \DataTypeTok{unsigned} \DataTypeTok{int}\NormalTok{ N}\OperatorTok{,} \DataTypeTok{unsigned} \DataTypeTok{int}\NormalTok{ i}\OperatorTok{);}
\DataTypeTok{unsigned} \DataTypeTok{int}\OperatorTok{\&}\NormalTok{ RTCRayN\_flags}\OperatorTok{(}\NormalTok{RTCRayN}\OperatorTok{*}\NormalTok{ ray}\OperatorTok{,} \DataTypeTok{unsigned} \DataTypeTok{int}\NormalTok{ N}\OperatorTok{,} \DataTypeTok{unsigned} \DataTypeTok{int}\NormalTok{ i}\OperatorTok{);}
\end{Highlighting}
\end{Shaded}

\subsubsection{DESCRIPTION}\label{description}

When the ray packet size is not known at compile time (e.g.~when Embree
returns a ray packet in the \texttt{RTCFilterFuncN} callback function),
Embree uses the \texttt{RTCRayN} type for ray packets. These ray packets
can only have sizes of 1, 4, 8, or 16. No other packet size will be
used.

You can either implement different special code paths for each of these
possible packet sizes and cast the ray to the appropriate ray packet
type, or implement one general code path that uses the
\texttt{RTCRayN\_XXX} helper functions to access the ray packet
components.

These helper functions get a pointer to the ray packet (\texttt{ray}
argument), the packet size (\texttt{N} argument), and returns a
reference to a component (e.g.~x-component of origin) of the the i-th
ray of the packet (\texttt{i} argument).

\subsubsection{EXIT STATUS}\label{exit-status}

\subsubsection{SEE ALSO}\label{see-also}

\hyperref[rtchitn]{RTCHitN}

\pagebreak

\section{RTCHitN}\label{rtchitn}

\subsubsection{NAME}\label{name}

\begin{Shaded}
\begin{Highlighting}[]
\NormalTok{RTCHitN }\OperatorTok{{-}}\NormalTok{ hit packet of runtime size}
\end{Highlighting}
\end{Shaded}

\subsubsection{SYNOPSIS}\label{synopsis}

\begin{Shaded}
\begin{Highlighting}[]
\PreprocessorTok{\#include }\ImportTok{\textless{}embree4/rtcore.h\textgreater{}}

\KeywordTok{struct}\NormalTok{ HitN}\OperatorTok{;}

\DataTypeTok{float}\OperatorTok{\&}\NormalTok{ RTCHitN\_Ng\_x}\OperatorTok{(}\NormalTok{RTCHitN}\OperatorTok{*}\NormalTok{ hit}\OperatorTok{,} \DataTypeTok{unsigned} \DataTypeTok{int}\NormalTok{ N}\OperatorTok{,} \DataTypeTok{unsigned} \DataTypeTok{int}\NormalTok{ i}\OperatorTok{);}
\DataTypeTok{float}\OperatorTok{\&}\NormalTok{ RTCHitN\_Ng\_y}\OperatorTok{(}\NormalTok{RTCHitN}\OperatorTok{*}\NormalTok{ hit}\OperatorTok{,} \DataTypeTok{unsigned} \DataTypeTok{int}\NormalTok{ N}\OperatorTok{,} \DataTypeTok{unsigned} \DataTypeTok{int}\NormalTok{ i}\OperatorTok{);}
\DataTypeTok{float}\OperatorTok{\&}\NormalTok{ RTCHitN\_Ng\_z}\OperatorTok{(}\NormalTok{RTCHitN}\OperatorTok{*}\NormalTok{ hit}\OperatorTok{,} \DataTypeTok{unsigned} \DataTypeTok{int}\NormalTok{ N}\OperatorTok{,} \DataTypeTok{unsigned} \DataTypeTok{int}\NormalTok{ i}\OperatorTok{);}

\DataTypeTok{float}\OperatorTok{\&}\NormalTok{ RTCHitN\_u}\OperatorTok{(}\NormalTok{RTCHitN}\OperatorTok{*}\NormalTok{ hit}\OperatorTok{,} \DataTypeTok{unsigned} \DataTypeTok{int}\NormalTok{ N}\OperatorTok{,} \DataTypeTok{unsigned} \DataTypeTok{int}\NormalTok{ i}\OperatorTok{);}
\DataTypeTok{float}\OperatorTok{\&}\NormalTok{ RTCHitN\_v}\OperatorTok{(}\NormalTok{RTCHitN}\OperatorTok{*}\NormalTok{ hit}\OperatorTok{,} \DataTypeTok{unsigned} \DataTypeTok{int}\NormalTok{ N}\OperatorTok{,} \DataTypeTok{unsigned} \DataTypeTok{int}\NormalTok{ i}\OperatorTok{);}

\DataTypeTok{unsigned}\OperatorTok{\&}\NormalTok{ RTCHitN\_primID}\OperatorTok{(}\NormalTok{RTCHitN}\OperatorTok{*}\NormalTok{ hit}\OperatorTok{,} \DataTypeTok{unsigned} \DataTypeTok{int}\NormalTok{ N}\OperatorTok{,} \DataTypeTok{unsigned} \DataTypeTok{int}\NormalTok{ i}\OperatorTok{);}
\DataTypeTok{unsigned}\OperatorTok{\&}\NormalTok{ RTCHitN\_geomID}\OperatorTok{(}\NormalTok{RTCHitN}\OperatorTok{*}\NormalTok{ hit}\OperatorTok{,} \DataTypeTok{unsigned} \DataTypeTok{int}\NormalTok{ N}\OperatorTok{,} \DataTypeTok{unsigned} \DataTypeTok{int}\NormalTok{ i}\OperatorTok{);}
\DataTypeTok{unsigned}\OperatorTok{\&}\NormalTok{ RTCHitN\_instID}\OperatorTok{(}\NormalTok{RTCHitN}\OperatorTok{*}\NormalTok{ hit}\OperatorTok{,} \DataTypeTok{unsigned} \DataTypeTok{int}\NormalTok{ N}\OperatorTok{,} \DataTypeTok{unsigned} \DataTypeTok{int}\NormalTok{ i}\OperatorTok{,} \DataTypeTok{unsigned} \DataTypeTok{int}\NormalTok{ level}\OperatorTok{);}
\end{Highlighting}
\end{Shaded}

\subsubsection{DESCRIPTION}\label{description}

When the hit packet size is not known at compile time (e.g.~when Embree
returns a hit packet in the \texttt{RTCFilterFuncN} callback function),
Embree uses the \texttt{RTCHitN} type for hit packets. These hit packets
can only have sizes of 1, 4, 8, or 16. No other packet size will be
used.

You can either implement different special code paths for each of these
possible packet sizes and cast the hit to the appropriate hit packet
type, or implement one general code path that uses the
\texttt{RTCHitN\_XXX} helper functions to access hit packet components.

These helper functions get a pointer to the hit packet (\texttt{hit}
argument), the packet size (\texttt{N} argument), and returns a
reference to a component (e.g.~x component of \texttt{Ng}) of the the
i-th hit of the packet (\texttt{i} argument).

\subsubsection{EXIT STATUS}\label{exit-status}

\subsubsection{SEE ALSO}\label{see-also}

\hyperref[rtcrayn]{RTCRayN}

\pagebreak

\section{RTCRayHitN}\label{rtcrayhitn}

\subsubsection{NAME}\label{name}

\begin{Shaded}
\begin{Highlighting}[]
\NormalTok{RTCRayHitN }\OperatorTok{{-}}\NormalTok{ combined ray}\OperatorTok{/}\NormalTok{hit packet of runtime size}
\end{Highlighting}
\end{Shaded}

\subsubsection{SYNOPSIS}\label{synopsis}

\begin{Shaded}
\begin{Highlighting}[]
\PreprocessorTok{\#include }\ImportTok{\textless{}embree4/rtcore\_ray.h\textgreater{}}

\KeywordTok{struct}\NormalTok{ RTCRayHitN}\OperatorTok{;}

\KeywordTok{struct}\NormalTok{ RTCRayN}\OperatorTok{*}\NormalTok{ RTCRayHitN\_RayN}\OperatorTok{(}\KeywordTok{struct}\NormalTok{ RTCRayHitN}\OperatorTok{*}\NormalTok{ rayhit}\OperatorTok{,} \DataTypeTok{unsigned} \DataTypeTok{int}\NormalTok{ N}\OperatorTok{);}
\KeywordTok{struct}\NormalTok{ RTCHitN}\OperatorTok{*}\NormalTok{ RTCRayHitN\_HitN}\OperatorTok{(}\KeywordTok{struct}\NormalTok{ RTCRayHitN}\OperatorTok{*}\NormalTok{ rayhit}\OperatorTok{,} \DataTypeTok{unsigned} \DataTypeTok{int}\NormalTok{ N}\OperatorTok{);}
\end{Highlighting}
\end{Shaded}

\subsubsection{DESCRIPTION}\label{description}

When the packet size of a ray/hit structure is not known at compile time
(e.g.~when Embree returns a ray/hit packet in the
\texttt{RTCIntersectFunctionN} callback function), Embree uses the
\texttt{RTCRayHitN} type for ray packets. These ray/hit packets can only
have sizes of 1, 4, 8, or 16. No other packet size will be used.

You can either implement different special code paths for each of these
possible packet sizes and cast the ray/hit to the appropriate ray/hit
packet type, or extract the \texttt{RTCRayN} and \texttt{RTCHitN}
components using the \texttt{rtcGetRayN} and \texttt{rtcGetHitN} helper
functions and use the \texttt{RTCRayN\_XXX} and \texttt{RTCHitN\_XXX}
functions to access the ray and hit parts of the structure.

\subsubsection{EXIT STATUS}\label{exit-status}

\subsubsection{SEE ALSO}\label{see-also}

\hyperref[rtchitn]{RTCHitN}

\pagebreak

\section{RTCFeatureFlags}\label{rtcfeatureflags}

\subsubsection{NAME}\label{name}

\begin{Shaded}
\begin{Highlighting}[]
\NormalTok{RTCFeatureFlags }\OperatorTok{{-}}\NormalTok{ specifies features to enable}
  \ControlFlowTok{for}\NormalTok{ ray queries}
\end{Highlighting}
\end{Shaded}

\subsubsection{SYNOPSIS}\label{synopsis}

\begin{Shaded}
\begin{Highlighting}[]
\PreprocessorTok{\#include }\ImportTok{\textless{}embree4/rtcore\_ray.h\textgreater{}}

\KeywordTok{enum}\NormalTok{ RTCFeatureFlags}
\OperatorTok{\{}
\NormalTok{  RTC\_FEATURE\_FLAG\_NONE }\OperatorTok{=} \DecValTok{0}\OperatorTok{,}
  
\NormalTok{  RTC\_FEATURE\_FLAG\_MOTION\_BLUR }\OperatorTok{=} \DecValTok{1} \OperatorTok{\textless{}\textless{}} \DecValTok{0}\OperatorTok{,}

\NormalTok{  RTC\_FEATURE\_FLAG\_TRIANGLE }\OperatorTok{=} \DecValTok{1} \OperatorTok{\textless{}\textless{}} \DecValTok{1}\OperatorTok{,}
\NormalTok{  RTC\_FEATURE\_FLAG\_QUAD }\OperatorTok{=} \DecValTok{1} \OperatorTok{\textless{}\textless{}} \DecValTok{2}\OperatorTok{,}
\NormalTok{  RTC\_FEATURE\_FLAG\_GRID }\OperatorTok{=} \DecValTok{1} \OperatorTok{\textless{}\textless{}} \DecValTok{3}\OperatorTok{,}
\NormalTok{  RTC\_FEATURE\_FLAG\_SUBDIVISION }\OperatorTok{=} \DecValTok{1} \OperatorTok{\textless{}\textless{}} \DecValTok{4}\OperatorTok{,}
\NormalTok{  RTC\_FEATURE\_FLAG\_POINT }\OperatorTok{=} \OperatorTok{...} \OperatorTok{,}
\NormalTok{  RTC\_FEATURE\_FLAG\_CURVES }\OperatorTok{=} \OperatorTok{...} \OperatorTok{,}
 
\NormalTok{  RTC\_FEATURE\_FLAG\_CONE\_LINEAR\_CURVE }\OperatorTok{=} \DecValTok{1} \OperatorTok{\textless{}\textless{}} \DecValTok{5}\OperatorTok{,}
\NormalTok{  RTC\_FEATURE\_FLAG\_ROUND\_LINEAR\_CURVE  }\OperatorTok{=} \DecValTok{1} \OperatorTok{\textless{}\textless{}} \DecValTok{6}\OperatorTok{,}
\NormalTok{  RTC\_FEATURE\_FLAG\_FLAT\_LINEAR\_CURVE }\OperatorTok{=} \DecValTok{1} \OperatorTok{\textless{}\textless{}} \DecValTok{7}\OperatorTok{,}

\NormalTok{  RTC\_FEATURE\_FLAG\_ROUND\_BEZIER\_CURVE }\OperatorTok{=} \DecValTok{1} \OperatorTok{\textless{}\textless{}} \DecValTok{8}\OperatorTok{,}
\NormalTok{  RTC\_FEATURE\_FLAG\_FLAT\_BEZIER\_CURVE }\OperatorTok{=} \DecValTok{1} \OperatorTok{\textless{}\textless{}} \DecValTok{9}\OperatorTok{,}
\NormalTok{  RTC\_FEATURE\_FLAG\_NORMAL\_ORIENTED\_BEZIER\_CURVE }\OperatorTok{=} \DecValTok{1} \OperatorTok{\textless{}\textless{}} \DecValTok{10}\OperatorTok{,}

\NormalTok{  RTC\_FEATURE\_FLAG\_ROUND\_BSPLINE\_CURVE }\OperatorTok{=} \DecValTok{1} \OperatorTok{\textless{}\textless{}} \DecValTok{11}\OperatorTok{,}
\NormalTok{  RTC\_FEATURE\_FLAG\_FLAT\_BSPLINE\_CURVE }\OperatorTok{=} \DecValTok{1} \OperatorTok{\textless{}\textless{}} \DecValTok{12}\OperatorTok{,}
\NormalTok{  RTC\_FEATURE\_FLAG\_NORMAL\_ORIENTED\_BSPLINE\_CURVE }\OperatorTok{=} \DecValTok{1} \OperatorTok{\textless{}\textless{}} \DecValTok{13}\OperatorTok{,}

\NormalTok{  RTC\_FEATURE\_FLAG\_ROUND\_HERMITE\_CURVE }\OperatorTok{=} \DecValTok{1} \OperatorTok{\textless{}\textless{}} \DecValTok{14}\OperatorTok{,}
\NormalTok{  RTC\_FEATURE\_FLAG\_FLAT\_HERMITE\_CURVE }\OperatorTok{=} \DecValTok{1} \OperatorTok{\textless{}\textless{}} \DecValTok{15}\OperatorTok{,}
\NormalTok{  RTC\_FEATURE\_FLAG\_NORMAL\_ORIENTED\_HERMITE\_CURVE }\OperatorTok{=} \DecValTok{1} \OperatorTok{\textless{}\textless{}} \DecValTok{16}\OperatorTok{,}

\NormalTok{  RTC\_FEATURE\_FLAG\_ROUND\_CATMULL\_ROM\_CURVE }\OperatorTok{=} \DecValTok{1} \OperatorTok{\textless{}\textless{}} \DecValTok{17}\OperatorTok{,}
\NormalTok{  RTC\_FEATURE\_FLAG\_FLAT\_CATMULL\_ROM\_CURVE }\OperatorTok{=} \DecValTok{1} \OperatorTok{\textless{}\textless{}} \DecValTok{18}\OperatorTok{,}
\NormalTok{  RTC\_FEATURE\_FLAG\_NORMAL\_ORIENTED\_CATMULL\_ROM\_CURVE }\OperatorTok{=} \DecValTok{1} \OperatorTok{\textless{}\textless{}} \DecValTok{19}\OperatorTok{,}

\NormalTok{  RTC\_FEATURE\_FLAG\_SPHERE\_POINT }\OperatorTok{=} \DecValTok{1} \OperatorTok{\textless{}\textless{}} \DecValTok{20}\OperatorTok{,}
\NormalTok{  RTC\_FEATURE\_FLAG\_DISC\_POINT }\OperatorTok{=} \DecValTok{1} \OperatorTok{\textless{}\textless{}} \DecValTok{21}\OperatorTok{,}
\NormalTok{  RTC\_FEATURE\_FLAG\_ORIENTED\_DISC\_POINT }\OperatorTok{=} \DecValTok{1} \OperatorTok{\textless{}\textless{}} \DecValTok{22}\OperatorTok{,}

\NormalTok{  RTC\_FEATURE\_FLAG\_ROUND\_CURVES }\OperatorTok{=} \OperatorTok{...} \OperatorTok{,}
\NormalTok{  RTC\_FEATURE\_FLAG\_FLAT\_CURVES }\OperatorTok{=} \OperatorTok{...} \OperatorTok{,}
\NormalTok{  RTC\_FEATURE\_FLAG\_NORMAL\_ORIENTED\_CURVES }\OperatorTok{=} \OperatorTok{...} \OperatorTok{,}
  
\NormalTok{  RTC\_FEATURE\_FLAG\_LINEAR\_CURVES }\OperatorTok{=} \OperatorTok{...} \OperatorTok{,}
\NormalTok{  RTC\_FEATURE\_FLAG\_BEZIER\_CURVES }\OperatorTok{=} \OperatorTok{...} \OperatorTok{,}
\NormalTok{  RTC\_FEATURE\_FLAG\_BSPLINE\_CURVES }\OperatorTok{=} \OperatorTok{...} \OperatorTok{,}
\NormalTok{  RTC\_FEATURE\_FLAG\_HERMITE\_CURVES }\OperatorTok{=} \OperatorTok{...} \OperatorTok{,}
  
\NormalTok{  RTC\_FEATURE\_FLAG\_INSTANCE }\OperatorTok{=} \DecValTok{1} \OperatorTok{\textless{}\textless{}} \DecValTok{23}\OperatorTok{,}

\NormalTok{  RTC\_FEATURE\_FLAG\_FILTER\_FUNCTION\_IN\_ARGUMENTS }\OperatorTok{=} \DecValTok{1} \OperatorTok{\textless{}\textless{}} \DecValTok{24}\OperatorTok{,}
\NormalTok{  RTC\_FEATURE\_FLAG\_FILTER\_FUNCTION\_IN\_GEOMETRY }\OperatorTok{=} \DecValTok{1} \OperatorTok{\textless{}\textless{}} \DecValTok{25}\OperatorTok{,}
\NormalTok{  RTC\_FEATURE\_FLAG\_FILTER\_FUNCTION }\OperatorTok{=} \OperatorTok{...} \OperatorTok{,}

\NormalTok{  RTC\_FEATURE\_FLAG\_USER\_GEOMETRY\_CALLBACK\_IN\_ARGUMENTS }\OperatorTok{=} \DecValTok{1} \OperatorTok{\textless{}\textless{}} \DecValTok{26}\OperatorTok{,}
\NormalTok{  RTC\_FEATURE\_FLAG\_USER\_GEOMETRY\_CALLBACK\_IN\_GEOMETRY }\OperatorTok{=} \DecValTok{1} \OperatorTok{\textless{}\textless{}} \DecValTok{27}\OperatorTok{,}
\NormalTok{  RTC\_FEATURE\_FLAG\_USER\_GEOMETRY }\OperatorTok{=} \OperatorTok{...} \OperatorTok{,}

\NormalTok{  RTC\_FEATURE\_FLAG\_32\_BIT\_RAY\_MASK }\OperatorTok{=} \DecValTok{1} \OperatorTok{\textless{}\textless{}} \DecValTok{28}\OperatorTok{,}

\NormalTok{  RTC\_FEATURE\_FLAG\_ALL }\OperatorTok{=} \BaseNTok{0xffffffff}
\OperatorTok{\};}
\end{Highlighting}
\end{Shaded}

\subsubsection{DESCRIPTION}\label{description}

The \texttt{RTCFeatureFlags} enum specify a bit mask to enable specific
ray tracing features for ray query operations. The feature flags are
passed to the \texttt{rtcIntersect1/4/8/16} and
\texttt{rtcOccluded1/4/8/16} functions through the
\texttt{RTCIntersectArguments} and \texttt{RTCOccludedArguments}
structures. Only a ray tracing feature whose bit is enabled in the
feature mask can get used. If a feature bit is not set, the behaviour is
undefined, thus the feature may work or not. To enable multiple features
the respective features have to get combined using a bitwise \texttt{OR}
operation.

The purpose of feature flags is to reduce code size on the GPU by
enabling just the features required to render the scene. On the CPU
there is no need to use feature flags, and the default of all features
enabled (\texttt{RTC\_FEATURE\_FLAG\_ALL}) can just be kept.

The following features can get enabled using feature flags:

\begin{itemize}
\item
  RTC\_FEATURE\_FLAG\_MOTION\_BLUR: Enables motion blur for all geometry
  types.
\item
  RTC\_FEATURE\_FLAG\_TRIANGLE: Enables triangle geometries
  (RTC\_GEOMETRY\_TYPE\_TRIANGLE).
\item
  RTC\_FEATURE\_FLAG\_QUAD: Enables quad geometries
  (RTC\_GEOMETRY\_TYPE\_QUAD).
\item
  RTC\_FEATURE\_FLAG\_GRID: Enables grid geometries
  (RTC\_GEOMETRY\_TYPE\_GRID).
\item
  RTC\_FEATURE\_FLAG\_SUBDIVISION: Enables subdivision geometries
  (RTC\_GEOMETRY\_TYPE\_SUBDIVISION).
\item
  RTC\_FEATURE\_FLAG\_POINT: Enables all point geometry types
  (RTC\_GEOMETRY\_TYPE\_XXX\_POINT)
\item
  RTC\_FEATURE\_FLAG\_CURVES: Enables all curve geometry types
  (RTC\_GEOMETRY\_TYPE\_XXX\_YYY\_CURVE)
\item
  RTC\_FEATURE\_FLAG\_ROUND\_CURVES: Enables all round curves
  (RTC\_GEOMETRY\_TYPE\_ROUND\_XXX\_CURVE).
\item
  RTC\_FEATURE\_FLAG\_FLAT\_CURVES: Enables all flat curves
  (RTC\_GEOMETRY\_TYPE\_FLAT\_XXX\_CURVE).
\item
  RTC\_FEATURE\_FLAG\_NORMAL\_ORIENTED\_CURVES: Enables all normal
  oriented curves (RTC\_GEOMETRY\_TYPE\_NORMAL\_ORIENTED\_XXX\_CURVE).
\item
  RTC\_FEATURE\_FLAG\_LINEAR\_CURVES: Enables all linear curves
  (RTC\_GEOMETRY\_TYPE\_XXX\_LINEAR\_CURVE).
\item
  RTC\_FEATURE\_FLAG\_BEZIER\_CURVES: Enables all Bézier curves
  (RTC\_GEOMETRY\_TYPE\_XXX\_BEZIER\_CURVE).
\item
  RTC\_FEATURE\_FLAG\_BSPLINE\_CURVES: Enables all B-spline curves
  (RTC\_GEOMETRY\_TYPE\_XXX\_BSPLINE\_CURVE).
\item
  RTC\_FEATURE\_FLAG\_HERMITE\_CURVES: Enables all Hermite curves
  (RTC\_GEOMETRY\_TYPE\_XXX\_HERMITE\_CURVE).
\item
  RTC\_FEATURE\_FLAG\_CONE\_LINEAR\_CURVE: Enables cone geometry type
  (RTC\_GEOMETRY\_TYPE\_CONE\_LINEAR\_CURVE).
\item
  RTC\_FEATURE\_FLAG\_ROUND\_LINEAR\_CURVE: Enables round linear curves
  (RTC\_GEOMETRY\_TYPE\_ROUND\_LINEAR\_CURVE).
\item
  RTC\_FEATURE\_FLAG\_FLAT\_LINEAR\_CURVE: Enables flat linear curves
  (RTC\_GEOMETRY\_TYPE\_FLAT\_LINEAR\_CURVE).
\item
  RTC\_FEATURE\_FLAG\_ROUND\_BEZIER\_CURVE: Enables round Bézier curves
  (RTC\_GEOMETRY\_TYPE\_ROUND\_BEZIER\_CURVE).
\item
  RTC\_FEATURE\_FLAG\_FLAT\_BEZIER\_CURVE: Enables flat Bézier curves
  (RTC\_GEOMETRY\_TYPE\_FLAT\_BEZIER\_CURVE).
\item
  RTC\_FEATURE\_FLAG\_NORMAL\_ORIENTED\_BEZIER\_CURVE: Enables normal
  oriented Bézier curves
  (RTC\_GEOMETRY\_TYPE\_NORMAL\_ORIENTED\_BEZIER\_CURVE).
\item
  RTC\_FEATURE\_FLAG\_ROUND\_BSPLINE\_CURVE: Enables round B-spline
  curves (RTC\_GEOMETRY\_TYPE\_ROUND\_BSPLINE\_CURVE).
\item
  RTC\_FEATURE\_FLAG\_FLAT\_BSPLINE\_CURVE: Enables flat B-spline curves
  (RTC\_GEOMETRY\_TYPE\_FLAT\_BSPLINE\_CURVE).
\item
  RTC\_FEATURE\_FLAG\_NORMAL\_ORIENTED\_BSPLINE\_CURVE: Enables normal
  oriented B-spline curves
  (RTC\_GEOMETRY\_TYPE\_NORMAL\_ORIENTED\_BSPLINE\_CURVE).
\item
  RTC\_FEATURE\_FLAG\_ROUND\_HERMITE\_CURVE: Enables round Hermite
  curves (RTC\_GEOMETRY\_TYPE\_ROUND\_HERMITE\_CURVE).
\item
  RTC\_FEATURE\_FLAG\_FLAT\_HERMITE\_CURVE: Enables flat Hermite curves
  (RTC\_GEOMETRY\_TYPE\_FLAT\_HERMITE\_CURVE).
\item
  RTC\_FEATURE\_FLAG\_NORMAL\_ORIENTED\_HERMITE\_CURVE: Enables normal
  oriented Hermite curves
  (RTC\_GEOMETRY\_TYPE\_NORMAL\_ORIENTED\_HERMITE\_CURVE).
\item
  RTC\_FEATURE\_FLAG\_ROUND\_CATMULL\_ROM\_CURVE: Enables round Catmull
  Rom curves (RTC\_GEOMETRY\_TYPE\_ROUND\_CATMULL\_ROM\_CURVE).
\item
  RTC\_FEATURE\_FLAG\_FLAT\_CATMULL\_ROM\_CURVE: Enables flat Catmull
  Rom curves (RTC\_GEOMETRY\_TYPE\_FLAT\_CATMULL\_ROM\_CURVE).
\item
  RTC\_FEATURE\_FLAG\_NORMAL\_ORIENTED\_CATMULL\_ROM\_CURVE: Enables
  normal oriented Catmull Rom curves
  (RTC\_GEOMETRY\_TYPE\_NORMAL\_ORIENTED\_CATMULL\_ROM\_CURVE).
\item
  RTC\_FEATURE\_FLAG\_SPHERE\_POINT: Enables sphere geometry type
  (RTC\_GEOMETRY\_TYPE\_SPHERE\_POINT).
\item
  RTC\_FEATURE\_FLAG\_DISC\_POINT: Enables disc geometry type
  (RTC\_GEOMETRY\_TYPE\_DISC\_POINT).
\item
  RTC\_FEATURE\_FLAG\_ORIENTED\_DISC\_POINT: Enables oriented disc
  geometry types (RTC\_GEOMETRY\_TYPE\_ORIENTED\_DISC\_POINT).
\item
  RTC\_FEATURE\_FLAG\_INSTANCE: Enables instance geometries
  (RTC\_GEOMETRY\_TYPE\_INSTANCE).
\item
  RTC\_FEATURE\_FLAG\_FILTER\_FUNCTION\_IN\_ARGUMENTS: Enables filter
  functions passed through intersect arguments.
\item
  RTC\_FEATURE\_FLAG\_FILTER\_FUNCTION\_IN\_GEOMETRY: Enable filter
  functions passed through geometry.
\item
  RTC\_FEATURE\_FLAG\_FILTER\_FUNCTION: Enables filter functions
  (argument and geometry version).
\item
  RTC\_FEATURE\_FLAG\_USER\_GEOMETRY\_CALLBACK\_IN\_ARGUMENTS: Enables
  RTC\_GEOMETRY\_TYPE\_USER with function pointer passed through
  intersect arguments.
\item
  RTC\_FEATURE\_FLAG\_USER\_GEOMETRY\_CALLBACK\_IN\_GEOMETRY: Enables
  RTC\_GEOMETRY\_TYPE\_USER with function pointer passed through
  geometry object.
\item
  RTC\_FEATURE\_FLAG\_USER\_GEOMETRY: Enables RTC\_GEOMETRY\_TYPE\_USER
  geometries (both argument and geometry callback versions).
\item
  RTC\_FEATURE\_FLAG\_32\_BIT\_RAY\_MASK: Enables full 32 bit ray masks.
  If not used, only the lower 7 bits in the ray mask are handled
  correctly.
\item
  RTC\_FEATURE\_FLAG\_ALL: Enables all features (default).
\end{itemize}

\subsubsection{EXIT STATUS}\label{exit-status}

\subsubsection{SEE ALSO}\label{see-also}

\hyperref[rtcintersect1]{rtcIntersect1},
\hyperref[rtcintersect4816]{rtcIntersect4/8/16},
\hyperref[rtcoccluded1]{rtcOccluded1},
\hyperref[rtcoccluded4816]{rtcOccluded4/8/16},

\pagebreak

\section{rtcInitIntersectArguments}\label{rtcinitintersectarguments}

\subsubsection{NAME}\label{name}

\begin{Shaded}
\begin{Highlighting}[]
\NormalTok{rtcInitIntersectArguments }\OperatorTok{{-}}\NormalTok{ initializes the intersect arguments }\KeywordTok{struct}
\end{Highlighting}
\end{Shaded}

\subsubsection{SYNOPSIS}\label{synopsis}

\begin{Shaded}
\begin{Highlighting}[]
\PreprocessorTok{\#include }\ImportTok{\textless{}embree4/rtcore.h\textgreater{}}

\KeywordTok{enum}\NormalTok{ RTCRayQueryFlags}
\OperatorTok{\{}
\NormalTok{  RTC\_RAY\_QUERY\_FLAG\_NONE}\OperatorTok{,}
\NormalTok{  RTC\_RAY\_QUERY\_FLAG\_INCOHERENT}\OperatorTok{,}
\NormalTok{  RTC\_RAY\_QUERY\_FLAG\_COHERENT}\OperatorTok{,}
\NormalTok{  RTC\_RAY\_QUERY\_FLAG\_INVOKE\_ARGUMENT\_FILTER}
\OperatorTok{\};}

\KeywordTok{struct}\NormalTok{ RTCIntersectArguments}
\OperatorTok{\{}
  \KeywordTok{enum}\NormalTok{ RTCRayQueryFlags flags}\OperatorTok{;}
  \KeywordTok{enum}\NormalTok{ RTCFeatureFlags feature\_mask}\OperatorTok{;}
  \KeywordTok{struct}\NormalTok{ RTCRayQueryContext}\OperatorTok{*}\NormalTok{ context}\OperatorTok{;}
\NormalTok{  RTCFilterFunctionN filter}\OperatorTok{;}
\NormalTok{  RTCIntersectFunctionN intersect}\OperatorTok{;}
\PreprocessorTok{\#if RTC\_MIN\_WIDTH}
  \DataTypeTok{float}\NormalTok{ minWidthDistanceFactor}\OperatorTok{;}
\PreprocessorTok{\#endif}
\OperatorTok{\};}

\DataTypeTok{void}\NormalTok{ rtcInitIntersectArguments}\OperatorTok{(}
  \KeywordTok{struct}\NormalTok{ RTCIntersectArguments}\OperatorTok{*}\NormalTok{ args}
\OperatorTok{);}
\end{Highlighting}
\end{Shaded}

\subsubsection{DESCRIPTION}\label{description}

The \texttt{rtcInitIntersectArguments} function initializes the optional
argument struct that can get passed to the \texttt{rtcIntersect1/4/8/16}
functions to default values. The arguments struct needs to get used for
more advanced Embree features as described here.

The \texttt{flags} member can get used to enable special traversal mode.
Using the \texttt{RTC\_RAY\_QUERY\_FLAG\_INCOHERENT} flag uses an
optimized traversal algorithm for incoherent rays (default), while
\texttt{RTC\_RAY\_QUERY\_FLAG\_COHERENT} uses an optimized traversal
algorithm for coherent rays (e.g.~primary camera rays).

The \texttt{feature\_mask} member should get used in SYCL to just enable
ray tracing features required to render a given scene. Please see
section \hyperref[rtcfeatureflags]{RTCFeatureFlags} for a more detailed
description.

The \texttt{context} member can get used to pass an optional
intersection context. It is guaranteed that the pointer to the context
passed to a ray query is directly passed to all callback functions. This
way it is possible to attach arbitrary data to the end of the context,
such as a per-ray payload. Please note that the ray pointer is not
guaranteed to be passed to the callback functions, thus reading
additional data from the ray pointer passed to callbacks is not
possible. See section
\hyperref[rtcinitrayquerycontext]{rtcInitRayQueryContext} for more
details.

The \texttt{filter} member specifies a filter function to invoke for
each encountered hit. The support for the argument filter function must
be enabled for a scene by using the
\texttt{RTC\_SCENE\_FLAG\_FILTER\_FUNCTION\_IN\_ARGUMENTS} scene flag.
In case of instancing this feature has to get enabled also for each
instantiated scene.

The argument filter function is invoked for each geometry for which it
got explicitely enabled using the
\texttt{rtcSetGeometryEnableFilterFunctionFromArguments} function. The
invokation of the argument filter function can also get enfored for each
geometry using the
\texttt{RTC\_RAY\_QUERY\_FLAG\_INVOKE\_ARGUMENT\_FILTER} ray query flag.
This argument filter function is invoked as a second filter stage after
the per-geometry filter function is invoked. Only rays that passed the
first filter stage are valid in this second filter stage. Having such a
per ray-query filter function can be useful to implement modifications
of the behavior of the query, such as collecting all hits or
accumulating transparencies.

The \texttt{intersect} member specifies the user geometry callback to
get invoked for each user geometry encountered during traversal. The
user geometry callback specified this way has preference over the one
specified inside the geometry.

The \texttt{minWidthDistanceFactor} value controls the target size of
the curve radii when the min-width feature is enabled. Please see the
\hyperref[rtcsetgeometrymaxradiusscale]{rtcSetGeometryMaxRadiusScale}
function for more details on the min-width feature.

\subsubsection{EXIT STATUS}\label{exit-status}

No error code is set by this function.

\subsubsection{SEE ALSO}\label{see-also}

\hyperref[rtcintersect1]{rtcIntersect1},
\hyperref[rtcintersect4816]{rtcIntersect4/8/16},
\hyperref[rtcfeatureflags]{RTCFeatureFlags},
\hyperref[rtcinitrayquerycontext]{rtcInitRayQueryContext},
\hyperref[rtc_geometry_type_user]{RTC\_GEOMETRY\_TYPE\_USER},
\hyperref[rtcsetgeometrymaxradiusscale]{rtcSetGeometryMaxRadiusScale}

\pagebreak

\section{rtcInitOccludedArguments}\label{rtcinitoccludedarguments}

\subsubsection{NAME}\label{name}

\begin{Shaded}
\begin{Highlighting}[]
\NormalTok{rtcInitOccludedArguments }\OperatorTok{{-}}\NormalTok{ initializes the occluded arguments }\KeywordTok{struct}
\end{Highlighting}
\end{Shaded}

\subsubsection{SYNOPSIS}\label{synopsis}

\begin{Shaded}
\begin{Highlighting}[]
\PreprocessorTok{\#include }\ImportTok{\textless{}embree4/rtcore.h\textgreater{}}

\KeywordTok{enum}\NormalTok{ RTCRayQueryFlags}
\OperatorTok{\{}
\NormalTok{  RTC\_RAY\_QUERY\_FLAG\_NONE}\OperatorTok{,}
\NormalTok{  RTC\_RAY\_QUERY\_FLAG\_INCOHERENT}\OperatorTok{,}
\NormalTok{  RTC\_RAY\_QUERY\_FLAG\_COHERENT}\OperatorTok{,}
\NormalTok{  RTC\_RAY\_QUERY\_FLAG\_INVOKE\_ARGUMENT\_FILTER}
\OperatorTok{\};}

\KeywordTok{struct}\NormalTok{ RTCOccludedArguments}
\OperatorTok{\{}
  \KeywordTok{enum}\NormalTok{ RTCRayQueryFlags flags}\OperatorTok{;}
  \KeywordTok{enum}\NormalTok{ RTCFeatureFlags feature\_mask}\OperatorTok{;}
  \KeywordTok{struct}\NormalTok{ RTCRayQueryContext}\OperatorTok{*}\NormalTok{ context}\OperatorTok{;}
\NormalTok{  RTCFilterFunctionN filter}\OperatorTok{;}
\NormalTok{  RTCOccludedFunctionN intersect}\OperatorTok{;}
\PreprocessorTok{\#if RTC\_MIN\_WIDTH}
  \DataTypeTok{float}\NormalTok{ minWidthDistanceFactor}\OperatorTok{;}
\PreprocessorTok{\#endif}
\OperatorTok{\};}

\DataTypeTok{void}\NormalTok{ rtcInitOccludedArguments}\OperatorTok{(}
  \KeywordTok{struct}\NormalTok{ RTCOccludedArguments}\OperatorTok{*}\NormalTok{ args}
\OperatorTok{);}
\end{Highlighting}
\end{Shaded}

\subsubsection{DESCRIPTION}\label{description}

The \texttt{rtcInitOccludedArguments} function initializes the optional
argument struct that can get passed to the \texttt{rtcOccluded1/4/8/16}
functions to default values. The arguments struct needs to get used for
more advanced Embree features as described here.

The \texttt{flags} member can get used to enable special traversal mode.
Using the \texttt{RTC\_RAY\_QUERY\_FLAG\_INCOHERENT} flag uses an
optimized traversal algorithm for incoherent rays (default), while
\texttt{RTC\_RAY\_QUERY\_FLAG\_COHERENT} uses an optimized traversal
algorithm for coherent rays (e.g.~primary camera rays).

The \texttt{feature\_mask} member should get used in SYCL to just enable
ray tracing features required to render a given scene. Please see
section \hyperref[rtcfeatureflags]{RTCFeatureFlags} for a more detailed
description.

The \texttt{context} member can get used to pass an optional
intersection context. It is guaranteed that the pointer to the context
passed to a ray query is directly passed to all callback functions. This
way it is possible to attach arbitrary data to the end of the context,
such as a per-ray payload. Please note that the ray pointer is not
guaranteed to be passed to the callback functions, thus reading
additional data from the ray pointer passed to callbacks is not
possible. See section
\hyperref[rtcinitrayquerycontext]{rtcInitRayQueryContext} for more
details.

The \texttt{filter} member specifies a filter function to invoked for
each encountered hit. The support for the argument filter function must
be enabled for a scene by using the
\texttt{RTC\_SCENE\_FLAG\_FILTER\_FUNCTION\_IN\_ARGUMENTS} scene flag.
In case of instancing this feature has to get enabled also for each
instantiated scene.

The argument filter function is invoked for each geometry for which it
got explicitely enabled using the
\texttt{rtcSetGeometryEnableFilterFunctionFromArguments} function. The
invokation of the argument filter function can also get enfored for each
geometry using the
\texttt{RTC\_RAY\_QUERY\_FLAG\_INVOKE\_ARGUMENT\_FILTER} ray query flag.
This argument filter function is invoked as a second filter stage after
the per-geometry filter function is invoked. Only rays that passed the
first filter stage are valid in this second filter stage. Having such a
per ray-query filter function can be useful to implement modifications
of the behavior of the query, such as collecting all hits or
accumulating transparencies.

The \texttt{intersect} member specifies the user geometry callback to
get invoked for each user geometry encountered during traversal. The
user geometry callback specified this way has preference over the one
specified inside the geometry.

The \texttt{minWidthDistanceFactor} value controls the target size of
the curve radii when the min-width feature is enabled. Please see the
\hyperref[rtcsetgeometrymaxradiusscale]{rtcSetGeometryMaxRadiusScale}
function for more details on the min-width feature.

\subsubsection{EXIT STATUS}\label{exit-status}

No error code is set by this function.

\subsubsection{SEE ALSO}\label{see-also}

\hyperref[rtcoccluded1]{rtcOccluded1},
\hyperref[rtcoccluded4816]{rtcOccluded4/8/16},
\hyperref[rtcfeatureflags]{RTCFeatureFlags},
\hyperref[rtcinitrayquerycontext]{rtcInitRayQueryContext},
\hyperref[rtc_geometry_type_user]{RTC\_GEOMETRY\_TYPE\_USER},
\hyperref[rtcsetgeometrymaxradiusscale]{rtcSetGeometryMaxRadiusScale}

\pagebreak

\section{rtcInitRayQueryContext}\label{rtcinitrayquerycontext}

\subsubsection{NAME}\label{name}

\begin{Shaded}
\begin{Highlighting}[]
\NormalTok{rtcInitRayQueryContext }\OperatorTok{{-}}\NormalTok{ initializes the ray query context}
\end{Highlighting}
\end{Shaded}

\subsubsection{SYNOPSIS}\label{synopsis}

\begin{Shaded}
\begin{Highlighting}[]
\PreprocessorTok{\#include }\ImportTok{\textless{}embree4/rtcore.h\textgreater{}}

\KeywordTok{struct}\NormalTok{ RTCRayQueryContext}
\OperatorTok{\{}
  \PreprocessorTok{\#if RTC\_MAX\_INSTANCE\_LEVEL\_COUNT \textgreater{} 1}
    \DataTypeTok{unsigned} \DataTypeTok{int}\NormalTok{ instStackSize}\OperatorTok{;}
  \PreprocessorTok{\#endif}
  
  \DataTypeTok{unsigned} \DataTypeTok{int}\NormalTok{ instID}\OperatorTok{[}\NormalTok{RTC\_MAX\_INSTANCE\_LEVEL\_COUNT}\OperatorTok{];}
\OperatorTok{\};}

\DataTypeTok{void}\NormalTok{ rtcInitRayQueryContext}\OperatorTok{(}
  \KeywordTok{struct}\NormalTok{ RTCRayQueryContext}\OperatorTok{*}\NormalTok{ context}
\OperatorTok{);}
\end{Highlighting}
\end{Shaded}

\subsubsection{DESCRIPTION}\label{description}

The \texttt{rtcInitRayQueryContext} function initializes the
intersection context to default values and should be called to
initialize every ray query context.

It is guaranteed that the pointer to the ray query context
(\texttt{RTCRayQueryContext} type) is passed to the registered callback
functions. This way it is possible to attach arbitrary data to the end
of the ray query context, such as a per-ray payload.

Inside the user geometry callback the ray query context can get used to
access the \texttt{instID} stack to know which instance the user
geometry object resides.

If not ray query context is specified when tracing a ray, a default
context is used.

\subsubsection{EXIT STATUS}\label{exit-status}

No error code is set by this function.

\subsubsection{SEE ALSO}\label{see-also}

\hyperref[rtcintersect1]{rtcIntersect1},
\hyperref[rtcintersect4816]{rtcIntersect4/8/16},
\hyperref[rtcoccluded1]{rtcOccluded1},
\hyperref[rtcoccluded4816]{rtcOccluded4/8/16}

\pagebreak

\section{rtcIntersect1}\label{rtcintersect1}

\subsubsection{NAME}\label{name}

\begin{Shaded}
\begin{Highlighting}[]
\NormalTok{rtcIntersect1 }\OperatorTok{{-}}\NormalTok{ finds the closest hit }\ControlFlowTok{for}\NormalTok{ a single ray}
\end{Highlighting}
\end{Shaded}

\subsubsection{SYNOPSIS}\label{synopsis}

\begin{Shaded}
\begin{Highlighting}[]
\PreprocessorTok{\#include }\ImportTok{\textless{}embree4/rtcore.h\textgreater{}}

\DataTypeTok{void}\NormalTok{ rtcIntersect1}\OperatorTok{(}
\NormalTok{  RTCScene scene}\OperatorTok{,}
  \KeywordTok{struct}\NormalTok{ RTCRayHit}\OperatorTok{*}\NormalTok{ rayhit}
  \KeywordTok{struct}\NormalTok{ RTCIntersectArguments}\OperatorTok{*}\NormalTok{ args }\OperatorTok{=}\NormalTok{ NULL}
\OperatorTok{);}
\end{Highlighting}
\end{Shaded}

\subsubsection{DESCRIPTION}\label{description}

The \texttt{rtcIntersect1} function finds the closest hit of a single
ray (\texttt{rayhit} argument) with the scene (\texttt{scene} argument).
The provided ray/hit structure contains the ray to intersect and some
hit output fields that are filled when a hit is found. The passed
optional arguments struct (\texttt{args} argument) can get used for
advanced use cases, see section
\hyperref[rtcinitintersectarguments]{rtcInitIntersectArguments} for more
details.

To trace a ray, the user has to initialize the ray origin (\texttt{org}
ray member), ray direction (\texttt{dir} ray member), ray segment
(\texttt{tnear}, \texttt{tfar} ray members), ray mask (\texttt{mask} ray
member), and set the ray flags to \texttt{0} (\texttt{flags} ray
member). The ray time (\texttt{time} ray member) must be initialized to
a value in the range \${[}0, 1{]}. The ray segment has to be in the
range \([0, \infty]\), thus ranges that start behind the ray origin are
not valid, but ranges can reach to infinity. See Section
\hyperref[rtcray]{RTCRay} for the ray layout description.

The geometry ID (\texttt{geomID} hit member) of the hit data must be
initialized to \texttt{RTC\_INVALID\_GEOMETRY\_ID} (-1).

When no intersection is found, the ray/hit data is not updated. When an
intersection is found, the hit distance is written into the
\texttt{tfar} member of the ray and all hit data is set, such as
unnormalized geometry normal in object space (\texttt{Ng} hit member),
local hit coordinates (\texttt{u}, \texttt{v} hit member), instance ID
stack (\texttt{instID} hit member), geometry ID (\texttt{geomID} hit
member), and primitive ID (\texttt{primID} hit member). See Section
\hyperref[rtchit]{RTCHit} for the hit layout description.

If the instance ID stack has a prefix of values not equal to
\texttt{RTC\_INVALID\_GEOMETRY\_ID}, the instance ID on each level
corresponds to the geometry ID of the hit instance of the higher-level
scene, the geometry ID corresponds to the hit geometry inside the hit
instanced scene, and the primitive ID corresponds to the n-th primitive
of that geometry.

If level 0 of the instance ID stack is equal to
\texttt{RTC\_INVALID\_GEOMETRY\_ID}, the geometry ID corresponds to the
hit geometry inside the top-level scene, and the primitive ID
corresponds to the n-th primitive of that geometry.

The implementation makes no guarantees that primitives whose hit
distance is exactly at (or very close to) \texttt{tnear} or
\texttt{tfar} are hit or missed. If you want to exclude intersections at
\texttt{tnear} just pass a slightly enlarged \texttt{tnear}, and if you
want to include intersections at \texttt{tfar} pass a slightly enlarged
\texttt{tfar}.

The ray pointer passed to callback functions is not guaranteed to be
identical to the original ray provided. To extend the ray with
additional data to be accessed in callback functions, use the ray query
context. See section
\hyperref[rtcinitrayquerycontext]{rtcInitRayQueryContext} for more
details.

The ray/hit structure must be aligned to 16 bytes.

\subsubsection{EXIT STATUS}\label{exit-status}

For performance reasons this function does not do any error checks, thus
will not set any error flags on failure.

\subsubsection{SEE ALSO}\label{see-also}

\hyperref[rtcoccluded1]{rtcOccluded1},
\hyperref[rtcintersect4816]{rtcIntersect4/8/16},
\hyperref[rtcrayhit]{RTCRayHit},
\hyperref[rtcinitintersectarguments]{rtcInitIntersectArguments}

\pagebreak

\section{rtcOccluded1}\label{rtcoccluded1}

\subsubsection{NAME}\label{name}

\begin{Shaded}
\begin{Highlighting}[]
\NormalTok{rtcOccluded1 }\OperatorTok{{-}}\NormalTok{ finds any hit }\ControlFlowTok{for}\NormalTok{ a single ray}
\end{Highlighting}
\end{Shaded}

\subsubsection{SYNOPSIS}\label{synopsis}

\begin{Shaded}
\begin{Highlighting}[]
\PreprocessorTok{\#include }\ImportTok{\textless{}embree4/rtcore.h\textgreater{}}

\DataTypeTok{void}\NormalTok{ rtcOccluded1}\OperatorTok{(}
\NormalTok{  RTCScene scene}\OperatorTok{,}
  \KeywordTok{struct}\NormalTok{ RTCRay}\OperatorTok{*}\NormalTok{ ray}\OperatorTok{,}
  \KeywordTok{struct}\NormalTok{ RTCOccludedArguments}\OperatorTok{*}\NormalTok{ args }\OperatorTok{=}\NormalTok{ NULL}
\OperatorTok{);}
\end{Highlighting}
\end{Shaded}

\subsubsection{DESCRIPTION}\label{description}

The \texttt{rtcOccluded1} function checks for a single ray (\texttt{ray}
argument) whether there is any hit with the scene (\texttt{scene}
argument). The passed optional arguments struct (\texttt{args} argument)
can get used for advanced use cases, see section
\hyperref[rtcinitoccludedarguments]{rtcInitOccludedArguments} for more
details.

To trace a ray, the user must initialize the ray origin (\texttt{org}
ray member), ray direction (\texttt{dir} ray member), ray segment
(\texttt{tnear}, \texttt{tfar} ray members), ray mask (\texttt{mask} ray
member), and must set the ray flags to \texttt{0} (\texttt{flags} ray
member). The ray time (\texttt{time} ray member) must be initialized to
a value in the range \([0, 1]\). The ray segment must be in the range
\([0, \infty]\), thus ranges that start behind the ray origin are not
valid, but ranges can reach to infinity. See Section
\hyperref[rtcray]{RTCRay} for the ray layout description.

When no intersection is found, the ray data is not updated. In case a
hit was found, the \texttt{tfar} component of the ray is set to
\texttt{-inf}.

The implementation makes no guarantees that primitives whose hit
distance is exactly at (or very close to) \texttt{tnear} or
\texttt{tfar} are hit or missed. If you want to exclude intersections at
\texttt{tnear} just pass a slightly enlarged \texttt{tnear}, and if you
want to include intersections at \texttt{tfar} pass a slightly enlarged
\texttt{tfar}.

The ray pointer passed to callback functions is not guaranteed to be
identical to the original ray provided. To extend the ray with
additional data to be accessed in callback functions, use the ray query
context. See section
\hyperref[rtcinitrayquerycontext]{rtcInitRayQueryContext} for more
details.

The ray must be aligned to 16 bytes.

\subsubsection{EXIT STATUS}\label{exit-status}

For performance reasons this function does not do any error checks, thus
will not set any error flags on failure.

\subsubsection{SEE ALSO}\label{see-also}

\hyperref[rtcintersect1]{rtcIntersect1},
\hyperref[rtcoccluded4816]{rtcOccluded4/8/16},
\hyperref[rtcray]{RTCRay},
\hyperref[rtcinitoccludedarguments]{rtcInitOccludedArguments}

\pagebreak

\section{rtcIntersect4/8/16}\label{rtcintersect4816}

\subsubsection{NAME}\label{name}

\begin{Shaded}
\begin{Highlighting}[]
\NormalTok{rtcIntersect4}\OperatorTok{/}\DecValTok{8}\OperatorTok{/}\DecValTok{16} \OperatorTok{{-}}\NormalTok{ finds the closest hits }\ControlFlowTok{for}\NormalTok{ a ray packet}
\end{Highlighting}
\end{Shaded}

\subsubsection{SYNOPSIS}\label{synopsis}

\begin{Shaded}
\begin{Highlighting}[]
\PreprocessorTok{\#include }\ImportTok{\textless{}embree4/rtcore.h\textgreater{}}

\DataTypeTok{void}\NormalTok{ rtcIntersect4}\OperatorTok{(}
  \AttributeTok{const} \DataTypeTok{int}\OperatorTok{*}\NormalTok{ valid}\OperatorTok{,}
\NormalTok{  RTCScene scene}\OperatorTok{,}
  \KeywordTok{struct}\NormalTok{ RTCRayHit4}\OperatorTok{*}\NormalTok{ rayhit}\OperatorTok{,}
  \KeywordTok{struct}\NormalTok{ RTCIntersectArguments}\OperatorTok{*}\NormalTok{ args }\OperatorTok{=}\NormalTok{ NULL}
\OperatorTok{);}

\DataTypeTok{void}\NormalTok{ rtcIntersect8}\OperatorTok{(}
  \AttributeTok{const} \DataTypeTok{int}\OperatorTok{*}\NormalTok{ valid}\OperatorTok{,}
\NormalTok{  RTCScene scene}\OperatorTok{,}
  \KeywordTok{struct}\NormalTok{ RTCRayHit8}\OperatorTok{*}\NormalTok{ rayhit}\OperatorTok{,}
  \KeywordTok{struct}\NormalTok{ RTCIntersectArguments}\OperatorTok{*}\NormalTok{ args }\OperatorTok{=}\NormalTok{ NULL}
\OperatorTok{);}

\DataTypeTok{void}\NormalTok{ rtcIntersect16}\OperatorTok{(}
  \AttributeTok{const} \DataTypeTok{int}\OperatorTok{*}\NormalTok{ valid}\OperatorTok{,}
\NormalTok{  RTCScene scene}\OperatorTok{,}
  \KeywordTok{struct}\NormalTok{ RTCRayHit16}\OperatorTok{*}\NormalTok{ rayhit}\OperatorTok{,}
  \KeywordTok{struct}\NormalTok{ RTCIntersectArguments}\OperatorTok{*}\NormalTok{ args }\OperatorTok{=}\NormalTok{ NULL}
\OperatorTok{);}
\end{Highlighting}
\end{Shaded}

\subsubsection{DESCRIPTION}\label{description}

The \texttt{rtcIntersect4/8/16} functions finds the closest hits for a
ray packet of size 4, 8, or 16 (\texttt{rayhit} argument) with the scene
(\texttt{scene} argument). The ray/hit input contains a ray packet and
hit packet. The passed optional arguments struct (\texttt{args}
argument) are used to pass additional arguments for advanced features.
See Section \hyperref[rtcintersect1]{rtcIntersect1} for more details and
a description of how to set up and trace rays.

A ray valid mask must be provided (\texttt{valid} argument) which stores
one 32-bit integer (\texttt{-1} means valid and \texttt{0} invalid) per
ray in the packet. Only active rays are processed, and hit data of
inactive rays is not changed.

The ray pointer passed to callback functions is not guaranteed to be
identical to the original ray provided. To extend the ray with
additional data to be accessed in callback functions, use the ray query
context. See section
\hyperref[rtcinitrayquerycontext]{rtcInitRayQueryContext} for more
details.

For \texttt{rtcIntersect4} the ray packet must be aligned to 16 bytes,
for \texttt{rtcIntersect8} the alignment must be 32 bytes, and for
\texttt{rtcIntersect16} the alignment must be 64 bytes.

The \texttt{rtcIntersect4}, \texttt{rtcIntersect8} and
\texttt{rtcIntersect16} functions may change the ray packet size and ray
order when calling back into filter functions or user geometry
callbacks. Under some conditions the application can assume packets to
stay intakt, which can determined by querying the
\texttt{RTC\_DEVICE\_PROPERTY\_NATIVE\_RAY4\_SUPPORTED},
\texttt{RTC\_DEVICE\_PROPERTY\_NATIVE\_RAY8\_SUPPORTED},
\texttt{RTC\_DEVICE\_PROPERTY\_NATIVE\_RAY16\_SUPPORTED} properties
through the \texttt{rtcGetDeviceProperty} function. See
\hyperref[rtcgetdeviceproperty]{rtcGetDeviceProperty} for more
information.

\subsubsection{EXIT STATUS}\label{exit-status}

For performance reasons this function does not do any error checks, thus
will not set any error flags on failure.

\subsubsection{SEE ALSO}\label{see-also}

\hyperref[rtcintersect1]{rtcIntersect1},
\hyperref[rtcoccluded4816]{rtcOccluded4/8/16},
\hyperref[rtcinitintersectarguments]{rtcInitIntersectArguments}

\pagebreak

\section{rtcOccluded4/8/16}\label{rtcoccluded4816}

\subsubsection{NAME}\label{name}

\begin{Shaded}
\begin{Highlighting}[]
\NormalTok{rtcOccluded4}\OperatorTok{/}\DecValTok{8}\OperatorTok{/}\DecValTok{16} \OperatorTok{{-}}\NormalTok{ finds any hits }\ControlFlowTok{for}\NormalTok{ a ray packet}
\end{Highlighting}
\end{Shaded}

\subsubsection{SYNOPSIS}\label{synopsis}

\begin{Shaded}
\begin{Highlighting}[]
\PreprocessorTok{\#include }\ImportTok{\textless{}embree4/rtcore.h\textgreater{}}

\DataTypeTok{void}\NormalTok{ rtcOccluded4}\OperatorTok{(}
  \AttributeTok{const} \DataTypeTok{int}\OperatorTok{*}\NormalTok{ valid}\OperatorTok{,}
\NormalTok{  RTCScene scene}\OperatorTok{,}
  \KeywordTok{struct}\NormalTok{ RTCRay4}\OperatorTok{*}\NormalTok{ ray}\OperatorTok{,}
  \KeywordTok{struct}\NormalTok{ RTCOccludedArguments}\OperatorTok{*}\NormalTok{ args }\OperatorTok{=}\NormalTok{ NULL}
\OperatorTok{);}

\DataTypeTok{void}\NormalTok{ rtcOccluded8}\OperatorTok{(}
  \AttributeTok{const} \DataTypeTok{int}\OperatorTok{*}\NormalTok{ valid}\OperatorTok{,}
\NormalTok{  RTCScene scene}\OperatorTok{,}
  \KeywordTok{struct}\NormalTok{ RTCRay8}\OperatorTok{*}\NormalTok{ ray}\OperatorTok{,}
  \KeywordTok{struct}\NormalTok{ RTCOccludedArguments}\OperatorTok{*}\NormalTok{ args }\OperatorTok{=}\NormalTok{ NULL}
\OperatorTok{);}

\DataTypeTok{void}\NormalTok{ rtcOccluded16}\OperatorTok{(}
  \AttributeTok{const} \DataTypeTok{int}\OperatorTok{*}\NormalTok{ valid}\OperatorTok{,}
\NormalTok{  RTCScene scene}\OperatorTok{,}
  \KeywordTok{struct}\NormalTok{ RTCRay16}\OperatorTok{*}\NormalTok{ ray}\OperatorTok{,}
  \KeywordTok{struct}\NormalTok{ RTCOccludedArguments}\OperatorTok{*}\NormalTok{ args }\OperatorTok{=}\NormalTok{ NULL}
\OperatorTok{);}
\end{Highlighting}
\end{Shaded}

\subsubsection{DESCRIPTION}\label{description}

The \texttt{rtcOccluded4/8/16} functions checks for each active ray of
the ray packet of size 4, 8, or 16 (\texttt{ray} argument) whether there
is any hit with the scene (\texttt{scene} argument). The passed optional
arguments struct (\texttt{args} argument) can get used for advanced use
cases, see section
\hyperref[rtcinitoccludedarguments]{rtcInitOccludedArguments} for more
details. See Section \hyperref[rtcoccluded1]{rtcOccluded1} for more
details and a description of how to set up and trace occlusion rays.

A ray valid mask must be provided (\texttt{valid} argument) which stores
one 32-bit integer (\texttt{-1} means valid and \texttt{0} invalid) per
ray in the packet. Only active rays are processed, and hit data of
inactive rays is not changed.

The ray pointer passed to callback functions is not guaranteed to be
identical to the original ray provided. To extend the ray with
additional data to be accessed in callback functions, use the ray query
context. See section
\hyperref[rtcinitrayquerycontext]{rtcInitRayQueryContext} for more
details.

For \texttt{rtcOccluded4} the ray packet must be aligned to 16 bytes,
for \texttt{rtcOccluded8} the alignment must be 32 bytes, and for
\texttt{rtcOccluded16} the alignment must be 64 bytes.

The \texttt{rtcOccluded4}, \texttt{rtcOccluded8} and
\texttt{rtcOccluded16} functions may change the ray packet size and ray
order when calling back into intersect filter functions or user geometry
callbacks. Under some conditions the application can assume packets to
stay intakt, which can determined by querying the
\texttt{RTC\_DEVICE\_PROPERTY\_NATIVE\_RAY4\_SUPPORTED},
\texttt{RTC\_DEVICE\_PROPERTY\_NATIVE\_RAY8\_SUPPORTED},
\texttt{RTC\_DEVICE\_PROPERTY\_NATIVE\_RAY16\_SUPPORTED} properties
through the \texttt{rtcGetDeviceProperty} function. See
\hyperref[rtcgetdeviceproperty]{rtcGetDeviceProperty} for more
information.

\subsubsection{EXIT STATUS}\label{exit-status}

For performance reasons this function does not do any error checks, thus
will not set any error flags on failure.

\subsubsection{SEE ALSO}\label{see-also}

\hyperref[rtcoccluded1]{rtcOccluded1},
\hyperref[rtcintersect4816]{rtcIntersect4/8/16},
\hyperref[rtcinitoccludedarguments]{rtcInitOccludedArguments}

\pagebreak

\section{rtcForwardIntersect1}\label{rtcforwardintersect1}

\subsubsection{NAME}\label{name}

\begin{Shaded}
\begin{Highlighting}[]
\NormalTok{rtcForwardIntersect1 }\OperatorTok{{-}}\NormalTok{ forwards a single ray to }\KeywordTok{new}\NormalTok{ scene}
\NormalTok{  from user geometry callback}
\end{Highlighting}
\end{Shaded}

\subsubsection{SYNOPSIS}\label{synopsis}

\begin{Shaded}
\begin{Highlighting}[]
\PreprocessorTok{\#include }\ImportTok{\textless{}embree4/rtcore.h\textgreater{}}

\DataTypeTok{void}\NormalTok{ rtcForwardIntersect1}\OperatorTok{(}
  \AttributeTok{const} \KeywordTok{struct}\NormalTok{ RTCIntersectFunctionNArguments}\OperatorTok{*}\NormalTok{ args}\OperatorTok{,}
\NormalTok{  RTCScene scene}\OperatorTok{,}
  \KeywordTok{struct}\NormalTok{ RTCRay}\OperatorTok{*}\NormalTok{ ray}\OperatorTok{,}
  \DataTypeTok{unsigned} \DataTypeTok{int}\NormalTok{ instID}
\OperatorTok{);}
\end{Highlighting}
\end{Shaded}

\subsubsection{DESCRIPTION}\label{description}

The \texttt{rtcForwardIntersect1} function forwards the traversal of a
transformed ray (\texttt{ray} argument) into a scene (\texttt{scene}
argument) from a user geometry callback. The function can only get
invoked from a user geometry callback for a ray traversal initiated with
the \texttt{rtcIntersect1} function. The callback arguments structure of
the callback invokation has to get passed to the ray forwarding
(\texttt{args} argument). The user geometry callback should instantly
terminate after invoking the \texttt{rtcForwardIntersect1} function.

Only the ray origin and ray direction members of the ray argument are
used for forwarding, all additional ray properties are inherited from
the initial ray traversal invokation of \texttt{rtcIntersect1}.

The implementation of the \texttt{rtcForwardIntersect1} function
recursively continues the ray traversal into the specified scene and
pushes the provided instance ID (\texttt{instID} argument) to the
instance ID stack. Hit information is updated into the ray hit structure
passed to the original \texttt{rtcIntersect1} invokation.

This function can get used to implement user defined instancing using
user geometries, e.g.~by transforming the ray in a special way, and/or
selecting between different scenes to instantiate.

When using Embree on the CPU it is possible to recursively invoke
\texttt{rtcIntersect1} directly from a user geometry callback. However,
when SYCL is used, recursively tracing rays is not directly supported,
and the \texttt{rtcForwardIntersect1} function must be used.

The ray structure must be aligned to 16 bytes.

\subsubsection{EXIT STATUS}\label{exit-status}

For performance reasons this function does not do any error checks, thus
will not set any error flags on failure.

\subsubsection{SEE ALSO}\label{see-also}

\hyperref[rtcintersect1]{rtcIntersect1}, \hyperref[rtcray]{RTCRay}

\pagebreak

\section{rtcForwardOccluded1}\label{rtcforwardoccluded1}

\subsubsection{NAME}\label{name}

\begin{Shaded}
\begin{Highlighting}[]
\NormalTok{rtcForwardOccluded1 }\OperatorTok{{-}}\NormalTok{ forwards a single ray to }\KeywordTok{new}\NormalTok{ scene}
\NormalTok{  from user geometry callback}
\end{Highlighting}
\end{Shaded}

\subsubsection{SYNOPSIS}\label{synopsis}

\begin{Shaded}
\begin{Highlighting}[]
\PreprocessorTok{\#include }\ImportTok{\textless{}embree4/rtcore.h\textgreater{}}

\DataTypeTok{void}\NormalTok{ rtcForwardOccluded1}\OperatorTok{(}
  \AttributeTok{const} \KeywordTok{struct}\NormalTok{ RTCOccludedFunctionNArguments}\OperatorTok{*}\NormalTok{ args}\OperatorTok{,}
\NormalTok{  RTCScene scene}\OperatorTok{,}
  \KeywordTok{struct}\NormalTok{ RTCRay}\OperatorTok{*}\NormalTok{ ray}\OperatorTok{,}
  \DataTypeTok{unsigned} \DataTypeTok{int}\NormalTok{ instID}
\OperatorTok{);}
\end{Highlighting}
\end{Shaded}

\subsubsection{DESCRIPTION}\label{description}

The \texttt{rtcForwardOccluded1} function forwards the traversal of a
transformed ray (\texttt{ray} argument) into a scene (\texttt{scene}
argument) from a user geometry callback. The function can only get
invoked from a user geometry callback for a ray traversal initiated with
the \texttt{rtcOccluded1} function. The callback arguments structure of
the callback invokation has to get passed to the ray forwarding
(\texttt{args} argument). The user geometry callback should instantly
terminate after invoking the \texttt{rtcForwardOccluded1} function.

Only the ray origin and ray direction members of the ray argument are
used for forwarding, all additional ray properties are inherited from
the initial ray traversal invokation of \texttt{rtcOccluded1}.

The implementation of the \texttt{rtcForwardOccluded1} function
recursively continues the ray traversal into the specified scene and
pushes the provided instance ID (\texttt{instID} argument) to the
instance ID stack. Hit information is updated into the ray structure
passed to the original \texttt{rtcOccluded1} invokation.

This function can get used to implement user defined instancing using
user geometries, e.g.~by transforming the ray in a special way, and/or
selecting between different scenes to instantiate.

When using Embree on the CPU it is possible to recursively invoke
\texttt{rtcOccluded1} directly from a user geometry callback. However,
when SYCL is used, recursively tracing rays is not directly supported,
and the \texttt{rtcForwardOccluded1} function must be used.

The ray structure must be aligned to 16 bytes.

\subsubsection{EXIT STATUS}\label{exit-status}

For performance reasons this function does not do any error checks, thus
will not set any error flags on failure.

\subsubsection{SEE ALSO}\label{see-also}

\hyperref[rtcoccluded1]{rtcOccluded1}, \hyperref[rtcray]{RTCRay}

\pagebreak

\section{rtcForwardIntersect4/8/16}\label{rtcforwardintersect4816}

\subsubsection{NAME}\label{name}

\begin{Shaded}
\begin{Highlighting}[]
\NormalTok{rtcForwardIntersect4}\OperatorTok{/}\DecValTok{8}\OperatorTok{/}\DecValTok{16} \OperatorTok{{-}}\NormalTok{ forwards a ray packet to }\KeywordTok{new}\NormalTok{ scene}
\NormalTok{  from user geometry callback}
\end{Highlighting}
\end{Shaded}

\subsubsection{SYNOPSIS}\label{synopsis}

\begin{Shaded}
\begin{Highlighting}[]
\PreprocessorTok{\#include }\ImportTok{\textless{}embree4/rtcore.h\textgreater{}}

\DataTypeTok{void}\NormalTok{ rtcForwardIntersect4}\OperatorTok{(}
  \DataTypeTok{void} \DataTypeTok{int}\OperatorTok{*}\NormalTok{ valid}\OperatorTok{,}
  \AttributeTok{const} \KeywordTok{struct}\NormalTok{ RTCIntersectFunctionNArguments}\OperatorTok{*}\NormalTok{ args}\OperatorTok{,}
\NormalTok{  RTCScene scene}\OperatorTok{,}
  \KeywordTok{struct}\NormalTok{ RTCRay4}\OperatorTok{*}\NormalTok{ ray}\OperatorTok{,}
  \DataTypeTok{unsigned} \DataTypeTok{int}\NormalTok{ instID}
\OperatorTok{);}

\DataTypeTok{void}\NormalTok{ rtcForwardIntersect4}\OperatorTok{(}
  \DataTypeTok{void} \DataTypeTok{int}\OperatorTok{*}\NormalTok{ valid}\OperatorTok{,}
  \AttributeTok{const} \KeywordTok{struct}\NormalTok{ RTCIntersectFunctionNArguments}\OperatorTok{*}\NormalTok{ args}\OperatorTok{,}
\NormalTok{  RTCScene scene}\OperatorTok{,}
  \KeywordTok{struct}\NormalTok{ RTCRay4}\OperatorTok{*}\NormalTok{ ray}\OperatorTok{,}
  \DataTypeTok{unsigned} \DataTypeTok{int}\NormalTok{ instID}
\OperatorTok{);}

\DataTypeTok{void}\NormalTok{ rtcForwardIntersect16}\OperatorTok{(}
  \DataTypeTok{void} \DataTypeTok{int}\OperatorTok{*}\NormalTok{ valid}\OperatorTok{,}
  \AttributeTok{const} \KeywordTok{struct}\NormalTok{ RTCIntersectFunctionNArguments}\OperatorTok{*}\NormalTok{ args}\OperatorTok{,}
\NormalTok{  RTCScene scene}\OperatorTok{,}
  \KeywordTok{struct}\NormalTok{ RTCRay16}\OperatorTok{*}\NormalTok{ ray}\OperatorTok{,}
  \DataTypeTok{unsigned} \DataTypeTok{int}\NormalTok{ instID}
\OperatorTok{);}
\end{Highlighting}
\end{Shaded}

\subsubsection{DESCRIPTION}\label{description}

The \texttt{rtcForwardIntersect4/8/16} functions forward the traversal
of a transformed ray packet (\texttt{ray} argument) into a scene
(\texttt{scene} argument) from a user geometry callback. The function
can only get invoked from a user geometry callback for a ray traversal
initiated with the \texttt{rtcIntersect4/8/16} function. The callback
arguments structure of the callback invokation has to get passed to the
ray forwarding (\texttt{args} argument). The user geometry callback
should instantly terminate after invoking the
\texttt{rtcForwardIntersect4/8/16} function.

Only the ray origin and ray direction members of the ray argument are
used for forwarding, all additional ray properties are inherited from
the initial ray traversal invokation of \texttt{rtcIntersect4/8/16}.

The implementation of the \texttt{rtcForwardIntersect4/8/16} function
recursively continues the ray traversal into the specified scene and
pushes the provided instance ID (\texttt{instID} argument) to the
instance ID stack. Hit information is updated into the ray hit structure
passed to the original \texttt{rtcIntersect4/8/16} invokation.

This function can get used to implement user defined instancing using
user geometries, e.g.~by transforming the ray in a special way, and/or
selecting between different scenes to instantiate.

When using Embree on the CPU it is possible to recursively invoke
\texttt{rtcIntersect4/8/16} directly from a user geometry callback.
However, when SYCL is used, recursively tracing rays is not directly
supported, and the \texttt{rtcForwardIntersect4/8/16} function must be
used.

For \texttt{rtcForwardIntersect4} the ray packet must be aligned to 16
bytes, for \texttt{rtcForwardIntersect8} the alignment must be 32 bytes,
and for \texttt{rtcForwardIntersect16} the alignment must be 64 bytes.

\subsubsection{EXIT STATUS}\label{exit-status}

For performance reasons this function does not do any error checks, thus
will not set any error flags on failure.

\subsubsection{SEE ALSO}\label{see-also}

\hyperref[rtcintersect4816]{rtcIntersect4/8/16}

\pagebreak

\section{rtcForwardOccluded4/8/16}\label{rtcforwardoccluded4816}

\subsubsection{NAME}\label{name}

\begin{Shaded}
\begin{Highlighting}[]
\NormalTok{rtcForwardOccluded4}\OperatorTok{/}\DecValTok{8}\OperatorTok{/}\DecValTok{16} \OperatorTok{{-}}\NormalTok{ forwards a ray packet to }\KeywordTok{new}\NormalTok{ scene}
\NormalTok{  from user geometry callback}
\end{Highlighting}
\end{Shaded}

\subsubsection{SYNOPSIS}\label{synopsis}

\begin{Shaded}
\begin{Highlighting}[]
\PreprocessorTok{\#include }\ImportTok{\textless{}embree4/rtcore.h\textgreater{}}

\DataTypeTok{void}\NormalTok{ rtcForwardOccluded4}\OperatorTok{(}
  \DataTypeTok{void} \DataTypeTok{int}\OperatorTok{*}\NormalTok{ valid}\OperatorTok{,}
  \AttributeTok{const} \KeywordTok{struct}\NormalTok{ RTCOccludedFunctionNArguments}\OperatorTok{*}\NormalTok{ args}\OperatorTok{,}
\NormalTok{  RTCScene scene}\OperatorTok{,}
  \KeywordTok{struct}\NormalTok{ RTCRay4}\OperatorTok{*}\NormalTok{ ray}\OperatorTok{,}
  \DataTypeTok{unsigned} \DataTypeTok{int}\NormalTok{ instID}
\OperatorTok{);}

\DataTypeTok{void}\NormalTok{ rtcForwardOccluded4}\OperatorTok{(}
  \DataTypeTok{void} \DataTypeTok{int}\OperatorTok{*}\NormalTok{ valid}\OperatorTok{,}
  \AttributeTok{const} \KeywordTok{struct}\NormalTok{ RTCOccludedFunctionNArguments}\OperatorTok{*}\NormalTok{ args}\OperatorTok{,}
\NormalTok{  RTCScene scene}\OperatorTok{,}
  \KeywordTok{struct}\NormalTok{ RTCRay4}\OperatorTok{*}\NormalTok{ ray}\OperatorTok{,}
  \DataTypeTok{unsigned} \DataTypeTok{int}\NormalTok{ instID}
\OperatorTok{);}

\DataTypeTok{void}\NormalTok{ rtcForwardOccluded16}\OperatorTok{(}
  \DataTypeTok{void} \DataTypeTok{int}\OperatorTok{*}\NormalTok{ valid}\OperatorTok{,}
  \AttributeTok{const} \KeywordTok{struct}\NormalTok{ RTCOccludedFunctionNArguments}\OperatorTok{*}\NormalTok{ args}\OperatorTok{,}
\NormalTok{  RTCScene scene}\OperatorTok{,}
  \KeywordTok{struct}\NormalTok{ RTCRay16}\OperatorTok{*}\NormalTok{ ray}\OperatorTok{,}
  \DataTypeTok{unsigned} \DataTypeTok{int}\NormalTok{ instID}
\OperatorTok{);}
\end{Highlighting}
\end{Shaded}

\subsubsection{DESCRIPTION}\label{description}

The \texttt{rtcForwardOccluded4/8/16} functions forward the traversal of
a transformed ray packet (\texttt{ray} argument) into a scene
(\texttt{scene} argument) from a user geometry callback. The function
can only get invoked from a user geometry callback for a ray traversal
initiated with the \texttt{rtcOccluded4/8/16} function. The callback
arguments structure of the callback invokation has to get passed to the
ray forwarding (\texttt{args} argument). The user geometry callback
should instantly terminate after invoking the
\texttt{rtcForwardOccluded4/8/16} function.

Only the ray origin and ray direction members of the ray argument are
used for forwarding, all additional ray properties are inherited from
the initial ray traversal invokation of \texttt{rtcOccluded4/8/16}.

The implementation of the \texttt{rtcForwardOccluded4/8/16} function
recursively continues the ray traversal into the specified scene and
pushes the provided instance ID (\texttt{instID} argument) to the
instance ID stack. Hit information is updated into the ray structure
passed to the original \texttt{rtcOccluded4/8/16} invokation.

This function can get used to implement user defined instancing using
user geometries, e.g.~by transforming the ray in a special way, and/or
selecting between different scenes to instantiate.

When using Embree on the CPU it is possible to recursively invoke
\texttt{rtcOccluded4/8/16} directly from a user geometry callback.
However, when SYCL is used, recursively tracing rays is not directly
supported, and the \texttt{rtcForwardOccluded4/8/16} function must be
used.

For \texttt{rtcForwardOccluded4} the ray packet must be aligned to 16
bytes, for \texttt{rtcForwardOccluded8} the alignment must be 32 bytes,
and for \texttt{rtcForwardOccluded16} the alignment must be 64 bytes.

\subsubsection{EXIT STATUS}\label{exit-status}

For performance reasons this function does not do any error checks, thus
will not set any error flags on failure.

\subsubsection{SEE ALSO}\label{see-also}

\hyperref[rtcoccluded4816]{rtcOccluded4/8/16}

\pagebreak

\section{rtcInitPointQueryContext}\label{rtcinitpointquerycontext}

\subsubsection{NAME}\label{name}

\begin{Shaded}
\begin{Highlighting}[]
\NormalTok{rtcInitPointQueryContext }\OperatorTok{{-}}\NormalTok{ initializes the context information }\OperatorTok{(}\NormalTok{e}\OperatorTok{.}\NormalTok{g}\OperatorTok{.}
\NormalTok{  stack of }\OperatorTok{(}\NormalTok{multilevel}\OperatorTok{{-})}\NormalTok{instance transformations}\OperatorTok{)} \ControlFlowTok{for}\NormalTok{ point queries}
\end{Highlighting}
\end{Shaded}

\subsubsection{SYNOPSIS}\label{synopsis}

\begin{Shaded}
\begin{Highlighting}[]
\PreprocessorTok{\#include }\ImportTok{\textless{}embree4/rtcore.h\textgreater{}}

\KeywordTok{struct}\NormalTok{ RTC\_ALIGN}\OperatorTok{(}\DecValTok{16}\OperatorTok{)}\NormalTok{ RTCPointQueryContext}
\OperatorTok{\{}
  \CommentTok{// accumulated 4x4 column major matrices from world to instance space.}
  \DataTypeTok{float}\NormalTok{ world2inst}\OperatorTok{[}\NormalTok{RTC\_MAX\_INSTANCE\_LEVEL\_COUNT}\OperatorTok{][}\DecValTok{16}\OperatorTok{];}
  
  \CommentTok{// accumulated 4x4 column major matrices from instance to world space.}
  \DataTypeTok{float}\NormalTok{ inst2world}\OperatorTok{[}\NormalTok{RTC\_MAX\_INSTANCE\_LEVEL\_COUNT}\OperatorTok{][}\DecValTok{16}\OperatorTok{];}

  \CommentTok{// instance ids.}
  \DataTypeTok{unsigned} \DataTypeTok{int}\NormalTok{ instID}\OperatorTok{[}\NormalTok{RTC\_MAX\_INSTANCE\_LEVEL\_COUNT}\OperatorTok{];}
  
  \CommentTok{// number of instances currently on the stack.}
  \DataTypeTok{unsigned} \DataTypeTok{int}\NormalTok{ instStackSize}\OperatorTok{;}
\OperatorTok{\};}

\DataTypeTok{void}\NormalTok{ rtcInitPointQueryContext}\OperatorTok{(}
  \KeywordTok{struct}\NormalTok{ RTCPointQueryContext}\OperatorTok{*}\NormalTok{ context}
\OperatorTok{);}
\end{Highlighting}
\end{Shaded}

\subsubsection{DESCRIPTION}\label{description}

A stack (\texttt{RTCPointQueryContext} type) which stores the IDs and
instance transformations during a BVH traversal for a point query. The
transformations are assumed to be affine transformations (3×3 matrix
plus translation) and therefore the last column is ignored (see
\hyperref[rtc_geometry_type_instance]{RTC\_GEOMETRY\_TYPE\_INSTANCE} for
details).

The \texttt{rtcInitPointContext} function initializes the context to
default values and should be called for initialization.

The context will be passed as an argument to the point query callback
function (see
\hyperref[rtcsetgeometrypointqueryfunction]{rtcSetGeometryPointQueryFunction})
and should be used to pass instance information down the instancing
chain for user defined instancing (see tutorial {[}ClosestPoint{]} for a
reference implementation of point queries with user defined instancing).

The context is an necessary argument to
\hyperref[rtcpointquery]{rtcPointQuery} and Embree internally uses the
topmost instance transformation of the stack to transform the point
query into instance space.

\subsubsection{EXIT STATUS}\label{exit-status}

No error code is set by this function.

\subsubsection{SEE ALSO}\label{see-also}

\hyperref[rtcpointquery]{rtcPointQuery},
\hyperref[rtcsetgeometrypointqueryfunction]{rtcSetGeometryPointQueryFunction}

\pagebreak

\section{rtcPointQuery}\label{rtcpointquery}

\subsubsection{NAME}\label{name}

\begin{Shaded}
\begin{Highlighting}[]
\NormalTok{rtcPointQuery }\OperatorTok{{-}}\NormalTok{ traverses the BVH with a point query object}
\end{Highlighting}
\end{Shaded}

\subsubsection{SYNOPSIS}\label{synopsis}

\begin{Shaded}
\begin{Highlighting}[]
\PreprocessorTok{\#include }\ImportTok{\textless{}embree4/rtcore.h\textgreater{}}

\KeywordTok{struct}\NormalTok{ RTC\_ALIGN}\OperatorTok{(}\DecValTok{16}\OperatorTok{)}\NormalTok{ RTCPointQuery}
\OperatorTok{\{}
  \CommentTok{// location of the query}
  \DataTypeTok{float}\NormalTok{ x}\OperatorTok{;}
  \DataTypeTok{float}\NormalTok{ y}\OperatorTok{;}
  \DataTypeTok{float}\NormalTok{ z}\OperatorTok{;}

  \CommentTok{// radius and time of the query}
  \DataTypeTok{float}\NormalTok{ radius}\OperatorTok{;}
  \DataTypeTok{float}\NormalTok{ time}\OperatorTok{;}
\OperatorTok{\};}

\DataTypeTok{void}\NormalTok{ rtcPointQuery}\OperatorTok{(}
\NormalTok{  RTCScene scene}\OperatorTok{,}
  \KeywordTok{struct}\NormalTok{ RTCPointQuery}\OperatorTok{*}\NormalTok{ query}\OperatorTok{,}
  \KeywordTok{struct}\NormalTok{ RTCPointQueryContext}\OperatorTok{*}\NormalTok{ context}\OperatorTok{,}
  \KeywordTok{struct}\NormalTok{ RTCPointQueryFunction}\OperatorTok{*}\NormalTok{ queryFunc}\OperatorTok{,}
  \DataTypeTok{void}\OperatorTok{*}\NormalTok{ userPtr}
\OperatorTok{);}
\end{Highlighting}
\end{Shaded}

\subsubsection{DESCRIPTION}\label{description}

The \texttt{rtcPointQuery} function traverses the BVH using a
\texttt{RTCPointQuery} object (\texttt{query} argument) and calls a user
defined callback function (e.g \texttt{queryFunc} argument) for each
primitive of the scene (\texttt{scene} argument) that intersects the
query domain.

The user has to initialize the query location (\texttt{x}, \texttt{y}
and \texttt{z} member) and query radius in the range \([0, \infty]\). If
the scene contains motion blur geometries, also the query time
(\texttt{time} member) must be initialized to a value in the range
\([0, 1]\).

Further, a \texttt{RTCPointQueryContext} (\texttt{context} argument)
must be created and initialized. It contains ID and transformation
information of the instancing hierarchy if (multilevel-)instancing is
used. See \hyperref[rtcinitpointquerycontext]{rtcInitPointQueryContext}
for further information.

For every primitive that intersects the query domain, the callback
function (\texttt{queryFunc} argument) is called, in which distance
computations to the primitive can be implemented. The user will be
provided with the primID and geomID of the according primitive, however,
the geometry information (e.g.~triangle index and vertex data) has to be
determined manually. The \texttt{userPtr} argument can be used to input
geometry data of the scene or output results of the point query
(e.g.~closest point currently found on surface geometry (see tutorial
{[}ClosestPoint{]})).

The parameter \texttt{queryFunc} is optional and can be NULL, in which
case the callback function is not invoked. However, a callback function
can still get attached to a specific \texttt{RTCGeometry} object using
\hyperref[rtcsetgeometrypointqueryfunction]{rtcSetGeometryPointQueryFunction}.
If a callback function is attached to a geometry and (a potentially
different) callback function is passed as an argument to
\texttt{rtcPointQuery}, both functions are called for the primitives of
the according geometries.

The query radius can be decreased inside the callback function, which
allows to efficiently cull parts of the scene during BVH traversal.
Increasing the query radius and modifying time or location of the query
will result in undefined behaviour.

The callback function will be called for all primitives in a leaf node
of the BVH even if the primitive is outside the query domain, since
Embree does not gather geometry information of primitives internally.

Point queries can be used with (multilevel)-instancing. However, care
has to be taken when the instance transformation contains anisotropic
scaling or sheering. In these cases distance computations have to be
performed in world space to ensure correctness and the ellipsoidal query
domain (in instance space) will be approximated with its axis aligned
bounding box internally. Therefore, the callback function might be
invoked even for primitives in inner BVH nodes that do not intersect the
query domain. See
\hyperref[rtcsetgeometrypointqueryfunction]{rtcSetGeometryPointQueryFunction}
for details.

The point query structure must be aligned to 16 bytes.

\subsubsection{SUPPORTED PRIMITIVES}\label{supported-primitives}

Currently, all primitive types are supported by the point query API
except of points (see
\hyperref[rtc_geometry_type_point]{RTC\_GEOMETRY\_TYPE\_POINT}), curves
(see \hyperref[rtc_geometry_type_curve]{RTC\_GEOMETRY\_TYPE\_CURVE}) and
sudivision surfaces (see {[}RTC\_GEOMETRY\_SUBDIVISION{]}).

\subsubsection{EXIT STATUS}\label{exit-status}

For performance reasons this function does not do any error checks, thus
will not set any error flags on failure.

\subsubsection{SEE ALSO}\label{see-also}

\hyperref[rtcsetgeometrypointqueryfunction]{rtcSetGeometryPointQueryFunction},
\hyperref[rtcinitpointquerycontext]{rtcInitPointQueryContext}

\pagebreak

\section{rtcCollide}\label{rtccollide}

\subsubsection{NAME}\label{name}

\begin{Shaded}
\begin{Highlighting}[]
\NormalTok{rtcCollide }\OperatorTok{{-}}\NormalTok{ intersects one BVH with another}
\end{Highlighting}
\end{Shaded}

\subsubsection{SYNOPSIS}\label{synopsis}

\begin{Shaded}
\begin{Highlighting}[]
\PreprocessorTok{\#include }\ImportTok{\textless{}embree4/rtcore.h\textgreater{}}

\KeywordTok{struct}\NormalTok{ RTCCollision }\OperatorTok{\{}
  \DataTypeTok{unsigned} \DataTypeTok{int}\NormalTok{ geomID0}\OperatorTok{,}\NormalTok{ primID0}\OperatorTok{;}
  \DataTypeTok{unsigned} \DataTypeTok{int}\NormalTok{ geomID1}\OperatorTok{,}\NormalTok{ primID1}\OperatorTok{;}
\OperatorTok{\};}

\KeywordTok{typedef} \DataTypeTok{void} \OperatorTok{(*}\NormalTok{RTCCollideFunc}\OperatorTok{)} \OperatorTok{(}
  \DataTypeTok{void}\OperatorTok{*}\NormalTok{ userPtr}\OperatorTok{,}
\NormalTok{  RTCCollision}\OperatorTok{*}\NormalTok{ collisions}\OperatorTok{,}
  \DataTypeTok{size\_t}\NormalTok{ num\_collisions}\OperatorTok{);}

\DataTypeTok{void}\NormalTok{ rtcCollide }\OperatorTok{(}
\NormalTok{    RTCScene hscene0}\OperatorTok{,} 
\NormalTok{    RTCScene hscene1}\OperatorTok{,} 
\NormalTok{    RTCCollideFunc callback}\OperatorTok{,} 
    \DataTypeTok{void}\OperatorTok{*}\NormalTok{ userPtr}
\OperatorTok{);}
\end{Highlighting}
\end{Shaded}

\subsubsection{DESCRIPTION}\label{description}

The \texttt{rtcCollide} function intersects the BVH of \texttt{hscene0}
with the BVH of scene \texttt{hscene1} and calls a user defined callback
function (e.g \texttt{callback} argument) for each pair of intersecting
primitives between the two scenes. A user defined data pointer
(\texttt{userPtr} argument) can also be passed in.

For every pair of primitives that may intersect each other, the callback
function (\texttt{callback} argument) is called. The user will be
provided with the primID's and geomID's of multiple potentially
intersecting primitive pairs. Currently, only scene entirely composed of
user geometries are supported, thus the user is expected to implement a
primitive/primitive intersection to filter out false positives in the
callback function. The \texttt{userPtr} argument can be used to input
geometry data of the scene or output results of the intersection query.

\subsubsection{SUPPORTED PRIMITIVES}\label{supported-primitives}

Currently, the only supported type is the user geometry type (see
\hyperref[rtc_geometry_type_user]{RTC\_GEOMETRY\_TYPE\_USER}).

\subsubsection{EXIT STATUS}\label{exit-status}

On failure an error code is set that can be queried using
\texttt{rtcGetDeviceError}.

\subsubsection{SEE ALSO}\label{see-also}

\pagebreak

\section{rtcNewBVH}\label{rtcnewbvh}

\subsubsection{NAME}\label{name}

\begin{Shaded}
\begin{Highlighting}[]
\NormalTok{rtcNewBVH }\OperatorTok{{-}}\NormalTok{ creates a }\KeywordTok{new}\NormalTok{ BVH object}
\end{Highlighting}
\end{Shaded}

\subsubsection{SYNOPSIS}\label{synopsis}

\begin{Shaded}
\begin{Highlighting}[]
\PreprocessorTok{\#include }\ImportTok{\textless{}embree4/rtcore.h\textgreater{}}

\NormalTok{RTCBVH rtcNewBVH}\OperatorTok{(}\NormalTok{RTCDevice device}\OperatorTok{);}
\end{Highlighting}
\end{Shaded}

\subsubsection{DESCRIPTION}\label{description}

This function creates a new BVH object and returns a handle to this BVH.
The BVH object is reference counted with an initial reference count of
1. The handle can be released using the \texttt{rtcReleaseBVH} API call.

The BVH object can be used to build a BVH in a user-specified format
over user-specified primitives. See the documentation of the
\texttt{rtcBuildBVH} call for more details.

\subsubsection{EXIT STATUS}\label{exit-status}

On failure an error code is set that can be queried using
\texttt{rtcGetDeviceError}.

\subsubsection{SEE ALSO}\label{see-also}

\hyperref[rtcretainbvh]{rtcRetainBVH},
\hyperref[rtcreleasebvh]{rtcReleaseBVH},
\hyperref[rtcbuildbvh]{rtcBuildBVH}

\pagebreak

\section{rtcRetainBVH}\label{rtcretainbvh}

\subsubsection{NAME}\label{name}

\begin{Shaded}
\begin{Highlighting}[]
\NormalTok{rtcRetainBVH }\OperatorTok{{-}}\NormalTok{ increments the BVH reference count}
\end{Highlighting}
\end{Shaded}

\subsubsection{SYNOPSIS}\label{synopsis}

\begin{Shaded}
\begin{Highlighting}[]
\PreprocessorTok{\#include }\ImportTok{\textless{}embree4/rtcore.h\textgreater{}}

\DataTypeTok{void}\NormalTok{ rtcRetainBVH}\OperatorTok{(}\NormalTok{RTCBVH bvh}\OperatorTok{);}
\end{Highlighting}
\end{Shaded}

\subsubsection{DESCRIPTION}\label{description}

BVH objects are reference counted. The \texttt{rtcRetainBVH} function
increments the reference count of the passed BVH object (\texttt{bvh}
argument). This function together with \texttt{rtcReleaseBVH} allows to
use the internal reference counting in a C++ wrapper class to handle the
ownership of the object.

\subsubsection{EXIT STATUS}\label{exit-status}

On failure an error code is set that can be queried using
\texttt{rtcGetDeviceError}.

\subsubsection{SEE ALSO}\label{see-also}

\hyperref[rtcnewbvh]{rtcNewBVH}, \hyperref[rtcreleasebvh]{rtcReleaseBVH}

\pagebreak

\section{rtcReleaseBVH}\label{rtcreleasebvh}

\subsubsection{NAME}\label{name}

\begin{Shaded}
\begin{Highlighting}[]
\NormalTok{rtcReleaseBVH }\OperatorTok{{-}}\NormalTok{ decrements the BVH reference count}
\end{Highlighting}
\end{Shaded}

\subsubsection{SYNOPSIS}\label{synopsis}

\begin{Shaded}
\begin{Highlighting}[]
\PreprocessorTok{\#include }\ImportTok{\textless{}embree4/rtcore.h\textgreater{}}

\DataTypeTok{void}\NormalTok{ rtcReleaseBVH}\OperatorTok{(}\NormalTok{RTCBVH bvh}\OperatorTok{);}
\end{Highlighting}
\end{Shaded}

\subsubsection{DESCRIPTION}\label{description}

BVH objects are reference counted. The \texttt{rtcReleaseBVH} function
decrements the reference count of the passed BVH object (\texttt{bvh}
argument). When the reference count falls to 0, the BVH gets destroyed.

\subsubsection{EXIT STATUS}\label{exit-status}

On failure an error code is set that can be queried using
\texttt{rtcGetDeviceError}.

\subsubsection{SEE ALSO}\label{see-also}

\hyperref[rtcnewbvh]{rtcNewBVH}, \hyperref[rtcretainbvh]{rtcRetainBVH}

\pagebreak

\section{rtcBuildBVH}\label{rtcbuildbvh}

\subsubsection{NAME}\label{name}

\begin{Shaded}
\begin{Highlighting}[]
\NormalTok{rtcBuildBVH }\OperatorTok{{-}}\NormalTok{ builds a BVH}
\end{Highlighting}
\end{Shaded}

\subsubsection{SYNOPSIS}\label{synopsis}

\begin{Shaded}
\begin{Highlighting}[]
\PreprocessorTok{\#include }\ImportTok{\textless{}embree4/rtcore.h\textgreater{}}

\KeywordTok{struct}\NormalTok{ RTC\_ALIGN}\OperatorTok{(}\DecValTok{32}\OperatorTok{)}\NormalTok{ RTCBuildPrimitive}
\OperatorTok{\{}
  \DataTypeTok{float}\NormalTok{ lower\_x}\OperatorTok{,}\NormalTok{ lower\_y}\OperatorTok{,}\NormalTok{ lower\_z}\OperatorTok{;} 
  \DataTypeTok{unsigned} \DataTypeTok{int}\NormalTok{ geomID}\OperatorTok{;}
  \DataTypeTok{float}\NormalTok{ upper\_x}\OperatorTok{,}\NormalTok{ upper\_y}\OperatorTok{,}\NormalTok{ upper\_z}\OperatorTok{;}
  \DataTypeTok{unsigned} \DataTypeTok{int}\NormalTok{ primID}\OperatorTok{;}
\OperatorTok{\};}

\KeywordTok{typedef} \DataTypeTok{void}\OperatorTok{*} \OperatorTok{(*}\NormalTok{RTCCreateNodeFunction}\OperatorTok{)} \OperatorTok{(}
\NormalTok{  RTCThreadLocalAllocator allocator}\OperatorTok{,}
  \DataTypeTok{unsigned} \DataTypeTok{int}\NormalTok{ childCount}\OperatorTok{,}
  \DataTypeTok{void}\OperatorTok{*}\NormalTok{ userPtr}
\OperatorTok{);}

\KeywordTok{typedef} \DataTypeTok{void} \OperatorTok{(*}\NormalTok{RTCSetNodeChildrenFunction}\OperatorTok{)} \OperatorTok{(}
  \DataTypeTok{void}\OperatorTok{*}\NormalTok{ nodePtr}\OperatorTok{,}
  \DataTypeTok{void}\OperatorTok{**}\NormalTok{ children}\OperatorTok{,}
  \DataTypeTok{unsigned} \DataTypeTok{int}\NormalTok{ childCount}\OperatorTok{,}
  \DataTypeTok{void}\OperatorTok{*}\NormalTok{ userPtr}
\OperatorTok{);}

\KeywordTok{typedef} \DataTypeTok{void} \OperatorTok{(*}\NormalTok{RTCSetNodeBoundsFunction}\OperatorTok{)} \OperatorTok{(}
  \DataTypeTok{void}\OperatorTok{*}\NormalTok{ nodePtr}\OperatorTok{,}
  \AttributeTok{const} \KeywordTok{struct}\NormalTok{ RTCBounds}\OperatorTok{**}\NormalTok{ bounds}\OperatorTok{,}
  \DataTypeTok{unsigned} \DataTypeTok{int}\NormalTok{ childCount}\OperatorTok{,}
  \DataTypeTok{void}\OperatorTok{*}\NormalTok{ userPtr}
\OperatorTok{);}

\KeywordTok{typedef} \DataTypeTok{void}\OperatorTok{*} \OperatorTok{(*}\NormalTok{RTCCreateLeafFunction}\OperatorTok{)} \OperatorTok{(}
\NormalTok{  RTCThreadLocalAllocator allocator}\OperatorTok{,}
  \AttributeTok{const} \KeywordTok{struct}\NormalTok{ RTCBuildPrimitive}\OperatorTok{*}\NormalTok{ primitives}\OperatorTok{,}
  \DataTypeTok{size\_t}\NormalTok{ primitiveCount}\OperatorTok{,}
  \DataTypeTok{void}\OperatorTok{*}\NormalTok{ userPtr}
\OperatorTok{);}

\KeywordTok{typedef} \DataTypeTok{void} \OperatorTok{(*}\NormalTok{RTCSplitPrimitiveFunction}\OperatorTok{)} \OperatorTok{(}
  \AttributeTok{const} \KeywordTok{struct}\NormalTok{ RTCBuildPrimitive}\OperatorTok{*}\NormalTok{ primitive}\OperatorTok{,}
  \DataTypeTok{unsigned} \DataTypeTok{int}\NormalTok{ dimension}\OperatorTok{,}
  \DataTypeTok{float}\NormalTok{ position}\OperatorTok{,}
  \KeywordTok{struct}\NormalTok{ RTCBounds}\OperatorTok{*}\NormalTok{ leftBounds}\OperatorTok{,}
  \KeywordTok{struct}\NormalTok{ RTCBounds}\OperatorTok{*}\NormalTok{ rightBounds}\OperatorTok{,}
  \DataTypeTok{void}\OperatorTok{*}\NormalTok{ userPtr}
\OperatorTok{);}

\KeywordTok{typedef} \DataTypeTok{bool} \OperatorTok{(*}\NormalTok{RTCProgressMonitorFunction}\OperatorTok{)(}
  \DataTypeTok{void}\OperatorTok{*}\NormalTok{ userPtr}\OperatorTok{,} \DataTypeTok{double}\NormalTok{ n}
\OperatorTok{);}

\KeywordTok{enum}\NormalTok{ RTCBuildFlags}
\OperatorTok{\{}
\NormalTok{  RTC\_BUILD\_FLAG\_NONE}\OperatorTok{,}
\NormalTok{  RTC\_BUILD\_FLAG\_DYNAMIC}
\OperatorTok{\};}

\KeywordTok{struct}\NormalTok{ RTCBuildArguments}
\OperatorTok{\{}
  \DataTypeTok{size\_t}\NormalTok{ byteSize}\OperatorTok{;}

  \KeywordTok{enum}\NormalTok{ RTCBuildQuality buildQuality}\OperatorTok{;}
  \KeywordTok{enum}\NormalTok{ RTCBuildFlags buildFlags}\OperatorTok{;}
  \DataTypeTok{unsigned} \DataTypeTok{int}\NormalTok{ maxBranchingFactor}\OperatorTok{;}
  \DataTypeTok{unsigned} \DataTypeTok{int}\NormalTok{ maxDepth}\OperatorTok{;}
  \DataTypeTok{unsigned} \DataTypeTok{int}\NormalTok{ sahBlockSize}\OperatorTok{;}
  \DataTypeTok{unsigned} \DataTypeTok{int}\NormalTok{ minLeafSize}\OperatorTok{;}
  \DataTypeTok{unsigned} \DataTypeTok{int}\NormalTok{ maxLeafSize}\OperatorTok{;}
  \DataTypeTok{float}\NormalTok{ traversalCost}\OperatorTok{;}
  \DataTypeTok{float}\NormalTok{ intersectionCost}\OperatorTok{;}

\NormalTok{  RTCBVH bvh}\OperatorTok{;}
  \KeywordTok{struct}\NormalTok{ RTCBuildPrimitive}\OperatorTok{*}\NormalTok{ primitives}\OperatorTok{;}
  \DataTypeTok{size\_t}\NormalTok{ primitiveCount}\OperatorTok{;}
  \DataTypeTok{size\_t}\NormalTok{ primitiveArrayCapacity}\OperatorTok{;}
  
\NormalTok{  RTCCreateNodeFunction createNode}\OperatorTok{;}
\NormalTok{  RTCSetNodeChildrenFunction setNodeChildren}\OperatorTok{;}
\NormalTok{  RTCSetNodeBoundsFunction setNodeBounds}\OperatorTok{;}
\NormalTok{  RTCCreateLeafFunction createLeaf}\OperatorTok{;}
\NormalTok{  RTCSplitPrimitiveFunction splitPrimitive}\OperatorTok{;}
\NormalTok{  RTCProgressMonitorFunction buildProgress}\OperatorTok{;}
  \DataTypeTok{void}\OperatorTok{*}\NormalTok{ userPtr}\OperatorTok{;}
\OperatorTok{\};}

\KeywordTok{struct}\NormalTok{ RTCBuildArguments rtcDefaultBuildArguments}\OperatorTok{();}

\DataTypeTok{void}\OperatorTok{*}\NormalTok{ rtcBuildBVH}\OperatorTok{(}
  \AttributeTok{const} \KeywordTok{struct}\NormalTok{ RTCBuildArguments}\OperatorTok{*}\NormalTok{ args}
\OperatorTok{);}
\end{Highlighting}
\end{Shaded}

\subsubsection{DESCRIPTION}\label{description}

The \texttt{rtcBuildBVH} function can be used to build a BVH in a
user-defined format over arbitrary primitives. All arguments to the
function are provided through the \texttt{RTCBuildArguments} structure.
The first member of that structure must be set to the size of the
structure in bytes (\texttt{bytesSize} member) which allows future
extensions of the structure. It is recommended to initialize the build
arguments structure using the \texttt{rtcDefaultBuildArguments}
function.

The \texttt{rtcBuildBVH} function gets passed the BVH to build
(\texttt{bvh} member), the array of primitives (\texttt{primitives}
member), the capacity of that array (\texttt{primitiveArrayCapacity}
member), the number of primitives stored inside the array
(\texttt{primitiveCount} member), callback function pointers, and a
user-defined pointer (\texttt{userPtr} member) that is passed to all
callback functions when invoked. The \texttt{primitives} array can be
freed by the application after the BVH is built. All callback functions
are typically called from multiple threads, thus their implementation
must be thread-safe.

Four callback functions must be registered, which are invoked during
build to create BVH nodes (\texttt{createNode} member), to set the
pointers to all children (\texttt{setNodeChildren} member), to set the
bounding boxes of all children (\texttt{setNodeBounds} member), and to
create a leaf node (\texttt{createLeaf} member).

The function pointer to the primitive split function
(\texttt{splitPrimitive} member) may be \texttt{NULL}, however, then no
spatial splitting in high quality mode is possible. The function pointer
used to report the build progress (\texttt{buildProgress} member) is
optional and may also be \texttt{NULL}.

Further, some build settings are passed to configure the BVH build.
Using the build quality settings (\texttt{buildQuality} member), one can
select between a faster, low quality build which is good for dynamic
scenes, and a standard quality build for static scenes. One can also
specify the desired maximum branching factor of the BVH
(\texttt{maxBranchingFactor} member), the maximum depth the BVH should
have (\texttt{maxDepth} member), the block size for the SAH heuristic
(\texttt{sahBlockSize} member), the minimum and maximum leaf size
(\texttt{minLeafSize} and \texttt{maxLeafSize} member), and the
estimated costs of one traversal step and one primitive intersection
(\texttt{traversalCost} and \texttt{intersectionCost} members). When
enabling the \texttt{RTC\_BUILD\_FLAG\_DYNAMIC} build flags
(\texttt{buildFlags} member), re-build performance for dynamic scenes is
improved at the cost of higher memory requirements.

To spatially split primitives in high quality mode, the builder needs
extra space at the end of the build primitive array to store split
primitives. The total capacity of the build primitive array is passed
using the \texttt{primitiveArrayCapacity} member, and should be about
twice the number of primitives when using spatial splits.

The \texttt{RTCCreateNodeFunc} and \texttt{RTCCreateLeafFunc} callbacks
are passed a thread local allocator object that should be used for fast
allocation of nodes using the \texttt{rtcThreadLocalAlloc} function. We
strongly recommend using this allocation mechanism, as alternative
approaches like standard \texttt{malloc} can be over 10× slower. The
allocator object passed to the create callbacks may be used only inside
the current thread. Memory allocated using \texttt{rtcThreadLocalAlloc}
is automatically freed when the \texttt{RTCBVH} object is deleted. If
you use your own memory allocation scheme you have to free the memory
yourself when the \texttt{RTCBVH} object is no longer used.

The \texttt{RTCCreateNodeFunc} callback additionally gets the number of
children for this node in the range from 2 to
\texttt{maxBranchingFactor} (\texttt{childCount} argument).

The \texttt{RTCSetNodeChildFunc} callback function gets a pointer to the
node as input (\texttt{nodePtr} argument), an array of pointers to the
children (\texttt{childPtrs} argument), and the size of this array
(\texttt{childCount} argument).

The \texttt{RTCSetNodeBoundsFunc} callback function gets a pointer to
the node as input (\texttt{nodePtr} argument), an array of pointers to
the bounding boxes of the children (\texttt{bounds} argument), and the
size of this array (\texttt{childCount} argument).

The \texttt{RTCCreateLeafFunc} callback additionally gets an array of
primitives as input (\texttt{primitives} argument), and the size of this
array (\texttt{primitiveCount} argument). The callback should read the
\texttt{geomID} and \texttt{primID} members from the passed primitives
to construct the leaf.

The \texttt{RTCSplitPrimitiveFunc} callback is invoked in high quality
mode to split a primitive (\texttt{primitive} argument) at the specified
position (\texttt{position} argument) and dimension (\texttt{dimension}
argument). The callback should return bounds of the clipped left and
right parts of the primitive (\texttt{leftBounds} and
\texttt{rightBounds} arguments).

The \texttt{RTCProgressMonitorFunction} callback function is called with
the estimated completion rate \texttt{n} in the range \([0,1]\).
Returning \texttt{true} from the callback lets the build continue;
returning \texttt{false} cancels the build.

\subsubsection{EXIT STATUS}\label{exit-status}

On failure an error code is set that can be queried using
\texttt{rtcGetDeviceError}.

\subsubsection{SEE ALSO}\label{see-also}

\hyperref[rtcnewbvh]{rtcNewBVH}

\pagebreak

\section{RTCQuaternionDecomposition}\label{rtcquaterniondecomposition}

\subsubsection{NAME}\label{name}

\begin{Shaded}
\begin{Highlighting}[]
\NormalTok{RTCQuaternionDecomposition }\OperatorTok{{-}}\NormalTok{ structure that represents a quaternion}
\NormalTok{  decomposition of an affine transformation}
\end{Highlighting}
\end{Shaded}

\subsubsection{SYNOPSIS}\label{synopsis}

\begin{Shaded}
\begin{Highlighting}[]
\KeywordTok{struct}\NormalTok{ RTCQuaternionDecomposition}
\OperatorTok{\{}
  \DataTypeTok{float}\NormalTok{ scale\_x}\OperatorTok{,}\NormalTok{ scale\_y}\OperatorTok{,}\NormalTok{ scale\_z}\OperatorTok{;}
  \DataTypeTok{float}\NormalTok{ skew\_xy}\OperatorTok{,}\NormalTok{ skew\_xz}\OperatorTok{,}\NormalTok{ skew\_yz}\OperatorTok{;}
  \DataTypeTok{float}\NormalTok{ shift\_x}\OperatorTok{,}\NormalTok{ shift\_y}\OperatorTok{,}\NormalTok{ shift\_z}\OperatorTok{;}
  \DataTypeTok{float}\NormalTok{ quaternion\_r}\OperatorTok{,}\NormalTok{ quaternion\_i}\OperatorTok{,}\NormalTok{ quaternion\_j}\OperatorTok{,}\NormalTok{ quaternion\_k}\OperatorTok{;}
  \DataTypeTok{float}\NormalTok{ translation\_x}\OperatorTok{,}\NormalTok{ translation\_y}\OperatorTok{,}\NormalTok{ translation\_z}\OperatorTok{;}
\OperatorTok{\};}
\end{Highlighting}
\end{Shaded}

\subsubsection{DESCRIPTION}\label{description}

The struct \texttt{RTCQuaternionDecomposition} represents an affine
transformation decomposed into three parts. An upper triangular
scaling/skew/shift matrix

\[
S = \left( \begin{array}{cccc}
scale_x & skew_{xy} & skew_{xz} & shift_x \\ 
0 & scale_y & skew_{yz} & shift_y \\ 
0 & 0 & scale_z & shift_z \\ 
0 & 0 & 0 & 1 \\ 
\end{array} \right),
\]

a translation matrix

\[
T = \left( \begin{array}{cccc}
1 & 0 & 0 & translation_x \\ 
0 & 1 & 0 & translation_y \\ 
0 & 0 & 1 & translation_z \\ 
0 & 0 & 0 & 1 \\ 
\end{array} \right),
\]

and a rotation matrix \(R\), represented as a quaternion

\(quaternion_r + quaternion_i \ \mathbf{i} + quaternion_j \ \mathbf{i} + quaternion_k \ \mathbf{k}\)

where \(\mathbf{i}\), \(\mathbf{j}\) \(\mathbf{k}\) are the imaginary
quaternion units. The passed quaternion will be normalized internally.

The affine transformation matrix corresponding to a
\texttt{RTCQuaternionDecomposition} is \(TRS\) and a point
\(p = (p_x, p_y, p_z, 1)^T\) will be transformed as
\[p' = T \ R \ S \ p.\]

The functions \texttt{rtcInitQuaternionDecomposition},
\texttt{rtcQuaternionDecompositionSetQuaternion},
\texttt{rtcQuaternionDecompositionSetScale},
\texttt{rtcQuaternionDecompositionSetSkew},
\texttt{rtcQuaternionDecompositionSetShift}, and
\texttt{rtcQuaternionDecompositionSetTranslation} allow to set the
fields of the structure more conveniently.

\subsubsection{EXIT STATUS}\label{exit-status}

No error code is set by this function.

\subsubsection{SEE ALSO}\label{see-also}

\hyperref[rtcsetgeometrytransformquaternion]{rtcSetGeometryTransformQuaternion},
\hyperref[rtcinitquaterniondecomposition]{rtcInitQuaternionDecomposition}

\pagebreak

\section{rtcInitQuaternionDecomposition}\label{rtcinitquaterniondecomposition}

\subsubsection{NAME}\label{name}

\begin{Shaded}
\begin{Highlighting}[]
\NormalTok{rtcInitQuaternionDecomposition }\OperatorTok{{-}}\NormalTok{ initializes quaternion decomposition}
\end{Highlighting}
\end{Shaded}

\subsubsection{SYNOPSIS}\label{synopsis}

\begin{Shaded}
\begin{Highlighting}[]
\DataTypeTok{void}\NormalTok{ rtcInitQuaternionDecomposition}\OperatorTok{(}
  \KeywordTok{struct}\NormalTok{ RTCQuaternionDecomposition}\OperatorTok{*}\NormalTok{ qd}
\OperatorTok{);}
\end{Highlighting}
\end{Shaded}

\subsubsection{DESCRIPTION}\label{description}

The \texttt{rtcInitQuaternionDecomposition} function initializes a
\texttt{RTCQuaternionDecomposition} structure to represent an identity
transformation.

\subsubsection{EXIT STATUS}\label{exit-status}

No error code is set by this function.

\subsubsection{SEE ALSO}\label{see-also}

\hyperref[rtcsetgeometrytransformquaternion]{rtcSetGeometryTransformQuaternion},
\hyperref[rtcquaterniondecomposition]{RTCQuaternionDecomposition}

\pagebreak

\chapter{CPU Performance
Recommendations}\label{cpu-performance-recommendations}

\section{MXCSR control and status
register}\label{mxcsr-control-and-status-register}

It is strongly recommended to have the \texttt{Flush\ to\ Zero} and
\texttt{Denormals\ are\ Zero} mode of the MXCSR control and status
register enabled for each thread before calling the
\texttt{rtcIntersect}-type and \texttt{rtcOccluded}-type functions.
Otherwise, under some circumstances special handling of denormalized
floating point numbers can significantly reduce application and Embree
performance. When using Embree together with the Intel® Threading
Building Blocks, it is sufficient to execute the following code at the
beginning of the application main thread (before the creation of the
\texttt{tbb::task\_scheduler\_init} object):

\begin{Shaded}
\begin{Highlighting}[]
\PreprocessorTok{\#include }\ImportTok{\textless{}xmmintrin.h\textgreater{}}
\PreprocessorTok{\#include }\ImportTok{\textless{}pmmintrin.h\textgreater{}}
\OperatorTok{...}
\NormalTok{\_MM\_SET\_FLUSH\_ZERO\_MODE}\OperatorTok{(}\NormalTok{\_MM\_FLUSH\_ZERO\_ON}\OperatorTok{);}
\NormalTok{\_MM\_SET\_DENORMALS\_ZERO\_MODE}\OperatorTok{(}\NormalTok{\_MM\_DENORMALS\_ZERO\_ON}\OperatorTok{);}
\end{Highlighting}
\end{Shaded}

If using a different tasking system, make sure each rendering thread has
the proper mode set.

\section{Thread Creation and Affinity
Settings}\label{thread-creation-and-affinity-settings}

Tasking systems like TBB create worker threads on demand, which will add
a runtime overhead for the very first \texttt{rtcCommitScene} call. In
case you want to benchmark the scene build time, you should start the
threads at application startup. You can let Embree start TBB threads by
passing \texttt{start\_threads=1} to the \texttt{cfg} parameter of
\texttt{rtcNewDevice}.

On machines with a high thread count (e.g.~dual-socket Xeon or Xeon Phi
machines), affinitizing TBB worker threads increases build and rendering
performance. You can let Embree affinitize TBB worker threads by passing
\texttt{set\_affinity=1} to the \texttt{cfg} parameter of
\texttt{rtcNewDevice}. By default, threads are not affinitized by Embree
with the exception of Xeon Phi Processors where they are affinitized by
default.

All Embree tutorials automatically start and affinitize TBB worker
threads by passing \texttt{start\_threads=1,set\_affinity=1} to
\texttt{rtcNewDevice}.

\section{Fast Coherent Rays}\label{fast-coherent-rays}

For getting the highest performance for highly coherent rays, e.g.
primary or hard shadow rays, it is recommended to use packets with
setting the \texttt{RTC\_RAY\_QUERY\_FLAG\_COHERENT} flag in the
\texttt{RTCIntersectArguments} struct passed to the
\texttt{rtcIntersect}/\texttt{rtcOccluded} calls. The rays inside each
packet should be grouped as coherent as possible.

\section{Huge Page Support}\label{huge-page-support}

It is recommended to use huge pages under Linux to increase rendering
performance. Embree supports 2MB huge pages under Windows, Linux, and
macOS. Under Linux huge page support is enabled by default, and under
Windows and macOS disabled by default. Huge page support can be enabled
in Embree by passing \texttt{hugepages=1} to \texttt{rtcNewDevice} or
disabled by passing \texttt{hugepages=0} to \texttt{rtcNewDevice}.

We recommend using 2MB huge pages with Embree under Linux as this
improves ray tracing performance by about 5-10\%. Under Windows using
huge pages requires the application to run in elevated mode which is a
security issue, thus likely not an option for most use cases. Under
macOS huge pages are rarely available as memory tends to get quickly
fragmented, thus we do not recommend using huge pages on macOS.

\subsection{Huge Pages under Linux}\label{huge-pages-under-linux}

Linux supports transparent huge pages and explicit huge pages. To enable
transparent huge page support under Linux, execute the following as
root:

\begin{Shaded}
\begin{Highlighting}[]
\NormalTok{echo always }\OperatorTok{\textgreater{}} \OperatorTok{/}\NormalTok{sys}\OperatorTok{/}\NormalTok{kernel}\OperatorTok{/}\NormalTok{mm}\OperatorTok{/}\NormalTok{transparent\_hugepage}\OperatorTok{/}\NormalTok{enabled}
\end{Highlighting}
\end{Shaded}

When transparent huge pages are enabled, the kernel tries to merge 4KB
pages to 2MB pages when possible as a background job. Many Linux
distributions have transparent huge pages enabled by default. See the
following webpage for more information on
\href{https://www.kernel.org/doc/Documentation/vm/transhuge.txt}{transparent
huge pages under Linux}. In this mode each application, including your
rendering application based on Embree, will automatically tend to use
huge pages.

Using transparent huge pages, the transitioning from 4KB to 2MB pages
might take some time. For that reason Embree also supports allocating
2MB pages directly when a huge page pool is configured. Such a pool can
be configured by writing some number of huge pages to allocate to
\texttt{/proc/sys/vm/nr\_overcommit\_hugepages} as root user. E.g. to
configure 2GB of address space for huge page allocation, execute the
following as root:

\begin{Shaded}
\begin{Highlighting}[]
\NormalTok{echo }\DecValTok{1000} \OperatorTok{\textgreater{}} \OperatorTok{/}\NormalTok{proc}\OperatorTok{/}\NormalTok{sys}\OperatorTok{/}\NormalTok{vm}\OperatorTok{/}\NormalTok{nr\_overcommit\_hugepages}
\end{Highlighting}
\end{Shaded}

See the following webpage for more information on
\href{https://www.kernel.org/doc/Documentation/vm/hugetlbpage.txt}{huge
pages under Linux}.

\subsection{Huge Pages under Windows}\label{huge-pages-under-windows}

To use huge pages under Windows, the current user must have the ``Lock
pages in memory'' (SeLockMemoryPrivilege) assigned. This can be
configured through the ``Local Security Policy'' application, by adding
a user to ``Local Policies'' -\textgreater{} ``User Rights Assignment''
-\textgreater{} ``Lock pages in memory''. You have to log out and in
again for this change to take effect.

Further, your application must be executed as an elevated process (``Run
as administrator'') and the ``SeLockMemoryPrivilege'' must be explicitly
enabled by your application. Example code on how to enable this
privilege can be found in the ``common/sys/alloc.cpp'' file of Embree.
Alternatively, Embree will try to enable this privilege when passing
\texttt{enable\_selockmemoryprivilege=1} to \texttt{rtcNewDevice}.
Further, huge pages should be enabled in Embree by passing
\texttt{hugepages=1} to \texttt{rtcNewDevice}.

When the system has been running for a while, physical memory gets
fragmented, which can slow down the allocation of huge pages
significantly under Windows.

\subsection{Huge Pages under macOS}\label{huge-pages-under-macos}

To use huge pages under macOS you have to pass \texttt{hugepages=1} to
\texttt{rtcNewDevice} to enable that feature in Embree.

When the system has been running for a while, physical memory gets
quickly fragmented, and causes huge page allocations to fail. For this
reason, huge pages are not very useful under macOS in practice.

\section{Avoid store-to-load forwarding issues with single
rays}\label{avoid-store-to-load-forwarding-issues-with-single-rays}

We recommend to use a single SSE store to set up the \texttt{org} and
\texttt{tnear} components, and a single SSE store to set up the
\texttt{dir} and \texttt{time} components of a single ray
(\texttt{RTCRay} type). Storing these values using scalar stores causes
a store-to-load forwarding penalty because Embree is reading these
components using SSE loads later on.

\pagebreak

\chapter{GPU Performance
Recommendations}\label{gpu-performance-recommendations}

\section{Low Code Complexity}\label{low-code-complexity}

As a general rule try to keep code complexity low, to avoid spill code
generation. To achieve this we recommend splitting your renderer into
separate kernels instead of using a single Uber kernel invokation.

Code can further get reduced by using SYCL specialization constants to
just enable rendering features required to render a given scene.

\section{Feature Flags}\label{feature-flags}

Use SYCL specialization constants and the feature flags (see section
\hyperref[rtcfeatureflags]{RTCFeatureFlags}) of the
\texttt{rtcIntersect1} and \texttt{rtcOccluded1} calls to JIT compile
minimal code. The passed feature flags should just contain features
required to render the current scene. If JIT compile times are an issue,
reduce the number of feature masks used and use JIT caching (see section
\hyperref[sycl-jit-caching]{SYCL JIT caching}).

\section{Inline Indirect Calls}\label{inline-indirect-calls}

Attaching user geometry and intersection filter callbacks to the
geometries of the scene is not supported in SYCL for performance
reasons.

Instead directly pass the user geometry and intersection filter callback
functions through the \texttt{RTCIntersectArguments} (and
\texttt{RTCOccludedArguments}) struct to \texttt{rtcIntersect1} (and
\texttt{rtcOccluded1}) API functions as in the following example:

\begin{Shaded}
\begin{Highlighting}[]
\NormalTok{RTC\_SYCL\_INDIRECTLY\_CALLABLE }\DataTypeTok{void}\NormalTok{ intersectionFilter}\OperatorTok{(}
  \AttributeTok{const}\NormalTok{ RTCFilterFunctionNArguments}\OperatorTok{*}\NormalTok{ args}
\OperatorTok{)} \OperatorTok{\{} \OperatorTok{...} \OperatorTok{\}}

\NormalTok{RTCIntersectArguments args}\OperatorTok{;}
\NormalTok{rtcInitIntersectArguments}\OperatorTok{(\&}\NormalTok{args}\OperatorTok{);}
\NormalTok{args}\OperatorTok{.}\NormalTok{filter }\OperatorTok{=}\NormalTok{ intersectionFilter}\OperatorTok{;}

\NormalTok{rtcIntersect1}\OperatorTok{(}\NormalTok{scene}\OperatorTok{,\&}\NormalTok{ray}\OperatorTok{,\&}\NormalTok{args}\OperatorTok{);}
\end{Highlighting}
\end{Shaded}

If the callback function is directly passed that way, the SYCL compiler
can inline the indirect call, which gives a huge performance benefit. Do
\emph{not} read a function pointer form some memory location and pass it
to \texttt{rtcIntersect1} (and \texttt{rtcOccluded1}) as this will also
prevent inlining.

\section{7 Bit Ray Mask}\label{bit-ray-mask}

Use just the lower 7 bits of the ray and geometry mask if possible, even
though Embree supports 32 bit ray masks for geometry masking. On the CPU
using any of the 32 bits yields the same performance, but the ray
tracing hardware only supports an 8 bit mask, thus Embree has to emulate
32 bit masking if used. For that reason the lower 7 mask bits are
hardware accelerated and fast, while the mask bits 7-31 require some
software intervention and using them reduces performance. To turn on 32
bit ray masks use the RTC\_FEATURE\_FLAG\_32\_BIT\_RAY\_MASK (see
section \hyperref[rtcfeatureflags]{RTCFeatureFlags}).

\section{Limit Motion Blur Motions}\label{limit-motion-blur-motions}

The motion blur implementation on SYCL has some limitations regarding
supported motion. Primitive motion should be maximally as large as a
small multiple of the primitive size, otherwise performance can degrade
a lot. If detailed geometry moves fast, best put the geometry into an
instance, and apply motion blur to the instance itself, which
efficiently allows larger motions. As a fallback, problematic scenes can
always still get rendered robustly on the CPU.

\section{Generic Pointers}\label{generic-pointers}

Embree uses standard C++ pointers in its implementation. SYCL might not
be able to detect the memory space these pointers refer to and has to
treat them as generic pointers which are not performing optimal. The
DPC++ compiler has advanced optimizations to infer the proper address
space to avoid usage of generic pointers.

However, if you still encounter the following warning during ahead of
time compilation of SYCL kernels, then loads from generic pointer are
present:

\begin{Shaded}
\begin{Highlighting}[]
\NormalTok{warning}\OperatorTok{:}\NormalTok{ Adding XX occurrences of additional control flow due to presence}
\NormalTok{         of generic address space operations in function YYY}\OperatorTok{.}
\end{Highlighting}
\end{Shaded}

To work around this issue we recommend:

\begin{itemize}
\item
  Do not use local memory inside kernels that trace rays. In this case
  the DPC++ compiler knows that no local memory pointer can exist and
  will optimize generic loads. As this is typically the case for
  renderers, generic pointer will typically not cause issues.
\item
  Indirectly callable functions may still cause problems, even if your
  kernel does not use local memory. Thus best use SYCL pointers like
  sycl::global\_ptr and sycl::private\_ptr in indirectly callable
  functions to avoid generic address space usage.
\item
  You can also enforce usage of global pointers using the following
  DPC++ compile flags:
  \texttt{-cl-intel-force-global-mem-allocation\ \ \ -cl-intel-no-local-to-generic}.
\end{itemize}

\chapter{Embree Tutorials}\label{embree-tutorials}

Embree comes with a set of tutorials aimed at helping users understand
how Embree can be used and extended. There is a very basic minimal that
can be compiled as both C and C++, which should get new users started
quickly. All other tutorials exist in an Intel® ISPC and C++ version to
demonstrate the two versions of the API. Look for files named
\texttt{tutorialname\_device.ispc} for the Intel® ISPC implementation of
the tutorial, and files named \texttt{tutorialname\_device.cpp} for the
single ray C++ version of the tutorial. To start the C++ version use the
\texttt{tutorialname} executables, to start the Intel® ISPC version use
the \texttt{tutorialname\_ispc} executables. All tutorials can print
available command line options using the \texttt{-\/-help} command line
parameter.

For all tutorials except minimal, you can select an initial camera using
the \texttt{-\/-vp} (camera position), \texttt{-\/-vi} (camera look-at
point), \texttt{-\/-vu} (camera up vector), and \texttt{-\/-fov}
(vertical field of view) command line parameters:

\begin{Shaded}
\begin{Highlighting}[]
\OperatorTok{./}\NormalTok{triangle\_geometry }\OperatorTok{{-}{-}}\NormalTok{vp }\DecValTok{10} \DecValTok{10} \DecValTok{10} \OperatorTok{{-}{-}}\NormalTok{vi }\DecValTok{0} \DecValTok{0} \DecValTok{0}
\end{Highlighting}
\end{Shaded}

You can select the initial window size using the \texttt{-\/-size}
command line parameter, or start the tutorials in full screen using the
\texttt{-\/-fullscreen} parameter:

\begin{Shaded}
\begin{Highlighting}[]
\OperatorTok{./}\NormalTok{triangle\_geometry }\OperatorTok{{-}{-}}\NormalTok{size }\DecValTok{1024} \DecValTok{1024}
\OperatorTok{./}\NormalTok{triangle\_geometry }\OperatorTok{{-}{-}}\NormalTok{fullscreen}
\end{Highlighting}
\end{Shaded}

The initialization string for the Embree device (\texttt{rtcNewDevice}
call) can be passed to the ray tracing core through the
\texttt{-\/-rtcore} command line parameter, e.g.:

\begin{Shaded}
\begin{Highlighting}[]
\OperatorTok{./}\NormalTok{triangle\_geometry }\OperatorTok{{-}{-}}\NormalTok{rtcore verbose}\OperatorTok{=}\DecValTok{2}\OperatorTok{,}\NormalTok{threads}\OperatorTok{=}\DecValTok{1}
\end{Highlighting}
\end{Shaded}

The navigation in the interactive display mode follows the camera orbit
model, where the camera revolves around the current center of interest.
With the left mouse button you can rotate around the center of interest
(the point initially set with \texttt{-\/-vi}). Holding Control pressed
while clicking the left mouse button rotates the camera around its
location. You can also use the arrow keys for navigation.

You can use the following keys:

\begin{description}
\tightlist
\item[F1]
Default shading
\item[F2]
Gray EyeLight shading
\item[F3]
Traces occlusion rays only.
\item[F4]
UV Coordinate visualization
\item[F5]
Geometry normal visualization
\item[F6]
Geometry ID visualization
\item[F7]
Geometry ID and Primitive ID visualization
\item[F8]
Simple shading with 16 rays per pixel for benchmarking.
\item[F9]
Switches to render cost visualization. Pressing again reduces
brightness.
\item[F10]
Switches to render cost visualization. Pressing again increases
brightness.
\item[f]
Enters or leaves full screen mode.
\item[c]
Prints camera parameters.
\item[ESC]
Exits the tutorial.
\item[q]
Exits the tutorial.
\end{description}

\section{Minimal}\label{minimal}

This tutorial is designed to get new users started with Embree. It can
be compiled as both C and C++. It demonstrates how to initialize a
device and scene, and how to intersect rays with the scene. There is no
image output to keep the tutorial as simple as possible.

\href{https://github.com/embree/embree/blob/master/tutorials/minimal/minimal.cpp}{Source
Code}

\section{Triangle Geometry}\label{triangle-geometry}

\href{https://github.com/embree/embree/blob/master/tutorials/triangle_geometry/triangle_geometry_device.cpp}{\includegraphics{images/triangle_geometry.jpg}}

This tutorial demonstrates the creation of a static cube and ground
plane using triangle meshes. It also demonstrates the use of the
\texttt{rtcIntersect1} and \texttt{rtcOccluded1} functions to render
primary visibility and hard shadows. The cube sides are colored based on
the ID of the hit primitive.

\href{https://github.com/embree/embree/blob/master/tutorials/triangle_geometry/triangle_geometry_device.cpp}{Source
Code}

\section{Dynamic Scene}\label{dynamic-scene}

\href{https://github.com/embree/embree/blob/master/tutorials/dynamic_scene/dynamic_scene_device.cpp}{\includegraphics{images/dynamic_scene.jpg}}

This tutorial demonstrates the creation of a dynamic scene, consisting
of several deforming spheres. Half of the spheres use the
\texttt{RTC\_BUILD\_QUALITY\_REFIT} geometry build quality, which allows
Embree to use a refitting strategy for these spheres, the other half
uses the \texttt{RTC\_BUILD\_QUALITY\_LOW} geometry build quality,
causing a high performance rebuild of their spatial data structure each
frame. The spheres are colored based on the ID of the hit sphere
geometry.

\href{https://github.com/embree/embree/blob/master/tutorials/dynamic_scene/dynamic_scene_device.cpp}{Source
Code}

\section{Multi Scene Geometry}\label{multi-scene-geometry}

\href{https://github.com/embree/embree/blob/master/tutorials/multiscene_geometry/multiscene_geometry_device.cpp}{\includegraphics{images/dynamic_scene.jpg}}

This tutorial demonstrates the creation of multiple scenes sharing the
same geometry objects. Here, three scenes are built. One with all the
dynamic spheres of the Dynamic Scene test and two others each with half.
The ground plane is shared by all three scenes. The space bar is used to
cycle the scene chosen for rendering.

\href{https://github.com/embree/embree/blob/master/tutorials/multiscene_geometry/multiscene_geometry_device.cpp}{Source
Code}

\section{User Geometry}\label{user-geometry}

\href{https://github.com/embree/embree/blob/master/tutorials/user_geometry/user_geometry_device.cpp}{\includegraphics{images/user_geometry.jpg}}

This tutorial shows the use of user-defined geometry, to re-implement
instancing, and to add analytic spheres. A two-level scene is created,
with a triangle mesh as ground plane, and several user geometries that
instance other scenes with a small number of spheres of different kinds.
The spheres are colored using the instance ID and geometry ID of the hit
sphere, to demonstrate how the same geometry instanced in different ways
can be distinguished.

\href{https://github.com/embree/embree/blob/master/tutorials/user_geometry/user_geometry_device.cpp}{Source
Code}

\section{Viewer}\label{viewer}

\href{https://github.com/embree/embree/blob/master/tutorials/viewer/viewer_device.cpp}{
\includegraphics{images/viewer.jpg}}

This tutorial demonstrates a simple OBJ viewer that traces primary
visibility rays only. A scene consisting of multiple meshes is created,
each mesh sharing the index and vertex buffer with the application. It
also demonstrates how to support additional per-vertex data, such as
shading normals.

You need to specify an OBJ file at the command line for this tutorial to
work:

\begin{Shaded}
\begin{Highlighting}[]
\OperatorTok{./}\NormalTok{viewer }\OperatorTok{{-}}\NormalTok{i model}\OperatorTok{.}\NormalTok{obj}
\end{Highlighting}
\end{Shaded}

\href{https://github.com/embree/embree/blob/master/tutorials/viewer/viewer_device.cpp}{Source
Code}

\section{Intersection Filter}\label{intersection-filter}

\href{https://github.com/embree/embree/blob/master/tutorials/intersection_filter/intersection_filter_device.cpp}{\includegraphics{images/intersection_filter.jpg}}

This tutorial demonstrates the use of filter callback functions to
efficiently implement transparent objects. The filter function used for
primary rays lets the ray pass through the geometry if it is entirely
transparent. Otherwise, the shading loop handles the transparency
properly, by potentially shooting secondary rays. The filter function
used for shadow rays accumulates the transparency of all surfaces along
the ray, and terminates traversal if an opaque occluder is hit.

\href{https://github.com/embree/embree/blob/master/tutorials/intersection_filter/intersection_filter_device.cpp}{Source
Code}

\section{Instanced Geometry}\label{instanced-geometry}

\href{https://github.com/embree/embree/blob/master/tutorials/instanced_geometry/instanced_geometry_device.cpp}{\includegraphics{images/instanced_geometry.jpg}}

This tutorial demonstrates the in-build instancing feature of Embree, by
instancing a number of other scenes built from triangulated spheres. The
spheres are again colored using the instance ID and geometry ID of the
hit sphere, to demonstrate how the same geometry instanced in different
ways can be distinguished.

\href{https://github.com/embree/embree/blob/master/tutorials/instanced_geometry/instanced_geometry_device.cpp}{Source
Code}

\section{Multi Level Instancing}\label{multi-level-instancing}

\href{https://github.com/embree/embree/blob/master/tutorials/multi_instanced_geometry/multi_instanced_geometry_device.cpp}{\includegraphics{images/multi_level_instancing.jpg}}

This tutorial demonstrates multi-level instancing, i.e., nesting
instances into instances. To enable the tutorial, set the compile-time
variable \texttt{EMBREE\_MAX\_INSTANCE\_LEVEL\_COUNT} to a value other
than the default 1. This variable is available in the code as
\texttt{RTC\_MAX\_INSTANCE\_LEVEL\_COUNT}.

The renderer uses a basic path tracing approach, and the image will
progressively refine over time. There are two levels of instances in
this scene: multiple instances of the same tree nest instances of a
twig. Intersections on up to \texttt{RTC\_MAX\_INSTANCE\_LEVEL\_COUNT}
nested levels of instances work out of the box. Users may obtain the
\emph{instance ID stack} for a given hitpoint from the \texttt{instID}
member. During shading, the instance ID stack is used to accumulate
normal transformation matrices for each hit. The tutorial visualizes
transformed normals as colors.

\href{https://github.com/embree/embree/blob/master/tutorials/multi_instanced_geometry/multi_instanced_geometry_device.cpp}{Source
Code}

\section{Path Tracer}\label{path-tracer}

\href{https://github.com/embree/embree/blob/master/tutorials/pathtracer/pathtracer_device.cpp}{
\includegraphics{images/pathtracer.jpg}}

This tutorial is a simple path tracer, based on the viewer tutorial.

You need to specify an OBJ file and light source at the command line for
this tutorial to work:

\begin{Shaded}
\begin{Highlighting}[]
\OperatorTok{./}\NormalTok{pathtracer }\OperatorTok{{-}}\NormalTok{i model}\OperatorTok{.}\NormalTok{obj }\OperatorTok{{-}{-}}\NormalTok{ambientlight }\DecValTok{1} \DecValTok{1} \DecValTok{1}
\end{Highlighting}
\end{Shaded}

As example models we provide the ``Austrian Imperial Crown'' model by
\href{http://www.loramel.net}{Martin Lubich} and the ``Asian Dragon''
model from the
\href{http://graphics.stanford.edu/data/3Dscanrep/}{Stanford 3D Scanning
Repository}.

\href{https://github.com/embree/models/releases/download/release/crown.zip}{crown.zip}

\href{https://github.com/embree/models/releases/download/release/asian_dragon.zip}{asian\_dragon.zip}

To render these models execute the following:

\begin{Shaded}
\begin{Highlighting}[]
\OperatorTok{./}\NormalTok{pathtracer }\OperatorTok{{-}}\NormalTok{c crown}\OperatorTok{/}\NormalTok{crown}\OperatorTok{.}\NormalTok{ecs}
\OperatorTok{./}\NormalTok{pathtracer }\OperatorTok{{-}}\NormalTok{c asian\_dragon}\OperatorTok{/}\NormalTok{asian\_dragon}\OperatorTok{.}\NormalTok{ecs}
\end{Highlighting}
\end{Shaded}

\href{https://github.com/embree/embree/blob/master/tutorials/pathtracer/pathtracer_device.cpp}{Source
Code}

\section{Hair}\label{hair}

\href{https://github.com/embree/embree/blob/master/tutorials/hair_geometry/hair_geometry_device.cpp}{\includegraphics{images/hair_geometry.jpg}}

This tutorial demonstrates the use of the hair geometry to render a
hairball.

\href{https://github.com/embree/embree/blob/master/tutorials/hair_geometry/hair_geometry_device.cpp}{Source
Code}

\section{Curve Geometry}\label{curve-geometry}

\href{https://github.com/embree/embree/blob/master/tutorials/curve_geometry/curve_geometry_device.cpp}{\includegraphics{images/curve_geometry.jpg}}

This tutorial demonstrates the use of the Linear Basis, B-Spline, and
Catmull-Rom curve geometries.

\href{https://github.com/embree/embree/blob/master/tutorials/curve_geometry/curve_geometry_device.cpp}{Source
Code}

\section{Subdivision Geometry}\label{subdivision-geometry}

\href{https://github.com/embree/embree/blob/master/tutorials/subdivision_geometry/subdivision_geometry_device.cpp}{\includegraphics{images/subdivision_geometry.jpg}}

This tutorial demonstrates the use of Catmull-Clark subdivision
surfaces.

\href{https://github.com/embree/embree/blob/master/tutorials/subdivision_geometry/subdivision_geometry_device.cpp}{Source
Code}

\section{Displacement Geometry}\label{displacement-geometry}

\href{https://github.com/embree/embree/blob/master/tutorials/displacement_geometry/displacement_geometry_device.cpp}{\includegraphics{images/displacement_geometry.jpg}}

This tutorial demonstrates the use of Catmull-Clark subdivision surfaces
with procedural displacement mapping using a constant edge tessellation
level.

\href{https://github.com/embree/embree/blob/master/tutorials/displacement_geometry/displacement_geometry_device.cpp}{Source
Code}

\section{Grid Geometry}\label{grid-geometry}

\href{https://github.com/embree/embree/tree/master/tutorials/grid_geometry}{\includegraphics{images/grid_geometry.jpg}}

This tutorial demonstrates the use of the memory efficient grid
primitive to handle highly tessellated and displaced geometry.

\href{https://github.com/embree/embree/tree/master/tutorials/grid_geometry}{Source
Code}

\section{Point Geometry}\label{point-geometry}

\href{https://github.com/embree/embree/blob/master/tutorials/point_geometry/point_geometry_device.cpp}{\includegraphics{images/point_geometry.jpg}}

This tutorial demonstrates the use of the three representations of point
geometry.

\href{https://github.com/embree/embree/blob/master/tutorials/point_geometry/point_geometry_device.cpp}{Source
Code}

\section{Motion Blur Geometry}\label{motion-blur-geometry}

\href{https://github.com/embree/embree/blob/master/tutorials/motion_blur_geometry/motion_blur_geometry_device.cpp}{\includegraphics{images/motion_blur_geometry.jpg}}

This tutorial demonstrates rendering of motion blur using the
multi-segment motion blur feature. Shown is motion blur of a triangle
mesh, quad mesh, subdivision surface, line segments, hair geometry,
Bézier curves, instantiated triangle mesh where the instance moves,
instantiated quad mesh where the instance and the quads move, and user
geometry.

The number of time steps used can be configured using the
\texttt{-\/-time-steps\ \textless{}int\textgreater{}} and
\texttt{-\/-time-steps2\ \textless{}int\textgreater{}} command line
parameters, and the geometry can be rendered at a specific time using
the the \texttt{-\/-time\ \textless{}float\textgreater{}} command line
parameter.

\href{https://github.com/embree/embree/blob/master/tutorials/motion_blur_geometry/motion_blur_geometry_device.cpp}{Source
Code}

\section{Quaternion Motion Blur}\label{quaternion-motion-blur}

\href{https://github.com/embree/embree/blob/master/tutorials/quaternion_motion_blur/quaternion_motion_blur_device.cpp}{\includegraphics{images/quaternion_motion_blur.jpg}}

This tutorial demonstrates rendering of motion blur using quaternion
interpolation. Shown is motion blur using spherical linear interpolation
of the rotational component of the instance transformation on the left
and simple linear interpolation of the instance transformation on the
right. The number of time steps can be modified as well.

\href{https://github.com/embree/embree/blob/master/tutorials/quaternion_motion_blur/quaternion_motion_blur_device.cpp}{Source
Code}

\section{Interpolation}\label{interpolation}

\href{https://github.com/embree/embree/blob/master/tutorials/interpolation/interpolation_device.cpp}{\includegraphics{images/interpolation.jpg}}

This tutorial demonstrates interpolation of user-defined per-vertex
data.

\href{https://github.com/embree/embree/blob/master/tutorials/interpolation/interpolation_device.cpp}{Source
Code}

\section{Closest Point}\label{closest-point}

\href{https://github.com/embree/embree/blob/master/tutorials/closest_point/closest_point_device.cpp}{\includegraphics{images/closest_point.jpg}}

This tutorial demonstrates a use-case of the point query API. The scene
consists of a simple collection of objects that are instanced and for
several point in the scene (red points) the closest point on the
surfaces of the scene are computed (white points). The closest point
functionality is implemented for Embree internal and for user-defined
instancing. The tutorial also illustrates how to handle instance
transformations that are not similarity transforms.

\href{https://github.com/embree/embree/blob/master/tutorials/closest_point/closest_point_device.cpp}{Source
Code}

\section{Voronoi}\label{voronoi}

\href{https://github.com/embree/embree/blob/master/tutorials/voronoi/voronoi_device.cpp}{
\includegraphics{images/voronoi.jpg}}

This tutorial demonstrates how to implement nearest neighbour lookups
using the point query API. Several colored points are located on a plane
and the corresponding voroni regions are illustrated.

\href{https://github.com/embree/embree/blob/master/tutorials/voronoi/voronoi_device.cpp}{Source
Code}

\section{Collision Detection}\label{collision-detection-1}

\href{https://github.com/embree/embree/blob/master/tutorials/collide/collide_device.cpp}{\includegraphics{images/collide.jpg}}

This tutorial demonstrates how to implement collision detection using
the collide API. A simple cloth solver is setup to collide with a
sphere.

The cloth can be reset with the \texttt{space} bar. The sim stepped once
with \texttt{n} and continuous simulation started and paused with
\texttt{p}.

\href{https://github.com/embree/embree/blob/master/tutorials/collide/collide_device.cpp}{Source
Code}

\section{BVH Builder}\label{bvh-builder}

This tutorial demonstrates how to use the templated hierarchy builders
of Embree to build a bounding volume hierarchy with a user-defined
memory layout using a high-quality SAH builder using spatial splits, a
standard SAH builder, and a very fast Morton builder.

\href{https://github.com/embree/embree/blob/master/tutorials/bvh_builder/bvh_builder_device.cpp}{Source
Code}

\section{BVH Access}\label{bvh-access}

This tutorial demonstrates how to access the internal triangle
acceleration structure build by Embree. Please be aware that the
internal Embree data structures might change between Embree updates.

\href{https://github.com/embree/embree/blob/master/tutorials/bvh_access/bvh_access.cpp}{Source
Code}

\section{Find Embree}\label{find-embree}

This tutorial demonstrates how to use the \texttt{FIND\_PACKAGE} CMake
feature to use an installed Embree. Under Linux and macOS the tutorial
finds the Embree installation automatically, under Windows the
\texttt{embree\_DIR} CMake variable must be set to the following folder
of the Embree installation:
\texttt{C:\textbackslash{}Program\ Files\textbackslash{}Intel\textbackslash{}Embree3}.

\href{https://github.com/embree/embree/blob/master/tutorials/find_embree/CMakeLists.txt}{Source
Code}

\section{Next Hit}\label{next-hit}

This tutorial demonstrates how to robustly enumerate all hits along the
ray using multiple ray queries and an intersection filter function. To
improve performance, the tutorial also supports collecting the next N
hits in a single ray query.

\href{https://github.com/embree/embree/blob/master/tutorials/next_hit/next_hit_device.cpp}{Source
Code}
